\documentclass[11pt, reqno]{amsart}
\usepackage[utf8]{inputenc}
\usepackage{amsmath}
\usepackage{amsthm}
\usepackage{amsfonts}
\usepackage{amssymb}
\usepackage{mathrsfs, mathtools, mathdots}
\usepackage{diagbox}
\DeclarePairedDelimiter \floor {\lfloor} {\rfloor}
\DeclareMathOperator{\Tr}{Tr}

\usepackage[top=27mm,bottom=27mm,left=27mm,right=27mm]{geometry}
\newtheorem{theorem}{Theorem}[section]
\newtheorem{conjecture}[theorem]{Conjecture}
\newtheorem{corollary}[theorem]{Corollary}
\newtheorem{problem}{Problem}
\newtheorem*{theorem*}{Theorem}
\newtheorem*{remark*}{Remark}
\newtheorem*{problem*}{Problem}
\newtheorem*{lemma*}{Lemma}
\newtheorem*{conjecture*}{Conjecture}
\newtheorem{lemma}[theorem]{Lemma}
\newtheorem{prop}[theorem]{Proposition}


\begin{document}

\title[Transcendence of special values of modular functions]{A note on transcendence of special values \\ of functions related to modularity}

\author[T. Bhowmik]{Tapas Bhowmik}

\author[S. Pathak]{Siddhi Pathak}

\address{Department of Mathematics, University of South Carolina, 1523 Greene Street, LeConte College, Columbia, South Carolina, USA 29208.}

\address{Chennai Mathematical Institute, H-1 SIPCOT IT Park, Siruseri, Kelambakkam, Tamil Nadu, India 603103.}

\email{tbhowmik@email.sc.edu}

\email{siddhi@cmi.ac.in}

\subjclass[2010]{11J91, 11F03, 11F11}

\keywords{Values of modular functions, values of modular forms, Nesterenko's theorem, Schneider's theorem}

\thanks{Research of the second author was partially supported by an INSPIRE Faculty fellowship.}



\begin{abstract}
In this note, we study the arithmetic nature of values of modular functions, meromorphic modular forms and meromorphic quasi-modular forms with respect to arbitrary congruence subgroups, that have algebraic Fourier coefficients. This approach unifies many of the known results, and leads to generalizations of the theorems of Schneider, Nesterenko and others. 
\end{abstract}

\maketitle

\section{\bf Introduction}
\bigskip

The study of the arithmetic nature of values of special transcendental functions at algebraic arguments has been a well-established theme in number theory. Continuing in this spirit, this note focuses on the transcendental nature and algebraic independence of values of functions arising in the modular world, such as modular functions, modular forms and quasi-modular forms (see Section 2 for definition of the functions appearing below). The genesis of this study can be traced back to a 1937 theorem of Schneider \cite{schneider}, namely,
\begin{theorem}[Schneider]\label{schneider}
    If $\tau\in \mathbb{H}$ is algebraic but not imaginary quadratic, then $j(\tau)$ is transcendental.
\end{theorem}
It is known from the theory of complex multiplication that if $\tau\in\mathbb{H}$ generates an imaginary quadratic field ($\tau$ is a CM point), then $j(\tau)$ is an algebraic number. Therefore, Schneider's theorem translates to the statement: if $\tau$ is algebraic, then $j(\tau)$ is algebraic if and only if $\tau$ is CM. A further conjecture by Mahler \cite{mahler} in this regard, proved by Barr\'e-Sirieix, Diaz, Gramain and Philibert \cite{b-d-g-p} states that
\begin{theorem}[Barr\'e-Sirieix, Diaz, Gramain, Philibert]\label{mahl}
For any $\tau\in \mathbb{H}$, at least one of the two numbers $e^{2\pi i\tau}$ and $j(\tau)$ is transcendental.
\end{theorem}
This can be derived as a consequence of a remarkable theorem of Nesterenko \cite{nesterenko} with several applications.
\begin{theorem}[Nesterenko] \label{nest}
If $\tau\in \mathbb{H}$, then at least three of the numbers
$$
e^{2 \pi i \tau},\quad E_2(\tau),\quad E_4(\tau),\quad E_6(\tau)
$$
are algebraically independent over $\overline{\mathbb{Q}}$.
\end{theorem}

Although the above theorems are about specific functions of ``level $1$", it is the aim of this note to highlight that they are sufficient to deduce the corresponding results for functions associated with arbitrary congruence subgroups. The authors believe that this fact may be known to experts, but is not well-documented. Often, the congruence subgroup in question is restricted to be the group $\Gamma_0(N)$. The results in this paper apply to functions satisfying appropriate modularity properties with respect to $\Gamma(N)$, and hence, arbitrary congruence subgroups.\\

A detailed investigation of algebraic independence of values of modular forms and quasi-modular forms was carried out by S. Gun, M. R. Murty and P. Rath \cite{r-m-g} in 2011. Their results on values of modular forms were further elaborated upon for higher level in \cite{hamieh-murty} by A. Hamieh and M. R. Murty\footnote{A small correction in their statement of Theorem 1.2 is required. The conclusion should read as ${(\pi/\omega_{\tau})}^k \, L_q(k,\chi)$ is algebraic, without the $L(1-k,\chi)$ term.}. These theorems will follow from our discussion later. In the context of quasimodular forms, it was proven independently by Gun, Murty and Rath \cite[Theorem 7]{r-m-g} and C. Y. Chang \cite{chang} that
\begin{theorem}[Gun-Murty-Rath and Chang]\label{GMR-Chang}
    If $f$ is a quasi-modular form of non-zero weight for $SL_2(\mathbb{Z})$ with algebraic Fourier coefficients and $\tau\in\mathbb{H}$ is such that $j(\tau)\in\overline{\mathbb{Q}}$, then either $f(\tau) = 0$ or $f(\tau)$ is transcendental.  
\end{theorem}
The above statement is also proved for quasimodular forms with respect to $\Gamma_0(N)$ in \cite{r-m-g}.\\


Another instance of investigation is a recent paper of D. Jeon, S.-Y. Kang and C. H. Kim \cite[Theorem 2.4]{jeon-kang-kim}, where they prove the following. Let $N \in \mathbb{N}$, $\mathfrak{g}:=\mathfrak{g}_0(N) = $ genus of $X_0(N)$, the modular curve obtained as the quotient of the extended upper half plane by $\Gamma_0(N)$. Suppose that $\mathfrak{g} > 0$ and $m \geq \mathfrak{g}+1 $ is an integer. Let $\mathfrak{f}_{N,m}$ denote the unique modular function with respect to $\Gamma_0(N)$ constructed in \cite{j-k-k-2} such that $\mathfrak{f}_{N,m}$ is holomorphic on the upper half plane and
\begin{equation*}
    \mathfrak{f}_{N,m}(q) = \frac{1}{q^m} + \sum_{l=1}^\mathfrak{g} a_N(m,-l) \frac{1}{q^l} + O(q),
\end{equation*}
with the coefficients of powers of $q$ being algebraic numbers. Then, they show the following.
\begin{theorem}[Jeon, Kang, Kim]\label{jjk-thm}
   Let $f$ be a non-zero meromorphic modular form with respect to $\Gamma_0(N)$ with algebraic Fourier coefficients. If $\tau$ is either a zero or a pole of $f$, then $\mathfrak{f}_{N,m}(\tau)$ is algebraic for all $m$. 
\end{theorem}
As a corollary, they deduce that any zero or pole of $f$ should be either CM or transcendental.\\

In this note, we first give an exposition of the algebraic structure of modular functions of higher level, following Shimura \cite{shimura}. Building upon this and using Schneider's theorem, we prove
\begin{theorem}\label{mdf-values-thm}
Let $g$ be a non-constant modular function with respect to a congruence subgroup $\Gamma$ of level $N$, with algebraic Fourier coefficients at $i\infty$. 
\begin{enumerate}
    \item[(a)] If $\tau\in\mathbb{H}$ is either a zero or a pole of $g$, then $j(\tau)$ is algebraic, and hence, $\tau$ is either CM or transcendental.
    \item[(b)] If $\tau\in\mathbb{H}$ is such that $\tau$ is not a pole of $g$, then $j(\tau)\in\overline{\mathbb{Q}}\Leftrightarrow g(\tau)\in\overline{\mathbb{Q}}$. Thus, at least one of $g(\tau)$ and $e^{2\pi i\tau}$ is transcendental for any $\tau \in \mathbb{H}$.
\end{enumerate}
\end{theorem}


In the context of meromorphic modular forms, we establish a generalization of Theorem \ref{jjk-thm} to arbitrary modular forms and arbitrary modular functions. 
\begin{theorem}\label{zero-pole-value}
Let $f$ be a non-constant meromorphic modular form of weight $k \in \mathbb{Z}$ with respect to a congruence subgroup of level $N$ and $g$ be a non-constant modular function with respect to a congruence subgroup of level $M$. Suppose that both $f$ and $g$ have algebraic Fourier coefficients at $i\infty$. Let $\tau \in \mathbb{H}$ be a zero or a pole of $f$. Then either $g(z)$ has a pole at $z=\tau$ or $g(\tau)$ is algebraic.
\end{theorem}


Furthermore, we generalize the theorems in \cite{r-m-g} and \cite{hamieh-murty} to the setting of meromorphic modular forms. In particular, we show the following. 
\begin{theorem}\label{general-mero-values-thm}
Let $f$ be a non-constant meromorphic modular form with respect to a congruence subgroup $\Gamma$ of level $N$. Suppose that $f$ has algebraic Fourier coefficients at $i\infty$. Suppose that $\tau$ is not a pole of $f$. 
  \begin{enumerate}
      \item[(a)] If $\tau\in\mathbb{H}$ is such that $e^{2\pi i\tau}$ is algebraic, then $f(\tau)$ is transcendental.
      \item[(b)] If $\tau\in\mathbb{H}$ is such that $j(\tau)\in \overline{\mathbb{Q}}$, then there exists a transcendental number $\omega_\tau$ which depends only on $\tau$ and is $\overline{\mathbb{Q}}$-linearly independent with $\pi$, such that $\left(\dfrac{\pi}{\omega_\tau}\right)^k f(\tau)\in \overline{\mathbb{Q}}$. Therefore, $f(\tau)$ is either zero or transcendental.
  \end{enumerate}  
\end{theorem}

With regard to quasi-modular forms, we extend the previously known results to meromorphic quasi-modular forms and prove
\begin{theorem}\label{quasi-mdf-thm}
    Let $\widetilde{f}$ be a non-constant meromorphic quasi-modular form with depth $p \geq 1$, with respect to a congruence subgroup $\Gamma$ of level $N$. Suppose that $\widetilde{f}$ has algebraic Fourier coefficients at $i\infty$ and that $\tau$ is not a pole of $\widetilde{f}$. 
  \begin{enumerate}
      \item[(a)] If $\tau\in\mathbb{H}$ is such that $e^{2\pi i\tau}$ is algebraic, then $\widetilde{f}(\tau)$ is transcendental.
      \item[(b)] Let 
      \begin{equation*}
          \widetilde{f} = \sum_{r=0}^p f_{r} \, E_2^r \qquad \text{ with } \qquad f_{r} \in M^m_{k-2r, N, \overline{\mathbb{Q}}}.
      \end{equation*}
      Here $M^m_{j,N,\overline{\mathbb{Q}}}$ denotes the space of meromorphic modular forms of weight $j$, level $N$ and algebraic Fourier coefficients at $i\infty$. If $\tau\in\mathbb{H}$ is such that $j(\tau)\in \overline{\mathbb{Q}}$, then 
      \begin{equation*}
          \widetilde{f}(\tau) = 0 \qquad \iff \qquad f_r(\tau) = 0 \text{  for all  } 0 \leq r \leq p.
      \end{equation*}
      Moreover, if $\widetilde{f}(\tau) \neq 0$, then $\widetilde{f}(\tau)$ is transcendental.
  \end{enumerate}  
\end{theorem}
This generalizes Theorem \ref{GMR-Chang} as well as \cite[Theorem 1.8]{balu-gun}.  \\

Let $\mathcal{Z}_{j, \overline{\mathbb{Q}}}:= \left\{ \tau \in \mathbb{H} \,: \, j(\tau) \in \overline{\mathbb{Q}} \right\}$. From the above results, we deduce the following interesting corollary.
\begin{corollary}
Let $\mathcal{Z}_{\text{mdfn}}$ be the set of zeros and poles of modular \emph{functions} of arbitrary level with algebraic Fourier coefficients, $\mathcal{Z}_{\text{mdfrm}}$ be the set of zeros and poles of meromorphic modular forms of arbitrary level with algebraic Fourier coefficients and $\mathcal{Z}_{\text{quasi-mdf}}$ be the set of zeros and poles of meromorphic quasi-modular forms of arbitrary level with algebraic Fourier coefficients. Then
\begin{equation*}
    \mathcal{Z}_{\text{quasi-mdf}} \subseteq \mathcal{Z}_{\text{mdfrm}} \subseteq \mathcal{Z}_{\text{mdfn}} \subseteq \mathcal{Z}_{j, \overline{\mathbb{Q}}}.
\end{equation*}
In particular, zeros and poles of quasi-modular forms, modular forms and modular functions are either CM or transcendental.
\end{corollary}


Finally, we have the following generalization of Theorem \ref{nest}. 
\begin{theorem}\label{tr-deg-thm}
Let $f$ be a non-constant meromorphic modular form of weight $k \in \mathbb{Z}$ with respect to a congruence subgroup of level $N$, $g$ be a non-constant modular function with respect to a congruence subgroup of level $M$ and $\widetilde{f}$ be a non-constant meromorphic quasi-modular function of depth at least $1$, with respect to a subgroup of level $\widetilde{N}$. Suppose that $f$, $g$ and $\widetilde{f}$ have algebraic Fourier coefficients at $i\infty$. If $\tau\in\mathbb{H}$ is such that $f(\tau)\neq 0$, $\widetilde{f}(\tau)\neq 0$ and $\tau$ is not a pole of $f$, $g$ and $\widetilde{f}$, then 
    $$
\operatorname{trdeg}_{\,\mathbb{Q}}\,  \,\overline{\mathbb{Q}}\left( e^{2\pi i\tau},\, f(\tau),\, g(\tau),\, \widetilde f(\tau)\right)\geq 3. 
    $$
\end{theorem}
This theorem is in the same spirit as \cite[Theorem 1.2]{wang}, where W. Wang considers the algebraic independence of the values of three algebraically independent quasi-modular forms. However, Theorem \ref{tr-deg-thm} allows one to also compare values of modular functions with those of quasi-modular forms. It can be shown that a modular function, a modular form of non-zero weight and a quasi-modular form of positive depth are algebraically independent. We include a proof of this assertion in Theorem \ref{algebraic-indep-fn-thm} for completeness. \\

\noindent We also remark that in Theorem \ref{zero-pole-value} and Theorem \ref{tr-deg-thm}, one can replace a meromorphic modular form by a half-integer weight modular form with algebraic Fourier coefficients since the square of a half-integer weight modular form is an integer weight modular form.

\bigskip

\section{\bf Preliminaries}
\bigskip

The aim of this section is to study the algebraic structure of the field of modular functions and to record the required results from transcendental number theory. For the sake of completeness and clarity of exposition, a brief account of the proofs is included, and appropriate references are given.


\subsection{Modular and quasi-modular forms}
For each $N\in \mathbb{N}$, $$
\Gamma(N):=\left\{  \left( 
\begin{array}{cc}
  a & b \\
  c & d
\end{array}
\right)\in SL_2(\mathbb{Z}):\left( 
\begin{array}{cc}
  a & b \\
  c & d
\end{array}
\right)\equiv \left( 
\begin{array}{cc}
  1 & 0 \\
  0 & 1
\end{array}
\right) (\operatorname{mod}N) \right\},
$$ 
with $\Gamma(1)=SL_2(\mathbb{Z})$. A subgroup $\Gamma$ of $SL_2(\mathbb{Z})$ is said to be congruence subgroup if there exists $N\in \mathbb{N}$ such that $\Gamma(N)\subseteq \Gamma$. The smallest such $N$ is said to be the level of $\Gamma$.\\

A \emph{holomorphic modular form} of integer weight $k \geq 0$ with respect to a congruence subgroup $\Gamma$ is a holomorphic function on the upper half-plane $\mathbb{H}$ which satisfies 
\begin{enumerate}
    \item[(i)] $f \big|_k \gamma = f$ for all $\gamma \in \Gamma$
    \item[(ii)] $f \big|_k \alpha$ is holomorphic at $i\infty$ for all $\alpha \in SL_2(\mathbb{Z})$.
\end{enumerate}
The function $f$ is said to be a \emph{weakly holomorphic modular form} if $f \big|_k \alpha$ is allowed to have poles at $i\infty$, that is, $f$ is \emph{meromorphic} at the cusps. More generally, $f$ is called a \emph{meromorphic modular form} if it is meromorphic on $\mathbb{H}$ and also at cusps. \\

We will say that \emph{$f$ has algebraic Fourier coefficients} if the Fourier coefficients of $f \big|_k \alpha$ at $i \infty$ are algebraic numbers. Denote the space of all holomorphic, weakly holomorphic and meromorphic modular forms, with algebraic Fourier coefficients by $M_{k,\overline{\mathbb{Q}}}(\Gamma)$, $M_{k,\overline{\mathbb{Q}}}^w(\Gamma)$ and $M_{k,\overline{\mathbb{Q}}}^m(\Gamma)$ respectively. Clearly, $M_{k,\overline{\mathbb{Q}}}(\Gamma)\subset M_{k,\overline{\mathbb{Q}}}^w(\Gamma)\subset M_{k,\overline{\mathbb{Q}}}^m(\Gamma)$. \\

For even integer $k\geq 2$, define the normalized Eisenstein series of weight $k$ for $S L_2(\mathbb{Z})$ by
$$
E_{k}(\tau) = 1-\frac{2 k}{B_{k}} \sum_{n=1}^{\infty} \sigma_{k-1}(n) q^{n}, \text{   where }q=e^{2\pi i\tau} \text{ and } \sigma_s(n)=\sum_{\substack{d|n\\ d>0}}d^s.
$$ 
Here, $B_k$ is the $k$-th Bernoulli number. Define
\begin{equation*}
    \Delta(\tau):=\frac{E_4(\tau)^3-E_6(\tau)^2}{1728}.
\end{equation*}
For $k\geq 4$, the function $E_k\in M_{k,\overline{\mathbb{Q}}}(S L_2(\mathbb{Z}))$ and $\Delta \in M_{12,\overline{\mathbb{Q}}}(S L_2(\mathbb{Z}))$. But $E_2(\tau)$ is not a modular form (see \cite{murty-dewar-graves}, Chapter $5$), as
\begin{equation*}
    E_2 \left( \frac{-1}{\tau} \right) = {\tau}^2 E_2(\tau) + \frac{6}{i\pi} \tau.
\end{equation*}
This motivates the definition of a quasi-modular form, of which there are several equivalent formulations. We use the following characterization, which was established for holomorphic quasi-modular forms in \cite[Proposition 20]{1-2-3} and meromorphic quasi-modular forms in \cite[Theorem 4.2]{wang-zhang}.
\begin{theorem}\label{quasistructure}
Every meromorphic quasi-modular form for a congruence subgroup $\Gamma$ is a polynomial in $E_2$ with modular coefficients. More precisely, if $\widetilde{f}$ is a meromorphic quasi-modular form of weight $k$ and depth $p$ with respect to $\Gamma$, then $\widetilde{f}$ can be uniquely written as $\widetilde{f} = \sum_{r=0}^p  f_r\, E_2^r$, where $f_r$ is a meromorphic modular form with respect to $\Gamma$ of weight $k-2r$ for all $0\leq r\leq p$ and $f_p\neq 0$. 
\end{theorem}
A quasi-modular form $\widetilde{f}$ is said to have algebraic Fourier coefficients if all the modular coefficients in the above expression of $f$ have algebraic Fourier coefficients. \\

\subsection{The Weierstrass $\wp$-function}
 Let $L=\omega_1 \mathbb{Z} \,  \oplus \, \omega_2 \mathbb{Z}$ be a two-dimensional lattice in $\mathbb{C}$ with $\omega_1/\omega_2 \in \mathbb{H}$. The Weierstrass $\wp$-function associated to $L$, given by
\begin{equation}\label{wpfunction}
\wp (z)=\frac{1}{z^2}+\sum_{\substack{\omega\in L\\
\omega\neq 0}}\left(\frac{1}{(z-\omega)^2}-\frac{1}{\omega^2} \right), \quad \text{for }z\in \mathbb{C}\setminus L,
\end{equation}
defines a meromorphic function on $\mathbb{C}$. It satisfies the differential equation
\begin{equation} \label{dell}
 \wp'(z)^2=4\wp(z)^3-g_2(L)\wp(z)-g_3(L),   
\end{equation}
where 
\begin{equation*}
    g_2(L)=60\, \sum_{\substack{\omega \in L\\ \omega\neq 0}} \frac{1}{\omega^{4}} \qquad \text{and} \qquad g_3(L)=140\, \sum_{\substack{\omega \in L\\ \omega\neq 0}} \frac{1}{\omega^{6}}.
\end{equation*}
For the lattice $L_\tau=\tau\,\mathbb{Z}\oplus \mathbb{Z}$ with $\tau\in\mathbb{H}$, 
\begin{equation*}
    g_2(L_\tau)=\frac{4\pi^4}{3}E_4(\tau) \qquad \text{and} \qquad g_3(L_\tau)=\frac{8\pi^6}{27}E_6(\tau).
\end{equation*}
Define the discriminant of a lattice, $$\Delta_0(L):=g_2(L)^3-27g_3(L)^2\neq 0, \text{ for any two dimensional lattice }L.$$ In particular, we have $\Delta_0(L_\tau)=(2\pi)^{12}\,\Delta(\tau)$ for all $\tau\in\mathbb{H}$.\\

The Weierstrass zeta-function associated to $L=\omega_1\,\mathbb{Z}\oplus \omega_2\,\mathbb{Z}$ is defined by
$$
\zeta_{L}(z)=\frac{1}{z}+\sum_{\substack{\omega\in L\\ \omega\neq 0}}\left(\frac{1}{z-\omega}+\frac{1}{\omega}+\frac{z}{\omega^2} \right)\quad \text{for } z\in \mathbb{C}\setminus L.
$$
Note that $\zeta_{L}'(z)=-\wp_L(z)$ is a periodic function with each point of $L$ as a period. Hence, the functions $$\zeta_{L}(z+\omega_1)-\zeta_{L}(z) \text{  and  } \zeta_{L}(z+\omega_2)-\zeta_{L}(z)$$ are constants. These constants are denoted by $\eta_1(L)$ and $ \eta_2(L)$ respectively and are called \emph{quasi-periods}. Moreover, for $\omega_1 = \tau$ and $\omega_2 = 1$, it is known that 
\begin{equation}\label{quasi-period}
\eta_2(L_\tau)=G_2(\tau)=\frac{\pi^2}{3}\, E_2(\tau) \text{ 
  for all }\tau\in\mathbb{H}. 
\end{equation}
The reader is referred to (\cite{lang}, chapter $18.3$) for proof of these results.\\

The Weierstrass $\wp$-function satisfies the following addition formula (\cite{r-m}, chapter $11$).
 \begin{theorem} \label{addf}
 For $z_1, z_2\in \mathbb{C}$ such that $z_1\pm z_2\notin L$ we have
$$
\wp\left(z_1+z_2\right)=-\wp\left(z_1\right)-\wp\left(z_2\right)+\frac{1}{4}\left(\frac{\wp^{\prime}\left(z_1\right)-\wp^{\prime}\left(z_2\right)}{\wp\left(z_1\right)-\wp\left(z_2\right)}\right)^2 .
$$
\end{theorem}
Using the addition formula of the Weierstrass $\wp$-function, Schneider \cite{schneider.} proved that 
\begin{theorem}\label{tranw}
    Let $L=\omega_1\,\mathbb{Z}\oplus\omega_2\,\mathbb{Z}$ be a lattice such that both $g_2(L)$, $g_3(L)$ are algebraic. If $\alpha$ is an algebraic number with $\alpha\notin L$, then $\wp(\alpha)$ is transcendental.
\end{theorem}
The addition formula also implies the following important proposition (see \cite{r-m}). Two generators $\omega_1$ and $\omega_2$ of a lattice $L$ are said to be primitive if both have minimal absolute value among all generators of $L$.
\begin{prop}\label{tranp}
Let $L=\omega_1\,\mathbb{Z}\oplus\omega_2\,\mathbb{Z}$ be such that both $g_2(L)$ and $g_3(L)$ are algebraic. Assume that $\omega_1$ and $\omega_2$ are primitive generators of $L$. Then, for any natural number $n>1$, the numbers $\wp(\frac{\omega_1}{n})$ and $\wp(\frac{\omega_2}{n})$ are algebraic. Moreover, any non-zero period of $L$ is necessarily transcendental.\\
\end{prop}

\subsection{Modular functions}

A meromorphic function $g$ on $\mathbb{H}$ is said to be a modular function if it satisfies
\begin{equation}
g\left(\frac{a\,\tau+b}{c\,\tau+d}\right)=g(\tau)\quad\text{for all }
\gamma=\left(\begin{smallmatrix}
   a  & b \\
   c  & d
\end{smallmatrix} \right)\in \Gamma,
\end{equation}
that is, $g \big|_0 \gamma = g$, and is also meromorphic at all the cusps. In particular, we call a modular function on the congruence subgroup $\Gamma(N)$ to a be \emph{modular function of level $N$}. Note that if $g$ is a modular function with respect to $\Gamma$, which is of level $N$, then $g$ is a modular function of level $N$.\\

An example of modular function of level one is given by 
$$
j(\tau):=\frac{E_4^3(\tau)}{\Delta(\tau)} \text{ for all } \tau\in\mathbb{H},
$$
which has the following Fourier expansion at $i\infty$
$$
j(\tau)=\frac{1}{q}+744+196884q+\cdots,$$
where $q=e^{2\pi i\tau}$. \\


More specifically, for any modular function (or form) $f$, define 
\begin{equation*}
    \mathbb{Q}(f) := \mathbb{Q} \left( \left\{ \text{Fourier coefficients of $f$ at all $\Gamma$-in-equivalent cusps} \right\} \right)
\end{equation*}
The $j$-function is a canonical example of a modular function of level one, and governs properties of modular functions of higher levels as well. This is made precise in the following series of propositions.
\begin{prop} \label{level1}
Let $g$ be a non-constant modular function of level one. Let $\mathcal{F}$ denote the standard fundamental domain for the action of $S L_2(\mathbb{Z})$ on $\mathbb{H}$. Suppose that the poles of $g$ in $\mathcal{F}$ are $\tau_1,\tau_2,\cdots,\tau_m$. Then $g(\tau)$ is a rational function in $j(\tau)$ with coefficients in the field $\mathbb{Q}(g) \left( j(\tau_1),\cdots,j(\tau_m) \right)$.
\end{prop}
\begin{proof}
By the compactness of $\mathcal{F} \cup i\infty$, we know that $g$ has only finitely many poles in $\mathcal{F}$. Consider the function 
$$
h(\tau):=\left(\prod_{j=1}^m \left(j(\tau)-j(\tau_j)\right)^{-\operatorname{ord}_{\tau_j}(g)}\right)g(\tau),
$$
where $\operatorname{ord}_{\tau_j}(g) = - $ order of the pole of $g$ at $\tau_j$.  Then $h$ is holomorphic on $\mathbb{H}$.\\

Suppose $h$ has pole of order $M \geq 1$ at $i\infty$. Then the Fourier expansion of
$h(\tau)$ at $i\infty$ has the form
$$
h(\tau)=\sum_{n=-M}^\infty c_n q^n, \, \text{ where } c_{-M}\neq 0.
$$
Note that the modular function $h(\tau)-c_{-M}\,j(\tau)^M$ is holomorphic on $\mathbb{H}$ and its Fourier expansion starts with at most a polar term of order $M-1$. Iterating this process, we can subtract a polynomial in $j(\tau)$ to get a holomorphic modular function that vanishes at $i\infty$, and hence is identically zero. Thus, $h(\tau)$ is a polynomial in $j(\tau)$ over $\mathbb{Q}(g)$, and so $g(\tau)$ is a rational function of $j(\tau)$ over $\mathbb{Q}(g)(j(\tau_1),\cdots,j(\tau_m))$.
\end{proof}

One can also conclude the following important fact from the above proof.
\begin{corollary}\label{rationalfunctions}
If $g$ is a modular function of level one which is holomorphic on $\mathbb{H}$ with a pole of order $M$ at $i\infty$, then $g(\tau)$ is a polynomial in $j(\tau)$ of degree $M$ with coefficients in $\mathbb{Q}(g)$.
\end{corollary}

The $j$-function is sufficient to `generate' all higher level modular functions as well. This is proved below.
\begin{theorem} \label{mdft}
Let $g$ be a modular function with respect to a congruence subgroup $\Gamma$ and let $\tau_1,\, \tau_2, \cdots, \, \tau_m$ be the poles of $g$ in $\mathcal{F}$. Set
$$
K:=\mathbb{Q}(g)\left(j(\tau_1), \, j(\tau_2), \cdots, \, j(\tau_m) \right).
$$
Then there exists a monic polynomial $P_g(X)\in K(j)[X]$ such that $P_g(g)=0$. 
\end{theorem}
\begin{proof}
Let $[S L_2(\mathbb{Z}):\Gamma]=r$ and $\{ \gamma_1=I,\gamma_2,\cdots,\gamma_r\}$ be a complete set of right coset representatives so that $$
S L_2(\mathbb{Z})=\bigsqcup_{i=1}^r\Gamma \gamma_i.
$$
For all $1\leq i \leq r$, define the functions $$g_i(\tau):=g(\gamma_i\tau) \quad \text{for } \tau \in \mathbb{H}.
$$ Each $g_i$ is independent of the choice of coset representatives as $g$ is $\Gamma$-invariant. Moreover, each $g_i$ is a modular function with respect to $\Gamma$ and the Fourier expansion of $g_i$ at $i \infty$ is precisely the expansion of $g$ at the cusp $\gamma_i(i\infty)$. For any $\gamma\in S L_2(\mathbb{Z})$, we have $g_i(\gamma \tau)=g(\gamma_i\gamma\tau)=g_j(\tau)$ for some $j$ with $1 \leq j \leq r$ such that $\gamma_i\gamma\in \Gamma\gamma_j$, i.e., the set $\{g_1,g_2,...,g_r\}$ gets permuted under the action of $S L_2(\mathbb{Z}).$ This observation implies that any elementary symmetric polynomial in $g_1,g_2,\cdots,g_r$ is a modular function of level one, and is in $K(j)$ by Theorem \ref{level1}. Note that the polynomial $$
P(X)=\prod_{i=1}^r(X-g_i)$$is satisfied by $g$ as $g=g_1$ and has coefficients that are elementary symmetric polynomials in $g_1,g_2,\cdots,g_r$. This proves the theorem.
\end{proof}



\begin{corollary}\label{useful}
Let $g$ be a modular function with respect to a congruence subgroup $\Gamma$ which is holomorphic on $\mathbb{H}$ and $K = \mathbb{Q}(g)$. Then the monic polynomial $P_g(X)$ satisfied by $g$ has coefficients in $K[j][X]$.
\end{corollary}
\begin{proof}
The coefficients of $P_g(X)$ constructed in the proof above are modular functions of level one, holomorphic on $\mathbb{H}$ with Fourier coefficients in $K$. The result now follows from Corollary \ref{rationalfunctions}. 
\end{proof}

\begin{remark*}
An important point to note here is that if $g$ is a modular function with respect $\Gamma$ that has algebraic Fourier coefficients, Theorem \ref{mdft} does not guarantee that the coefficients of $P_g$ are algebraic, unless $g$ is holomorphic in $\mathbb{H}$. 
% For general modular functions, one needs to delve deeper into the structure of the field generated by modular functions of higher level in order to assert that a modular function of higher level with algebraic Fourier coefficients satisfies a polynomial over $\overline{\mathbb{Q}}(j)$.
\end{remark*}


To that effect, we now study the structure of the field of modular functions of a fixed level $N>1$. For any field $\mathbb{F} \subseteq \mathbb{C}$, let
\begin{equation*}
    \mathcal{F}_{N, \mathbb{F}} := \left\{ \text{Modular functions of level $N$ whose Fourier coefficients at $i \infty$ are in $\mathbb{F}$} \right\}.
\end{equation*}
Theorem \ref{mdft} implies that $\mathcal{F}_{N,\mathbb{C}}$ is an algebraic extension of $\mathbb{C}(j)$. 

% Furthermore, we have
% \begin{prop}\label{fnc}
% The field $\mathcal{F}_{N,\mathbb{C}}$ is a finite Galois extension of $\mathbb{C}(j)$ and the Galois group is isomorphic to $S L_2(\mathbb{Z})/\{\pm I\}\Gamma(N)$.
% \end{prop}
% \begin{proof}
% From Theorem \ref{mdft}, we know that every $g\in \mathcal{F}_{N,\mathbb{C}}$ satisfies $P_g(X)$. The other roots of $P_g(X)$ are the modular functions $g_i$, which are also of level $N$. Hence $\mathcal{F}_{N,\mathbb{C}}/\mathbb{C}(j)$ is a Galois extension. To find the Galois group, note that any element $\gamma\in S L_2(\mathbb{Z})$ defines an automophism of $\mathcal{F}_{N,\mathbb{C}}$ by $g \mapsto g\circ\gamma $, where $(g\circ\gamma )(\tau):=g(\gamma \tau)$ for all $\tau\in\mathbb{H}$. Moreover, $\gamma$ fixes $\mathbb{C}(j)$ as $j$ is modular function of level $1$. Define a map $$
% \begin{aligned}
% \Psi : \,\,  & S L_2(\mathbb{Z}) \longrightarrow \text{Gal}(\mathcal{F}_{N,\mathbb{C}}/\mathbb{C}(j))\\
% & \quad\gamma \longmapsto \Psi_\gamma : \mathcal{F}_{N,\mathbb{C}} \longrightarrow \mathcal{F}_{N, \mathbb{C}}\\
% &\quad \quad\quad\quad \quad \quad h \longmapsto h\circ\gamma.
% \end{aligned}
% $$
% Then $\Psi$ is a group homomorphism with kernel $\{\pm I\}\Gamma(N)$. By the primitive element theorem, $\mathcal{F}_{N,\mathbb{C}}=\mathbb{C}(\widetilde g)$ for some $\widetilde{g} \in \mathcal{F}_{N,\mathbb{C}}$. Note that any $\sigma\in Gal(\mathcal{F}_{N,\mathbb{C}}/\mathbb{C}(j)) $ sends $\widetilde{g}$ to one of its conjugates, which is of the form $\widetilde g_i=\widetilde g\circ\gamma_i$. Then we get
% $$
% \sigma(h)=h\circ\gamma_i\quad \forall\ h\in \mathcal{F}_{N,\mathbb{C}},
% $$
% and so $\sigma=\Psi_{\gamma_i}$. This shows that $\Psi$ is onto and hence, $\text{Gal}(\mathcal{F}_N/\mathbb{C}(j)) \simeq SL_2(\mathbb{Z})/\{\pm I\}\Gamma(N).$
% \end{proof}






% For $N>1$, let $F_{N,\overline{\mathbb{Q}}}$ denote the field of modular functions of level $N$ with algebraic Fourier coefficients at all the $\Gamma(N)$-in-equivalent cusps. Theorem \ref{mdft} implies that $F_{N,\overline{\mathbb{Q}}}$ is an algebraic extension of $\overline{\mathbb{Q}}(j)$.


Following \cite[Chapter 6]{shimura}, we consider explicit modular function of level $N$ whose Fourier coefficients have good rationality properties. Let $\boldsymbol{a}=(a_1,a_2)\in \frac{1}{N}\mathbb{Z}^2 \,\setminus \,\mathbb{Z}^2$, consider the function
$$
f_{\boldsymbol{a}}(\tau):=\frac{g_2(L_\tau)\,g_3(L_\tau)}{\Delta_0(L_\tau)}\, \wp_{L_\tau} (a_1\tau+a_2),
$$
which is holomorphic on $\mathbb{H}$. The following properties can be checked routinely, and we leave the proof to the reader.
\begin{prop}
Let $S$ denotes the set $\{\left(\frac{r}{N},\frac{s}{N}\right): 0\leq r,s\leq N-1 \text{ and }(r,s)\neq(0,0)\}$. Let $\boldsymbol{a}$, $\boldsymbol{b} \in \frac{1}{N}\mathbb{Z}^2 \,\setminus \,\mathbb{Z}^2$ and $f_{\boldsymbol{a}}$ be as defined above. For $\gamma = \big(\begin{smallmatrix}
p & q\\
r & s
\end{smallmatrix}\big)$ and $\boldsymbol{a} = (a_1,a_2)$, let ${\boldsymbol{a\gamma}=(pa_1+ra_2, \, qa_1+sa_2)}$. Then
\begin{enumerate}
    \item[(a)] $f_{\boldsymbol{a}} = f_{\boldsymbol{b}}$ if and only if $\boldsymbol{a} \equiv \boldsymbol{b} \bmod \mathbb{Z}^2 $,
    \item[(b)] for $\gamma \in SL_2(\mathbb{Z})$, $f_{\boldsymbol{a}} (\gamma \tau) = f_{\boldsymbol{a\gamma}}(\tau)$ for all $\tau \in \mathbb{H}$.
\end{enumerate}
Therefore, all elements in $\{f_{\boldsymbol{a}}: \boldsymbol{a} \in S\}$ satisfy modularity with respect to $\Gamma(N)$.
\end{prop}

In order to conclude that the functions $f_{\boldsymbol{a}}$ are modular functions, we study their Fourier expansions at the cusp. But first, we need to understand the behaviour of the Weierstrass $\wp$-function.
\begin{lemma*}
Fix $\tau \in \mathbb{H}$ and denote $\wp_{\tau} = \wp_{L_{\tau}}$. For $z \not\in L_{\tau}$, we have
\begin{equation*}\label{weifou}
    (2\pi i)^{-2}\wp_\tau(z)=\frac{1}{12}+\frac{1}{\xi^{-1}-2+\xi}+\sum_{n=1}^\infty c_n q^n, 
\end{equation*}
where 
$$
\begin{aligned}
  \xi=e^{2\pi i z},\, q=e^{2\pi i\tau} \text{ and }
    c_n=\sum_{d|n}d\left(\xi^{-d}-2+\xi^d\right)\ \forall\ n\geq 1.
\end{aligned}
$$
\end{lemma*}
\begin{proof}
    From the definition of the $\wp$-function, we have 
$$
\begin{aligned}
\wp_\tau(z)&=\frac{1}{z^2}+\sum_{\substack{(m,n)\in \mathbb{Z}^2\\ (m,n)\neq (0,0)}}\left[\frac{1}{(z-m\tau-n)^2}-\frac{1}{(m\tau+n)^2} \right]\\
&=\frac{1}{z^2}+\sum_{n\in \mathbb{Z}\setminus \{0\}}\frac{1}{(z+n)^2}-\sum_{n\in \mathbb{Z}\setminus \{0\}}\frac{1}{n^2}+\sum_{m\neq 0}\sum_{n\in \mathbb{Z}}\frac{1}{(z-m\tau-n)^2}-\sum_{m\neq 0}\sum_{n\in \mathbb{Z}}\frac{1}{(m\tau+n)^2}\\
&=-2\zeta(2)+\sum_{n\in\mathbb{Z}}\frac{1}{(z+n)^2}-2\sum_{m=1}^\infty\sum_{n\in \mathbb{Z}}\frac{1}{(m\tau+n)^2}\\
& \quad \quad \quad \quad \quad \quad \quad \quad \quad \quad \quad +\sum_{m=1}^\infty\sum_{n\in \mathbb{Z}}\left[\frac{1}{(-z+m\tau+n)^2}+\frac{1}{(z+m\tau+n)^2} \right].
\end{aligned}
$$
Recall the Lipschitz summation formula, which states 
\begin{equation*} \label{lipc}
\sum_{n \in \mathbb{Z}} \frac{1}{(z+n)^{k}}=\frac{(-2 \pi i)^{k}}{(k-1) !} \sum_{n=1}^{\infty} n^{k-1} q^{n}.
\end{equation*}
For a proof, see \cite[Theorem 4.2.2]{murty-dewar-graves}. Applying this, we obtain
$$
\begin{aligned}
& \wp_\tau(z)=-\frac{\pi^2}{3}+(-2\pi i)^2\sum_{n=1}^\infty ne^{2\pi i n z}\\
& \quad \quad \quad \quad \quad\quad\quad\quad +(-2\pi i)^2\sum_{m=1}^\infty\sum_{n=1}^\infty n\left[e^{2\pi i n(-z+m\tau)}+e^{2\pi i n(z+m\tau)}-2e^{2\pi i m n \tau} \right] \\
&\Rightarrow (2\pi i)^{-2}\wp_\tau(z)=\frac{1}{12}+\sum_{n=1}^\infty ne^{2\pi i n z}+\sum_{m,n=1}^\infty nq^{mn}\left[e^{2\pi i n z}+e^{-2 \pi i n z}-2 \right]
  \end{aligned}
$$
This implies the result.
\end{proof}

Using the above expansion, the interpretation of $g_2(L_{\tau)}$, $g_3(L_{\tau})$ and $\Delta_{0}(L_{\tau})$ in terms of Eisenstein series, we get for $\boldsymbol{a}= \left( \frac{r}{N}, \frac{s}{N} \right) \in S$,
\begin{align}
& f_{\boldsymbol{a}}(\tau) \nonumber \\
& =\frac{-1}{2592}\cdot\frac{E_4(\tau)E_6(\tau)}{\Delta(\tau)}\left[\frac{1}{12} -\sum_{n=1}^\infty ne^{\frac{2\pi i n s}{N}} q_N^{n r}+\sum_{m,n=1}^\infty n q^{m n}\left[e^{\frac{2\pi i n s}{N}}q_N^{n r}+e^{\frac{-2\pi i n s}{N}}q_N^{-n r}-2 \right]\right]. \label{foufa}
\end{align}
Here $q_N = e^{2 \pi i \tau/N}$. Since the Fourier series of $\Delta(\tau)$ begins with $q=q_N^N$, the Fourier expansion of $f_{\boldsymbol{a}}(\tau)$  begins with a rational multiple of $q_N^{-N}$. Thus, $f_{\boldsymbol{a}}$ has pole of order $N$ at $i\infty$ for all $\boldsymbol{a}\in S$. If $s$ is any other cusp, then there exists $\gamma\in S L_2(\mathbb{Z})$ such that $s=\gamma(i\infty)$. Since $f_{\boldsymbol{a}}(\gamma\tau)=f_{\boldsymbol{a\gamma}}(\tau)$, which has pole of order $N$ at $i\infty$, we conclude that $f_{\boldsymbol{a}}$ is meromorphic at all cusps, with pole of order $N$. Thus, $f_{\boldsymbol{a}}$ is a modular function of level $N$ for all $ \boldsymbol{a} \in S$. \\

This helps us to conclude the following.
\begin{theorem}\label{primiext}
For all $\boldsymbol{a} \in S$, the Fourier coefficients of $f_{\boldsymbol{a}}$ with respect to all cusps belong to $\mathbb{Q}(\mu_N)$, where $\mu_N=e^{2\pi i/N}$. 
\end{theorem}
\begin{proof}
   Recall that Fourier coefficients of $E_4(\tau), E_6(\tau)$ and $\Delta(\tau)$ are integers. Hence, from \eqref{foufa}, it follows that the Fourier coefficients of $f_{\boldsymbol{a}}$ at $i\infty$ lie in $\mathbb{Q}(\mu_N)$ for all $\boldsymbol{a}\in S$. If $s$ is any other cusp, then there exists $\gamma\in S L_2(\mathbb{Z})$ such $s=\gamma(i\infty)$. But $f_{\boldsymbol{a}}(\gamma\tau)=f_{\boldsymbol{a\gamma}}(\tau)$ also has Fourier coefficients in $\mathbb{Q}(\mu_N)$ with respect to $i\infty$. This completes the proof.
\end{proof}

These modular functions, together with the $j$-function serve to generate all modular functions of level $N$. That is,
\begin{theorem}\label{mdft..}
We have $\mathcal{F}_{N,\overline{\mathbb{Q}}}=\overline{\mathbb{Q}}\left( j, \, \{ f_{\boldsymbol{a}} |{\boldsymbol{a}}\in S\} \right)$ and $\mathcal{F}_{N,\overline{\mathbb{Q}}}$ is a finite Galois extension of $\overline{\mathbb{Q}}(j)$.
\end{theorem}
\begin{proof}
% Let $E_{N,\mathbb{C}}$ denote the field $\mathbb{C}\left(j, \, \{ f_{\boldsymbol{a}} | {\boldsymbol{a}}\in S\}\right)$. Then we have containment $\mathbb{C}(j)\subseteq E_{N,\mathbb{C}} \subseteq \mathcal{F}_{N,\mathbb{C}}.$ To complete the proof, it is enough to show $[\mathcal{F}_{N,\mathbb{C}}:E_{N,\mathbb{C}}]=1$, i.e., any element of $\text{Gal}(\mathcal{F}_{N,\mathbb{C}}/\mathbb{C}(j))$ that fixes $E_{N,\mathbb{C}}$ must be trivial. Recall that $\text{Gal}(\mathcal{F}_{N,\mathbb{C}}/\mathbb{C}(j))=S L_2(\mathbb{Z})/\{\pm I\}\Gamma(N)$. Let $\gamma=\big(\begin{smallmatrix}
% p & q\\
% r & s\end{smallmatrix}\big)\in S L_2(\mathbb{Z}) $ fix $E_{N,\mathbb{C}}$, that is, $$f_{\boldsymbol{a}}\circ\gamma=f_{\boldsymbol{a}}\quad \text{for all }{\boldsymbol{a}}\in S.$$ Equivalently,
% $$
%  \quad f_{\boldsymbol{a\gamma}}=f_{\boldsymbol{a}} \quad \text{for all } {\boldsymbol{a}} \in S
% $$
% and hence 
% $$
%  {\boldsymbol{a}\gamma}\equiv \pm {\boldsymbol{a}}\pmod{\mathbb{Z}^2} \quad \text{ for all } {\boldsymbol{a}}\in S.
% $$
% In particular, taking two choices of ${\boldsymbol{a}}= (\frac{1}{N},0)$ and $(0,\frac{1}{N})$ we conclude $\gamma \in \{\pm I\}\Gamma(N).$ This completes the proof. \\

Let $E_{N,\overline{\mathbb{Q}}}:= \overline{\mathbb{Q}}\left(j, \, \{ f_{\boldsymbol{a}} | {\boldsymbol{a}}\in S\}\right)$. Then $E_{N, \overline{\mathbb{Q}}}$ is a Galois extension of $\overline{\mathbb{Q}}(j)$. Indeed, for each ${\boldsymbol{a}}\in S$, the modular function $f_{\boldsymbol{a}}$ is holomorphic on $\mathbb{H}$ with algebraic Fourier coefficients at all cusps. Hence, by Corollary \ref{useful} we get a polynomial $P(X)\in\overline{\mathbb{Q}}(j)[X]$, which is satisfied by $f_{\boldsymbol{a}}$. Thus, $E_{N,\overline{\mathbb{Q}}} \, / \, \overline{\mathbb{Q}}(j)$ is an algebraic extension. Moreover, this extension is normal because each conjugate of $f_{\boldsymbol{a}}$ is of the form $f_{\boldsymbol{a\gamma_i}}$, which is equal to $f_{\boldsymbol{b}}$ for some $\boldsymbol{b}\in S$. Since $E_{N,\overline{\mathbb{Q}}}$ is obtained from $\overline{\mathbb{Q}}(j)$ by adjoining finite number of elements, $E_{N,\overline{\mathbb{Q}}} \, / \, \overline{\mathbb{Q}}(j)$ is a finite Galois extension.\\

Now we show that $E_{N,\overline{\mathbb{Q}}}= \mathcal{F}_{N,\overline{\mathbb{Q}}}$. To begin with, observe that $\mathbb{C}$ and $\mathcal{F}_{N,\overline{\mathbb{Q}}}$ are linearly disjoint over $\overline{\mathbb{Q}}$. This can be seen as follows. Suppose $\{c_1,c_2,\ldots,c_r\}\subset \mathbb{C}$ is an arbitrary $\overline{\mathbb{Q}}$-linearly independent subset of $\mathbb{C}$. If there exists $g_i \in \mathcal{F}_{N,\overline{\mathbb{Q}}}$ such that $\sum_{i=1}^r g_i(\tau) \, c_i = 0$ for all $\tau \in \mathbb{H}$ with
$$g_i(\tau)=\sum_n d_{i n} \, \, q_N^n, \qquad d_{in} \in \overline{\mathbb{Q}} \quad \text{for }0\leq i\leq r,$$
then $$\sum_{i=1}^rc_i\,d_{i n}=0\quad \text{ for all } n.$$ The linear independence of $c_i$ over $\overline{\mathbb{Q}}$ implies that all $d_{i n}=0$ for all $0\leq i\leq r$, and hence, $f_1=f_2=\cdots=f_r=0$. Thus, we have $E_{N,\overline{\mathbb{Q}}}\subseteq \mathcal{F}_{N,\overline{\mathbb{Q}}} \subseteq \mathbb{C}E_{N,\overline{\mathbb{Q}}}.$ Suppose there exists $f\in \mathcal{F}_{N,\overline{\mathbb{Q}}}\setminus E_{N,\overline{\mathbb{Q}}}$. Since $f \in \mathbb{C}E_{N,\overline{\mathbb{Q}}}$, we get a $\overline{\mathbb{Q}}$-linearly independent subset $\{f_1,f_2,\ldots,f_m\}\subseteq E_{N,\overline{\mathbb{Q}}}$ such that 
\begin{equation}\label{rep.}
  f=\sum_{i=1}^m \alpha_i f_i, \quad \alpha_i\in \mathbb{C},  
\end{equation}
with at least one of the $\alpha_i \in \mathbb{C} \setminus \overline{\mathbb{Q}}$. Since $E_{N,\overline{\mathbb{Q}}}\subseteq \mathcal{F}_{N,\overline{\mathbb{Q}}}$, the set $\{f, \,f_1,f_2,\ldots,f_m\}$ is $\overline{\mathbb{Q}}$-linearly independent subset of $\mathcal{F}_{N,\overline{\mathbb{Q}}}$ and hence, $\mathbb{C}$-linearly independent. This contradicts \eqref{rep.}. Therefore, $E_{N,\overline{\mathbb{Q}}}= \mathcal{F}_{N,\overline{\mathbb{Q}}}$ and $\mathcal{F}_{N,\overline{\mathbb{Q}}} \, / \, \overline{\mathbb{Q}}(j)$ is a finite Galois extension.
\end{proof}

Thus, Theorem \ref{primiext} and Theorem \ref{mdft..} imply the following crucial fact.
\begin{corollary}\label{cor-main-mdf}
Let $g \in \mathcal{F}_{N,\overline{\mathbb{Q}}}$. Then the Fourier coefficients of $g$ with respect to all cusps are algebraic numbers and $g$ satisfies a polynomial over $\overline{\mathbb{Q}}(j)$. 
\end{corollary}


\subsection{Algebraic independence of modular, quasi-modular forms and modular functions} 

The aim of the discussion below is to establish the algebraic independence (as functions) of the three functions arising from modular considerations - modular functions, modular forms and quasi-modular forms. For the basic theory of quasi-modular forms, we refer the reader to \cite{royer}. To begin with, we prove the following lemma.
\begin{lemma}\label{indep}
A sum of meromorphic quasi-modular forms of distinct weights is not identically zero, unless, each of these quasi-modular forms is identically zero.
\end{lemma}
\begin{proof}
Suppose $f_1, f_2,\ldots,f_r$ are non-identically zero meromorphic quasi-modular forms of weights $k_1<k_2<\cdots<k_r$ with respect to the congruence subgroups $\Gamma_1,\Gamma_2,\ldots,\Gamma_r$ of level $N_1,N_2,\ldots,N_r$ respectively. Let $p$ be the greatest depth of these quasi-modular forms. Suppose that for each $i\in\{1,2,\ldots,r\}$, they have the following transformation formulae:
$$
\begin{aligned}
&f_i\big|_{k_i}\left(\begin{smallmatrix}
a & b\\
c & d
\end{smallmatrix}\right)(\tau)=\sum_{j=0}^p f_{i,j}(\tau)\left(\frac{c}{c\tau+d}\right)^j\\
\text{i.e.,}\quad & f_i\left(\frac{a\tau+b}{c\tau+d}\right)=\sum_{j=0}^p f_{i,j}(\tau)\,c^j\, (c\tau+d)^{k_i-j}
 \end{aligned}
$$
for every $\left(\begin{smallmatrix}
a & b\\
c & d
\end{smallmatrix}\right)\in \Gamma_i$ and $\tau\in\mathbb{H}$. Here the functions $f_{i,0},f_{i,1},\ldots, f_{i,p}$ are the components of $f_i$, and in particular, $f_{i,0}=f_i$ for all $1\leq i\leq r$. Then we consider $N:=\prod_{i=1}^r N_i$ and $S:=\{bN^2+1: b\in\mathbb{N}\}$ such that for each $b'=Nb^2+1\in S$, the matrix $\left(\begin{smallmatrix}
1 & bN\\
N & b'
\end{smallmatrix}\right)\in \cap_{i=1}^r\Gamma_i$. For all such matrices we have 
$$
f_i\left(\frac{\tau+bN}{N\tau+b'}\right)=\sum_{j=0}^p f_{i,j}(\tau)\,N^j\, \left(N\tau+b'\right)^{k_i-j} \quad \text{  for all } 1\leq i\leq r.
$$
Suppose that $\sum_{i=1}^r f_i=0$. Then for $\tau\in\mathbb{H}$ and $b'\in S$, we have 
$$
\sum_{i=1}^r f_i\left(\frac{\tau+bN}{N\tau+b'}\right)=0.
$$
 From the above transformation formula for each $f_i$, we obtain
\begin{equation}\label{1}
\sum_{i=1}^r\sum_{j=0}^p f_{i,j}(\tau)\, N^j\, (N\tau+b')^{k_i-j}=0.
 \end{equation}
For a fixed $\tau\in\mathbb{H}$, multiplying \eqref{1} by $(N\tau+b')^{p'}$, where $p'=2\cdot\max\{p,|k_1|,\ldots,|k_r|\} $,  we get $$
\begin{aligned}
0&=\sum_{i=1}^r\sum_{j=0}^p f_{i,j}(\tau)\, N^j\, (N\tau+b')^{k_i+p'-j}\\
&=\sum_{i=1}^r N^{k_i+p'}\sum_{j=0}^pf_{i,j}(\tau)\left(\tau+\frac{b'}{N}\right)^{k_i+p'-j},
\end{aligned}
$$
which holds for each $b'\in S$. This shows that for any fixed $\tau\in \mathbb{H}$, the polynomial 
$$
P(X)=\sum_{i=1}^r N^{k_i+p'}\sum_{j=0}^pf_{i,j}(\tau)\left(\tau+X\right)^{k_i+p'-j}\in \mathbb{C}[X]
$$
has infinitely many roots, and hence $P(X)$ is identically zero. Note that even if any of the $k_j$'s are negative, the maximum exponent of $X$ occurring in $P(X)$ is $p'+ k_r$. Thus, the leading coefficient of $P(X)$ is $N^{k_r+p'}\, f_{r,0}(\tau)$. Since $P(X)=0$, we get $f_{r,0}(\tau)=f_r(\tau)=0$. This is true for any $\tau\in\mathbb{H}$. Hence, $f_r=0$, which is a contradiction, proving the lemma.  
\end{proof}

We are now ready to prove the main theorem in this context.
\begin{theorem}\label{algebraic-indep-fn-thm}
Let $f$ a non-constant meromorphic modular form of non-zero weight. Let $g$ and $\widetilde{f}$ be a modular function and a meromorphic quasi-modular form of positive depth and non-zero weight, respectively. Assume that $f$, $g$ and $\widetilde{f}$ are of arbitrary level. Then the functions $f$, $g$ and $\widetilde{f}$ are algebraically independent. 
\end{theorem}
\begin{proof}
Let $f$ be of weight $m$ and level $N$, $g$ be of level $M$ and $\widetilde{f}$ be of weight $n$ and level $\widetilde{N}$. Suppose that $P\in \mathbb{C}[X,Y,Z]$ is such that $P(f,g,\widetilde{f})=0$. Recall that a product of a modular form of weight $m$ and level $N$ and a quasimodular form of weight $n$ and level $\widetilde{N}$ is a quasimodular form of weight $mn$ and level $N\widetilde{N}$. Denote $\mathbb{N}_0 = \mathbb{Z}_{\geq 0}$. By grouping the terms of $P(f,g,\widetilde{f})$ of the same weight, we can rewrite it as
\begin{equation}\label{2}
P(f,g,\widetilde{f})=\sum_{k=-K'}^K\,\sum_{\substack{{(r,s,t)\in\mathbb{N}_{0}^3}\\ {mr+nt=k}}} p_{r,s,t} \, \, f^r\,g^s\,{\widetilde{f}}^t , \qquad \text{with   } p_{r,s,t} \in \mathbb{C}.
 \end{equation} 
As $P(X,Y,Z)$ is a polynomial, $p_{r,s,t}=0$ for all but finitely many $r$, $s$ and $t$. As noted above, for each $k\in \{-K',\ldots,K\}$ the inner sum in \eqref{2} is a meromorphic quasi-modular form of weight $k$ and level $NM\widetilde{N}$. Thus, Lemma \ref{indep} implies that
\begin{equation}\label{3}
\sum_{\substack{{(r,s,t)\in\mathbb{N}_0^3}\\ {mr+nt=k}}} p_{r,s,t} \,\, f^r\,g^s\,{\widetilde{f}}^t=0.
\end{equation}
 for each $-K'\leq k\leq K$. \\
 
\noindent Fix a $k\in\{-K',\ldots,K\} $ and consider the corresponding inner sum from \eqref{3}. If there exists an integer $t_0\neq 0$ such that $p_{r,s,t_0}\neq 0$ for at least one tuple $(r,s,t_0)\in\mathbb{N}_0^3$ satisfying $mr+nt_0=k$, then the term on the left in \eqref{3}
is a meromorphic quasi-modular form of depth at least $t_0$. But the uniqueness of depth implies that the term on the left in \eqref{3} is of depth $0$. Hence, the coefficient $p_{r,s,t}=0$ when $t\neq 0$. If $t=0$, then $mr=k$, or in other words, $r = k/m \in \mathbb{Z}$. Thus, for all $0 \neq k\in\{-K',\ldots, K\}$ such that $\frac{k}{m}\in\mathbb{Z}$, the relation in \eqref{3} takes the form  
\begin{equation}\label{5}
\sum_{s \in \mathbb{N}_0} p_{\frac{k}{m},s,0} \,\,g^s = 0,
\end{equation}
as $f$ is not identically zero. If $k=0$, then clearly, $r=0$ and we get
\begin{equation}\label{4}
    \sum_{s \in \mathbb{N}_0} p_{0,s,0} \, \, g^s=0.
\end{equation}
Since $g$ is a non-constant modular function, $g$ must have either a zero or a pole in $\mathbb{H}$. Say that the order of $g$ at $\tau_0$ is $b \neq 0$. Then any function $\sum_{r=0}^R c_r\, g^r$ with not all $c_r = 0$ will have order $bR \neq 0$ at $\tau_0$. Thus, any combination of the form  $\sum_{r=0}^R c_r\, g^r$ cannot be identically zero. Thus, relations \eqref{5} and \eqref{4} imply that $p_{0,s,0}=p_{\frac{k}{m},s,0}=0$ for all $k$ and $s$. Hence $P=0$ and the theorem is proved.
\end{proof}

\bigskip

\section{\bf Proof of Main Results}
\bigskip

\subsection{Values of modular functions when $j(\tau) \not\in \overline{\mathbb{Q}}$}
Before proceeding with the proof of the main theorems, we establish the following intermediary observation.
\begin{lemma} \label{mdfz}
    Let $g\in  \mathcal{F}_{N,\overline{\mathbb{Q}}}$. If $\tau\in \mathbb{H}$ such that $j(\tau)\notin\overline{\mathbb{Q}}$, then $g(\tau)$ is non-zero.
\end{lemma}
\begin{proof}
By Theorem \ref{mdft..}, we know that $g$ satisfies a non-trivial polynomial over $\overline{\mathbb{Q}}(j)$. On clearing denominators, we can assume that $g$ satisfies the irreducible polynomial
$$
P(X)= \sum_{r=0}^m c_r(j) \, X^r =\sum_{r=0}^m \bigg( \sum_{s=0}^{d_r}c_{r s}\,j^s \bigg) \,X^r \in \overline{\mathbb{Q}}[j][X], \quad m=\deg P(X), \, \, d_r=\deg c_r(j),
$$
Note that $c_0(j) \in \overline{\mathbb{Q}}[j]$ is not the identically zero polynomial. Now suppose $\tau \in \mathbb{H}$ is such that $j(\tau)$ is transcendental. If $g(\tau)$ is zero, then we get that $c_0(j(\tau))=0$ which contradicts the transcendence of $j(\tau)$. This completes the proof. 
\end{proof}

\begin{proof}[\bf Proof of Theorem \ref{mdf-values-thm}(a)]
If $g(\tau)=0$, then it follows from Lemma \ref{mdfz} that $j(\tau)\in \overline{\mathbb{Q}}$. If $\tau$ is a pole of $g$, we consider the modular function
$$
h(z):=\frac{1}{g(z)}\in \mathcal{F}_{N,\overline{\mathbb{Q}}},
$$
which vanishes at $z=\tau$, and so $j(\tau)\in\overline{\mathbb{Q}}$ by Lemma \ref{mdfz}.    
\end{proof}

The proof of Theorem \ref{mdf-values-thm}(b) is established in two parts. We first show that if $j(\tau)$ is transcendental, then $g(\tau)$ is as well. The converse implication is proved later.
\begin{prop}\label{tranmf}
 Let $g\in \mathcal{F}_{N,\overline{\mathbb{Q}}}$ and $\tau\in \mathbb{H}$ be such that it is not a pole of $g$. If $j(\tau)$ is transcendental, then $g(\tau)$ is transcendental.
\end{prop}
\begin{proof}
Suppose that $g(\tau)\in\overline{\mathbb{Q}}$. Then the modular function $$h(z):=g(z)-g(\tau)$$ belongs to $ \mathcal{F}_{N,\overline{\mathbb{Q}}}$ and vanishes at $\tau$. This contradicts Lemma \ref{mdfz}. Therefore, $g(\tau)$ is transcendental.
\end{proof}


\subsection{Values of modular functions when $j(\tau)\in\overline{\mathbb{Q}}$}
From Theorem \ref{mdft..}, we have $$ \mathcal{F}_{N,\overline{\mathbb{Q}}}=  \overline{\mathbb{Q}}\left(\{j,f_{\boldsymbol{a}} : {\boldsymbol{a}}\in S\}\right), \text{ where }S=\left\{\left(\frac{r}{N},\frac{s}{N}\right): 0\leq r, s\leq N-1 \text{ and }(r,s)\neq (0,0)\right\}.$$
Hence, for any $g\in \mathcal{F}_{N,\overline{\mathbb{Q}}}$, the algebraic nature of $g(\tau)$ is determined by the numbers in the set $\{f_{\boldsymbol{a}}(\tau):{\boldsymbol{a}}\in S\}$.\\

For $N=1$, $\mathcal{F}_{1,\overline{\mathbb{Q}}} = \overline{\mathbb{Q}}(j)$. Hence, for any $g \in \mathcal{F}_{1,\overline{\mathbb{Q}}}$, $j(\tau) \in \overline{\mathbb{Q}}$ implies that $g(\tau) \in \overline{\mathbb{Q}}$. For $N>1$, we have to study the nature of the values $f_{\boldsymbol{a}}(\tau)$ for all ${\boldsymbol{a}}\in S$. We start by proving the following important lemma.
\begin{lemma}\label{omegatau}
 If $\tau\in \mathbb{H}$ such that $j(\tau)$ is algebraic, then there exists a unique transcendental number (up to an algebraic multiples) $\omega_\tau$ for which $g_2(\omega_\tau L_\tau)$ and $g_3(\omega_\tau L_\tau)$ are both algebraic numbers.   
\end{lemma}
\begin{proof}
Note that for any $\alpha \in\mathbb{C}^\times $,
\begin{equation}\label{omegatau.}
    j(\tau)=\frac{E_4(\tau)^3}{\Delta(\tau)}=\frac{1728\,  g_2(L_\tau)^3}{\Delta_0(L_\tau)}=\frac{1728\, g_2\left(\alpha L_\tau\right)^3}{\Delta_0\left(\alpha L_\tau\right)},
\end{equation}
because of the homogeneity properties of the $g_2$, $g_3$ functions. Here $L_\tau=\tau\,\mathbb{Z}\oplus\mathbb{Z}$ and $\alpha L_\tau= \alpha \, \tau\,\mathbb{Z}\oplus\alpha \, \mathbb{Z}$. If we choose $\omega_\tau$ such that $\omega_\tau^4=g_2(L_\tau)$, then $g_2(\omega_\tau L_\tau)=1$. Given that $j(\tau)$ is algebraic, from \eqref{omegatau.}, we get that $\Delta_0(\omega_\tau L_\tau)$ is also algebraic. Thus, the numbers $g_2(\omega_\tau L_\tau)$ and $g_3(\omega_\tau L_\tau)$ are both algebraic. Moreover, $\omega_\tau$ is a period of $\omega_\tau L_\tau$ and both $g_2(\omega_\tau L_\tau)$ and $g_3(\omega_\tau L_\tau)$ are algebraic. By Corollary \ref{tranp} we conclude that $\omega_\tau$ is transcendental.\\

To prove uniqueness, consider an arbitrary $\omega_\tau'\in\mathbb{C}^\times$ such that both $g_2(\omega_\tau' L_\tau)$ and $g_3(\omega_\tau' L_\tau)$ are algebraic. By homogeneity of $g_2$, we have 
$$
 \quad g_2(\omega_\tau' L_\tau)= (\omega_\tau')^{-4}\,g_2(L_\tau)=:\beta \in \overline{\mathbb{Q} }.
$$
Then, we get $g_2(L_\tau)=\beta \cdot (\omega_\tau')^{4}$ and hence $\omega_\tau^4=\beta \cdot (\omega_\tau')^{4}.$
This completes the proof that $\omega_\tau$ is unique up to algebraic multiples.
\end{proof}

Now recall that
$$
f_{\boldsymbol{a}}(\tau):=\frac{g_2(L_\tau)\,g_3(L_\tau)}{\Delta_0(L_\tau)}\,\wp_{L_\tau}(a_1\tau+a_2),
$$ 
where $L_\tau=\tau\,\mathbb{Z}\oplus\mathbb{Z}$ and ${\boldsymbol{a}}=(a_1,a_2)\in S$. Using homogeneity properties of the functions involved, we can rewrite this as follows:
$$
f_{\boldsymbol{a}}(\tau)=\frac{g_2(\omega_\tau L_\tau)\, g_3(\omega_\tau L_\tau)}{\Delta_0(\omega_\tau L_\tau)}\,\wp_{\omega_\tau L_\tau}(a_1\omega_\tau\tau+a_2\omega_\tau).
$$
\begin{lemma}\label{algfa}
    If $\tau\in\mathbb{H}$ is such that $j(\tau)\in\overline{\mathbb{Q}}$, then $f_{\boldsymbol{a}}(\tau)\in\overline{\mathbb{Q}}$ for all ${\boldsymbol{a}}\in S$.
\end{lemma}
\begin{proof}
    Let ${\boldsymbol{a}}=\left( \frac{r}{N},\frac{s}{N} \right)\in S$. Then we have
$$
f_{\boldsymbol{a}}(\tau)=\frac{g_2(\omega_\tau L_\tau)\, g_3(\omega_\tau L_\tau)}{\Delta_0(\omega_\tau L_\tau)}\,\wp_{\omega_\tau L_\tau}\left(\frac{r\,\omega_\tau\,\tau}{N}+\frac{s\,\omega_\tau}{N}\right),
$$
where $\omega_\tau$ is chosen as in Lemma \ref{omegatau} so that $g_2(\omega_\tau L_\tau),\, g_3(\omega_\tau L_\tau)\in\overline{\mathbb{Q}}$.\\

Recall that both $\wp_{\omega_\tau L_\tau}\left(\frac{\omega_\tau\, \tau}{N} \right)$ and $\wp_{\omega_\tau L_\tau}\left(\frac{\omega_\tau}{N} \right)$  are algebraic by Corollary \ref{tranp}. Moreover, by the addition formula for Weierstrass $\wp$-function in Theorem \ref{addf}, we get that $$\wp_{\omega_\tau L_\tau}\left(\frac{r\,\omega_\tau\, \tau}{N}+\frac{s\,\omega_\tau}{N}\right)\in\overline{\mathbb{Q}}$$ for all ${\boldsymbol{a}}=\left(\frac{r}{N},\frac{s}{N}\right)\in S$. Thus, the number $f_{\boldsymbol{a}}(\tau)$ is algebraic for all ${\boldsymbol{a}}\in S$.
\end{proof}
\begin{prop}\label{algemf}
Let $g\in \mathcal{F}_{N,\overline{\mathbb{{Q}}}}$ and $\tau\in\mathbb{H}$ such that $j(\tau)\in \overline{\mathbb{Q}}$. If $\tau$ is not a pole of $g$, then $g(\tau)$ is algebraic.
\end{prop}
\begin{proof}
 From Lemma \ref{algfa}, we have   $f_{\boldsymbol{a}}(\tau)\in\overline{\mathbb{Q}}$ for all ${\boldsymbol{a}}\in S$. Since 
 $$g \in \mathcal{F}_{N,\overline{\mathbb{Q}}}=  \overline{\mathbb{Q}}\left(\{j,f_{\boldsymbol{a}} |{\boldsymbol{a}}\in S\}\right)$$
 by Theorem \ref{mdft..}, $g(\tau)$ is also algebraic. 
\end{proof}
\bigskip

Propositions \ref{tranmf} and \ref{algemf} together complete the proof of Theorem \ref{mdf-values-thm}(b). We now prove the generalization of \cite[Theorem 2.4]{jeon-kang-kim}.
\begin{proof}[\bf Proof of Theorem \ref{zero-pole-value}]
Let $f$ be a meromorphic modular form of weight $k \in \mathbb{Z}$. Define 
   $$
h(\tau):=\frac{f(\tau)^{12}}{\Delta(\tau)^k}\quad \forall \tau\in\mathbb{H}.
   $$
The function $h \in \mathcal{F}_{N,\overline{\mathbb{Q}}}$ because the Fourier coefficients of $f$ and $\Delta$ are algebraic. Moreover, any $\tau\in\mathbb{H}$ is a zero or a pole of $h$ if and only if it is a zero or a pole of $f$ because $\Delta$ is non-vanishing and holomorphic on $\mathbb{H}$. Thus, if $\tau\in\mathbb{H}$ is a zero or pole of $f$, then $h$ is non-constant and by Theorem \ref{mdf-values-thm}(a), we deduce that $j(\tau)\in\overline{\mathbb{Q}}$, and by Theorem \ref{mdf-values-thm}(b), that either $g(z)$ has a pole at $z = \tau$ or that $g(\tau)\in\overline{\mathbb{Q}}$.
\end{proof}
\bigskip

\subsection{Values of modular forms} Recall that for any integer $k$ and $N$, we denote $M_{k,N,\overline{\mathbb{Q}}}$, $M_{k,N,\overline{\mathbb{Q}}}^w$ and $M_{k,N,\overline{\mathbb{Q}}}^m$ to be the set of holomorphic, weakly holomorphic and meromorphic modular forms respectively of level $N$ with algebraic Fourier coefficients at $i\infty$. By Corollary \ref{cor-main-mdf}, all elements of these sets have algebraic Fourier coefficients with respect to all cusps as well.\\

\begin{proof}[Proof of Theorem \ref{general-mero-values-thm}(a) ]
Let $f \in M^m_{k,N,\overline{\mathbb{Q}}}$. Consider the modular function 
\begin{equation*}
    g(\tau) := \frac{{f(\tau)}^{12}}{\Delta(\tau)^k}.
\end{equation*}
If $g$ is a non-zero constant, then $f(\tau) = c \Delta(\tau)^{k/12}$ for some $c \in \overline{\mathbb{Q}}^{\times}$ and all $\tau \in \mathbb{H}$. Since $1728 \Delta(\tau) = E_4^3(\tau) - E_6^2(\tau)$, and $E_4(\tau)$ is algebraically independent with $E_6(\tau)$ for $\tau \in \mathbb{H}$ such that $e^{2 \pi i \tau} \in \overline{\mathbb{Q}}$ by Nesterenko's theorem \ref{nest}, we deduce that $f(\tau)$ is transcendental. \\

Now suppose that $g$ is non-constant. Since $f \in M_{k,N,\overline{\mathbb{Q}}}$, $g \in \mathcal{F}_{N,\overline{\mathbb{Q}}}$. Note that for $\tau_0 \in \mathbb{H}$ such that $e^{2 \pi i \tau_0}$ is algebraic, $j(\tau_0)$ is transcendental by Theorem \ref{mahl}. Such a $\tau_0$ cannot be a pole of $f$. For if $\tau_0$ is a pole of $f$, then $\tau_0$ is a pole of $g$ and by Theorem \ref{mdf-values-thm}(a), $j(\tau_0)$ would be algebraic, leading to a contradiction. As $\tau_0$ is not a pole of $g$, Theorem \ref{mdf-values-thm}(b) implies that $g(\tau_0)$ is transcendental. Moreover, by Theorem \ref{mdft..}, $g(\tau_0)$ is algebraic over $\overline{\mathbb{Q}}(j(\tau_0))$. By Nesterenko's theorem, 
\begin{equation*}
    \Delta(\tau_0) = \frac{E_4^3(\tau_0) - E_6^2(\tau_0)}{1728} \qquad \text{ and } \qquad j(\tau_0) = \frac{E_4^3(\tau_0)}{\Delta(\tau_0)}
\end{equation*}
are algebraically independent. Therefore,
\begin{equation*}
    \text{trdeg}_{\,\mathbb{Q}} \, \, \overline{\mathbb{Q}} \left ( g(\tau_0), \, \Delta(\tau_0) \right) = \text{trdeg}_{\,\mathbb{Q}}\, \, \overline{\mathbb{Q}} \left ( j(\tau_0), \, \Delta(\tau_0) \right) = 2.
\end{equation*}
Hence, $f(\tau_0)^{12} = g(\tau_0) \, \Delta(\tau_0)^k$ is transcendental, proving the claim.
\end{proof}








% We first establish a lemma which is crucial to the proof of Theorem \ref{general-mero-values}.

% \begin{lemma}\label{det}
% Let $n$ and $m$ be positive integers such that $n \geq 2m$, then the determinant of the matrix $A = [a_{ij}]_{(m+1) \times (m+1)}$ with $a_{ij} = {n+1-j \choose i-1}$
% %     $$
% % A=\left[\begin{array}{cccccc}
% %     {{n}\choose{0}} & {{n-1}\choose{0}} & {{n-2}\choose{0}} & \ldots & {{n-m+1}\choose{0}} & {{n-m}\choose{0}} \\\\
% %     {{n}\choose{1}} & {{n-1}\choose{1}} & {{n-2}\choose{1}} & \ldots & {{n-m+1}\choose{1}} & {{n-m}\choose{1}}\\\\
% %     {{n}\choose{2}} & {{n-1}\choose{2}} & {{n-2}\choose{2}} & \ldots & {{n-m+1}\choose{2}} & {{n-m}\choose{2}}\\\\
% %     \vdots & \vdots & \vdots & \ldots & \vdots & \vdots \\\\
% %     {{n}\choose{m}} & {{n-1}\choose{m}} & {{n-2}\choose{m}} & \ldots & {{n-m+1}\choose{m}} & {{n-m}\choose{m}}
% % \end{array} \right]_{(m+1)\times (m+1)}
% %   $$ 
%   is given by $(-1)^{\frac{m(m+1)}{2}}$.
% \end{lemma}
% \begin{proof}
% We will use the identity ${n\choose{i}} - {n-1\choose{i}} ={n-1\choose{i-1}}$ for $1 \leq j \leq m$. On applying the column operation $C_j \mapsto C_j - C_{j+1}$ successively for $1 \leq j \leq m$, the matrix $A$ transforms into 
% $$
% A'=\left[\begin{array}{cccccc}
%     0 & 0 & 0 & \ldots & 0 & {{n-m}\choose{0}} \\\\
%     {{n-1}\choose{0}} & {{n-2}\choose{0}} & {{n-3}\choose{0}} & \ldots & {{n-m}\choose{0}} & {{n-m}\choose{1}}\\\\
%     {{n-1}\choose{1}} & {{n-2}\choose{1}} & {{n-3}\choose{1}} & \ldots & {{n-m}\choose{1}} & {{n-m}\choose{2}}\\\\
%     \vdots & \vdots & \vdots & \ldots & \vdots & \vdots \\\\
%     {{n-1}\choose{m-1}} & {{n-2}\choose{m-1}} & {{n-3}\choose{m-1}} & \ldots & {{n-m}\choose{m-1}} & {{n-m}\choose{m}}
% \end{array} \right]_{(m+1)\times (m+1)}.
% $$
% We can now repeat the same column operations as above, repeatedly, to reduce the matrix to a lower triangular form 
%  $$
% B=\left[\begin{array}{cccccc}
%     0 & 0 & 0 & \ldots & 0 & {{n-m}\choose{0}} \\\\
%     0 & 0 & 0 & \ldots & {{n-m}\choose{0}} & *\\\\
%     0 & 0 & {{n-m}\choose{0}} & \iddots & * & *\\\\
%     \vdots & \iddots & \vdots & \ldots & \vdots & \vdots \\\\
%     {{n-m}\choose{0}} & * & * & \ldots & * & *
% \end{array} \right]_{(m+1)\times (m+1)}.
%  $$
% Since $B$ is obtained from $A$ by elementary column operations that do not change the determinant, we conclude that $\det(A)=\det(B)=(-1)^{\frac{m(m+1)}{2}}$.
% \end{proof}



% \begin{prop}\label{algem'}
%    Let $f$ be a non-constant meromorphic modular form with algebraic Fourier coefficients. If $\tau\in\mathbb{H}$ such that $e^{2\pi i\tau}$ is algebraic, then $f(\tau)$ is transcendental. 
% \end{prop}
% \begin{proof}[Proof of Theorem \ref{general-mero-values}(a)]
% Consider the modular function
% $$g(\tau):=\frac{f(\tau)^{12}}{\Delta(\tau)^k}\quad \forall\, \tau\in\mathbb{H},$$
% which lies in $\mathcal{F}_{N,\overline{\mathbb{Q}}}$. By Theorem \ref{mdft..}, we know that $g$ has the minimal polynomial $P[X]\in \overline{\mathbb{Q}}(j)[X]$, and we can clear denominators of the coefficients of $P[X]$ to get a polynomial 
% $$P(X)=\sum_{r=0}^m c_r(j)\,X^r\in \overline{\mathbb{Q}}[j][X],\quad \deg P(X)=m$$
% such that
% $P(g(\tau))=0$ for any $\tau\in \mathbb{H}$. Here, we have
% $$
% c_r(j)=\sum_{s=0}^{d_r}c_{r,s}\,j^s\in \overline{\mathbb{Q}}[j]\quad \forall\ 0 \leq r\leq m
% $$
% and $d_r=\deg c_r(j)$ so that $c_{r,d_r}\neq 0$ for all $0\leq r\leq m$. Moreover, Corollary \ref{useful} shows that if $f$ is holomorphic on $\mathbb{H}$, then the minimal polynomial $P[X]$ of $g$ is monic and belongs to $\overline{\mathbb{Q}}[j][X]$, and hence $c_m(j)=1$.\\

% For any $\tau\in\mathbb{H}$, we have $$P\left(\frac{f(\tau)^{12}}{\Delta(\tau)^k}\right)=0,$$
% and hence for any positive integer $l$, 
% $$
% \begin{aligned}
% & \Delta(\tau)^l\, P\left(\frac{f(\tau)^{12}}{\Delta(\tau)^k}\right)=0\\
% \Rightarrow &\sum_{r=0}^m\sum_{s=0}^{d_r}c_{r,s}\,j(\tau)^s f(\tau)^{12 r}\Delta(\tau)^{l-k r}=0\\
% \Rightarrow &\sum_{r=0}^m\sum_{s=0}^{d_r}c_{r,s}\,E_4(\tau)^{3 s}f(\tau)^{12\, r}\Delta(\tau)^{l-k r-s}=0.
% \end{aligned}
% $$
% We choose $l> k m +2\cdot\operatorname{max}\{d_0,d_1,\ldots,d_m\}$ so that the powers of $\Delta(\tau)$ are always positive integers. Then we put $$\Delta(\tau)=\frac{E_4(\tau)^3-E_6(\tau)^2}{1728}$$ to obtain \begin{equation} \label{trane}
% \sum_{r=0}^m \sum_{s=0}^{d_r} \sum_{t=0}^{l-k r-s} \frac{(-1)^t c_{r, s}}{1728^{l-k r-s}}\left(\begin{array}{c}
% l-k\,r-s \\
% t
% \end{array}\right) f(\tau)^{12\, r} E_4(\tau)^{3(l-k\, r-t)} E_6(\tau)^{2\, t}=0.   
% \end{equation}
% Now we consider $\tau\in \mathbb{H}$ such that $e^{2\pi i\tau}$ is algebraic. By Theorem \ref{schneider}, the number $j(\tau)$ is transcendental, and hence Theorem \ref{zero-pole} guarantees that $\tau$ is neither a zero nor a pole of $f$. Moreover, we claim that the number $f(\tau)$ is not algebraic. For otherwise, \eqref{trane} becomes a polynomial in two variables over $\overline{\mathbb{Q}}$, satisfied by $E_4(\tau)$  and $E_6(\tau)$. If we can prove that this is a non-zero polynomial of $E_4(\tau)$ and $E_6(\tau)$, then we would get a contradiction with Nesterenko's Theorem \ref{nest}. Hence, it would complete the proof that $f(\tau)$ is transcendental.\\

% To show that the left side of \eqref{trane} becomes a non-zero polynomial of $E_4(\tau)$ and $E_6(\tau)$ when $e^{2\pi i\tau}\in\overline{\mathbb{Q}}$, we consider the following cases:\\

% \emph{Case-$1$: } If $f\in M_{k,N,\overline{\mathbb{Q}}}\,\cup M_{k,N,\overline{\mathbb{Q}}}^w$, then $c_m(j)=1$, and hence the coefficient of $E_4(\tau)^{3(l-k m)}\,E_6(\tau)^{0}$ (which appears in \eqref{trane} only for $r=m,\, t=0$) is given by $$\frac{1}{1728^{l-k m}}\cdot 1\cdot f(\tau)^{12 m}\neq 0,$$
% which proves that the polynomial is non-zero.\\

% \emph{Case-$2$: } If $f\in M_{k,N,\overline{\mathbb{Q}}}^m\setminus \left(M_{k,N,\overline{\mathbb{Q}}}\,\cup M_{k,N,\overline{\mathbb{Q}}}^w\right)$, then the leading coefficient of $P[X]$ is given by $$c_m(j)=\sum_{s=0}^{d_m}c_{m,s}\,j^s.$$
% Moreover, it cannot be \emph{identically zero polynomial} because the degree of $P[X]$ is same as the degree of the minimal polynomial of $g$. Since $l>k\, m+2\, d_m$, the following terms
% $$
% E_4(\tau)^{3(l-k\, m-t)}\, E_6(\tau)^{2\, t} \quad \forall\, 0\leq t\leq d_m
% $$
% with coefficients
% $$
% \frac{(-1)^t\, f(\tau)^{12 m}}{1728^{l-km}}\sum_{s=0}^{d_m}\left(\begin{array}{c}
% l-k\,m-s \\
% t
% \end{array}\right) (1728)^s\, c_{m,s} \quad \forall\, 0\leq t\leq d_m
% $$
% appear in \eqref{trane}. We claim that at least one of these terms has non-zero coefficient. If not, we consider the following system of equations: $ A\, X=0$, where 
% $$
% A=\left[\begin{array}{ccccc}
%     {{l-km}\choose{0}} & {{l-km-1}\choose{0}} & {{l-km-2}\choose{0}} & \ldots & {{l-km-d_m}\choose{0}} \\\\
%     {{l-km}\choose{1}} & {{l-km-1}\choose{1}} & {{l-km-2}\choose{1}} & \ldots & {{l-km-d_m}\choose{1}}\\\\
%     {{l-km}\choose{2}} & {{l-km-1}\choose{2}} & {{l-km-2}\choose{2}} & \ldots & {{l-km-d_m}\choose{2}}\\\\
%     \vdots & \vdots & \vdots & \ldots & \vdots \\\\
%     {{l-km}\choose{d_m}} & {{l-km-1}\choose{d_m}} & {{l-km-2}\choose{d_m}} & \ldots & {{l-km-d_m}\choose{d_m}}
% \end{array} \right]_{(d_m+1)\times (d_m+1)}  X=\left[\begin{array}{c}
%    x_0  \\
%     x_1 \\
%     x_2 \\
%     \vdots\\
%     x_{d_m}
% \end{array}\right]_{d_m+1.}
% $$
% Then this system of equations is satisfied by the non-trivial vector $$\left[\begin{array}{c}
%    (1728)^0\, c_{m,0} \\
%     (1728)^1\, c_{m,1} \\
%     (1728)^2\, c_{m,2} \\
%     \vdots\\
%     (1728)^{d_m}\, c_{m,d_m}
% \end{array}\right]_{(d_m+1)\times 1}.$$ From Lemma \ref{det}, we get that $\det(A)\neq 0$, which contradicts the fact that the above system has only trivial solution. Thus, it completes the proof that the polynomial relation of $E_4(\tau)$ and $E_6(\tau)$ in \eqref{trane} is non-trivial.
% \end{proof}

We remark here that the proof of Theorem \ref{general-mero-values-thm}(a) only requires Nesterenko's theorem and the structure of modular functions of higher level. One also immediately deduces the following. 
\begin{prop}\label{algi-cor}
If $f\in M^m_{k,N,\overline{\mathbb{Q}}}$, then $f$ is algebraically dependent with $E_4$ and $E_6$ over $\overline{\mathbb{Q}}$. In particular, If $f\in M_{k,N,\overline{\mathbb{Q}}}^m$ and $\tau\in\mathbb{H}$ is not a pole of $f$, then $f(\tau)$ is algebraic over $\overline{\mathbb{Q}}(E_4(\tau),E_6(\tau))$. 
\end{prop}
\begin{proof}
Let $f \in M^m_{k,N,\overline{\mathbb{Q}}}$ and consider 
\begin{equation*}
    g (\tau) := \frac{f^{12}(\tau)}{\Delta^k(\tau)} \in \mathcal{F}_{N,\overline{\mathbb{Q}}}.
\end{equation*}
By Corollary \ref{cor-main-mdf}, there exists a polynomial $P(X) \in \overline{\mathbb{Q}}[j](X)$ such that $P(g) = 0$. More specifically,
$$\sum_{r=0}^m \sum_{s=0}^{d_r} c_{r,s} \, {j(\tau)}^s \,{g(\tau)}^r= 0 \qquad \text{for all} \qquad \tau \in \mathbb{H}.$$
Here $c_{r,s} \in \overline{\mathbb{Q}}$ for all $0 \leq s \leq d_r$ and $0 \leq r \leq m$. Multiplying by $\Delta^l(\tau)$ for any positive integer $l > km$ and substituting $j(\tau)$ and $\Delta(\tau)$ in terms of $E_4(\tau)$ and $E_6(\tau)$ gives
\begin{equation*} 
\sum_{r=0}^m \sum_{s=0}^{d_r} \sum_{t=0}^{l-k r-s} \frac{(-1)^t c_{r, s}}{1728^{l-k r-s}}\left(\begin{array}{c}
l-k\,r-s \\
t
\end{array}\right) f(\tau)^{12\, r} E_4(\tau)^{3(l-k\, r-t)} E_6(\tau)^{2\, t}=0.   
\end{equation*}
This proves the proposition.
\end{proof}


\noindent We now consider the complementary case, namely, points $\tau\in\mathbb{H}$ such that $j(\tau)\in\overline{\mathbb{Q}}$. 
\begin{proof}[Proof of Theorem \ref{general-mero-values-thm}(b)]
Fix a $\tau\in\mathbb{H}$ such that $j(\tau)\in\overline{\mathbb{Q}}$. From Lemma \ref{omegatau}, we get a transcendental number $\omega_\tau$ such that $g_2(\omega_\tau L_\tau)$ and $g_3(\omega_\tau L_\tau)$ are both algebraic. Moreover, we have the formulae
$$
\begin{aligned}
E_4(\tau)&=\frac{3}{4\pi^4}\, g_2(L_\tau)=\frac{3}{4}\left(\frac{\omega_\tau}{\pi} \right)^4 g_2(\omega_\tau L_\tau),\\ E_6(\tau)&=\frac{27}{8\pi^6}\,g_3(L_\tau)=\frac{27}{8}\left(\frac{\omega_\tau}{\pi}\right)^6g_3(\omega_\tau L_\tau).
\end{aligned}
$$
From Nesterenko's Theorem \ref{nest}, we know that at most one of $E_4(\tau)$ and $E_6(\tau)$ is algebraic. The above formulae imply that if ${\omega_{\tau}}/{\pi}\in\overline{\mathbb{Q}}$, then both $E_4(\tau)$ and $E_6(\tau)$ are algebraic, which is a contradiction. Hence, ${\omega_\tau}/{\pi}$ is a transcendental number. Besides, we have 
$$
\begin{aligned}
\Delta(\tau)=\frac{E_4(\tau)^3-E_6(\tau)^2}{1728}&=\frac{1}{4^6}\left(\frac{\omega_\tau}{\pi}\right)^{12}\bigg( g_2(\omega_\tau L_\tau)^3-27 g_3(\omega_\tau L_\tau)^2 \bigg)\\
&=\frac{1}{4^6}\left(\frac{\omega_\tau}{\pi}\right)^{12}\Delta_0(\omega_\tau L_\tau).  
\end{aligned}
$$
As $g_2(\omega_\tau L_\tau)$ and $g_3(\omega_\tau L_\tau)$ both are algebraic, the number $\Delta_0(\omega_\tau L_\tau)\in\overline{\mathbb{Q}}\setminus \{0\}$. Since ${\omega_\tau}/{\pi}$ is transcendental and $\Delta(\tau)$ is a non-zero algebraic multiple of $\left({\omega_\tau}/{\pi}\right)^{12}$, we deduce that $\Delta(\tau)$ is transcendental. \\

Consider the modular function $$g(\tau):=\frac{f(\tau)^{12}}{\Delta(\tau)^k}\quad \forall \tau\in\mathbb{H},$$ which lies in $ \mathcal{F}_{N,\overline{\mathbb{Q}}}$. Since $j(\tau)\in\overline{\mathbb{Q}}$, Theorem \ref{algemf} implies that $g(\tau)$ is algebraic, say $\alpha$. Thus, we get that
\begin{equation}\label{forms-value-j-alg}
f(\tau)^{12}= \alpha \, \cdot \, \Delta(\tau)^k=\frac{\alpha }{4^6} \, \cdot \, \Delta_0(\omega_\tau L_\tau)^k \, \left(\frac{\omega_\tau}{\pi} \right)^{12 k}. 
\end{equation}
This shows that if $f(\tau) \neq 0$, 
 then it is a non-zero algebraic multiple of $\left({\omega_\tau}/{\pi} \right)^{k}$ and hence, is transcendental.
\end{proof}
\bigskip


\subsection{Values of quasi-modular forms}
Let $\widetilde{M}_{k,\overline{\mathbb{Q}}}^{(p)}(\Gamma)$ denote the set of all meromorphic quasi-modular forms of weight $k$ and depth $p(>0)$ for $\Gamma$ with algebraic Fourier coefficients. We study their values at the points $\tau\in\mathbb{H}$, where exactly one of $e^{2\pi i\tau}$ and $j(\tau)$ is algebraic.\\

\begin{proof}[Proof of Theorem \ref{quasi-mdf-thm}(a)]
By Theorem \ref{quasistructure}, we can write $\widetilde{f}$ in the form $$\widetilde{f}=\sum_{r=0}^p f_r E_2^r,\quad f_r\in M_{k,N,\overline{\mathbb{Q}}}.$$ 
Suppose that both $e^{2\pi i\tau}, \widetilde f(\tau)\in \overline{\mathbb{Q}}$. By Corollary \ref{algi-cor}, we know that each number $f_r(\tau)$ is algebraically dependent with $E_4(\tau), E_6(\tau)$. Thus, we get that $E_2(\tau)$ is algebraic over $\overline{\mathbb{Q}}(E_4(\tau), E_6(\tau))$. This implies that
\begin{equation*}
    \text{trdeg}_{\,\mathbb{Q}} \, \,\overline{\mathbb{Q}} \bigg( e^{2 \pi i \tau}, E_2(\tau), E_4(\tau), E_6(\tau) \bigg) \leq 2,
\end{equation*}
contradicting Nesterenko's theorem \ref{nest}. This proves the claim.
\end{proof}


To study values of quasi-modular forms in the complementary case (i.e., for $j(\tau)\in\overline{\mathbb{Q}}$), we need the following lemma.
\begin{lemma}\label{peri}
 Let $\tau\in \mathbb{H}$ be such that $j(\tau)$ is algebraic and $\omega_{\tau}$ be the transcendental number determined in Lemma \ref{omegatau} such that $g_2(\omega_{\tau}L_{\tau})$ and $g_3(\omega_{\tau}L_{\tau)}$ are algebraic. Let $\eta_2(\omega_{\tau}) := \eta_2(\omega_{\tau} L_{\tau})$ be the quasi-period 
 \begin{equation*}
     \eta_2( \omega_{\tau} L_{\tau}) = \zeta_{\omega_{\tau}L_{\tau}}(z+\omega_{\tau}) - \zeta_{\omega_{\tau}L_{\tau}}(z).
 \end{equation*}
 Then $\frac{\omega_\tau}{\pi}$ and $\frac{\eta_2}{\pi}$ are algebraically independent over $\overline{\mathbb{Q}}$. 
\end{lemma}
\begin{proof}
From the definition of the Weierstrass zeta-function, one gets
 $   
{\zeta}_{\omega_\tau L_\tau}(z)= \frac{1}{\omega_\tau}\cdot \zeta_{L_\tau}\left(\frac{z}{\omega_\tau}\right).
$ 
Using the definition of a quasi-period and the identity \eqref{quasi-period}, we obtain
$$
\eta_2= \zeta_{\omega_\tau L_\tau}(z+\omega_\tau)-\zeta_{\omega_\tau L_\tau}(z)= \frac{1}{\omega_\tau}\,\eta_2(L_\tau)=\frac{1}{\omega_\tau}\,G_2(\tau)= \frac{1}{3}\cdot\frac{\pi^2}{ \omega_\tau}\,E_2(\tau).
$$
Thus, we get the following formulae 
$$
E_2(\tau)=3\,\frac{\omega_\tau}{\pi}\frac{\eta_2}{\pi}, \quad E_4(\tau)=\frac{3}{4}\left(\frac{\omega_\tau}{\pi} \right)^4 g_2(\omega_\tau L_\tau),\quad E_6(\tau)=\frac{27}{8}\left(\frac{\omega_\tau}{\pi}\right)^6g_3(\omega_\tau L_\tau).$$
The above formulae imply that $E_2(\tau)$, $E_4(\tau)$ and $E_6(\tau)$ are algebraic over $\overline{\mathbb{Q}}\left(\frac{\omega_\tau}{\pi}, \frac{\eta_2}{\pi}\right)$. Suppose that $\frac{\omega_\tau}{\pi}, \frac{\eta_2}{\pi}$ are algebraically dependent. This implies that $$ \text{trdeg}_{\,\mathbb{Q}} \, \,\overline{\mathbb{Q}} \bigg(  E_2(\tau), E_4(\tau), E_6(\tau) \bigg)=1,$$ which contradicts Nesterenko's Theorem \ref{nest}.
\end{proof}

Lemma \ref{peri} and Theorem \ref{quasistructure} together allow us to describe values of quasi-modular forms explicitly at the points $\tau\in\mathbb{H}$ where the $j$-function is algebraic.
\begin{proof}[Proof of Theorem \ref{quasi-mdf-thm}(b)]
From Theorem \ref{quasistructure}, we have the expression $$\widetilde{f}(\tau)=\sum_{r=0}^p f_r(\tau) E_2(\tau)^r,\quad \text{where } f_r\in M^m_{k,N,\overline{\mathbb{Q}}}.$$
Writing the value $f_r(\tau)$ as in \eqref{forms-value-j-alg} for each coefficient $f_r$, and using the above formula for $E_2$, we obtain
\begin{equation}\label{last}
\widetilde{f}(\tau)=\sum_{r=0}^p c_r\left(\frac{\omega_\tau}{\pi}\right)^{k-2 r}\left(\frac{\omega_\tau}{\pi}\cdot\frac{\eta_2}{\pi}\right)^r=\sum_{r=0}^p\, c_r\, \left(\frac{\omega_\tau}{\pi}\right)^{k-r}\, \left(\frac{\eta_2}{\pi}\right)^r,
 \end{equation}
where each $c_r$ is an algebraic number. If  $c_r \neq 0$ for some $r$ satisfying $0 \leq r \leq p$, then \eqref{last} gives a non-trivial algebraic relation among $\frac{\omega_\tau}{\pi}$ and $\frac{\eta_2}{\pi}$. This contradicts Lemma \ref{peri}. Therefore, the number $\widetilde{f}(\tau)$ is either zero, precisely when each $\widetilde{f}(\tau)\in\overline{\mathbb{Q}}\setminus \{0\}$, or is transcendental.     
\end{proof}


\bigskip

\subsection{Algebraic independence of special values}

\begin{proof}[Proof of Theorem \ref{tr-deg-thm}]
From Theorem \ref{mdft..}, we know that for any $\tau \in \mathbb{H}$ which is not a pole of $g$, $g(\tau)$ is algebraic over $\overline{\mathbb{Q}}(j(\tau))$. Moreover, Corollary \ref{algi-cor} gives that $f(\tau)$ is algebraic over $\overline{\mathbb{Q}}(E_4(\tau), E_6(\tau))$. Since $\widetilde{f}(\tau)=\sum_{r=0}^p f_r(\tau) E_2(\tau)^r$, we have that $\widetilde{f}(\tau)$ is algebraic over $${\overline{\mathbb{Q}}(E_2(\tau), E_4(\tau), E_6(\tau)).}$$ Hence, we get that
\begin{align*}
    \text{trdeg}_{\, \mathbb{Q}} \, \, \overline{\mathbb{Q}} \bigg(e^{2 \pi i \tau}, g(\tau), f(\tau), \widetilde{f}(\tau) \bigg) & = \text{trdeg}_{\, \mathbb{Q}} \, \, \overline{\mathbb{Q}} \bigg(e^{2 \pi i \tau}, j(\tau), E_2(\tau), E_4(\tau), E_6(\tau) \bigg) \\
    & = \text{trdeg}_{\, \mathbb{Q}} \, \, \overline{\mathbb{Q}} \bigg(e^{2 \pi i \tau}, E_2(\tau), E_4(\tau), E_6(\tau) \bigg) \geq 3,
\end{align*}
by Nesterenko's theorem \ref{nest}. This establishes the claim.
\end{proof}


\section*{Acknowledgment}
The authors thank Prof. Michel Waldschmidt and Dr. Veekesh Kumar for insightful suggestions on an earlier version of the paper. They are also grateful to the referee for helpful comments.




\begin{thebibliography}{0}

\bibitem{balu-gun}
R. Balasubramanian, S. Gun, \textit{On zeros of quasi-modular forms}, Journal of Number Theory, Vol. 132 (2012) 2228-2241.

\bibitem{b-d-g-p}
K. Barr\'e-Sirieix, G. Diaz, F. Gramain, G. Philibert, \textit{Une preuve de la conjecture de Mahler-Manin}, Invent. Math., 124 (1996), 1-9.


\bibitem{chang}
C. Y. Chang, \textit{Transcendence of special values of quasi-modular forms}, Forum Math. 24 (2012) 539-551.


\bibitem{stromberg}
H. Cohen and F. Stromberg, \textit{Modular Forms, A Classical Approach}, Graduate Studies In Mathematics 179. 

\bibitem{jeon-kang-kim}
D. Jeon, S.-Y. Kang, C.H. Kim, \textit{On values of weakly holomorphic modular functions at divisors of meromorphic modular forms}, J. Number Theory 239, 183-206.

\bibitem{j-k-k-2}
D. Jeon, S.-Y. Kang, C.H. Kim, \textit{Hecke system of harmonic Maass functions and its applications to modular curves of higher genera}, Ramanujan J., Vol. 62, (2023) 675-717.  


\bibitem{hamieh-murty}
Alia Hamieh and M. Ram Murty, {A note on $q$-analogues of Dirichlet L-functions}, International Journal of Number Theory, Vol. 12 (2016) 765-773.

\bibitem{lang}
Serge Lang, \textit{Elliptic Functions}, Addision-Wesley Publishing Company, Inc.

\bibitem{mahler}
Kurt Mahler, \textit{Remarks on a paper by W. Schwarz.}, Journal of Number Theory, 1 (1969) 512–521.

\bibitem{murty-dewar-graves}
M. R. Murty, M. Dewar and H. Graves, \textit{Problems in 
the Theory of 
Modular
Forms}, Hindustan Book Agency.

\bibitem{r-m}
M. Ram Murty and P. Rath, \textit{Transcendental Numbers} Springer-Verlag, (2014).

\bibitem{r-m-g}
S. Gun, M. Ram Murty, P. Rath, \textit{Algebraic independence of values of
modular forms}. Int. J. Number Theory 7(4), 1065–1074 (2011).


\bibitem{nesterenko}
Y. Nesterenko, \textit{Modular functions and transcendence questions}, Mat. Sb. 187 (9)  (1996) 65–96 (in Russian); translation in: Sb. Math. 187 (9) (1996) 1319–1348.

\bibitem{royer}
E. Royer, \textit{Quasimodular forms: an introduction}, Annales Math\'ematiques Blaise Pascal, Vol. 19 no. 2 (2012) 297-306.

\bibitem{schneider}
T. Schneider, \textit{Arithmetische Untersuchungen elliptischer Integrale}, Math. Ann. 113 (1937) 1–13.

\bibitem{schneider.}
T. Schneider, \textit{Ein Satz über ganzwertige Funktionen als Prinzip für Transzendenzbeweise}, Mathematische Annalen, 121 (1949): 131–140. 


\bibitem{shimura}
 G. Shimura, \textit{Introduction to the Arithmetic Theory of Automorphic Functions}, (Princeton University Press, $1971$).

\bibitem{wang-zhang}
W. Wang and H. Zhang, \textit{Meromorphic quasi-modular forms and their $L$-functions}, Journal of Number Theory, Vol. 241 (2022) 465-503.

% \bibitem{mero-quasi-mdf}
% W. Wang, H. Zhang, \textit{Meromorphic quasi-modular forms and their $L$-functions}, J. Number Theory ($2022$).

\bibitem{wang}
W. Wang, \textit{Algebraic independence of values of quasi-modular forms}, Journal of Number Theory, Vol. 255 (2024) 85-97.

\bibitem{1-2-3}
D. Zagier, \textit{Elliptic Modular Forms and Their Applications}, Lectures at a Summer School
in Nordfjordeid, Norway.


\end{thebibliography}

\end{document}


\documentclass[11pt, reqno]{amsart}
\usepackage[utf8]{inputenc}
\usepackage{amsmath}
\usepackage{amsthm}
\usepackage{amsfonts}
\usepackage{amssymb}
\usepackage{mathrsfs, mathtools, mathdots}
\usepackage{diagbox}
\DeclarePairedDelimiter \floor {\lfloor} {\rfloor}
\DeclareMathOperator{\Tr}{Tr}

\usepackage[top=35mm,bottom=35mm,left=35mm,right=35mm]{geometry}
\newtheorem{theorem}{Theorem}[section]
\newtheorem{conjecture}{Conjecture}
\newtheorem{corollary}{Corollary}
\newtheorem{problem}{Problem}
\newtheorem*{theorem*}{Theorem}
\newtheorem*{remark*}{Remark}
\newtheorem*{problem*}{Problem}
\newtheorem*{conjecture*}{Conjecture}
\newtheorem{lemma}[theorem]{Lemma}
\newtheorem{prop}[theorem]{Proposition}


\begin{document}

\title[Transcendence]{Modular functions and Transcendence}

\author{[Author]}



\subjclass[2010]{-------}

\keywords{Algebraic independence, Transcendence degree, modular forms}
\thanks{Research ----}
\begin{abstract}

\end{abstract}

\maketitle

\section{\bf Introduction}
\bigskip

In this note, we study the arithmetic nature of values of special functions arising from the modular world. The values of these functions on the upper half-plane $\mathbb{H}$ is a rich source of transcendental numbers. An earlier result regarding transcendence of values of 
 normalized Eisenstein series $E_4$ and $E_6$ was given by D. Bertrand \cite{bertrand}. He proved that for $\tau\in\mathbb{H}$, at least one of the numbers $E_4(\tau), E_6(\tau)$ is transcendental. Later, Chudnovsky \cite{chudnovsky} showed that if $\tau\in\mathbb{H}$, then at least two of the numbers $E_2(\tau), E_4(\tau), E_6(\tau)$ are algebraically independent over $\overline{\mathbb{Q}}$. In $1996$, Nesterenko made a significant advancement in this direction by proving the following theorem \cite{nesterenko}. 
\begin{theorem}[Nesterenko] \label{nest}
If $\tau\in \mathbb{H}$, then at least three of the numbers
$$
e^{2 \pi i \tau},\quad E_2(\tau),\quad E_4(\tau),\quad E_6(\tau)
$$
are algebraically independent over $\overline{\mathbb{Q}}$.
\end{theorem}
 An important consequence of Theorem \ref{nest} is the following result which was conjectured earlier by Mahler and proved in a joint work of Barr\'e-sirieix, Diaz, Gramain and Philibert \cite{b-d-g-p}. 
\begin{theorem}[Barr\'e-sirieix et al]\label{mahl}
For any $\tau\in \mathbb{H}$, at least one of the two numbers $e^{2\pi i\tau}$ and $j(\tau)$ is transcendental.
\end{theorem}
\bigskip

 In 2010, an extensive study on algebraic independencece of values of modular forms was done by M. R. Murty, P. Rath and S. Gun \cite{r-m-g}.  Here we mention some of their results.
 \begin{theorem}\label{algem}
     Let $f$ be a non-zero modular form of weight $k$ for $SL_2(\mathbb{Z})$. If the Fourier coefficients of $f$ are algebraic, then
     \begin{enumerate}
         \item any zero of $f$ is either CM or transcendental.
         \item $f(\tau)$ is algebraically independent with $e^{2\pi i\tau}$ except when $f(\tau)=0$.
     \end{enumerate}
 \end{theorem}
In another paper of A. Hamieh and M. R. Murty \cite{hamieh-murty}, they pointed out that similar results hold for higher level modular forms. 
\begin{theorem}\label{algem''}
  Let $f$ be a non-zero modular form of weight $k$ and level $N$ with algebraic Fourier coefficients.\begin{enumerate}
      \item If $\tau\in\mathbb{
      H
      }$ such that $e^{2\pi i\tau}$ is algebraic, then $f(\tau)$ is transcendental.
      \item If $\tau\in\mathbb{H}$ such that  such that $j(\tau)\in \overline{\mathbb{Q}}$, then there exists a transcendental number $\omega_\tau$ which depends only on $\tau$ and is $\overline{\mathbb{Q}}$-linearly independent with $\pi$ such that $\left(\frac{\pi}{\omega_\tau}\right)^k f(\tau)\in \overline{\mathbb{Q}}$. Therefore, $f(\tau)$ is either zero or transcendental.
  \end{enumerate}  
\end{theorem}
We shall restate Theorem \ref{algem''} and proof for weakly holomorphic modular forms also following the technique from \cite{hamieh-murty}.\\

Investigation on transcendence on special values of quasi-modular forms for $SL_2(\mathbb{Z})$ also has been carried out by C. Y. Chang \cite{quasi}.
 \begin{theorem}\label{changt}
 Let $f$ be a quasi-modular form of non-zero weight for $SL_2(\mathbb{Z})$ with algebraic Fourier coefficients. If $\tau\in\mathbb{H}$ such that $j(\tau)\in\overline{\mathbb{Q}}$ and $f(\tau)\neq 0$, then $f(\tau)$ is transcendental.    
 \end{theorem}
  In the paper \cite{hamieh-murty}, they combined Theorem \ref{algem''} and Theorem \ref{changt} with Theorem \ref{mahl} to deduce the following result.
\begin{theorem}\label{mg}
    Let $f$ be a non-zero modular form of level $N$ and weight $k$ with algebraic Fourier coefficients. If $\tau\in\mathbb{H}$ such that $f(\tau)\neq 0$, then at least two of 
    $$
e^{2\pi i\tau},\quad j(\tau),\quad f(\tau)
    $$
    are transcendental.
\end{theorem}
 
Here, we first study values of modular functions of arbitrary level with algebraic Fourier coefficients. Using these results we determine the transcendence of values of modular and quasi-modular forms for both holomorphic and weakly holomorphic cases. Finally, we investigate the transcendence degree of the field generated by values of these special functions.\\

Let $f$ and $g$ denote non-constant modular form (\emph{holomorphic or weakly holomorphic}) and modular function respectively of arbitrary level with algebraic Fourier coefficients. Then we have the following results.
\begin{theorem}\label{suma}
If $\tau\in\mathbb{H}$ such that $\tau$ is not a pole of $g$, then

(i) $g(\tau)\in\overline{\mathbb{Q}} \Leftrightarrow j(\tau)\in\overline{\mathbb{Q}}$.

(ii) At least one of $e^{2\pi i\tau}$ and $g(\tau)$ is transcendental.
\end{theorem}
\begin{theorem}\label{deg.}
    If $\tau\in\mathbb{H}$ such that $f(\tau)\neq 0$ and $\tau$ is not a pole of $g$, then the transcendence degree of the field $\overline{\mathbb{Q}}\left(e^{2\pi i\tau},\, f(\tau),\, g(\tau) \right)$ over $\overline{\mathbb{Q}}$ is at least two, that is
    $$
trdeg \,\overline{\mathbb{Q}}\left(e^{2\pi i\tau},\, f(\tau),\, g(\tau) \right) \geq 2.
$$  
\end{theorem}
We note that Theorem \ref{deg.} generalizes  Theorem \ref{mg}. Moreover, the above theorem extends to the following. 
\begin{theorem}\label{deg..}
 Let $\widetilde{f}$ be a holomorphic or weakly holomorphic quasi-modular form of positive depth and arbitrary level with algebraic Fourier coefficients. If $\tau\in\mathbb{H}$ such that $f(\tau)\neq 0$, $\widetilde{f}(\tau)\neq 0$ and $\tau$ is not a pole of $g$, then 
    $$
trdeg\,\overline{\mathbb{Q}}\left( e^{2\pi i\tau},\, f(\tau),\, g(\tau),\, \widetilde f(\tau)\right)\geq 3. 
    $$
\end{theorem}
\begin{remark*}
 The proof of Theorem \ref{deg..} mainly uses our knowledge on transcendence of values of modular functions (Theorem \ref{suma}) and Nesterenko's Theorem \ref{nest}. Again, if we choose 
$$
f=E_4, \, g=j,\, \widetilde f= E_6+E_4\, E_2,
$$
the statement of Theorem \ref{deg..} implies Nesterenko's theorem. Therefore, Theorem \ref{deg..} is essentially the Nesterenko's theorem in more general setup.   
\end{remark*}








\bigskip

\section{\bf Notations and Preliminary Results}
\subsection{Modular and quasi-modular forms}\quad\bigskip

 Recall that for each $N\in \mathbb{N}$, $$
\Gamma(N)=\left\{  \left( 
\begin{array}{cc}
  a & b \\
  c & d
\end{array}
\right)\in SL_2(\mathbb{Z}):\left( 
\begin{array}{cc}
  a & b \\
  c & d
\end{array}
\right)\equiv \left( 
\begin{array}{cc}
  1 & 0 \\
  0 & 1
\end{array}
\right) (\operatorname{mod}N) \right\}
$$ is called the \emph{principal congruence subgroup of level N}.
 Note that $\Gamma(1)=SL_2(\mathbb{Z})$. A subgroup $\Gamma$ of $SL_2(\mathbb{Z})$ is said to be congruence subgroup if there exists $N\in \mathbb{N}$ such that $\Gamma(N)\subseteq \Gamma$, and the smallest such $N$ is  called the level of $\Gamma$.\\

A \emph{holomorphic modular form} of integer weight $k(\geq 0)$ with respect to a congruence subgroup $\Gamma$ is a holomorphic function on the upper half-plane $\mathbb{H}$ which satisfies 
\begin{equation} \label{hmod}
f\left(\frac{a\,\tau+b}{c\,\tau+d}\right)=(c\,\tau+d)^kf(\tau)\quad\text{for all }
\gamma=\left(\begin{smallmatrix}
   a  & b \\
   c  & d
\end{smallmatrix} \right)\in \Gamma,\end{equation}
and also is \emph{holomorphic} at all cusps. The function $f$ is said to be a \emph{weakly holomorphic modular form} if it is \emph{meromorphic} at cusps. Moreover, $M_{k,\overline{\mathbb{Q}}}(\Gamma)$ and $M_{k,\overline{\mathbb{Q}}}^w(\Gamma)$ denote the set of all holomorphic and weakly modular forms respectively with algebraic Fourier coefficients. \\

For even integer $k\geq 2$, define the nomalized Eisenstein series of weight $k$ for $S L_2(\mathbb{Z})$ by
$$
E_{k}(\tau)=1-\frac{2 k}{B_{k}} \sum_{n=1}^{\infty} \sigma_{k-1}(n) q^{n}, \text{   where }q=e^{2\pi i\tau} \text{ and } \sigma_s(n)=\sum_{\substack{d|n\\ d>0}}d^s.
$$ 
Here, $B_k$ is called the $k$-th Bernoulli number. One can show that for $k\geq 4$, $E_k\in M_{k,\overline{\mathbb{Q}}}(S L_2(\mathbb{Z}))$ and $\Delta(\tau):=\frac{E_4(\tau)^3-E_6(\tau)^2}{1728}\in M_{12,\overline{\mathbb{Q}}}(S L_2(\mathbb{Z}))$. But $E_2(\tau)$ is not a modular form (see \cite{murty-dewar-graves}, Chapter $5$) as it fails to satisfy \eqref{hmod} for the matrix $\left(\begin{smallmatrix}
   0  & -1 \\
   1  & 0
\end{smallmatrix} \right)\in SL_2(\mathbb{Z})$.\\

There are several equivalent  definitions of \emph{quasi-modular} forms, but we use the following characterization.
\begin{theorem}\label{quasistructure}
Every quasi-modular form on $\Gamma$ is a polynomial in $E_2$ with modular coefficients. More precisely, if $f$ is a quasi-modular form of weight $k$ and depth $p$ on $\Gamma$, then $f$ can be uniquely written as $f=\sum_{r=0}^p g_r E_2^r$, where $g_r\in M_{k-2 r}(\Gamma)$ for all $0\leq r\leq p$ and $g_p\neq 0$. 
\end{theorem}

A quasi-modular form $f$ is said to be \emph{holomorphic or weakly holomorphic} depending on whether all modular coefficients of $f$ are holomorphic or weakly holomorphic modular forms.
\bigskip

\subsection{Weierstrass $\wp$-function and zeta-function}\quad\bigskip

The Weierstrass $\wp$-function associated with a two dimensional
lattice $L=\omega_1 \mathbb{Z} \,  \oplus \, \omega_2 \mathbb{Z}$ is defined as follows
\begin{equation}\label{wpfunction}
\wp_L (z)=\frac{1}{z^2}+\sum_{\substack{\omega\in L\\
\omega\neq 0}}\left(\frac{1}{(z-\omega)^2}-\frac{1}{\omega^2} \right), \quad \text{for }z\in \mathbb{C}\setminus L. 
\end{equation}
This is a meromorphic function on $\mathbb{C}$ with double poles at each lattice points. Moreover, it satisfies the following differential equation
\begin{equation} \label{dell}
 \wp'(z)^2=4\wp(z)^3-g_2(L)\wp(z)-g_3(L),   
\end{equation}
where $$g_2(L)=60\, \sum_{\substack{\omega \in L\\ \omega\neq 0}} \frac{1}{\omega^{4}} \text{ and } g_3(L)=140\, \sum_{\substack{\omega \in L\\ \omega\neq 0}} \frac{1}{\omega^{6}}.$$
In particular, if we take $L_\tau=\tau\,\mathbb{Z}\oplus \mathbb{Z}$ for $\tau\in\mathbb{H}$, then $$g_2(L_\tau)=60\,G_4(\tau)=\frac{4\pi^4}{3}E_4(\tau) \text{ and } g_3(L_\tau)=140\,G_6(\tau)=\frac{8\,\pi^6}{27}E_6(\tau).$$

Define the discriminant function by$$\Delta_0(L):=g_2(L)^3-27g_3(L)^2\neq 0, \text{ for any lattice }L.$$
 In particular, we have $\Delta_0(L_\tau)=(2\pi)^{12}\Delta(\tau)$ for all $\tau\in\mathbb{H}$.\\

 The Weierstrass zeta-function associated with lattice $L=\omega_1\,\mathbb{Z}\oplus \omega_2\,\mathbb{Z}$ is defined by
$$
Z_{L}(z)=\frac{1}{z}+\sum_{\substack{\omega\in L\\ \omega\neq 0}}\left(\frac{1}{z-\omega}+\frac{1}{\omega}+\frac{z}{\omega^2} \right)\quad \text{for } z\in \mathbb{C}\setminus L.
$$
 Note that $Z_{L}'(z)=-\wp_L(z)$ is a periodic function with each lattice point of $L$ as a period. Hence, the functions $$Z_{L}(z+\omega_1)-Z_{L}(z) \text{  and  } Z_{L}(z+\omega_2)-Z_{L}(z)$$ are constants. These constants are denoted by $\eta_1(L), \eta_2(L)$ respectively and called \emph{quasi-periods}. Moreover, it is known that \begin{equation}\label{quasi-period}
\eta_2(L_\tau)=G_2(\tau)=\frac{\pi^2}{3}\, E_2(\tau) \text{ 
  for all }\tau\in\mathbb{H}. 
\end{equation}
The reader is referred to (\cite{lang}, chapter $18.3$) for proof of these results.\\

The Weierstrass $\wp$-function satisfies the following addition formula (\cite{r-m}, chapter $11$).
 \begin{theorem} \label{addf}
 For $z_1, z_2\in \mathbb{C}$ such that $z_1\pm z_2\notin L$ we have
$$
\wp\left(z_1+z_2\right)=-\wp\left(z_1\right)-\wp\left(z_2\right)+\frac{1}{4}\left(\frac{\wp^{\prime}\left(z_1\right)-\wp^{\prime}\left(z_2\right)}{\wp\left(z_1\right)-\wp\left(z_2\right)}\right)^2 .
$$
When $z_1=z_2=z$, we have
$$
\wp(2 z)=-2 \wp(z)+\frac{1}{4}\left(\frac{\wp^{\prime \prime}(z)}{\wp^{\prime}(z)}\right)^2 .
$$
 \end{theorem}
By using this addition formula together with a remarkable theorem known as Schneider–Lang theorem, one can prove the following result (see \cite{r-m}, Theorem $11.2$).
\begin{theorem}\label{tranw}
    Let $L=\omega_1\,\mathbb{Z}\oplus\omega_2\,\mathbb{Z}$ be a lattice and both $g_2(L)$, $g_3(L)$ are algebraic. If $\alpha$ is an algebraic number such that $\alpha\notin L$, then $\wp(\alpha)$ is transcendental.
\end{theorem}
We have the following corollary of Theorem \ref{tranw}.
\begin{corollary}\label{tranp}
    Let the Weierstrass $\wp$ function be associated with $L=\omega_1\,\mathbb{Z}\oplus\omega_2\,\mathbb{Z}$ such that both $g_2(L)$ and $g_3(L)$ are algebraic. Then, for any natural number $n$, the numbers $\wp(\frac{\omega_1}{n})$ and $\wp(\frac{\omega_2}{n})$ are algebraic. Moreover, any non-zero period of $\wp(z)$ is necessarily transcendental.
\end{corollary}









































\subsection{Modular function}
\quad\bigskip

A meromorphic function $f$ on $\mathbb{H}$ is said to be a modular function if it satisfies
\begin{equation}
f\left(\frac{a\,\tau+b}{c\,\tau+d}\right)=f(\tau)\quad\text{for all }
\gamma=\left(\begin{smallmatrix}
   a  & b \\
   c  & d
\end{smallmatrix} \right)\in \Gamma\end{equation}
and also is meromorphic at cusps. In particular, a modular function on the congruence subgroup $\Gamma(N)$ is said to a be modular function of level $N$.\\

A canonical example of modular function of level one is given by 
$$
j(\tau):=\frac{E_4^3(\tau)}{\Delta(\tau)} \text{ for all } \tau\in\mathbb{H}.
$$
It has the following Fourier expansion
$$
j(\tau)=\frac{1}{q}+744+196884q+\cdots,$$
where $q=e^{2\pi i\tau}$. This is made precise in the following theorem.
\begin{theorem} \label{level1}
Let $f$ be a modular function of level one and let $$K := \mathbb{Q}(\text{Fourier coefficients of }f)\subset\mathbb{C}.$$ Suppose that the poles of $f$ in $\mathcal{F}$ are $\tau_1,\tau_2,\cdots,\tau_m$, where $\mathcal{F}$ denotes the standard fundamental domain for the action of $S L_2(\mathbb{Z})$ on $\mathbb{H}$. Then $f\in K(j(\tau_1),\cdots,j(\tau_m))(j)$ i.e., $f$ is a rational function in $j$ over $K(j(\tau_1),\cdots,j(\tau_m))$.
\end{theorem}
\begin{proof}
First, we consider the function 
$$
g(\tau):=\left(\prod_{\tau_0\in\mathcal{F}} \left(j(\tau)-j(\tau_0)\right)^{\operatorname{ord}_{\tau_0}f}\right)f(\tau),
$$
which is holomorphic on $\mathbb{H}$. Since $f$ is meromorphic at $i\infty$, there exists $c>0$ such that the region $\{z\in\mathcal{F} : \operatorname{Im}z>c\}\cup \{i\infty\}$ contains no pole of $f$ except possibly $i\infty$. Again, $f$ is meromorphic on $\mathbb{H}$ and the region $\{z\in \mathcal{F}: \operatorname{Im}z\leq c\}$ is compact, and so it contains only a finite number of poles of $f$. Thus, $f$ can have only finitely many poles in $\mathcal{F}$ and hence above product is finite.\\

Suppose $g$ has pole of order $M$ at $i\infty$. Then the Fourier expansion of
$g(\tau)$ at $i\infty$ has the form
$$
g(\tau)=\sum_{n=-M}^\infty c_n q^n, \, \text{ where } c_{-M}\neq 0.
$$
Note that the modular function $g(\tau)-c_{-M}j(\tau)^M$ is holomorphic on $\mathbb{H}$ and its $q$-series starts with at most a polar term of order $M-1$. Repeating this process, we can subtract a polynomial of $j(\tau)$ to get a holomorphic modular function that vanishes at $i\infty$, and hence identically zero. Thus, $g(\tau)$ is a polynomial of $j(\tau)$ over $K$ and so $f(\tau)$ is a rational function of $j(\tau)$ over $K(j(\tau_1),\cdots,j(\tau_m))$.
\end{proof}
The above proof proves the following :
\begin{corollary}\label{rationalfunctions}
If $f(\tau)$ is a modular function of level one which is holomorphic on $\mathbb{H}$ with the following $q$-expansion at $i\infty$
$$
f(\tau)=\sum_{n=-M}^\infty c_n q^n, \ c_{-M}\neq 0
$$
with $K=\mathbb{Q}(\text{Fourier coefficients of }f)$, then $f(\tau)$ is a polynomial in $j(\tau)$ of degree $M$ with coefficients in $K$.
\end{corollary}
\begin{theorem} \label{mdft}
Let $f$ be a modular function with respect to a congruence subgroup $\Gamma$ and
$$
K:=\mathbb{Q}\left(\text{Values of }j\text{ at poles of }f\text{ in }\mathcal{F}, \text{ Fourier coefficients of }f\text{ at all cusps}\right).
$$
Then there exists a monic polynomial $P(X)\in K(j)[X]$ such that $P(f)=0$. 
\end{theorem}
\begin{proof}
We construct a polynomial satisfied by $f$ with coefficients that are rational functions of $j(\tau)$. Theorem \ref{level1} suggests that these coefficients are modular functions of level $1$.\\

Let $[S L_2(\mathbb{Z}):\Gamma]=r$ and $\{ \gamma_1=I,\gamma_2,\cdots,\gamma_r\}$ be a complete set of right coset representatives so that $$
S L_2(\mathbb{Z})=\bigsqcup_{i=1}^r\Gamma \gamma_i.
$$
For all $1\leq i \leq r$, define the functions $$
f_i(\tau):=f(\gamma_i\tau) \quad \text{for } \tau \in \mathbb{H}.
$$ Clearly, each $f_i$ is independent of the choice of coset representatives as $f$ is $\Gamma$-invariant. Again, for any $\gamma\in S L_2(\mathbb{Z})$, we have $f_i(\gamma \tau)=f(\gamma_i\gamma\tau)=f_j(\tau)$ if $\gamma_i\gamma\in \Gamma\gamma_j$, i.e., the set $\{f_1,f_2,...,f_r\}$ gets permuted under the action of $S L_2(\mathbb{Z}).$ This observation implies that any elementary symmetric polynomial of $f_1,f_2,\cdots,f_r$ is a modular function with respect to $S L_2(\mathbb{Z}).$\\

Consider the polynomial $$
P(X)=\prod_{i=1}^r(X-f_i),$$ which is satisfied by $f$ as $f=f_1$. Note that the coefficients of $P$ are elementary symmetric polynomials in $f_1,f_2,\cdots,f_r$, and hence, modular functions with respect to $S L_2(\mathbb{Z})$ with Fourier coefficients in $K$. Thus by Theorem \ref{level1}, the coefficients of $P$ are in $K(j)$ and so $P(X)$ is the required polynomial.
\end{proof}
\begin{corollary}\label{useful}
 If the modular function $f$ is holomorphic on $\mathbb{H}$ and $$K := \mathbb{Q}(\text{Fourier coefficients of }f\text{ at all cusps}),$$
 then the monic polynomial in Theorem \ref{mdft} lies in $K[j][X]$.
\end{corollary}
\begin{proof}
   This is clear from above proof and Corollary \ref{rationalfunctions}, since coefficients of $P$ constructed above are modular functions of level one, holomorphic on $\mathbb{H}$ with Fourier coefficients in $K$.  
\end{proof}
\bigskip

For $N>1$ and $a=(a_1,a_2)\in N^{-1}\mathbb{Z}^2\setminus\mathbb{Z}^2$, consider the function
$$
f_{a}(\tau):=\frac{g_2(L_\tau)\,g_3(L_\tau)}{\Delta_0(L_\tau)}\,\wp_{L_\tau} (a_1\tau+a_2),
$$
which is holomorphic on $\mathbb{H}$. Moreover, It is known that $f_a$ is a modular function of level $N$ with algebraic Fourier coefficients with respect to $e^{\frac{2\pi i\tau}{N}}$. For an elaborate discussion on these functions, the reader is referred to (\cite{shimura}, Chapter 6).\\

Let $F_{N,\overline{\mathbb{Q}}}$ denote the field of modular functions of level $N$ with algebraic Fourier coefficients with respect to $e^{\frac{2\pi i\tau}{N}}$, and $S:=\{\left(\frac{r}{N},\frac{s}{N}\right): 0\leq r,s\leq N-1 \text{ and }(r,s)\neq(0,0)\}$. Thus, $f_a\in F_{N,\overline{\mathbb{Q}}}$ for all $a\in S$.
\begin{theorem}\label{mdft..}
$F_{N,\overline{\mathbb{Q}}}$ is a finite Galois extension of $\overline{\mathbb{Q}}(j)$ and $F_{N,\overline{\mathbb{Q}}}=\overline{\mathbb{Q}}(j,f_a |a\in S)$.
\end{theorem}
For an proof of Theorem \ref{mdft..}, we refer to (\cite{shimura}, Proposition $6.9.(1)$), where they consider modular functions with Fourier coefficients in $\mathbb{Q}\big(e^{\frac{2\pi i}{N}}\big)$. But, one can use the same argument for the field $\overline{\mathbb{Q}}$ instead of $\mathbb{Q}\big(e^{\frac{2\pi i}{N}}\big)$.
\bigskip

\section{\bf Proof of Main Results}
\bigskip

\subsection{Values of modular functions when $j(\tau)$ is transcendental}
\begin{lemma} \label{mdfz}
    Let $g$ be a non-zero modular function of level $N(\geq 1)$ with algebraic Fourier coefficients, i.e., $g\in  F_{N,\overline{\mathbb{Q}}}$. If $\tau\in \mathbb{H}$ such that $j(\tau)\notin\overline{\mathbb{Q}}$ and $\tau$ is not a pole of $g$, then $g(\tau)\neq 0$.
\end{lemma}
\begin{proof}
By Theorem \ref{mdft..}, we have an irreducible polynomial 
$$
P(X)=\sum_{r=0}^m c_r(j)\,X^r \in \overline{\mathbb{Q}}[j][X], \quad m=\deg P(X)
$$
such that
$$
P(g)=0.
$$
Here $$c_r(j)=\sum_{s=0}^{d_r}c_{r s}\,j^s\in \overline{\mathbb{Q}}[j]\quad \forall\ 0\leq r\leq m, \quad d_r=\deg c_r(j).$$
Also, $c_0(j)$ is not identically zero polynomial because $P(X)$ is irreducible over $\overline{\mathbb{Q}}[j]$. Thus, we get
\begin{equation}\label{funzero}
\sum_{r=0}^m\sum_{s=0}^{d_r}c_{r s}\,j(\tau)^s\, g(\tau)^r=0\quad \forall \tau\in\mathbb{H}.
\end{equation}
Now we consider $\tau\in\mathbb{H}$ such that $j(\tau)\notin\overline{\mathbb{Q}}$. We claim that $g(\tau)\neq 0$. If not, from \eqref{funzero} we get $c_0(j(\tau))=0$ and hence, $j(\tau)\in\overline{\mathbb{Q}}$, which contradicts our assumption. This completes the proof that $g(\tau)$ is non-zero. 
\end{proof}
The following theorem shows that the non-zero number $g(\tau)$ is indeed transcendental.
\begin{theorem}\label{tranmf}
 Let $g\in F_{N,\overline{\mathbb{Q}}}$ and $\tau\in \mathbb{H}$ such that it is not a pole of $g$. If $j(\tau)\notin\overline{\mathbb{Q}}$, then $g(\tau)$ is transcendental.
\end{theorem}
\begin{proof}
If possible, let $g(\tau)\in\overline{\mathbb{Q}}$. Then the following modular function $$h(z):=g(z)-g(\tau)\quad \forall z\in\mathbb{H}$$ belongs to $ F_{N,\overline{\mathbb{Q}}}$ but vanishes at $\tau$, which contradicts Lemma \ref{mdfz}. Therefore, $g(\tau)$ is a transcendental number.
\end{proof}
\begin{corollary}
If $g\in F_{N,\overline{\mathbb{Q}}}$ and $\tau_0$ is a zero or a pole of $g$, then $j(\tau_0)\in \overline{\mathbb{Q}}$. 
\end{corollary}
\begin{proof}
If $g(\tau_0)=0$, then Lemma \ref{mdfz} implies that $j(\tau_0)\in \overline{\mathbb{Q}}$. If $\tau_0$ is a pole of $g$, we consider the modular function
$$
h(\tau):=\frac{1}{g(\tau)}\in F_{N,\overline{\mathbb{Q}}},
$$
which satisfies $h(\tau_0)=0$, and hence $j(\tau_0)\in\overline{\mathbb{Q}}$ by Lemma \ref{mdfz}. 
\end{proof}
\subsection{Values of modular functions when $j(\tau)\in\overline{\mathbb{Q}}$}
\quad \bigskip

Here we study values of modular functions at the points on $\mathbb{H}$ where $j$-function takes algebraic values. From Theorem \ref{mdft..}, we have $$ F_{N,\overline{\mathbb{Q}}}=  \overline{\mathbb{Q}}\left(\{j,f_a |a\in S\}\right), \text{ where }S=\left\{\left(\frac{r}{N},\frac{s}{N}\right): 0\leq r, s\leq N-1 \text{ and }(r,s)\neq (0,0)\right\}.$$
Hence, for any $g\in F_{N,\overline{\mathbb{Q}}}$, the algebraic nature of $g(\tau)$ is completely determined by the numbers in the set $\{f_a(\tau):a\in S\}$.\\

Note that for $N=1$, the set $S$ is empty. This shows that for any $g\in F_{1,\overline{\mathbb{Q}}}$, the number $g(\tau)\in\overline{\mathbb{Q}}$ as $j(\tau)\in\overline{\mathbb{Q}}$. For $N>1$, we have to study the numbers $f_a(\tau)$ for all $a\in S$. We start by proving the following important lemma.\bigskip

\begin{lemma}\label{omegatau}
 If $\tau\in \mathbb{H}$ such that $j(\tau)$ is algebraic, then there exists a unique (up to an algebraic multiple) transcendental number $\omega_\tau$ for which $g_2(\omega_\tau L_\tau)$ and $g_3(\omega_\tau L_\tau)$ are both algebraic numbers.   
\end{lemma}
\begin{proof}
Note that for any $\omega_\tau\in\mathbb{C}^\times $,
\begin{equation}\label{omegatau.}
    j(\tau)=\frac{E_4(\tau)^3}{\Delta(\tau)}=\frac{1728\,  g_2(L_\tau)^3}{\Delta_0(L_\tau)}=\frac{1728\, g_2\left(\omega_\tau L_\tau\right)^3}{\Delta_0\left(\omega_\tau L_\tau\right)},
\end{equation}
where $L_\tau=\tau\,\mathbb{Z}\oplus\mathbb{Z}$ and $\omega_\tau L_\tau=\tau\,\omega_\tau\,\mathbb{Z}\oplus\omega_\tau\mathbb{Z}$. Here the last equality follows from the fact that \begin{equation}\label{homogenity}
g_2(\omega_\tau L_\tau)=\omega_\tau^{-4}g_2(L_\tau).
\end{equation} If we choose $\omega_\tau$ such that $\omega_\tau^4=g_2(L_\tau)$, then $g_2(\omega_\tau L_\tau)=1$. Given that $j(\tau)$ is algebraic, and so \eqref{omegatau.} shows that $\Delta_0(\omega_\tau L_\tau)$ is also algebraic. Thus, the numbers $g_2(\omega_\tau L_\tau)$ and $g_3(\omega_\tau L_\tau)$ both are algebraic. Moreover,  $\omega_\tau$ is a period of $\omega_\tau L_\tau$ and $g_2(\omega_\tau L_\tau),\, g_3(\omega_\tau L_\tau)\in\overline{\mathbb{Q}}$. By Corollary \ref{tranp} we conclude that $\omega_\tau$ is transcendental.\\

To prove uniqueness, consider an arbitrary $\omega_\tau'\in\mathbb{C}^\times$ such that both $g_2(\omega_\tau' L_\tau),\, g_3(\omega_\tau' L_\tau)\in\overline{\mathbb{Q}}$. From \eqref{homogenity} we have 
$$
 \quad g_2(\omega_\tau' L_\tau)= (\omega_\tau')^{-4}\,g_2(L_\tau)=\alpha, (\text{say}).
$$
Therefore, we get $g_2(L_\tau)=\alpha\cdot (\omega_\tau')^{4}$ and hence $\omega_\tau^4=\alpha\cdot (\omega_\tau')^{4}.$
This completes the proof that $\omega_\tau$ is unique up to an algebraic multiple.
\end{proof}

Recall that
$$
f_a(\tau):=\frac{g_2(L_\tau)\,g_3(L_\tau)}{\Delta_0(L_\tau)}\,\wp_{L_\tau}(a_1\tau+a_2),
$$ 
where $L_\tau=\tau\,\mathbb{Z}\oplus\mathbb{Z}$ and $a=(a_1,a_2)\in S$. We can rewrite it as follows:
$$
f_a(\tau)=\frac{g_2(\omega_\tau L_\tau)\, g_3(\omega_\tau L_\tau)}{\Delta_0(\omega_\tau L_\tau)}\,\wp_{\omega_\tau L_\tau}(a_1\omega_\tau\tau+a_2\omega_\tau).
$$
\begin{lemma}\label{algfa}
    If $\tau\in\mathbb{H}$ such that $j(\tau)\in\overline{\mathbb{Q}}$, then $f_a(\tau)\in\overline{\mathbb{Q}}$ for all $a\in S$.
\end{lemma}
\begin{proof}
    Let $a=(\frac{r}{N},\frac{s}{N})\in S$. Then we have
$$
f_a(\tau)=\frac{g_2(\omega_\tau L_\tau)\, g_3(\omega_\tau L_\tau)}{\Delta_0(\omega_\tau L_\tau)}\,\wp_{\omega_\tau L_\tau}\left(\frac{r\,\omega_\tau\,\tau}{N}+\frac{s\,\omega_\tau}{N}\right),
$$
where $\omega_\tau$ is chosen as in Lemma \ref{omegatau} so that $g_2(\omega_\tau L_\tau),\, g_3(\omega_\tau L_\tau)\in\overline{\mathbb{Q}}$.\\

Recall that both $\wp_{\omega_\tau L_\tau}\left(\frac{\omega_\tau\, \tau}{N} \right)$ and $\wp_{\omega_\tau L_\tau}\left(\frac{\omega_\tau}{N} \right)$  are algebraic by Corollary \ref{tranp}. Moreover, by the addition formula for Weierstrass $\wp$-function in Theorem \ref{addf}, we get that $$\wp_{\omega_\tau L_\tau}\left(\frac{r\,\omega_\tau\, \tau}{N}+\frac{s\,\omega_\tau}{N}\right)\in\overline{\mathbb{Q}}$$ for all $a=\left(\frac{r}{N},\frac{s}{N}\right)\in S$. Thus, the number $f_a(\tau)$ is algebraic for all $a\in S$.
\end{proof}
\begin{theorem}\label{algemf}
Let $g\in F_{N,\overline{\mathbb{{Q}}}}$ and $\tau\in\mathbb{H}$ such that $j(\tau)\in \overline{\mathbb{Q}}$. If $\tau$ is not a pole of $g$, then $g(\tau)$ is algebraic.
\end{theorem}
\begin{proof}
 From Lemma \ref{algfa}, we have   $f_a(\tau)\in\overline{\mathbb{Q}}$ for all $a\in S$. Then we use the following fact $$g\in  F_{N,\overline{\mathbb{Q}}}=  \overline{\mathbb{Q}}\left(\{j,f_a |a\in S\}\right)$$  to conclude that $g(\tau)$ is algebraic. 
\end{proof}
\bigskip

\begin{proof}[\bf Proof of Theorem \ref{suma}]
    \, Combining Theorem \ref{tranmf} with Theorem \ref{algemf} we get (i), and (ii) follows from Theorem \ref{mahl} and (i).
\end{proof}
\bigskip

\subsection{Values of modular forms}
\quad\bigskip

For positive integers $k$ and $N$, we denote $M_{k,N,\overline{\mathbb{Q}}}$ and $M_{k,N,\overline{\mathbb{Q}}}^w$ as the space of holomorphic and weakly holomorphic modular forms respectively with algebraic Fourier coefficients. Here, we restate Theorem \ref{algem''} and give a proof following the same technique from \cite{hamieh-murty}.
\begin{theorem}\label{algem'}
   Let $f\in M_{k,N,\overline{\mathbb{Q}}}\,\cup M_{k,N,\overline{\mathbb{Q}}}^w$. If $\tau\in\mathbb{H}$ such that $e^{2\pi i\tau}$ is algebraic, then $f(\tau)$ is transcendental. 
\end{theorem}
\begin{proof}
    We first construct the function
$$g(\tau):=\frac{f(\tau)^{12}}{\Delta(\tau)^k}\quad \forall \tau\in\mathbb{H},$$
which is holomorphic on $\mathbb{H}$ because $\Delta(\tau)$ is non-vanishing on $\mathbb{H}$.
Moreover, $g\in F_{N,\overline{\mathbb{Q}}}$ as Fourier coefficients of $f$ and $\Delta$ are algebraic.\\

From Theorem \ref{mdft}, we know that the minimal polynomial of $g$ is of the form: $$P(X)=\sum_{r=0}^m c_r(j)X^r\in \overline{\mathbb{Q}}(j)[X],\quad \deg P(X)=m.$$ Moreover, Corollary \ref{useful} shows that $P(X)$ is a monic polynomial and belongs to $\overline{\mathbb{Q}}[j][X]$ i.e., $$
\begin{aligned}
c_r(j)&=\sum_{s=0}^{d_r}c_{r,s}j^s\in \overline{\mathbb{Q}}[j]\quad \forall\ 0 \leq r\leq m-1  \quad \text{with }d_r=\deg c_r(j),\\  
   c_m(j)&=c_{m,0}=1.
\end{aligned}
$$ Here $c_0(j)$ cannot be the zero polynomial of $j$ because $P(X)$ is irreducible over $\overline{\mathbb{Q}}(j)$.\\

Now we have $$P\left(\frac{f(\tau)^{12}}{\Delta(\tau)^k}\right)=0\quad \text{for all } \tau\in\mathbb{H}.$$
For any positive integer $l$, we get
$$
\begin{aligned}
& \Delta(\tau)^l\, P\left(\frac{f(\tau)^{12}}{\Delta(\tau)^k}\right)=0\\
\Rightarrow &\sum_{r=0}^m\sum_{s=0}^{d_r}c_{r,s}\,j(\tau)^s f(\tau)^{12 r}\Delta(\tau)^{l-k r}=0\\
\Rightarrow &\sum_{r=0}^m\sum_{s=0}^{d_r}c_{r,s}\,E_4(\tau)^{3 s}f(\tau)^{12 r}\Delta(\tau)^{l-k r-s}=0.
\end{aligned}
$$
We choose $l> k m +\operatorname{max}\{d_0,d_1,\ldots,d_m\}$ so that the powers of $\Delta(\tau)$ are always positive integers. Then we use $$\Delta(\tau)=\frac{E_4(\tau)^3-E_6(\tau)^2}{1728}$$ to obtain \begin{equation} \label{trane}
\sum_{r=0}^m \sum_{s=0}^{d_r} \sum_{t=0}^{l-k r-s} \frac{(-1)^t c_{r, s}}{1728^{l-k r-s}}\left(\begin{array}{c}
l-k r-s \\
t
\end{array}\right) f(\tau)^{12 r} E_4(\tau)^{3(l-k r-t)} E_6(\tau)^{2 t}=0.   
\end{equation}
Note that the equation \eqref{trane} holds for all $\tau\in\mathbb{H}$. Now we consider $\tau\in \mathbb{H}$ such that $e^{2\pi i\tau}$ is algebraic. Then we have $f(\tau)\neq 0$, as otherwise $$c_0(j(\tau))=P\left(\frac{f(\tau)^{12}}{\Delta(\tau)^k}\right)=0.$$ This implies that $j(\tau)\in\overline{\mathbb{Q}}$, which contradicts Theorem \ref{mahl}. Moreover, we claim that the number $f(\tau)$ is not algebraic. For otherwise, \eqref{trane} becomes a polynomial in two variables over $\overline{\mathbb{Q}}$, satisfied by $E_4(\tau)$  and $E_6(\tau)$. This polynomial is non-zero because the coefficient of the term $E_4(\tau)^{3(l-k m)}\,E_6(\tau)^{0}$ (which appears in \eqref{trane} only for $r=m,\, t=0$) is given by $$\frac{1}{1728^{l-k m}}\cdot c_{m,0}\cdot f(\tau)^{12 m}\neq 0.$$ This shows that if $f(\tau)$ is algebraic, the equation \eqref{trane} gives  an algebraic dependence relation between $E_4(\tau)$ and $E_6(\tau)$, which contradicts Theorem \ref{nest}. Thus, for any $\tau\in \mathbb{H}$ such that $e^{2 \pi i\tau}$ is algebraic, the number $f(\tau)$ is transcendental.
\end{proof}
\begin{corollary}\label{algi}
If $f\in M_{k,N,\overline{\mathbb{Q}}}\,\cup M_{k,N,\overline{\mathbb{Q}}}^w$, then $f$ is algebraically dependent with $E_4$ and $E_6$ over $\overline{\mathbb{Q}}$.
\end{corollary}
\begin{proof}
    It directly follows from the equation \eqref{trane} because the left-hand side is a non-zero polynomial in $f, E_4$ and $E_6$ over $\overline{\mathbb{Q}}.$
\end{proof}
\bigskip

Now we consider the points $\tau\in\mathbb{H}$ such that $j(\tau)\in\overline{\mathbb{Q}}$. From Lemma \ref{omegatau}, we get a transcendental number $\omega_\tau$ such that $g_2(\omega_\tau L_\tau)$ and $g_3(\omega_\tau L_\tau)$ are both algebraic. Moreover, we have the formulae
$$
\begin{aligned}
E_4(\tau)&=\frac{3}{4\pi^4}\, g_2(L_\tau)=\frac{3}{4}\left(\frac{\omega_\tau}{\pi} \right)^4 g_2(\omega_\tau L_\tau),\\ E_6(\tau)&=\frac{27}{8\pi^6}\,g_3(L_\tau)=\frac{27}{8}\left(\frac{\omega_\tau}{\pi}\right)^6g_3(\omega_\tau L_\tau).
\end{aligned}
$$
From Nesterenko's Theorem \ref{nest}, we know that at most one of $E_4(\tau)$ and $E_6(\tau)$ is algebraic. The above formulae imply that if $\frac{\omega_\tau}{\pi}\in\overline{\mathbb{Q}}$, then the numbers $E_4(\tau),\, E_6(\tau)\in\overline{\mathbb{Q}}$, which is a contradiction. Hence, $\frac{\omega_\tau}{\pi}$ is a transcendental number. Besides, we have 
$$
\begin{aligned}
\Delta(\tau)=\frac{E_4(\tau)^3-E_6(\tau)^2}{1728}&=\frac{1}{4^6}\left(\frac{\omega_\tau}{\pi}\right)^{12}\left(g_2(\omega_\tau L_\tau)^3-27 g_3(\omega_\tau L_\tau)^2 \right)\\
&=\frac{1}{4^6}\left(\frac{\omega_\tau}{\pi}\right)^{12}\Delta_0(\omega_\tau L_\tau).  
\end{aligned}
$$
As $g_2(\omega_\tau L_\tau)$ and $g_3(\omega_\tau L_\tau)$ both are algebraic, the number $\Delta_0(\omega_\tau L_\tau)\in\overline{\mathbb{Q}}\setminus \{0\}$. Since $\frac{\omega_\tau}{\pi}$ is transcendental and $\Delta(\tau)$ is a non-zero algebraic multiple of $\left(\frac{\omega_\tau}{\pi}\right)^{12}$, we deduce that $\Delta(\tau)$ is transcendental. We use this observation to prove the following result.
\begin{theorem}\label{algem}
If $f\in M_{k,N,\overline{\mathbb{Q}}}\,\cup M_{k,N,\overline{\mathbb{Q}}}^w$ and $\tau\in\mathbb{H}$ such that $j(\tau)\in \overline{\mathbb{Q}}$, then there exists a transcendental number $\omega_\tau$ (which depends only on $\tau$ and is $\overline{\mathbb{Q}}$-linearly independent with $\pi$) such that $\left(\frac{\pi}{\omega_\tau}\right)^k f(\tau)\in \overline{\mathbb{Q}}$. Therefore, $f(\tau)$ is either zero or transcendental.
\end{theorem}
\begin{proof}
     We first construct the function $$g(\tau):=\frac{f(\tau)^{12}}{\Delta(\tau)^k}\quad \forall \tau\in\mathbb{H},$$ which lies in $ F_{N,\overline{\mathbb{Q}}}$ by the same argument as in Theorem \ref{algem'}. Since $j(\tau)\in\overline{\mathbb{Q}}$, Theorem \ref{algemf} implies that $g(\tau)$ is an algebraic number $c$, (say). Thus, we get  \begin{equation}\label{forms}
f(\tau)^{12}=c.\Delta(\tau)^k=\frac{c}{4^6}.\Delta_0(\omega_\tau L_\tau)^k\left(\frac{\omega_\tau}{\pi} \right)^{12 k}. 
\end{equation}
This shows that $f(\tau)$ is zero when $c=0$, and otherwise, it a non-zero algebraic multiple of $\left(\frac{\omega_\tau}{\pi} \right)^{k}$ and hence transcendental.
\end{proof}
\bigskip

\subsection{Proof of Theorem \ref{deg.}}\, \bigskip

 By theorem \ref{mdft..}, the numbers $g(\tau)$ and $j(\tau)$ are algebraically dependent over $\overline{\mathbb{Q}}$. Moreover, the number $j(\tau)$ is algebraic if and only if $g(\tau)\in\overline{\mathbb{Q}}$ by Corollary \ref{suma}. These facts together give the first equality below:
$$
\begin{aligned}
 trdeg\,\overline{\mathbb{Q}}\left(e^{2\pi i\tau},\, f(\tau),\, g(\tau) \right)&= trdeg \,\overline{\mathbb{Q}}\left(e^{2\pi i\tau},\, f(\tau),\, j(\tau) \right)\\
 &= trdeg \,\overline{\mathbb{Q}}\left(e^{2\pi i\tau},\, \frac{f(\tau)^{12}}{\Delta(\tau)^k},\, f(\tau),\, j(\tau) \right)\\
 &= trdeg \,\overline{\mathbb{Q}}\left(e^{2\pi i\tau},\, \Delta(\tau)^k,\, f(\tau),\, j(\tau) \right)\\
 &= trdeg \,\overline{\mathbb{Q}}\left(e^{2\pi i\tau},\, \Delta(\tau),\, f(\tau),\, j(\tau) \right) \\
 &= trdeg \,\overline{\mathbb{Q}}\left(e^{2\pi i\tau},\, E_4(\tau)^3,\, E_6(\tau)^2,\, f(\tau),\, j(\tau) \right)\\
 &=trdeg \,\overline{\mathbb{Q}}\left(e^{2\pi i\tau},\, E_4(\tau),\, E_6(\tau),\, f(\tau),\, j(\tau) \right)\\
 &= trdeg \,\overline{\mathbb{Q}}\left(e^{2\pi i\tau},\, E_4(\tau),\, E_6(\tau),\, f(\tau) \right)\\
 &= trdeg \,\overline{\mathbb{Q}}\left(e^{2\pi i\tau},\, E_4(\tau),\, E_6(\tau)\right).
\end{aligned}
$$ 
Here, the second equality follows from the fact that $f^{12}/\Delta^k\in F_{N,\overline{\mathbb{Q}}}$ and  $f(\tau)^{12}/\Delta(\tau)^k$ is algebraically dependent with $j(\tau)$ by Theorem \ref{mdft..}. The third and the seventh equality hold as the fields are identical.  
The final equality holds because $f(\tau)$ is algebraic over $\overline{\mathbb{Q}}\left(E_4(\tau), E_6(\tau)\right)$ by Corollary \ref{algi}. Then we apply Nesterenko's Theorem \ref{nest} on the right side of the final equality to conclude that the transcendence degree of the field $\overline{\mathbb{Q}}\left(e^{2\pi i\tau},\, f(\tau),\, g(\tau) \right)$ over $\overline{\mathbb{Q}}$ is at least two. 
\bigskip

\subsection{Values of quasi-modular forms}\,\bigskip

Let $\widetilde{M}_{k,\overline{\mathbb{Q}}}^{(p)}(\Gamma)$ denote the set of all \emph{holomorphic and weakly holomorphic} quasi-modular forms of weight $k$ and depth $p(>0)$ for $\Gamma$ with algebraic Fourier coefficients. We study their values at the points $\tau\in\mathbb{H}$ such that one of $e^{2\pi i\tau}$ and $j(\tau)$ is algebraic.
\begin{theorem}
Let $\Gamma$ be a congruence subgroup of level $N(\geq 1)$ and $\widetilde{f}\in \widetilde{M}_{k,\overline{\mathbb{Q}}}^{(p)}(\Gamma)$. If $\tau\in\mathbb{H}$ such that $e^{2\pi i\tau}$ is algebraic, then $f(\tau)$ is transcendental. 
\end{theorem}
\begin{proof}
   By Theorem \ref{quasistructure}, we can write $\widetilde{f}$ in the form $$\widetilde{f}=\sum_{r=0}^p g_r E_2^r,\quad g_r\in M_{k,N,\overline{\mathbb{Q}}}\,\cup M_{k,N,\overline{\mathbb{Q}}}^w.$$ From Corollary \ref{algi}, we know that each number $g_r(\tau)$ is algebraically dependent with $E_4(\tau), E_6(\tau)$. Therefore, we get
 $$ trdeg\,\overline{\mathbb{Q}}(g_0(\tau), g_2(\tau),\ldots, g_p(\tau), E_4(\tau), E_6(\tau))=trdeg\,\overline{\mathbb{Q}}( E_4(\tau), E_6(\tau)).$$
If possible, let both $e^{2\pi i\tau}, \widetilde f(\tau)\in \overline{\mathbb{Q}}$. Since 
$\widetilde f(\tau)=\sum_{r=0}^p g_r(\tau) E_2(\tau)^r$ and $g_r(\tau)\neq 0$ (by Theorem \ref{algem'}), the number $E_2(\tau)$ is algebraic over $\overline{\mathbb{Q}}\left(g_0(\tau), g_1(\tau),\ldots, g_p(\tau), E_4(\tau), E_6(\tau)\right)$ and so
$$
\begin{aligned}
&\quad trdeg\, \overline{\mathbb{Q}}\left(E_2(\tau), g_0(\tau), g_1(\tau),\ldots, g_p(\tau), E_4(\tau), E_6(\tau)\right)\\ &=trdeg\,\overline{\mathbb{Q}}(g_0(\tau), g_1(\tau),\ldots, g_p(\tau), E_4(\tau), E_6(\tau))\\ 
&=trdeg\,\overline{\mathbb{Q}}(E_4(\tau), E_6(\tau)).
\end{aligned}
$$
Since $e^{2\pi i\tau}$ is an algebraic number, the following field containment holds
$$
\overline{\mathbb{Q}}(e^{2\pi i\tau},\, E_2(\tau),\,E_4(\tau), E_6(\tau))\subseteq \overline{\mathbb{Q}}\left(E_2(\tau), g_0(\tau), g_1(\tau),\ldots, g_p(\tau), E_4(\tau), E_6(\tau)\right).
$$
Therefore, we get 
$$
\begin{aligned}
  trdeg\, \overline{\mathbb{Q}}(e^{2\pi i\tau},\, E_2(\tau),\,E_4(\tau), E_6(\tau))& \leq trdeg\,\overline{\mathbb{Q}}\left(E_2(\tau), g_0(\tau), g_1(\tau),\ldots, g_p(\tau), E_4(\tau), E_6(\tau)\right)\\
  &= trdeg\,\overline{\mathbb{Q}}(E_4(\tau), E_6(\tau))\\
  &\leq 2,
\end{aligned}
$$
which contradicts Nesterenko's Theorem \ref{nest}. Hence, $\widetilde{f}(\tau)$ is a transcendental number. 
\end{proof}
To study values of quasi-modular forms in the complementary case (i.e., for $j(\tau)\in\overline{\mathbb{Q}}$), we need the following lemma.
\begin{lemma}\label{peri}
 If $\tau\in \mathbb{H}$ such that $j(\tau)$ is algebraic, then there exists a unique (up to algebraic multiple) transcendental number $\omega_\tau$ for which $\frac{\omega_\tau}{\pi}$ and $\frac{\eta_2}{\pi}$ are algebraically independent over $\overline{\mathbb{Q}}$. Here, $\eta_2$ is the second quasi-period with respect to the lattice $\omega_\tau L_\tau$ 
\end{lemma}
\begin{proof}
    We have already proved existence of $\omega_\tau$ in Lemma \ref{omegatau}.  From the definition of the Weierstrass zeta-function, one gets
 $   
Z_{\omega_\tau L_\tau}(z)= \frac{1}{\omega_\tau}\cdot Z_{L_\tau}\left(\frac{z}{\omega_\tau}\right).
$ Then we use the definition of quasi-period and the identity \eqref{quasi-period} obtain
$$
\eta_2=Z_{\omega_\tau L_\tau}(z+\omega_\tau)-Z_{\omega_\tau L_\tau}(z)= \frac{1}{\omega_\tau}\,\eta_2(L_\tau)=\frac{1}{\omega_\tau}\,G_2(\tau)= \frac{1}{3}\cdot\frac{\pi^2}{ \omega_\tau}\,E_2(\tau).
$$
Thus, we get the following formulae 
$$
E_2(\tau)=3\,\frac{\omega_\tau}{\pi}\frac{\eta_2}{\pi}, \quad E_4(\tau)=\frac{3}{4}\left(\frac{\omega_\tau}{\pi} \right)^4 g_2(\omega_\tau L_\tau),\quad E_6(\tau)=\frac{27}{8}\left(\frac{\omega_\tau}{\pi}\right)^6g_3(\omega_\tau L_\tau).$$
The above formulae imply that $E_2(\tau)$, $E_4(\tau)$ and $E_6(\tau)$ are algebraic over $\overline{\mathbb{Q}}\left(\frac{\omega_\tau}{\pi}, \frac{\eta_2}{\pi}\right)$. If possible, let $\frac{\omega_\tau}{\pi}, \frac{\eta_2}{\pi}$ be algebraically dependent. Then we get that $trdeg\,\overline{\mathbb{Q}}(E_2(\tau),E_4(\tau),E_6(\tau))=1$, which contradicts Nesterenko's Theorem \ref{nest}. Hence, the numbers $\frac{\omega_\tau}{\pi}$ and $\frac{\eta_2}{\pi}$ are algebraically independent over $\overline{\mathbb{Q}}$.
\end{proof}
Lemma \ref{peri} and Theorem \ref{quasistructure} together allow us to describe values of quasi-modular forms explicitly at the points $\tau\in\mathbb{H}$ where $j$-function is algebraic.
\begin{theorem}\label{tranq}
 Let $\widetilde{f}\in \widetilde{M}_{k,\overline{\mathbb{Q}}}^{(p)}(\Gamma)$ and $\tau\in\mathbb{H}$ such that $j(\tau)\in\overline{\mathbb{Q}}$, then $f(\tau)$ is a polynomial in $\frac{\omega_\tau}{\pi}$ and $\frac{\eta_2}{\pi}$ with algebraic coefficients. Therefore, $f(\tau)$ is either zero or transcendental.   
\end{theorem}
\begin{proof}
    From Theorem \ref{quasistructure}, we have the expression $$\widetilde{f}(\tau)=\sum_{r=0}^p g_r(\tau) E_2(\tau)^r,\quad \text{where } g_r\in M_{k,N,\overline{\mathbb{Q}}}\,\cup M_{k,N,\overline{\mathbb{Q}}}^w.$$
Then we apply Theorem \ref{algem} for each coefficient $g_r$, and the above formula for $E_2$ such that 
\begin{equation}\label{last}
\widetilde{f}(\tau)=\sum_{r=0}^p c_r\left(\frac{\omega_\tau}{\pi}\right)^{k-2 r}\left(\frac{\omega_\tau}{\pi}\cdot\frac{\eta_2}{\pi}\right)^r=\sum_{r=0}^p\, c_r\, \left(\frac{\omega_\tau}{\pi}\right)^{k-r}\, \left(\frac{\eta_2}{\pi}\right)^r,
 \end{equation}
where each $c_r$ is an algebraic constant. If $\widetilde{f}(\tau)\in\overline{\mathbb{Q}}\setminus \{0\}$, then the above expression of $f(\tau)$ shows that $\frac{\omega_\tau}{\pi}$ and $\frac{\eta_2}{\pi}$ are algebraically dependent, which contradicts Lemma \ref{peri}. Therefore, the number $\widetilde{f}(\tau)$ is either zero or transcendental.     
\end{proof}
\bigskip

\subsection{Proof of Theorem \ref{deg..}}
\,\bigskip

By the same argument as in Theorem \ref{deg.}, we get
\begin{equation}\label{final.}
trdeg\,\overline{\mathbb{Q}}\left( e^{2\pi i\tau},\, f(\tau),\, g(\tau),\, \widetilde f(\tau)\right)= trdeg\,\overline{\mathbb{Q}}\left( e^{2\pi i\tau},\, E_4(\tau),\, E_6(\tau),\, \widetilde f(\tau)\right).
\end{equation}
Moreover, we can apply Theorem \ref{quasistructure} to express
$$\widetilde{f}(\tau)=\sum_{r=0}^p g_r(\tau) E_2(\tau)^r,\quad \text{where } g_r\in M_{k,N,\overline{\mathbb{Q}}}\,\cup M_{k,N,\overline{\mathbb{Q}}}^w.$$
Thus, we obtain 
  $$
\begin{aligned}
&trdeg\,\overline{\mathbb{Q}}\left( e^{2\pi i\tau},\, E_4(\tau),\, E_6(\tau),\, \widetilde f(\tau)\right)\\
  & =trdeg\,\overline{\mathbb{Q}}\left( e^{2\pi i\tau},\, E_4(\tau),\, E_6(\tau),\,\sum_{r=0}^p g_r(\tau)\, E_2(\tau)^r \right)\\
&=trdeg\,\overline{\mathbb{Q}}\left( e^{2\pi i\tau}, E_4(\tau), E_6(\tau), g_0(\tau),\ldots,g_p(\tau),\sum_{r=0}^p g_r(\tau)\, E_2(\tau)^r \right)\\
&=trdeg\,\overline{\mathbb{Q}}\left( e^{2\pi i\tau},\, E_4(\tau),\, E_6(\tau),\, g_0(\tau),\,\ldots,\,g_p(\tau),\,E_2(\tau) \right)\\
&=trdeg\,\overline{\mathbb{Q}}\left( e^{2\pi i\tau},\, E_4(\tau),\, E_6(\tau),\,E_2(\tau) \right),
\end{aligned}
 $$
where the second and fourth equalities follow from the fact that each coefficient $g_r(\tau)$ is algebraic over $\overline{\mathbb{Q}}\left(E_4(\tau), E_6(\tau)\right)$ by Corollary \ref{algi}. Moreover, Nesterenko's theorem \ref{nest} implies that the right side of the final equality is at least $3$, and then we use the identity \eqref{final.} to conclude the proof. 












\begin{thebibliography}{0}

\bibitem{ahlfors}
Lars V. Ahlfors, \textit{An Introduction to the Theory of Analytic 
Functions of One Complex Variable}, Third Edition, International series in pure and applied mathematics.


\bibitem{b-d-g-p}
K. Barr\'e-Sirieix, G. Diaz, F. Gramain, G. Philibert, \textit{Une preuve de la conjecture de Mahler-Manin}, Invent. Math., 124 (1996), 1-9.

\bibitem{bertrand}
D. Bertrand, \textit{S\'eries d’Eisenstein et transcendance.} Bull. Soc. Math. France
$104(3)$, $309–321 (1976)$

\bibitem{quasi}
    C. Y. Chang, \textit{Transcendence of special values of quasi-modular forms}, Forum Math. $24$ (2012) $539-551$.


\bibitem{chudnovsky}
    G.V. Chudnovsky, \textit{Algebraic independence of constants connected with the exponential and the elliptic functions}, Dokl.
Akad. Nauk Ukrain. SSR Ser. A $8$ $(1976)$ $698–701$.


\bibitem{stromberg}
H. Cohen and F. Stromberg, \textit{Modular Forms, A Classical Approach}, Graduate Studies In Matheatics $179$.

\bibitem{cohen}
P. Cohen, \textit{On the coefficients of the transformation polynomials for the elliptic modular function}, Math. Proc. Cambridge Phil. Soc. $95$, (1984), 389-402. 

\bibitem{hamieh-murty}
Alia Hamieh and M. Ram Murty, {A note on $q$-analogues of Dirichlet L-functions}, International Journal of Number Theory, Vol. 12 (2016) 765-773.

\bibitem{lang}
Serge Lang, \textit{Elliptic Functions}, Addision-Wesley Publishing Company, Inc.


\bibitem{murty-dewar-graves}
M. R. Murty, M. Dewar and H. Graves, \textit{Problems in 
the Theory of 
Modular
Forms}, Hindustan Book Agency.

\bibitem{r-m}
M. Ram Murty and P. Rath, \textit{Transcendental Numbers} (Springer-Verlag, $2014$).

\bibitem{r-m-g}
S. Gun, M. Ram Murty, P. Rath, \textit{Algebraic independence of values of
modular forms}. Int. J. Number Theory 7(4), 1065–1074 (2011).


\bibitem{nesterenko}
Y. Nesterenko, \textit{Modular functions and transcendence questions}, Mat. Sb. $187$ $(9)$ $(1996)$ $65–96$ (in Russian); translation in:
Sb. Math. $187$ $(9)$ $(1996)$ $1319–1348$.


\bibitem{schneider}
T. Schneider, \textit{Arithmetische Untersuchungen elliptischer Integrale}, Math. Ann.$ 113$ $(1973)$ $1–13$.

\bibitem{shimura}
 G. Shimura, \textit{Introduction to the arithmetic theory of automorphic functions}, Kano Memorial Lectures $1$.


\end{thebibliography}

\end{document}


\documentclass[11pt, reqno]{amsart}
\usepackage[utf8]{inputenc}
\usepackage{amsmath}
\usepackage{amsthm}
\usepackage{amsfonts}
\usepackage{amssymb}
\usepackage{mathrsfs, mathtools, mathdots}
\usepackage{diagbox}
\DeclarePairedDelimiter \floor {\lfloor} {\rfloor}
\DeclareMathOperator{\Tr}{Tr}

\usepackage[top=35mm,bottom=35mm,left=35mm,right=35mm]{geometry}
\newtheorem{theorem}{Theorem}[section]
\newtheorem{conjecture}{Conjecture}
\newtheorem{corollary}{Corollary}
\newtheorem{problem}{Problem}
\newtheorem*{theorem*}{Theorem}
\newtheorem*{remark*}{Remark}
\newtheorem*{problem*}{Problem}
\newtheorem*{conjecture*}{Conjecture}
\newtheorem{lemma}[theorem]{Lemma}
\newtheorem{prop}[theorem]{Proposition}


\begin{document}

\title[Transcendence]{Modular functions and Transcendence}

\author{[Author]}



\subjclass[2010]{-------}

\keywords{Algebraic independence, Transcendence degree, modular forms}
\thanks{Research ----}
\begin{abstract}
We study values of modular forms (possibly  meromorphic and weakly holomorphic) and modular functions of arbitrary level with algebraic Fourier coefficients. Then we investigate the transcendence degree of the field over $\overline{\mathbb{Q}}$ obtained by adjoining $e^{2\pi i\tau}$ and the values of a modular form, modular function and quasi-modular form at a point $\tau$ in the upper half-plane. As a consequence, we obtain a generalization of Nesterenko's Theorem in higher level. 
\end{abstract}

\maketitle

\section{\bf Introduction}
\bigskip

In this note, we investigate the arithmetic nature of values of special functions arising from the modular world. The values of these functions on the upper half-plane $\mathbb{H}$ is a rich source of transcendental numbers. An earlier result regarding transcendence of values of 
 normalized Eisenstein series $E_4$ and $E_6$ was given by D. Bertrand \cite{bertrand}. He proved that for $\tau\in\mathbb{H}$, at least one of the numbers $E_4(\tau), E_6(\tau)$ is transcendental. Later, Chudnovsky \cite{chudnovsky} showed that if $\tau\in\mathbb{H}$, then at least two of the numbers $E_2(\tau), E_4(\tau), E_6(\tau)$ are algebraically independent over $\overline{\mathbb{Q}}$. In $1996$, Nesterenko made a significant advancement in this direction by proving the following theorem \cite{nesterenko}. 
\begin{theorem}[Nesterenko] \label{nest}
If $\tau\in \mathbb{H}$, then at least three of the numbers
$$
e^{2 \pi i \tau},\quad E_2(\tau),\quad E_4(\tau),\quad E_6(\tau)
$$
are algebraically independent over $\overline{\mathbb{Q}}$.
\end{theorem}
 An important consequence of Theorem \ref{nest} is the following result which was conjectured earlier by Mahler and proved in a joint work of Barr\'e-sirieix, Diaz, Gramain and Philibert \cite{b-d-g-p}. 
\begin{theorem}[Barr\'e-sirieix et al]\label{mahl}
For any $\tau\in \mathbb{H}$, at least one of the two numbers $e^{2\pi i\tau}$ and $j(\tau)$ is transcendental.
\end{theorem}
It is known from the theory of complex multiplication that if $\tau\in\mathbb{H}$ is an imaginary quadratic irrational number (i.e. CM point), then $j(\tau)$ is an algebraic number. For non-CM algebraic points in $\mathbb{H}$, Schneider \cite{schneider} proved the following result.
\begin{theorem}[Schneider]\label{schneider}
    If $\tau\in \mathbb{H}$ is algebraic but not a CM point, then $j(\tau)$ is transcendental.
\end{theorem}

 In 2010, a study on algebraic independencece of values of modular forms was done by M. R. Murty, P. Rath and S. Gun \cite{r-m-g}. 
In another paper \cite{hamieh-murty}, Hamieh and Murty proved the following results for modular forms of arbitrary level. 
\begin{theorem}\label{algem''}
  Let $f$ be a non-zero modular form of weight $k$ and level $N$ with algebraic Fourier coefficients.\begin{enumerate}
      \item If $\tau\in\mathbb{
      H
      }$ such that $e^{2\pi i\tau}$ is algebraic, then $f(\tau)$ is transcendental.
      \item If $\tau\in\mathbb{H}$ such that  such that $j(\tau)\in \overline{\mathbb{Q}}$, then there exists a transcendental number $\omega_\tau$ which depends only on $\tau$ and is $\overline{\mathbb{Q}}$-linearly independent with $\pi$ such that $\left(\frac{\pi}{\omega_\tau}\right)^k f(\tau)\in \overline{\mathbb{Q}}$. Therefore, $f(\tau)$ is either zero or transcendental.
  \end{enumerate}  
\end{theorem}
We note that Theorem \ref{algem''} immediately extends for weakly holomorphic modular forms, but it requires a little work for meromorphic modular forms. We will restate Theorem \ref{algem''} and give a  proof for meromorphic modular forms (see Theorem \ref{algem'}).\\

In a recent paper (\cite{jeon-kang-kim}, Theorem $2.4$) Jeon, Kang and Kim studied values of a particular family of modular functions at divisors of meromorphic modular forms. As a consequence (\cite{jeon-kang-kim}, Corollary $2.5$), they proved that any zero or pole of a meromorphic modular with algebraic Fourier coefficients is either transcendental or imaginary quadratic. Indeed, These result follows from a more general correspondence between the transcendental nature of values of the $j$-function and arbitrary modular functions. We give a proof of this correspondence in Theorem \ref{suma}. As immediate consequence, one can state the classical result of Schneider (Theorem \ref{schneider}) and Theorem \ref{mahl} for arbitrary modular functions with algebraic Fourier coefficients.\\

Investigation on transcendence on special values of quasi-modular forms for $SL_2(\mathbb{Z})$ also has been carried out by C. Y. Chang. In \cite{quasi}, Chang proved the following theorem.
 \begin{theorem}\label{changt}
 Let $f$ be a quasi-modular form of non-zero weight for $SL_2(\mathbb{Z})$ with algebraic Fourier coefficients. If $\tau\in\mathbb{H}$ such that $j(\tau)\in\overline{\mathbb{Q}}$ and $f(\tau)\neq 0$, then $f(\tau)$ is transcendental.    
 \end{theorem}
  In (\cite{hamieh-murty}, Corollary $3.3$), Hamieh and Murty combined Theorem \ref{algem''} and Theorem \ref{changt} with Theorem \ref{mahl} to deduce the following result.
\begin{theorem}\label{mg}
    Let $f$ be a non-zero modular form of level $N$ and weight $k$ with algebraic Fourier coefficients. If $\tau\in\mathbb{H}$ such that $f(\tau)\neq 0$, then at least two of 
    $$
e^{2\pi i\tau},\quad j(\tau),\quad f(\tau)
    $$
    are transcendental.
\end{theorem}
In Theorem \ref{deg.}, we prove that at least two of these numbers are algebraically independent over $\overline{\mathbb{Q}}$.\\
 
Here, we first study values of modular functions of arbitrary level with algebraic Fourier coefficients. Using these results we determine the transcendence of values of modular (possibly weakly holomorphic or meromorphic) and quasi-modular forms of arbitrary level. Finally, we investigate the transcendence degree of the field generated by values of these functions.\\

Let $f$ and $g$ denote a non-constant element of $M_{k,N,\overline{\mathbb{Q}}}\,\cup M_{k,N,\overline{\mathbb{Q}}}^w\, \cup M_{k,N,\overline{\mathbb{Q}}}^m$ and $F_{N,\overline{\mathbb{Q}}}$ respectively with algebraic Fourier coefficients. Then we have the following results.
\begin{theorem}\label{suma}
If $\tau\in\mathbb{H}$ such that $\tau$ is not a pole of $g$, then

(i) $g(\tau)\in\overline{\mathbb{Q}} \Leftrightarrow j(\tau)\in\overline{\mathbb{Q}}$.

(ii) At least one of $e^{2\pi i\tau}$ and $g(\tau)$ is transcendental.
\end{theorem}
\begin{theorem}\label{deg.}
    If $\tau\in\mathbb{H}$ such that $f(\tau)\neq 0$ and $\tau$ is not a pole of $f$ and $g$, then the transcendence degree of the field $\overline{\mathbb{Q}}\left(e^{2\pi i\tau},\, f(\tau),\, g(\tau) \right)$ over $\overline{\mathbb{Q}}$ is at least two, that is
    $$
trdeg \,\overline{\mathbb{Q}}\left(e^{2\pi i\tau},\, f(\tau),\, g(\tau) \right) \geq 2.
$$  
\end{theorem}
We note that Theorem \ref{deg.} generalizes  Theorem \ref{mg}. Moreover, the above theorem extends to the following. 
\begin{theorem}\label{deg..}
 Let $\widetilde{f}$ be a quasi-modular form of positive depth and arbitrary level with algebraic Fourier coefficients. If $\tau\in\mathbb{H}$ such that $f(\tau)\neq 0$, $\widetilde{f}(\tau)\neq 0$ and $\tau$ is not a pole of $f$ and $g$, then 
    $$
trdeg\,\overline{\mathbb{Q}}\left( e^{2\pi i\tau},\, f(\tau),\, g(\tau),\, \widetilde f(\tau)\right)\geq 3. 
    $$
\end{theorem}
\begin{remark*}
 The proof of Theorem \ref{deg..} mainly uses the result on transcendence of values of modular functions (Theorem \ref{suma}) and Nesterenko's Theorem \ref{nest}. Again, if we choose 
$$
f=E_4, \, g=j,\, \widetilde f= E_6+E_4\, E_2,
$$
the statement of Theorem \ref{deg..} implies Nesterenko's theorem. Therefore, Theorem \ref{deg..} is essentially the Nesterenko's theorem in more general setup.   
\end{remark*}








\bigskip

\section{\bf Notations and Preliminary Results}
\subsection{Modular and quasi-modular forms}\quad\bigskip

 Recall that for each $N\in \mathbb{N}$, $$
\Gamma(N)=\left\{  \left( 
\begin{array}{cc}
  a & b \\
  c & d
\end{array}
\right)\in SL_2(\mathbb{Z}):\left( 
\begin{array}{cc}
  a & b \\
  c & d
\end{array}
\right)\equiv \left( 
\begin{array}{cc}
  1 & 0 \\
  0 & 1
\end{array}
\right) (\operatorname{mod}N) \right\}
$$ is called the \emph{principal congruence subgroup of level N}.
 Note that $\Gamma(1)=SL_2(\mathbb{Z})$. A subgroup $\Gamma$ of $SL_2(\mathbb{Z})$ is said to be congruence subgroup if there exists $N\in \mathbb{N}$ such that $\Gamma(N)\subseteq \Gamma$, and the smallest such $N$ is  called the level of $\Gamma$.\\

A \emph{holomorphic modular form} of integer weight $k(\geq 0)$ with respect to a congruence subgroup $\Gamma$ is a holomorphic function on the upper half-plane $\mathbb{H}$ which satisfies 
\begin{equation} \label{hmod}
f\left(\frac{a\,\tau+b}{c\,\tau+d}\right)=(c\,\tau+d)^kf(\tau)\quad\text{for all }
\gamma=\left(\begin{smallmatrix}
   a  & b \\
   c  & d
\end{smallmatrix} \right)\in \Gamma,\end{equation}
and also is \emph{holomorphic} at all cusps. The function $f$ is said to be a \emph{weakly holomorphic modular form} if it is \emph{holomorphic} on $\mathbb{H}$ but \emph{meromorphic} at cusps. In general, $f$ is called \emph{meromorphic modular form} if it is meromorphic on $\mathbb{H}$ and also at cusps. Moreover, $M_{k,\overline{\mathbb{Q}}}(\Gamma)$, $M_{k,\overline{\mathbb{Q}}}^w(\Gamma)$ and $M_{k,\overline{\mathbb{Q}}}^m(\Gamma)$ denote the set of all holomorphic, weakly holomorphic and meromorphic modular forms respectively with algebraic Fourier coefficients. \\

For even integer $k\geq 2$, define the nomalized Eisenstein series of weight $k$ for $S L_2(\mathbb{Z})$ by
$$
E_{k}(\tau)=1-\frac{2 k}{B_{k}} \sum_{n=1}^{\infty} \sigma_{k-1}(n) q^{n}, \text{   where }q=e^{2\pi i\tau} \text{ and } \sigma_s(n)=\sum_{\substack{d|n\\ d>0}}d^s.
$$ 
Here, $B_k$ is called the $k$-th Bernoulli number. One can show that for $k\geq 4$, the function $E_k\in M_{k,\overline{\mathbb{Q}}}(S L_2(\mathbb{Z}))$ and $\Delta(\tau):=\frac{E_4(\tau)^3-E_6(\tau)^2}{1728}\in M_{12,\overline{\mathbb{Q}}}(S L_2(\mathbb{Z}))$. But $E_2(\tau)$ is not a modular form (see \cite{murty-dewar-graves}, Chapter $5$) as it fails to satisfy \eqref{hmod} for the matrix $\left(\begin{smallmatrix}
   0  & -1 \\
   1  & 0
\end{smallmatrix} \right)\in SL_2(\mathbb{Z})$.\\

There are several equivalent  definitions of \emph{quasi-modular} forms, but we use the following characterization.
\begin{theorem}\label{quasistructure}
Every quasi-modular form for a congruence subgroup $\Gamma$ is a polynomial in $E_2$ with modular coefficients. More precisely, if $f$ is a quasi-modular form of weight $k$ and depth $p$ on $\Gamma$, then $f$ can be uniquely written as $f=\sum_{r=0}^p g_r E_2^r$, where $g_r\in M_{k-2 r}(\Gamma)$ for all $0\leq r\leq p$ and $g_p\neq 0$. 
\end{theorem}

A quasi-modular form $f$ is said to have algebraic Fourier coefficients if all modular coefficient $g_r$'s of $f$ have algebraic Fourier coefficients. 
\bigskip

\subsection{Weierstrass $\wp$-function and zeta-function}\quad\bigskip

The Weierstrass $\wp$-function associated with a two dimensional
lattice $L=\omega_1 \mathbb{Z} \,  \oplus \, \omega_2 \mathbb{Z}$ is defined as follows
\begin{equation}\label{wpfunction}
\wp_L (z)=\frac{1}{z^2}+\sum_{\substack{\omega\in L\\
\omega\neq 0}}\left(\frac{1}{(z-\omega)^2}-\frac{1}{\omega^2} \right), \quad \text{for }z\in \mathbb{C}\setminus L. 
\end{equation}
This is a meromorphic function on $\mathbb{C}$ with double poles at each lattice points. Moreover, it satisfies the following differential equation
\begin{equation} \label{dell}
 \wp'(z)^2=4\wp(z)^3-g_2(L)\wp(z)-g_3(L),   
\end{equation}
where $$g_2(L)=60\, \sum_{\substack{\omega \in L\\ \omega\neq 0}} \frac{1}{\omega^{4}} \text{ and } g_3(L)=140\, \sum_{\substack{\omega \in L\\ \omega\neq 0}} \frac{1}{\omega^{6}}.$$
In particular, if we take $L_\tau=\tau\,\mathbb{Z}\oplus \mathbb{Z}$ for $\tau\in\mathbb{H}$, then $$g_2(L_\tau)=60\,G_4(\tau)=\frac{4\pi^4}{3}E_4(\tau) \text{ and } g_3(L_\tau)=140\,G_6(\tau)=\frac{8\,\pi^6}{27}E_6(\tau).$$

Define the discriminant function by$$\Delta_0(L):=g_2(L)^3-27g_3(L)^2\neq 0, \text{ for any lattice }L.$$
 In particular, we have $\Delta_0(L_\tau)=(2\pi)^{12}\Delta(\tau)$ for all $\tau\in\mathbb{H}$.\\

 The Weierstrass zeta-function associated with lattice $L=\omega_1\,\mathbb{Z}\oplus \omega_2\,\mathbb{Z}$ is defined by
$$
Z_{L}(z)=\frac{1}{z}+\sum_{\substack{\omega\in L\\ \omega\neq 0}}\left(\frac{1}{z-\omega}+\frac{1}{\omega}+\frac{z}{\omega^2} \right)\quad \text{for } z\in \mathbb{C}\setminus L.
$$
 Note that $Z_{L}'(z)=-\wp_L(z)$ is a periodic function with each lattice point of $L$ as a period. Hence, the functions $$Z_{L}(z+\omega_1)-Z_{L}(z) \text{  and  } Z_{L}(z+\omega_2)-Z_{L}(z)$$ are constants. These constants are denoted by $\eta_1(L), \eta_2(L)$ respectively and called \emph{quasi-periods}. Moreover, it is known that \begin{equation}\label{quasi-period}
\eta_2(L_\tau)=G_2(\tau)=\frac{\pi^2}{3}\, E_2(\tau) \text{ 
  for all }\tau\in\mathbb{H}. 
\end{equation}
The reader is referred to (\cite{lang}, chapter $18.3$) for proof of these results.\\

The Weierstrass $\wp$-function satisfies the following addition formula (\cite{r-m}, chapter $11$).
 \begin{theorem} \label{addf}
 For $z_1, z_2\in \mathbb{C}$ such that $z_1\pm z_2\notin L$ we have
$$
\wp\left(z_1+z_2\right)=-\wp\left(z_1\right)-\wp\left(z_2\right)+\frac{1}{4}\left(\frac{\wp^{\prime}\left(z_1\right)-\wp^{\prime}\left(z_2\right)}{\wp\left(z_1\right)-\wp\left(z_2\right)}\right)^2 .
$$
When $z_1=z_2=z$, we have
$$
\wp(2 z)=-2 \wp(z)+\frac{1}{4}\left(\frac{\wp^{\prime \prime}(z)}{\wp^{\prime}(z)}\right)^2 .
$$
 \end{theorem}
By using this addition formula together with a remarkable theorem known as Schneider–Lang theorem, one can prove the following result (see \cite{r-m}, Theorem $11.2$).
\begin{theorem}\label{tranw}
    Let $L=\omega_1\,\mathbb{Z}\oplus\omega_2\,\mathbb{Z}$ be a lattice and both $g_2(L)$, $g_3(L)$ are algebraic. If $\alpha$ is an algebraic number such that $\alpha\notin L$, then $\wp(\alpha)$ is transcendental.
\end{theorem}
We have the following corollary of Theorem \ref{tranw}.
\begin{corollary}\label{tranp}
    Let the Weierstrass $\wp$ function be associated with $L=\omega_1\,\mathbb{Z}\oplus\omega_2\,\mathbb{Z}$ such that both $g_2(L)$ and $g_3(L)$ are algebraic. Then, for any natural number $n$, the numbers $\wp(\frac{\omega_1}{n})$ and $\wp(\frac{\omega_2}{n})$ are algebraic. Moreover, any non-zero period of $\wp(z)$ is necessarily transcendental.
\end{corollary}









































\subsection{Modular function}
\quad\bigskip

A meromorphic function $f$ on $\mathbb{H}$ is said to be a modular function if it satisfies
\begin{equation}
f\left(\frac{a\,\tau+b}{c\,\tau+d}\right)=f(\tau)\quad\text{for all }
\gamma=\left(\begin{smallmatrix}
   a  & b \\
   c  & d
\end{smallmatrix} \right)\in \Gamma\end{equation}
and also is meromorphic at cusps. In particular, a modular function on the congruence subgroup $\Gamma(N)$ is said to a be modular function of level $N$.\\

A canonical example of modular function of level one is given by 
$$
j(\tau):=\frac{E_4^3(\tau)}{\Delta(\tau)} \text{ for all } \tau\in\mathbb{H}.
$$
It has the following Fourier expansion
$$
j(\tau)=\frac{1}{q}+744+196884q+\cdots,$$
where $q=e^{2\pi i\tau}$. This is made precise in the following theorem.
\begin{theorem} \label{level1}
Let $f$ be a modular function of level one and let $$K := \mathbb{Q}(\text{Fourier coefficients of }f)\subset\mathbb{C}.$$ Suppose that the poles of $f$ in $\mathcal{F}$ are $\tau_1,\tau_2,\cdots,\tau_m$, where $\mathcal{F}$ denotes the standard fundamental domain for the action of $S L_2(\mathbb{Z})$ on $\mathbb{H}$. Then $f\in K(j(\tau_1),\cdots,j(\tau_m))(j)$ i.e., $f$ is a rational function in $j$ over $K(j(\tau_1),\cdots,j(\tau_m))$.
\end{theorem}
\begin{proof}
First, we consider the function 
$$
g(\tau):=\left(\prod_{\tau_0\in\mathcal{F}} \left(j(\tau)-j(\tau_0)\right)^{\operatorname{ord}_{\tau_0}f}\right)f(\tau),
$$
which is holomorphic on $\mathbb{H}$. Since $f$ is meromorphic at $i\infty$, there exists $c>0$ such that the region $\{z\in\mathcal{F} : \operatorname{Im}z>c\}\cup \{i\infty\}$ contains no pole of $f$ except possibly $i\infty$. Again, $f$ is meromorphic on $\mathbb{H}$ and the region $\{z\in \mathcal{F}: \operatorname{Im}z\leq c\}$ is compact, and so it contains only a finite number of poles of $f$. Thus, $f$ can have only finitely many poles in $\mathcal{F}$ and hence above product is finite.\\

Suppose $g$ has pole of order $M$ at $i\infty$. Then the Fourier expansion of
$g(\tau)$ at $i\infty$ has the form
$$
g(\tau)=\sum_{n=-M}^\infty c_n q^n, \, \text{ where } c_{-M}\neq 0.
$$
Note that the modular function $g(\tau)-c_{-M}j(\tau)^M$ is holomorphic on $\mathbb{H}$ and its $q$-series starts with at most a polar term of order $M-1$. Repeating this process, we can subtract a polynomial of $j(\tau)$ to get a holomorphic modular function that vanishes at $i\infty$, and hence identically zero. Thus, $g(\tau)$ is a polynomial of $j(\tau)$ over $K$ and so $f(\tau)$ is a rational function of $j(\tau)$ over $K(j(\tau_1),\cdots,j(\tau_m))$.
\end{proof}
The above proof proves the following :
\begin{corollary}\label{rationalfunctions}
If $f(\tau)$ is a modular function of level one which is holomorphic on $\mathbb{H}$ with the following $q$-expansion at $i\infty$
$$
f(\tau)=\sum_{n=-M}^\infty c_n q^n, \ c_{-M}\neq 0
$$
with $K=\mathbb{Q}(\text{Fourier coefficients of }f)$, then $f(\tau)$ is a polynomial in $j(\tau)$ of degree $M$ with coefficients in $K$.
\end{corollary}
\begin{theorem} \label{mdft}
Let $f$ be a modular function with respect to a congruence subgroup $\Gamma$ and
$$
K:=\mathbb{Q}\left(\text{Values of }j\text{ at poles of }f\text{ in }\mathcal{F}, \text{ Fourier coefficients of }f\text{ at all cusps}\right).
$$
Then there exists a monic polynomial $P(X)\in K(j)[X]$ such that $P(f)=0$. 
\end{theorem}
\begin{proof}
We construct a polynomial satisfied by $f$ with coefficients that are rational functions of $j(\tau)$. Theorem \ref{level1} suggests that these coefficients are modular functions of level $1$.\\

Let $[S L_2(\mathbb{Z}):\Gamma]=r$ and $\{ \gamma_1=I,\gamma_2,\cdots,\gamma_r\}$ be a complete set of right coset representatives so that $$
S L_2(\mathbb{Z})=\bigsqcup_{i=1}^r\Gamma \gamma_i.
$$
For all $1\leq i \leq r$, define the functions $$
f_i(\tau):=f(\gamma_i\tau) \quad \text{for } \tau \in \mathbb{H}.
$$ Clearly, each $f_i$ is independent of the choice of coset representatives as $f$ is $\Gamma$-invariant. Again, for any $\gamma\in S L_2(\mathbb{Z})$, we have $f_i(\gamma \tau)=f(\gamma_i\gamma\tau)=f_j(\tau)$ if $\gamma_i\gamma\in \Gamma\gamma_j$, i.e., the set $\{f_1,f_2,...,f_r\}$ gets permuted under the action of $S L_2(\mathbb{Z}).$ This observation implies that any elementary symmetric polynomial of $f_1,f_2,\cdots,f_r$ is a modular function with respect to $S L_2(\mathbb{Z}).$\\

Consider the polynomial $$
P(X)=\prod_{i=1}^r(X-f_i),$$ which is satisfied by $f$ as $f=f_1$. Note that the coefficients of $P$ are elementary symmetric polynomials in $f_1,f_2,\cdots,f_r$, and hence, modular functions with respect to $S L_2(\mathbb{Z})$ with Fourier coefficients in $K$. Thus by Theorem \ref{level1}, the coefficients of $P$ are in $K(j)$ and so $P(X)$ is the required polynomial.
\end{proof}
\begin{corollary}\label{useful}
 If the modular function $f$ is holomorphic on $\mathbb{H}$ and $$K := \mathbb{Q}(\text{Fourier coefficients of }f\text{ at all cusps}),$$
 then the monic polynomial in Theorem \ref{mdft} lies in $K[j][X]$.
\end{corollary}
\begin{proof}
   This is clear from above proof and Corollary \ref{rationalfunctions}, since coefficients of $P$ constructed above are modular functions of level one, holomorphic on $\mathbb{H}$ with Fourier coefficients in $K$.  
\end{proof}
\bigskip

For $N>1$ and $a=(a_1,a_2)\in N^{-1}\mathbb{Z}^2\setminus\mathbb{Z}^2$, consider the function
$$
f_{a}(\tau):=\frac{g_2(L_\tau)\,g_3(L_\tau)}{\Delta_0(L_\tau)}\,\wp_{L_\tau} (a_1\tau+a_2),
$$
which is holomorphic on $\mathbb{H}$. Moreover, It is known that $f_a$ is a modular function of level $N$ with algebraic Fourier coefficients with respect to $e^{\frac{2\pi i\tau}{N}}$. For an elaborate discussion on these functions, the reader is referred to (\cite{shimura}, Chapter 6).\\

Let $F_{N,\overline{\mathbb{Q}}}$ denote the field of modular functions of level $N$ with algebraic Fourier coefficients with respect to $e^{\frac{2\pi i\tau}{N}}$, and $S:=\{\left(\frac{r}{N},\frac{s}{N}\right): 0\leq r,s\leq N-1 \text{ and }(r,s)\neq(0,0)\}$. Thus, $f_a\in F_{N,\overline{\mathbb{Q}}}$ for all $a\in S$.
\begin{theorem}\label{mdft..}
$F_{N,\overline{\mathbb{Q}}}$ is a finite Galois extension of $\overline{\mathbb{Q}}(j)$ and $F_{N,\overline{\mathbb{Q}}}=\overline{\mathbb{Q}}(j,f_a |a\in S)$.
\end{theorem}
For a proof of Theorem \ref{mdft..}, we refer to (\cite{shimura}, Proposition $6.9.(1)$), where they consider modular functions with Fourier coefficients in $\mathbb{Q}\big(e^{\frac{2\pi i}{N}}\big)$. But, one can use the same argument for the field $\overline{\mathbb{Q}}$ instead of $\mathbb{Q}\big(e^{\frac{2\pi i}{N}}\big)$.
\bigskip

\section{\bf Proof of Main Results}
\bigskip

\subsection{Values of modular functions when $j(\tau)$ is transcendental}
\begin{lemma} \label{mdfz}
    Let $g$ be a non-zero modular function of level $N(\geq 1)$ with algebraic Fourier coefficients, i.e., $g\in  F_{N,\overline{\mathbb{Q}}}$. If $\tau\in \mathbb{H}$ such that $j(\tau)\notin\overline{\mathbb{Q}}$ and $\tau$ is not a pole of $g$, then $g(\tau)\neq 0$.
\end{lemma}
\begin{proof}
By Theorem \ref{mdft..}, we have an irreducible polynomial 
$$
P(X)=\sum_{r=0}^m c_r(j)\,X^r \in \overline{\mathbb{Q}}[j][X], \quad m=\deg P(X)
$$
such that
$$
P(g)=0.
$$
Here $$c_r(j)=\sum_{s=0}^{d_r}c_{r s}\,j^s\in \overline{\mathbb{Q}}[j]\quad \forall\ 0\leq r\leq m, \quad d_r=\deg c_r(j).$$
Also, $c_0(j)$ is not identically zero polynomial because $P(X)$ is irreducible over $\overline{\mathbb{Q}}[j]$. Thus, we get
\begin{equation}\label{funzero}
\sum_{r=0}^m\sum_{s=0}^{d_r}c_{r s}\,j(\tau)^s\, g(\tau)^r=0\quad \forall \tau\in\mathbb{H}.
\end{equation}
Now we consider $\tau\in\mathbb{H}$ such that $j(\tau)\notin\overline{\mathbb{Q}}$. We claim that $g(\tau)\neq 0$. If not, from \eqref{funzero} we get $c_0(j(\tau))=0$ and hence, $j(\tau)\in\overline{\mathbb{Q}}$, which contradicts our assumption. This completes the proof that $g(\tau)$ is non-zero. 
\end{proof}
The following theorem shows that the non-zero number $g(\tau)$ is indeed transcendental.
\begin{theorem}\label{tranmf}
 Let $g\in F_{N,\overline{\mathbb{Q}}}$ and $\tau\in \mathbb{H}$ such that it is not a pole of $g$. If $j(\tau)\notin\overline{\mathbb{Q}}$, then $g(\tau)$ is transcendental.
\end{theorem}
\begin{proof}
If possible, let $g(\tau)\in\overline{\mathbb{Q}}$. Then the following modular function $$h(z):=g(z)-g(\tau)\quad \forall z\in\mathbb{H}$$ belongs to $ F_{N,\overline{\mathbb{Q}}}$ but vanishes at $\tau$, which contradicts Lemma \ref{mdfz}. Therefore, $g(\tau)$ is a transcendental number.
\end{proof}
\begin{corollary}\label{zero-pole}
If $g\in F_{N,\overline{\mathbb{Q}}}$ and $\tau_0$ is a zero or a pole of $g$, then $j(\tau_0)\in \overline{\mathbb{Q}}$. 
\end{corollary}
\begin{proof}
If $g(\tau_0)=0$, then Lemma \ref{mdfz} implies that $j(\tau_0)\in \overline{\mathbb{Q}}$. If $\tau_0$ is a pole of $g$, we consider the modular function
$$
h(\tau):=\frac{1}{g(\tau)}\in F_{N,\overline{\mathbb{Q}}},
$$
which satisfies $h(\tau_0)=0$, and hence $j(\tau_0)\in\overline{\mathbb{Q}}$ by Lemma \ref{mdfz}. 
\end{proof}
\subsection{Values of modular functions when $j(\tau)\in\overline{\mathbb{Q}}$}
\quad \bigskip

Here we study values of modular functions at the points on $\mathbb{H}$ where $j$-function takes algebraic values. From Theorem \ref{mdft..}, we have $$ F_{N,\overline{\mathbb{Q}}}=  \overline{\mathbb{Q}}\left(\{j,f_a |a\in S\}\right), \text{ where }S=\left\{\left(\frac{r}{N},\frac{s}{N}\right): 0\leq r, s\leq N-1 \text{ and }(r,s)\neq (0,0)\right\}.$$
Hence, for any $g\in F_{N,\overline{\mathbb{Q}}}$, the algebraic nature of $g(\tau)$ is completely determined by the numbers in the set $\{f_a(\tau):a\in S\}$.\\

Note that for $N=1$, the set $S$ is empty. This shows that for any $g\in F_{1,\overline{\mathbb{Q}}}$, the number $g(\tau)\in\overline{\mathbb{Q}}$ as $j(\tau)\in\overline{\mathbb{Q}}$. For $N>1$, we have to study the numbers $f_a(\tau)$ for all $a\in S$. We start by proving the following important lemma.\bigskip

\begin{lemma}\label{omegatau}
 If $\tau\in \mathbb{H}$ such that $j(\tau)$ is algebraic, then there exists a unique (up to an algebraic multiple) transcendental number $\omega_\tau$ for which $g_2(\omega_\tau L_\tau)$ and $g_3(\omega_\tau L_\tau)$ are both algebraic numbers.   
\end{lemma}
\begin{proof}
Note that for any $\omega_\tau\in\mathbb{C}^\times $,
\begin{equation}\label{omegatau.}
    j(\tau)=\frac{E_4(\tau)^3}{\Delta(\tau)}=\frac{1728\,  g_2(L_\tau)^3}{\Delta_0(L_\tau)}=\frac{1728\, g_2\left(\omega_\tau L_\tau\right)^3}{\Delta_0\left(\omega_\tau L_\tau\right)},
\end{equation}
where $L_\tau=\tau\,\mathbb{Z}\oplus\mathbb{Z}$ and $\omega_\tau L_\tau=\tau\,\omega_\tau\,\mathbb{Z}\oplus\omega_\tau\mathbb{Z}$. Here the last equality follows from the fact that \begin{equation}\label{homogenity}
g_2(\omega_\tau L_\tau)=\omega_\tau^{-4}g_2(L_\tau).
\end{equation} If we choose $\omega_\tau$ such that $\omega_\tau^4=g_2(L_\tau)$, then $g_2(\omega_\tau L_\tau)=1$. Given that $j(\tau)$ is algebraic, and so \eqref{omegatau.} shows that $\Delta_0(\omega_\tau L_\tau)$ is also algebraic. Thus, the numbers $g_2(\omega_\tau L_\tau)$ and $g_3(\omega_\tau L_\tau)$ both are algebraic. Moreover,  $\omega_\tau$ is a period of $\omega_\tau L_\tau$ and both $g_2(\omega_\tau L_\tau)$ and $g_3(\omega_\tau L_\tau)$ are algebraic. By Corollary \ref{tranp} we conclude that $\omega_\tau$ is transcendental.\\

To prove uniqueness, consider an arbitrary $\omega_\tau'\in\mathbb{C}^\times$ such that both $g_2(\omega_\tau' L_\tau)$ and $g_3(\omega_\tau' L_\tau)$ are algebraic. From \eqref{homogenity} we have 
$$
 \quad g_2(\omega_\tau' L_\tau)= (\omega_\tau')^{-4}\,g_2(L_\tau)=\alpha, (\text{say}).
$$
Therefore, we get $g_2(L_\tau)=\alpha\cdot (\omega_\tau')^{4}$ and hence $\omega_\tau^4=\alpha\cdot (\omega_\tau')^{4}.$
This completes the proof that $\omega_\tau$ is unique up to an algebraic multiple.
\end{proof}

Recall that
$$
f_a(\tau):=\frac{g_2(L_\tau)\,g_3(L_\tau)}{\Delta_0(L_\tau)}\,\wp_{L_\tau}(a_1\tau+a_2),
$$ 
where $L_\tau=\tau\,\mathbb{Z}\oplus\mathbb{Z}$ and $a=(a_1,a_2)\in S$. We can rewrite it as follows:
$$
f_a(\tau)=\frac{g_2(\omega_\tau L_\tau)\, g_3(\omega_\tau L_\tau)}{\Delta_0(\omega_\tau L_\tau)}\,\wp_{\omega_\tau L_\tau}(a_1\omega_\tau\tau+a_2\omega_\tau).
$$
\begin{lemma}\label{algfa}
    If $\tau\in\mathbb{H}$ such that $j(\tau)\in\overline{\mathbb{Q}}$, then $f_a(\tau)\in\overline{\mathbb{Q}}$ for all $a\in S$.
\end{lemma}
\begin{proof}
    Let $a=(\frac{r}{N},\frac{s}{N})\in S$. Then we have
$$
f_a(\tau)=\frac{g_2(\omega_\tau L_\tau)\, g_3(\omega_\tau L_\tau)}{\Delta_0(\omega_\tau L_\tau)}\,\wp_{\omega_\tau L_\tau}\left(\frac{r\,\omega_\tau\,\tau}{N}+\frac{s\,\omega_\tau}{N}\right),
$$
where $\omega_\tau$ is chosen as in Lemma \ref{omegatau} so that $g_2(\omega_\tau L_\tau),\, g_3(\omega_\tau L_\tau)\in\overline{\mathbb{Q}}$.\\

Recall that both $\wp_{\omega_\tau L_\tau}\left(\frac{\omega_\tau\, \tau}{N} \right)$ and $\wp_{\omega_\tau L_\tau}\left(\frac{\omega_\tau}{N} \right)$  are algebraic by Corollary \ref{tranp}. Moreover, by the addition formula for Weierstrass $\wp$-function in Theorem \ref{addf}, we get that $$\wp_{\omega_\tau L_\tau}\left(\frac{r\,\omega_\tau\, \tau}{N}+\frac{s\,\omega_\tau}{N}\right)\in\overline{\mathbb{Q}}$$ for all $a=\left(\frac{r}{N},\frac{s}{N}\right)\in S$. Thus, the number $f_a(\tau)$ is algebraic for all $a\in S$.
\end{proof}
\begin{theorem}\label{algemf}
Let $g\in F_{N,\overline{\mathbb{{Q}}}}$ and $\tau\in\mathbb{H}$ such that $j(\tau)\in \overline{\mathbb{Q}}$. If $\tau$ is not a pole of $g$, then $g(\tau)$ is algebraic.
\end{theorem}
\begin{proof}
 From Lemma \ref{algfa}, we have   $f_a(\tau)\in\overline{\mathbb{Q}}$ for all $a\in S$. Then we use the following fact $$g\in  F_{N,\overline{\mathbb{Q}}}=  \overline{\mathbb{Q}}\left(\{j,f_a |a\in S\}\right)$$  to conclude that $g(\tau)$ is algebraic. 
\end{proof}
\bigskip

\begin{proof}[\bf Proof of Theorem \ref{suma}]
    \, Combining Theorem \ref{tranmf} with Theorem \ref{algemf} we get (i), and (ii) follows from Theorem \ref{mahl} and (i).
\end{proof}
\bigskip

\subsection{Values of modular forms}
\quad\bigskip

Recall that for positive integers $k$ and $N$, we denote $M_{k,N,\overline{\mathbb{Q}}}$, $M_{k,N,\overline{\mathbb{Q}}}^w$ and $M_{k,N,\overline{\mathbb{Q}}}^m$ as the set of holomorphic, weakly holomorphic and meromorphic modular forms respectively with algebraic Fourier coefficients. Here, we restate Theorem \ref{algem''} and give a proof following the same technique from \cite{hamieh-murty} and include the case for meromorphic modular forms. First, we prove the following lemma:
\begin{lemma}\label{det}
    If $n$ and $m$ are positive integers such that $n>2m$, then the determinant of the matrix
    $$
A=\left[\begin{array}{cccccc}
    {{n}\choose{0}} & {{n-1}\choose{0}} & {{n-2}\choose{0}} & \ldots & {{n-m+1}\choose{0}} & {{n-m}\choose{0}} \\\\
    {{n}\choose{1}} & {{n-1}\choose{1}} & {{n-2}\choose{1}} & \ldots & {{n-m+1}\choose{1}} & {{n-m}\choose{1}}\\\\
    {{n}\choose{2}} & {{n-1}\choose{2}} & {{n-2}\choose{2}} & \ldots & {{n-m+1}\choose{2}} & {{n-m}\choose{2}}\\\\
    \vdots & \vdots & \vdots & \ldots & \vdots & \vdots \\\\
    {{n}\choose{m}} & {{n-1}\choose{m}} & {{n-2}\choose{m}} & \ldots & {{n-m+1}\choose{m}} & {{n-m}\choose{m}}
\end{array} \right]_{(m+1)\times (m+1)}
  $$ 
  is given by $(-1)^{\frac{m(m+1)}{2}}$.
\end{lemma}
\begin{proof}
  To find the determinant of $A$, we apply elementary column operations successively on $A$. By using the identity ${n\choose{k}}={n-1\choose{k-1}}+{n-1\choose{k}}$, we first rewrite the first column as follows: 
$$C_1=\left[\begin{array}{c}
     {n\choose{0}} \\\\
     {n-1\choose{0}}+{n-1\choose{1}}\\\\
     {n-1\choose{1}}+{n-1\choose{2}}\\\\
     \vdots\\\\
     {n-1\choose{m-1}}+{n-1\choose{m}}
\end{array} \right]_{(m+1)\times 1},$$ and then subtract the second column $C_2$ from $C_1$ (i.e.,  $C_1'=C_1-C_2$) to get
$$C_1'=\left[\begin{array}{c}
     0 \\\\
     {n-1\choose{0}}\\\\
     {n-1\choose{1}}\\\\
     \vdots\\\\
     {n-1\choose{m-1}}
\end{array} \right]_{(m+1)\times 1}.$$
We continue this process for each column $C_i$ and and replace $C_i$ by $C_i-C_{i+1}$ for all $1\leq i\leq m$ to obtain the following matrix  
$$
A'=\left[\begin{array}{cccccc}
    0 & 0 & 0 & \ldots & 0 & {{n-m}\choose{0}} \\\\
    {{n-1}\choose{0}} & {{n-2}\choose{0}} & {{n-3}\choose{0}} & \ldots & {{n-m}\choose{0}} & {{n-m}\choose{1}}\\\\
    {{n-1}\choose{1}} & {{n-2}\choose{1}} & {{n-3}\choose{1}} & \ldots & {{n-m}\choose{1}} & {{n-m}\choose{2}}\\\\
    \vdots & \vdots & \vdots & \ldots & \vdots & \vdots \\\\
    {{n-1}\choose{m-1}} & {{n-2}\choose{m-1}} & {{n-3}\choose{m-1}} & \ldots & {{n-m}\choose{m-1}} & {{n-m}\choose{m}}
\end{array} \right]_{(m+1)\times (m+1)}.
$$
 Again, we start from the first column of $A'$ and repeat the same process to reduce the matrix in the form 
 $$
B=\left[\begin{array}{cccccc}
    0 & 0 & 0 & \ldots & 0 & {{n-m}\choose{0}} \\\\
    0 & 0 & 0 & \ldots & {{n-m}\choose{0}} & *\\\\
    0 & 0 & {{n-m}\choose{0}} & \iddots & * & *\\\\
    \vdots & \iddots & \vdots & \ldots & \vdots & \vdots \\\\
    {{n-m}\choose{0}} & * & * & \ldots & * & *
\end{array} \right]_{(m+1)\times (m+1)}.
 $$
 Since $B$ is obtained from $A$ by elementary column operations, we conclude that $\det(A)=\det(B)=(-1)^{\frac{m(m+1)}{2}}$.
\end{proof}





















\begin{theorem}\label{algem'}
   Let $f\in M_{k,N,\overline{\mathbb{Q}}}\,\cup M_{k,N,\overline{\mathbb{Q}}}^w\,\cup M_{k,N,\overline{\mathbb{Q}}}^m$. If $\tau\in\mathbb{H}$ such that $e^{2\pi i\tau}$ is algebraic, then $f(\tau)$ is transcendental. 
\end{theorem}
\begin{proof}
    We first construct the function
$$g(\tau):=\frac{f(\tau)^{12}}{\Delta(\tau)^k}\quad \forall \tau\in\mathbb{H},$$
which is holomorphic or meromorphic on $\mathbb{H}$ depending on $f$ because $\Delta$ is non-vanishing on $\mathbb{H}$.
Moreover, $g\in F_{N,\overline{\mathbb{Q}}}$ as Fourier coefficients of $f$ and $\Delta$ are algebraic.\\

From Theorem \ref{mdft}, we know that $g$ has the minimal polynomial $P[X]\in \overline{\mathbb{Q}}(j)[X]$, and if needed, we can clear denominators of the coefficients of $P[X]$ to get a polynomial (use the same notation)
$$P(X)=\sum_{r=0}^m c_r(j)\,X^r\in \overline{\mathbb{Q}}[j][X],\quad \deg P(X)=m$$
such that
$P(g(\tau))=0$ for any $\tau\in \mathbb{H}$. Here, we have
$$
c_r(j)=\sum_{s=0}^{d_r}c_{r,s}\,j^s\in \overline{\mathbb{Q}}[j]\quad \forall\ 0 \leq r\leq m
$$
and $d_r=\deg c_r(j)$ so that $c_{r,d_r}\neq 0$ for all $0\leq r\leq m$. Moreover, Corollary \ref{useful} shows that if $f$ is holomorphic on $\mathbb{H}$, then the minimal polynomial $P[X]$ of $g$ is monic and belongs to $\overline{\mathbb{Q}}[j][X]$, and hence $c_m(j)=c_{m,0}=1$.\\

For any $\tau\in\mathbb{H}$, we have $$P\left(\frac{f(\tau)^{12}}{\Delta(\tau)^k}\right)=0,$$
and hence for any positive integer $l$, 
$$
\begin{aligned}
& \Delta(\tau)^l\, P\left(\frac{f(\tau)^{12}}{\Delta(\tau)^k}\right)=0\\
\Rightarrow &\sum_{r=0}^m\sum_{s=0}^{d_r}c_{r,s}\,j(\tau)^s f(\tau)^{12 r}\Delta(\tau)^{l-k r}=0\\
\Rightarrow &\sum_{r=0}^m\sum_{s=0}^{d_r}c_{r,s}\,E_4(\tau)^{3 s}f(\tau)^{12\, r}\Delta(\tau)^{l-k r-s}=0.
\end{aligned}
$$
We choose $l> k m +2\cdot\operatorname{max}\{d_0,d_1,\ldots,d_m\}$ so that the powers of $\Delta(\tau)$ are always positive integers. Then we put $$\Delta(\tau)=\frac{E_4(\tau)^3-E_6(\tau)^2}{1728}$$ to obtain \begin{equation} \label{trane}
\sum_{r=0}^m \sum_{s=0}^{d_r} \sum_{t=0}^{l-k r-s} \frac{(-1)^t c_{r, s}}{1728^{l-k r-s}}\left(\begin{array}{c}
l-k\,r-s \\
t
\end{array}\right) f(\tau)^{12\, r} E_4(\tau)^{3(l-k\, r-t)} E_6(\tau)^{2\, t}=0.   
\end{equation}
Now we consider $\tau\in \mathbb{H}$ such that $e^{2\pi i\tau}$ is algebraic. By Theorem \ref{suma}, the number $j(\tau)$ is transcendental, and hence Corollary \ref{zero-pole} guarantees that $\tau$ is neither a zero nor a pole of $f$. Moreover, we claim that the number $f(\tau)$ is not algebraic. For otherwise, \eqref{trane} becomes a polynomial in two variables over $\overline{\mathbb{Q}}$, satisfied by $E_4(\tau)$  and $E_6(\tau)$. If we can prove that this is a non-zero polynomial of $E_4(\tau)$ and $E_6(\tau)$, then we would get a contradiction with Nesterenko's Theorem \ref{nest}. Hence, it would complete the proof that $f(\tau)$ is transcendental.\\

To show that the left side of \eqref{trane} becomes a non-zero polynomial of $E_4(\tau)$ and $E_6(\tau)$ when $e^{2\pi i\tau}\in\overline{\mathbb{Q}}$, we consider the following cases:\\

\emph{Case-$1$: } If $f\in M_{k,N,\overline{\mathbb{Q}}}\,\cup M_{k,N,\overline{\mathbb{Q}}}^w$, then $c_m(j)=c_{m,0}$, and hence the coefficient of $E_4(\tau)^{3(l-k m)}\,E_6(\tau)^{0}$ (which appears in \eqref{trane} only for $r=m,\, t=0$) is given by $$\frac{1}{1728^{l-k m}}\cdot c_{m,0}\cdot f(\tau)^{12 m}\neq 0,$$
which proves that the polynomial is non-zero.\\

\emph{Case-$2$: } If $f\in M_{k,N,\overline{\mathbb{Q}}}^m$, then the leading coefficient of $P[X]$ is given by $$c_m(j)=\sum_{s=0}^{d_m}c_{m,s}\,j^s.$$
Moreover, it cannot be \emph{identically zero polynomial} because the degree of $P[X]$ is same as the degree of the minimal polynomial of $g$. Since $l>k\, m+2\, d_m$, the following terms
$$
E_4(\tau)^{3(l-k\, m-t)}\, E_6(\tau)^{2\, t} \quad \forall\, 0\leq t\leq d_m
$$
with coefficients
$$
\frac{(-1)^t\, f(\tau)^{12\, m}}{1728^{l-k\,m}}\sum_{s=0}^{d_m}\left(\begin{array}{c}
l-k\,m-s \\
t
\end{array}\right) (1728)^s\, c_{m,s} \quad \forall\, 0\leq t\leq d_m
$$
appear in \eqref{trane}. We claim that at least one of these terms has non-zero coefficient. If not, we consider the following system of equations: $ A\, X=0$, where 
$$
A=\left[\begin{array}{ccccc}
    {{l-km}\choose{0}} & {{l-km-1}\choose{0}} & {{l-km-2}\choose{0}} & \ldots & {{l-km-d_m}\choose{0}} \\\\
    {{l-km}\choose{1}} & {{l-km-1}\choose{1}} & {{l-km-2}\choose{1}} & \ldots & {{l-km-d_m}\choose{1}}\\\\
    {{l-km}\choose{2}} & {{l-km-1}\choose{2}} & {{l-km-2}\choose{2}} & \ldots & {{l-km-d_m}\choose{2}}\\\\
    \vdots & \vdots & \vdots & \ldots & \vdots \\\\
    {{l-km}\choose{d_m}} & {{l-km-1}\choose{d_m}} & {{l-km-2}\choose{d_m}} & \ldots & {{l-km-d_m}\choose{d_m}}
\end{array} \right]_{(d_m+1)\times (d_m+1)}  X=\left[\begin{array}{c}
   x_0  \\
    x_1 \\
    x_2 \\
    \vdots\\
    x_{d_m}
\end{array}\right]_{d_m+1.}
$$
This system of equations is satisfied by the non-trivial vector $$\left[\begin{array}{c}
   (1728)^0\, c_{m,0} \\
    (1728)^1\, c_{m,1} \\
    (1728)^2\, c_{m,2} \\
    \vdots\\
    (1728)^{d_m}\, c_{m,d_m}
\end{array}\right].$$ From Lemma \ref{det}, we get that $\det(A)\neq 0$, which contradicts the fact that the above system has only trivial solution. Thus, it completes the proof that the polynomial relation of $E_4(\tau)$ and $E_6(\tau)$ in \eqref{trane} is non-trivial.
\end{proof}
\begin{corollary}\label{algi}
If $f\in M_{k,N,\overline{\mathbb{Q}}}\,\cup M_{k,N,\overline{\mathbb{Q}}}^w$, then $f$ is algebraically dependent with $E_4$ and $E_6$ over $\overline{\mathbb{Q}}$. In particular, If $f\in M_{k,N,\overline{\mathbb{Q}}}^m$ and $\tau\in\mathbb{H}$ is not a pole of $f$, then $f(\tau)$ is algebraic over $\overline{\mathbb{Q}}(E_4(\tau),E_6(\tau))$. 
\end{corollary}
\begin{proof}
    It directly follows from the equation \eqref{trane} because the left-hand side is a non-zero polynomial in $f, E_4$ and $E_6$ over $\overline{\mathbb{Q}}.$
\end{proof}
\bigskip

Now we consider the points $\tau\in\mathbb{H}$ such that $j(\tau)\in\overline{\mathbb{Q}}$. From Lemma \ref{omegatau}, we get a transcendental number $\omega_\tau$ such that $g_2(\omega_\tau L_\tau)$ and $g_3(\omega_\tau L_\tau)$ are both algebraic. Moreover, we have the formulae
$$
\begin{aligned}
E_4(\tau)&=\frac{3}{4\pi^4}\, g_2(L_\tau)=\frac{3}{4}\left(\frac{\omega_\tau}{\pi} \right)^4 g_2(\omega_\tau L_\tau),\\ E_6(\tau)&=\frac{27}{8\pi^6}\,g_3(L_\tau)=\frac{27}{8}\left(\frac{\omega_\tau}{\pi}\right)^6g_3(\omega_\tau L_\tau).
\end{aligned}
$$
From Nesterenko's Theorem \ref{nest}, we know that at most one of $E_4(\tau)$ and $E_6(\tau)$ is algebraic. The above formulae imply that if $\frac{\omega_\tau}{\pi}\in\overline{\mathbb{Q}}$, then both $E_4(\tau)$ and $E_6(\tau)$ are algebraic, which is a contradiction. Hence, $\frac{\omega_\tau}{\pi}$ is a transcendental number. Besides, we have 
$$
\begin{aligned}
\Delta(\tau)=\frac{E_4(\tau)^3-E_6(\tau)^2}{1728}&=\frac{1}{4^6}\left(\frac{\omega_\tau}{\pi}\right)^{12}\left(g_2(\omega_\tau L_\tau)^3-27 g_3(\omega_\tau L_\tau)^2 \right)\\
&=\frac{1}{4^6}\left(\frac{\omega_\tau}{\pi}\right)^{12}\Delta_0(\omega_\tau L_\tau).  
\end{aligned}
$$
As $g_2(\omega_\tau L_\tau)$ and $g_3(\omega_\tau L_\tau)$ both are algebraic, the number $\Delta_0(\omega_\tau L_\tau)\in\overline{\mathbb{Q}}\setminus \{0\}$. Since $\frac{\omega_\tau}{\pi}$ is transcendental and $\Delta(\tau)$ is a non-zero algebraic multiple of $\left(\frac{\omega_\tau}{\pi}\right)^{12}$, we deduce that $\Delta(\tau)$ is transcendental. We use this observation to prove the following result.
\begin{theorem}\label{algem}
If $f\in M_{k,N,\overline{\mathbb{Q}}}\,\cup M_{k,N,\overline{\mathbb{Q}}}^w\,\cup M_{k,N,\overline{\mathbb{Q}}}^m$ and $\tau\in\mathbb{H}$ such that $j(\tau)\in \overline{\mathbb{Q}}$, then there exists a transcendental number $\omega_\tau$ (which depends only on $\tau$ and is $\overline{\mathbb{Q}}$-linearly independent with $\pi$) such that $\left(\frac{\pi}{\omega_\tau}\right)^k f(\tau)\in \overline{\mathbb{Q}}$. Therefore, $f(\tau)$ is either zero or transcendental.
\end{theorem}
\begin{proof}
     We first construct the function $$g(\tau):=\frac{f(\tau)^{12}}{\Delta(\tau)^k}\quad \forall \tau\in\mathbb{H},$$ which lies in $ F_{N,\overline{\mathbb{Q}}}$ by the same argument as in Theorem \ref{algem'}. Since $j(\tau)\in\overline{\mathbb{Q}}$, Theorem \ref{algemf} implies that $g(\tau)$ is an algebraic number $c$, (say). Thus, we get  \begin{equation}\label{forms}
f(\tau)^{12}=c.\Delta(\tau)^k=\frac{c}{4^6}.\Delta_0(\omega_\tau L_\tau)^k\left(\frac{\omega_\tau}{\pi} \right)^{12 k}. 
\end{equation}
This shows that $f(\tau)$ is zero when $c=0$, and otherwise, it a non-zero algebraic multiple of $\left(\frac{\omega_\tau}{\pi} \right)^{k}$ and hence transcendental.
\end{proof}
\bigskip

\subsection{Proof of Theorem \ref{deg.}}\, \bigskip

 By theorem \ref{mdft..}, the numbers $g(\tau)$ and $j(\tau)$ are algebraically dependent over $\overline{\mathbb{Q}}$. Moreover, the number $j(\tau)$ is algebraic if and only if $g(\tau)\in\overline{\mathbb{Q}}$ by Corollary \ref{suma}. These facts together give the first equality below:
$$
\begin{aligned}
 trdeg\,\overline{\mathbb{Q}}\left(e^{2\pi i\tau},\, f(\tau),\, g(\tau) \right)&= trdeg \,\overline{\mathbb{Q}}\left(e^{2\pi i\tau},\, f(\tau),\, j(\tau) \right)\\
 &= trdeg \,\overline{\mathbb{Q}}\left(e^{2\pi i\tau},\, \frac{f(\tau)^{12}}{\Delta(\tau)^k},\, f(\tau),\, j(\tau) \right)\\
 &= trdeg \,\overline{\mathbb{Q}}\left(e^{2\pi i\tau},\, \Delta(\tau)^k,\, f(\tau),\, j(\tau) \right)\\
 &= trdeg \,\overline{\mathbb{Q}}\left(e^{2\pi i\tau},\, \Delta(\tau),\, f(\tau),\, j(\tau) \right) \\
 &= trdeg \,\overline{\mathbb{Q}}\left(e^{2\pi i\tau},\, E_4(\tau)^3,\, E_6(\tau)^2,\, f(\tau),\, j(\tau) \right)\\
 &=trdeg \,\overline{\mathbb{Q}}\left(e^{2\pi i\tau},\, E_4(\tau),\, E_6(\tau),\, f(\tau),\, j(\tau) \right)\\
 &= trdeg \,\overline{\mathbb{Q}}\left(e^{2\pi i\tau},\, E_4(\tau),\, E_6(\tau),\, f(\tau) \right)\\
 &= trdeg \,\overline{\mathbb{Q}}\left(e^{2\pi i\tau},\, E_4(\tau),\, E_6(\tau)\right).
\end{aligned}
$$ 
Here, the second equality follows from the fact that $f^{12}/\Delta^k\in F_{N,\overline{\mathbb{Q}}}$ and  $f(\tau)^{12}/\Delta(\tau)^k$ is algebraically dependent with $j(\tau)$ by Theorem \ref{mdft..}. The third and the seventh equality hold as the fields are identical. To prove the equality of fields in the third step, we need the assumption that $f(\tau)\neq 0$.  
The final equality holds because $f(\tau)$ is algebraic over $\overline{\mathbb{Q}}\left(E_4(\tau), E_6(\tau)\right)$ by Corollary \ref{algi}. Then we apply Nesterenko's Theorem \ref{nest} on the right side of the final equality to conclude that the transcendence degree of the field $\overline{\mathbb{Q}}\left(e^{2\pi i\tau},\, f(\tau),\, g(\tau) \right)$ over $\overline{\mathbb{Q}}$ is at least two. 
\bigskip

\subsection{Values of quasi-modular forms}\,\bigskip

Let $\widetilde{M}_{k,\overline{\mathbb{Q}}}^{(p)}(\Gamma)$ denote the set of all quasi-modular forms of weight $k$ and depth $p(>0)$ for $\Gamma$ with algebraic Fourier coefficients. We study their values at the points $\tau\in\mathbb{H}$, where one of $e^{2\pi i\tau}$ and $j(\tau)$ is algebraic.
\begin{theorem}
Let $\Gamma$ be a congruence subgroup of level $N(\geq 1)$ and $\widetilde{f}\in \widetilde{M}_{k,\overline{\mathbb{Q}}}^{(p)}(\Gamma)$. If $\tau\in\mathbb{H}$ such that $e^{2\pi i\tau}$ is algebraic, then $f(\tau)$ is transcendental. 
\end{theorem}
\begin{proof}
   By Theorem \ref{quasistructure}, we can write $\widetilde{f}$ in the form $$\widetilde{f}=\sum_{r=0}^p g_r E_2^r,\quad g_r\in M_{k,N,\overline{\mathbb{Q}}}.$$ From Corollary \ref{algi}, we know that each number $g_r(\tau)$ is algebraically dependent with $E_4(\tau), E_6(\tau)$. Therefore, we get
 $$ trdeg\,\overline{\mathbb{Q}}(g_0(\tau), g_2(\tau),\ldots, g_p(\tau), E_4(\tau), E_6(\tau))=trdeg\,\overline{\mathbb{Q}}( E_4(\tau), E_6(\tau)).$$
If possible, let both $e^{2\pi i\tau}, \widetilde f(\tau)\in \overline{\mathbb{Q}}$. Since 
$\widetilde f(\tau)=\sum_{r=0}^p g_r(\tau) E_2(\tau)^r$ and $g_r(\tau)\neq 0$ (by Theorem \ref{algem'}), the number $E_2(\tau)$ is algebraic over $\overline{\mathbb{Q}}\left(g_0(\tau), g_1(\tau),\ldots, g_p(\tau), E_4(\tau), E_6(\tau)\right)$ and so
$$
\begin{aligned}
&\quad trdeg\, \overline{\mathbb{Q}}\left(E_2(\tau), g_0(\tau), g_1(\tau),\ldots, g_p(\tau), E_4(\tau), E_6(\tau)\right)\\ &=trdeg\,\overline{\mathbb{Q}}(g_0(\tau), g_1(\tau),\ldots, g_p(\tau), E_4(\tau), E_6(\tau))\\ 
&=trdeg\,\overline{\mathbb{Q}}(E_4(\tau), E_6(\tau)).
\end{aligned}
$$
Since $e^{2\pi i\tau}$ is an algebraic number, the following field containment holds
$$
\overline{\mathbb{Q}}(e^{2\pi i\tau},\, E_2(\tau),\,E_4(\tau), E_6(\tau))\subseteq \overline{\mathbb{Q}}\left(E_2(\tau), g_0(\tau), g_1(\tau),\ldots, g_p(\tau), E_4(\tau), E_6(\tau)\right).
$$
Therefore, we get 
$$
\begin{aligned}
 &\quad trdeg\, \overline{\mathbb{Q}}(e^{2\pi i\tau},\, E_2(\tau),\,E_4(\tau),\, E_6(\tau))\\ &\leq trdeg\,\overline{\mathbb{Q}}\left(E_2(\tau),\, g_0(\tau),\, g_1(\tau),\ldots, g_p(\tau), E_4(\tau),\, E_6(\tau)\right)\\
  &= trdeg\,\overline{\mathbb{Q}}(E_4(\tau),\, E_6(\tau))\\
  &\leq 2,
\end{aligned}
$$
which contradicts Nesterenko's Theorem \ref{nest}. Hence, $\widetilde{f}(\tau)$ is a transcendental number. 
\end{proof}
To study values of quasi-modular forms in the complementary case (i.e., for $j(\tau)\in\overline{\mathbb{Q}}$), we need the following lemma.
\begin{lemma}\label{peri}
 If $\tau\in \mathbb{H}$ such that $j(\tau)$ is algebraic, then there exists a unique (up to algebraic multiple) transcendental number $\omega_\tau$ for which $\frac{\omega_\tau}{\pi}$ and $\frac{\eta_2}{\pi}$ are algebraically independent over $\overline{\mathbb{Q}}$. Here, $\eta_2$ is the second quasi-period with respect to the lattice $\omega_\tau L_\tau$ 
\end{lemma}
\begin{proof}
    We have already proved existence of $\omega_\tau$ in Lemma \ref{omegatau}.  From the definition of the Weierstrass zeta-function, one gets
 $   
Z_{\omega_\tau L_\tau}(z)= \frac{1}{\omega_\tau}\cdot Z_{L_\tau}\left(\frac{z}{\omega_\tau}\right).
$ Then we use the definition of quasi-period and the identity \eqref{quasi-period} obtain
$$
\eta_2=Z_{\omega_\tau L_\tau}(z+\omega_\tau)-Z_{\omega_\tau L_\tau}(z)= \frac{1}{\omega_\tau}\,\eta_2(L_\tau)=\frac{1}{\omega_\tau}\,G_2(\tau)= \frac{1}{3}\cdot\frac{\pi^2}{ \omega_\tau}\,E_2(\tau).
$$
Thus, we get the following formulae 
$$
E_2(\tau)=3\,\frac{\omega_\tau}{\pi}\frac{\eta_2}{\pi}, \quad E_4(\tau)=\frac{3}{4}\left(\frac{\omega_\tau}{\pi} \right)^4 g_2(\omega_\tau L_\tau),\quad E_6(\tau)=\frac{27}{8}\left(\frac{\omega_\tau}{\pi}\right)^6g_3(\omega_\tau L_\tau).$$
The above formulae imply that $E_2(\tau)$, $E_4(\tau)$ and $E_6(\tau)$ are algebraic over $\overline{\mathbb{Q}}\left(\frac{\omega_\tau}{\pi}, \frac{\eta_2}{\pi}\right)$. If possible, let $\frac{\omega_\tau}{\pi}, \frac{\eta_2}{\pi}$ be algebraically dependent. Then we get that $$trdeg\,\overline{\mathbb{Q}}(E_2(\tau),E_4(\tau),E_6(\tau))=1,$$ which contradicts Nesterenko's Theorem \ref{nest}. Hence, the numbers $\frac{\omega_\tau}{\pi}$ and $\frac{\eta_2}{\pi}$ are algebraically independent over $\overline{\mathbb{Q}}$.
\end{proof}
Lemma \ref{peri} and Theorem \ref{quasistructure} together allow us to describe values of quasi-modular forms explicitly at the points $\tau\in\mathbb{H}$ where $j$-function is algebraic.
\begin{theorem}\label{tranq}
 Let $\widetilde{f}\in \widetilde{M}_{k,\overline{\mathbb{Q}}}^{(p)}(\Gamma)$ and $\tau\in\mathbb{H}$ such that $j(\tau)\in\overline{\mathbb{Q}}$, then $f(\tau)$ is a polynomial in $\frac{\omega_\tau}{\pi}$ and $\frac{\eta_2}{\pi}$ with algebraic coefficients. Therefore, $f(\tau)$ is either zero or transcendental.   
\end{theorem}
\begin{proof}
    From Theorem \ref{quasistructure}, we have the expression $$\widetilde{f}(\tau)=\sum_{r=0}^p g_r(\tau) E_2(\tau)^r,\quad \text{where } g_r\in M_{k,N,\overline{\mathbb{Q}}}\,\cup M_{k,N,\overline{\mathbb{Q}}}^w.$$
Then we apply Theorem \ref{algem} for each coefficient $g_r$, and the above formula for $E_2$ such that 
\begin{equation}\label{last}
\widetilde{f}(\tau)=\sum_{r=0}^p c_r\left(\frac{\omega_\tau}{\pi}\right)^{k-2 r}\left(\frac{\omega_\tau}{\pi}\cdot\frac{\eta_2}{\pi}\right)^r=\sum_{r=0}^p\, c_r\, \left(\frac{\omega_\tau}{\pi}\right)^{k-r}\, \left(\frac{\eta_2}{\pi}\right)^r,
 \end{equation}
where each $c_r$ is an algebraic constant. If $\widetilde{f}(\tau)\in\overline{\mathbb{Q}}\setminus \{0\}$, then the above expression of $f(\tau)$ shows that $\frac{\omega_\tau}{\pi}$ and $\frac{\eta_2}{\pi}$ are algebraically dependent, which contradicts Lemma \ref{peri}. Therefore, the number $\widetilde{f}(\tau)$ is either zero or transcendental.     
\end{proof}
\bigskip

\subsection{Proof of Theorem \ref{deg..}}
\,\bigskip

By the same argument as in Theorem \ref{deg.}, we get
\begin{equation}\label{final.}
trdeg\,\overline{\mathbb{Q}}\left( e^{2\pi i\tau},\, f(\tau),\, g(\tau),\, \widetilde f(\tau)\right)= trdeg\,\overline{\mathbb{Q}}\left( e^{2\pi i\tau},\, E_4(\tau),\, E_6(\tau),\, \widetilde f(\tau)\right).
\end{equation}
Moreover, we can apply Theorem \ref{quasistructure} to write
$$\widetilde{f}(\tau)=\sum_{r=0}^p g_r(\tau) E_2(\tau)^r,\quad \text{where } g_r\in M_{k,N,\overline{\mathbb{Q}}}.$$
Thus, we obtain 
  $$
\begin{aligned}
&trdeg\,\overline{\mathbb{Q}}\left( e^{2\pi i\tau},\, E_4(\tau),\, E_6(\tau),\, \widetilde f(\tau)\right)\\
&=trdeg\,\overline{\mathbb{Q}}\left( e^{2\pi i\tau},\, E_4(\tau),\, E_6(\tau),\,\sum_{r=0}^p g_r(\tau)\, E_2(\tau)^r \right)\\
&=trdeg\,\overline{\mathbb{Q}}\left( e^{2\pi i\tau}, E_4(\tau), E_6(\tau), g_0(\tau),\ldots,g_p(\tau),\sum_{r=0}^p g_r(\tau)\, E_2(\tau)^r \right)\\
&=trdeg\,\overline{\mathbb{Q}}\left( e^{2\pi i\tau},\, E_4(\tau),\, E_6(\tau),\, g_0(\tau),\,\ldots,\,g_p(\tau),\,E_2(\tau) \right)\\
&=trdeg\,\overline{\mathbb{Q}}\left( e^{2\pi i\tau},\, E_4(\tau),\, E_6(\tau),\,E_2(\tau) \right),
\end{aligned}
 $$
where the second and fourth equalities follow from the fact that each coefficient $g_r(\tau)$ is algebraic over $\overline{\mathbb{Q}}\left(E_4(\tau), E_6(\tau)\right)$ by Corollary \ref{algi}. Moreover, Nesterenko's theorem \ref{nest} implies that the right side of the final equality is at least $3$, and then we use the identity \eqref{final.} to conclude the proof. 












\begin{thebibliography}{0}

\bibitem{b-d-g-p}
K. Barr\'e-Sirieix, G. Diaz, F. Gramain, G. Philibert, \textit{Une preuve de la conjecture de Mahler-Manin}, Invent. Math., 124 (1996), 1-9.

\bibitem{bertrand}
D. Bertrand, \textit{S\'eries d’Eisenstein et transcendance.} Bull. Soc. Math. France
$104(3)$, $309–321 (1976)$

\bibitem{quasi}
    C. Y. Chang, \textit{Transcendence of special values of quasi-modular forms}, Forum Math. $24$ (2012) $539-551$.


\bibitem{chudnovsky}
    G.V. Chudnovsky, \textit{Algebraic independence of constants connected with the exponential and the elliptic functions}, Dokl.
Akad. Nauk Ukrain. SSR Ser. A $8$ $(1976)$ $698–701$.


\bibitem{stromberg}
H. Cohen and F. Stromberg, \textit{Modular Forms, A Classical Approach}, Graduate Studies In Matheatics $179$.

\bibitem{cohen}
P. Cohen, \textit{On the coefficients of the transformation polynomials for the elliptic modular function}, Math. Proc. Cambridge Phil. Soc. $95$, (1984), 389-402. 

\bibitem{jeon-kang-kim}
D Jeon, SY Kang, CH Kim, \textit{On values of weakly holomorphic modular functions at divisors of meromorphic modular forms}, J. Number Theory 239, 183-206.

\bibitem{hamieh-murty}
Alia Hamieh and M. Ram Murty, {A note on $q$-analogues of Dirichlet L-functions}, International Journal of Number Theory, Vol. 12 (2016) 765-773.

\bibitem{lang}
Serge Lang, \textit{Elliptic Functions}, Addision-Wesley Publishing Company, Inc.


\bibitem{murty-dewar-graves}
M. R. Murty, M. Dewar and H. Graves, \textit{Problems in 
the Theory of 
Modular
Forms}, Hindustan Book Agency.

\bibitem{r-m}
M. Ram Murty and P. Rath, \textit{Transcendental Numbers} (Springer-Verlag, $2014$).

\bibitem{r-m-g}
S. Gun, M. Ram Murty, P. Rath, \textit{Algebraic independence of values of
modular forms}. Int. J. Number Theory 7(4), 1065–1074 (2011).


\bibitem{nesterenko}
Y. Nesterenko, \textit{Modular functions and transcendence questions}, Mat. Sb. $187$ $(9)$ $(1996)$ $65–96$ (in Russian); translation in:
Sb. Math. $187$ $(9)$ $(1996)$ $1319–1348$.


\bibitem{schneider}
T. Schneider, \textit{Arithmetische Untersuchungen elliptischer Integrale}, Math. Ann.$ 113$ $(1973)$ $1–13$.

\bibitem{shimura}
 G. Shimura, \textit{Introduction to the arithmetic theory of automorphic functions}, Kano Memorial Lectures $1$.


\end{thebibliography}

\end{document}


\documentclass[11pt, reqno]{amsart}
\usepackage[utf8]{inputenc}
\usepackage{amsmath}
\usepackage{amsthm}
\usepackage{amsfonts}
\usepackage{amssymb}
\usepackage{mathrsfs, mathtools, mathdots}
\usepackage{diagbox}
\DeclarePairedDelimiter \floor {\lfloor} {\rfloor}
\DeclareMathOperator{\Tr}{Tr}

\usepackage[top=35mm,bottom=35mm,left=35mm,right=35mm]{geometry}
\newtheorem{theorem}{Theorem}[section]
\newtheorem{conjecture}{Conjecture}
\newtheorem{corollary}{Corollary}
\newtheorem{problem}{Problem}
\newtheorem*{theorem*}{Theorem}
\newtheorem*{remark*}{Remark}
\newtheorem*{problem*}{Problem}
\newtheorem*{conjecture*}{Conjecture}
\newtheorem{lemma}[theorem]{Lemma}
\newtheorem{prop}[theorem]{Proposition}


\begin{document}

\title[Transcendence]{Modular functions and Transcendence}

\author{[Author]}



\subjclass[2010]{-------}

\keywords{Algebraic independence, Transcendence degree, modular forms}
\thanks{Research ----}
\begin{abstract}
We study values of meromorphic modular forms  and modular functions of arbitrary level with algebraic Fourier coefficients. Then we investigate the transcendence degree of the field over $\overline{\mathbb{Q}}$ obtained by adjoining $e^{2\pi i\tau}$ and the values of a modular form, modular function and quasi-modular form at a point $\tau$ in the upper half-plane. As a consequence, we obtain a generalization of Nesterenko's Theorem in higher level. 
\end{abstract}

\maketitle

\section{\bf Introduction}
\bigskip

In this note, we investigate the arithmetic nature of values of special functions arising from the modular world. The values of these functions on the upper half-plane $\mathbb{H}$ is a rich source of transcendental numbers. An earlier result regarding transcendence of values of 
 normalized Eisenstein series $E_4$ and $E_6$ was given by D. Bertrand \cite{bertrand}. He proved that for $\tau\in\mathbb{H}$, at least one of the numbers $E_4(\tau), E_6(\tau)$ is transcendental. Later, Chudnovsky \cite{chudnovsky} showed that if $\tau\in\mathbb{H}$, then at least two of the numbers $E_2(\tau), E_4(\tau)$ and $E_6(\tau)$ are algebraically independent over $\overline{\mathbb{Q}}$. In $1996$, Nesterenko made a significant advancement in this direction by proving the following result \cite{nesterenko}. 
\begin{theorem}[Nesterenko] \label{nest}
If $\tau\in \mathbb{H}$, then at least three of the numbers
$$
e^{2 \pi i \tau},\quad E_2(\tau),\quad E_4(\tau),\quad E_6(\tau)
$$
are algebraically independent over $\overline{\mathbb{Q}}$.
\end{theorem}
 An important consequence of Theorem \ref{nest} is the following result which was conjectured earlier by Mahler and proved in a joint work of Barr\'e-sirieix, Diaz, Gramain and Philibert \cite{b-d-g-p}. 
\begin{theorem}[Barr\'e-sirieix et al]\label{mahl}
For any $\tau\in \mathbb{H}$, at least one of the two numbers $e^{2\pi i\tau}$ and $j(\tau)$ is transcendental.
\end{theorem}
It is known from the theory of complex multiplication that if $\tau\in\mathbb{H}$ is a CM point (that is, an algebraic number which generates an imaginary quadratic field), then $j(\tau)$ is an algebraic number. For non-CM algebraic points in $\mathbb{H}$, Schneider \cite{schneider} proved the following result.
\begin{theorem}[Schneider]\label{schneider}
    If $\tau\in \mathbb{H}$ is algebraic but not a CM point, then $j(\tau)$ is transcendental.
\end{theorem}

 In 2010, a study on algebraic independencece of values of modular forms was done by M. R. Murty, P. Rath and S. Gun \cite{r-m-g}. 
In another paper \cite{hamieh-murty}, Hamieh and Murty proved the following results for modular forms of arbitrary level. 
\begin{theorem}\label{algem''}
  Let $f$ be a non-zero modular form of weight $k$ and level $N$ with algebraic Fourier coefficients.\begin{enumerate}
      \item If $\tau\in\mathbb{
      H
      }$ such that $e^{2\pi i\tau}$ is algebraic, then $f(\tau)$ is transcendental.
      \item If $\tau\in\mathbb{H}$ such that  such that $j(\tau)\in \overline{\mathbb{Q}}$, then there exists a transcendental number $\omega_\tau$ which depends only on $\tau$ and is $\overline{\mathbb{Q}}$-linearly independent with $\pi$ such that $\left(\frac{\pi}{\omega_\tau}\right)^k f(\tau)\in \overline{\mathbb{Q}}$. Therefore, $f(\tau)$ is either zero or transcendental.
  \end{enumerate}  
\end{theorem}
We note that Theorem \ref{algem''} immediately extends for weakly holomorphic modular forms, but it requires a little work for meromorphic modular forms. We will restate Theorem \ref{algem''} and give a  proof for meromorphic modular forms (see Theorem \ref{algem'}).\\

Investigation on transcendence on special values of quasi-modular forms for $SL_2(\mathbb{Z})$ also has been carried out by C. Y. Chang. In \cite{quasi}, Chang proved the following theorem.
 \begin{theorem}\label{changt}
 Let $f$ be a quasi-modular form of non-zero weight for $SL_2(\mathbb{Z})$ with algebraic Fourier coefficients. If $\tau\in\mathbb{H}$ such that $j(\tau)\in\overline{\mathbb{Q}}$ and $f(\tau)\neq 0$, then $f(\tau)$ is transcendental.    
 \end{theorem}
  In (\cite{hamieh-murty}, Corollary $3.3$), Hamieh and Murty combined Theorem \ref{algem''} and Theorem \ref{changt} with Theorem \ref{mahl} to deduce the following result.
\begin{theorem}\label{mg}
    Let $f$ be a non-zero modular form of level $N$ and weight $k$ with algebraic Fourier coefficients. If $\tau\in\mathbb{H}$ such that $f(\tau)\neq 0$, then at least two of 
    $$
e^{2\pi i\tau},\quad j(\tau),\quad f(\tau)
    $$
    are transcendental.
\end{theorem}
More generally, in Theorem \ref{deg.}, we will prove that at least two of these numbers are algebraically independent over $\overline{\mathbb{Q}}$, where $j$ can be replaced by arbitrary modular function with algebraic Fourier coefficients and $f$ can be taken as a meromorphic modular form.\\

In a recent paper (\cite{jeon-kang-kim}, Theorem $2.4$) Jeon, Kang and Kim studied values of a particular family of modular functions at divisors of meromorphic modular forms. As a consequence (\cite{jeon-kang-kim}, Corollary $2.5$), they proved that any zero or pole of a meromorphic modular with algebraic Fourier coefficients is either transcendental or imaginary quadratic. Indeed, these results follow from a more general correspondence between the transcendental nature of values of the $j$-function and arbitrary modular functions. We give a proof of this correspondence in Theorem \ref{suma}. As immediate consequence, one can generalize the classical results of Schneider (Theorem \ref{schneider}) and Barr\'e-sirieix et al (Theorem \ref{mahl}) for arbitrary modular functions with algebraic Fourier coefficients.\\
 
Here, we first study values of modular functions of arbitrary level with algebraic Fourier coefficients. Using these results we determine transcendental nature of values of meromorphic modular and and quasi-modular forms of arbitrary level. Finally, we investigate the transcendence degree of the field generated by values of these functions.\\

Let $f$ and $g$ denote a non-constant \emph{meromorphic modular form} and modular function respectively with algebraic Fourier coefficients. In particular, one can take $f$ as a holomorphic or weakly holomorphic modular form. Then we have the following results.
\begin{theorem}\label{zero-pole}
    If $\tau\in\mathbb{H}$ is a zero or pole of $g$, then $j(\tau)$ is an algebraic number.
\end{theorem}
For other point in $\mathbb{H}$, we get the following correspondence between values of $j$-function and arbitrary modular functions.
\begin{theorem}\label{suma}
If $\tau\in\mathbb{H}$ such that $\tau$ is not a pole of $g$, then

(i) $ j(\tau)\in\overline{\mathbb{Q}}\Leftrightarrow g(\tau)\in\overline{\mathbb{Q}}$.

(ii) At least one of $e^{2\pi i\tau}$ and $g(\tau)$ is transcendental.
\end{theorem}
Then we study values of arbitrary modular functions at any zero or pole of meromorphic modular forms with algebraic Fourier coefficients, which generalizes the result of Jeon, Kang and Kim (\cite{jeon-kang-kim}, Theorem $2.4$) as follows: 
\begin{theorem}\label{zero-pole-value}
    Let $f$ and $g$ denote a meromorphic modular form and function respectively as above. If $\tau\in\mathbb{H}$ is a zero or pole of $f$, then $g(\tau)$ is algebraic.
\end{theorem}
\begin{theorem}\label{deg.}
    If $\tau\in\mathbb{H}$ such that $f(\tau)\neq 0$ and $\tau$ is not a pole of $f$ and $g$, then the transcendence degree of the field $\overline{\mathbb{Q}}\left(e^{2\pi i\tau},\, f(\tau),\, g(\tau) \right)$ over $\overline{\mathbb{Q}}$ is at least two, that is
    $$
trdeg \,\overline{\mathbb{Q}}\left(e^{2\pi i\tau},\, f(\tau),\, g(\tau) \right) \geq 2.
$$  
\end{theorem}
We note that Theorem \ref{deg.} generalizes  Theorem \ref{mg}. Moreover, the above theorem extends to the following. 
\begin{theorem}\label{deg..}
 Let $\widetilde{f}$ be a quasi-modular form of positive depth and arbitrary level with algebraic Fourier coefficients. If $\tau\in\mathbb{H}$ such that $f(\tau)\neq 0$, $\widetilde{f}(\tau)\neq 0$ and $\tau$ is not a pole of $f$ and $g$, then 
    $$
trdeg\,\overline{\mathbb{Q}}\left( e^{2\pi i\tau},\, f(\tau),\, g(\tau),\, \widetilde f(\tau)\right)\geq 3. 
    $$
\end{theorem}
\begin{remark*}
 The proof of Theorem \ref{deg..} mainly uses the result on transcendence of values of modular functions (Theorem \ref{suma}) and Nesterenko's Theorem \ref{nest}. Again, if we choose 
$$
f=E_4, \, g=j,\, \widetilde f= E_6+E_4\, E_2,
$$
the statement of Theorem \ref{deg..} implies Nesterenko's theorem. Therefore, Theorem \ref{deg..} is essentially the Nesterenko's theorem in more general setup.   
\end{remark*}








\bigskip

\section{\bf Notations and Preliminary Results}
\subsection{Modular and quasi-modular forms}\quad\bigskip

 Recall that for each $N\in \mathbb{N}$, $$
\Gamma(N)=\left\{  \left( 
\begin{array}{cc}
  a & b \\
  c & d
\end{array}
\right)\in SL_2(\mathbb{Z}):\left( 
\begin{array}{cc}
  a & b \\
  c & d
\end{array}
\right)\equiv \left( 
\begin{array}{cc}
  1 & 0 \\
  0 & 1
\end{array}
\right) (\operatorname{mod}N) \right\}
$$ is called the \emph{principal congruence subgroup of level N}.
 Note that $\Gamma(1)=SL_2(\mathbb{Z})$. A subgroup $\Gamma$ of $SL_2(\mathbb{Z})$ is said to be congruence subgroup if there exists $N\in \mathbb{N}$ such that $\Gamma(N)\subseteq \Gamma$, and the smallest such $N$ is  called the level of $\Gamma$.\\

A \emph{holomorphic modular form} of integer weight $k(\geq 0)$ with respect to a congruence subgroup $\Gamma$ is a holomorphic function on the upper half-plane $\mathbb{H}$ which satisfies 
\begin{equation} \label{hmod}
f\left(\frac{a\,\tau+b}{c\,\tau+d}\right)=(c\,\tau+d)^kf(\tau)\quad\text{for all }
\gamma=\left(\begin{smallmatrix}
   a  & b \\
   c  & d
\end{smallmatrix} \right)\in \Gamma,\end{equation}
and also is \emph{holomorphic} at all cusps. The function $f$ is said to be \emph{weakly holomorphic modular form} if it is \emph{holomorphic} on $\mathbb{H}$ but \emph{meromorphic} at cusps. More generally, $f$ is called \emph{meromorphic modular form} if it is meromorphic on $\mathbb{H}$ and also at cusps. 
We denote the space of all holomorphic, weakly holomorphic and meromorphic modular forms by $M_{k,\overline{\mathbb{Q}}}(\Gamma)$, $M_{k,\overline{\mathbb{Q}}}^w(\Gamma)$ and $M_{k,\overline{\mathbb{Q}}}^m(\Gamma)$ respectively, where $M_{k,\overline{\mathbb{Q}}}(\Gamma)\subset M_{k,\overline{\mathbb{Q}}}^w(\Gamma)\subset M_{k,\overline{\mathbb{Q}}}^m(\Gamma)$. \\

For even integer $k\geq 2$, define the nomalized Eisenstein series of weight $k$ for $S L_2(\mathbb{Z})$ by
$$
E_{k}(\tau)=1-\frac{2 k}{B_{k}} \sum_{n=1}^{\infty} \sigma_{k-1}(n) q^{n}, \text{   where }q=e^{2\pi i\tau} \text{ and } \sigma_s(n)=\sum_{\substack{d|n\\ d>0}}d^s.
$$ 
Here, $B_k$ is called the $k$-th Bernoulli number. One can show that for $k\geq 4$, the function $E_k\in M_{k,\overline{\mathbb{Q}}}(S L_2(\mathbb{Z}))$ and $\Delta(\tau):=\frac{E_4(\tau)^3-E_6(\tau)^2}{1728}\in M_{12,\overline{\mathbb{Q}}}(S L_2(\mathbb{Z}))$. But $E_2(\tau)$ is not a modular form (see \cite{murty-dewar-graves}, Chapter $5$) because it fails to satisfy \eqref{hmod} for the matrix $\left(\begin{smallmatrix}
   0  & -1 \\
   1  & 0
\end{smallmatrix} \right)\in SL_2(\mathbb{Z})$.\\

There are several equivalent  definitions of \emph{quasi-modular forms}, but we use the following characterization (\cite{1-2-3}, Proposition $20$).
\begin{theorem}\label{quasistructure}
Every quasi-modular form for a congruence subgroup $\Gamma$ is a polynomial in $E_2$ with modular coefficients. More precisely, if $f$ is a quasi-modular form of weight $k$ and depth $p$ on $\Gamma$, then $f$ can be uniquely written as $f=\sum_{r=0}^p g_r\, E_2^r$, where $g_r\in M_{k-2 r}(\Gamma)$ for all $0\leq r\leq p$ and $g_p\neq 0$. 
\end{theorem}

A quasi-modular form $f$ is said to have algebraic Fourier coefficients if all modular coefficient $g_r$'s of $f$ have algebraic Fourier coefficients. 
\bigskip

\subsection{Weierstrass $\wp$-function and zeta-function}\quad\bigskip

The Weierstrass $\wp$-function associated with a two dimensional
lattice $L=\omega_1 \mathbb{Z} \,  \oplus \, \omega_2 \mathbb{Z}$ is defined as follows
\begin{equation}\label{wpfunction}
\wp (z)=\frac{1}{z^2}+\sum_{\substack{\omega\in L\\
\omega\neq 0}}\left(\frac{1}{(z-\omega)^2}-\frac{1}{\omega^2} \right), \quad \text{for }z\in \mathbb{C}\setminus L. 
\end{equation}
This is a meromorphic function on $\mathbb{C}$ with double poles at each lattice points. Moreover, it satisfies the following differential equation
\begin{equation} \label{dell}
 \wp'(z)^2=4\wp(z)^3-g_2(L)\wp(z)-g_3(L),   
\end{equation}
where $$g_2(L)=60\, \sum_{\substack{\omega \in L\\ \omega\neq 0}} \frac{1}{\omega^{4}} \text{ and } g_3(L)=140\, \sum_{\substack{\omega \in L\\ \omega\neq 0}} \frac{1}{\omega^{6}}.$$
In particular, if we take $L_\tau=\tau\,\mathbb{Z}\oplus \mathbb{Z}$ for $\tau\in\mathbb{H}$, then $$g_2(L_\tau)=\frac{4\pi^4}{3}E_4(\tau) \text{ and } g_3(L_\tau)=\frac{8\pi^6}{27}E_6(\tau).$$

Define the discriminant function by$$\Delta_0(L):=g_2(L)^3-27g_3(L)^2\neq 0, \text{ for any two dimensional lattice }L.$$
 In particular, we have $\Delta_0(L_\tau)=(2\pi)^{12}\,\Delta(\tau)$ for all $\tau\in\mathbb{H}$.\\

 The Weierstrass zeta-function associated with $L=\omega_1\,\mathbb{Z}\oplus \omega_2\,\mathbb{Z}$ is defined by
$$
Z_{L}(z)=\frac{1}{z}+\sum_{\substack{\omega\in L\\ \omega\neq 0}}\left(\frac{1}{z-\omega}+\frac{1}{\omega}+\frac{z}{\omega^2} \right)\quad \text{for } z\in \mathbb{C}\setminus L.
$$
 Note that $Z_{L}'(z)=-\wp_L(z)$ is a periodic function with each point of $L$ as a period. Hence, the functions $$Z_{L}(z+\omega_1)-Z_{L}(z) \text{  and  } Z_{L}(z+\omega_2)-Z_{L}(z)$$ are constants. These constants are denoted by $\eta_1(L), \eta_2(L)$ respectively and called \emph{quasi-periods}. Moreover, it is known that \begin{equation}\label{quasi-period}
\eta_2(L_\tau)=G_2(\tau)=\frac{\pi^2}{3}\, E_2(\tau) \text{ 
  for all }\tau\in\mathbb{H}. 
\end{equation}
The reader is referred to (\cite{lang}, chapter $18.3$) for proof of these results.\\

The Weierstrass $\wp$-function satisfies the following addition formula (\cite{r-m}, chapter $11$).
 \begin{theorem} \label{addf}
 For $z_1, z_2\in \mathbb{C}$ such that $z_1\pm z_2\notin L$ we have
$$
\wp\left(z_1+z_2\right)=-\wp\left(z_1\right)-\wp\left(z_2\right)+\frac{1}{4}\left(\frac{\wp^{\prime}\left(z_1\right)-\wp^{\prime}\left(z_2\right)}{\wp\left(z_1\right)-\wp\left(z_2\right)}\right)^2 .
$$
When $z_1=z_2=z$, we have
$$
\wp(2 z)=-2 \wp(z)+\frac{1}{4}\left(\frac{\wp^{\prime \prime}(z)}{\wp^{\prime}(z)}\right)^2 .
$$
 \end{theorem}
By using this addition formula together with a remarkable theorem known as Schneider–Lang theorem, one can prove the following result (see \cite{r-m}, Theorem $11.2$).
\begin{theorem}\label{tranw}
    Let $L=\omega_1\,\mathbb{Z}\oplus\omega_2\,\mathbb{Z}$ be a lattice and both $g_2(L)$, $g_3(L)$ are algebraic. If $\alpha$ is an algebraic number such that $\alpha\notin L$, then $\wp(\alpha)$ is transcendental.
\end{theorem}
We have the following corollary of Theorem \ref{tranw}.
\begin{corollary}\label{tranp}
    Let the Weierstrass $\wp$ function be associated with $L=\omega_1\,\mathbb{Z}\oplus\omega_2\,\mathbb{Z}$ such that both $g_2(L)$ and $g_3(L)$ are algebraic. Then, for any natural number $n$, the numbers $\wp(\frac{\omega_1}{n})$ and $\wp(\frac{\omega_2}{n})$ are algebraic. Moreover, any non-zero period of $\wp(z)$ is necessarily transcendental.
\end{corollary}









































\subsection{Modular function}
\quad\bigskip

A meromorphic function $f$ on $\mathbb{H}$ is said to be a modular function if it satisfies
\begin{equation}
f\left(\frac{a\,\tau+b}{c\,\tau+d}\right)=f(\tau)\quad\text{for all }
\gamma=\left(\begin{smallmatrix}
   a  & b \\
   c  & d
\end{smallmatrix} \right)\in \Gamma\end{equation}
and also is meromorphic at cusps. In particular, a modular function on the congruence subgroup $\Gamma(N)$ is said to a be modular function of level $N$.\\

A canonical example of modular function of level one is given by 
$$
j(\tau):=\frac{E_4^3(\tau)}{\Delta(\tau)} \text{ for all } \tau\in\mathbb{H}.
$$
It has the following Fourier expansion
$$
j(\tau)=\frac{1}{q}+744+196884q+\cdots,$$
where $q=e^{2\pi i\tau}$. This is made precise in the following theorem.
\begin{theorem} \label{level1}
Let $f$ be a modular function of level one and let $$K := \mathbb{Q}(\text{Fourier coefficients of }f)\subset\mathbb{C}.$$ Suppose that the poles of $f$ in $\mathcal{F}$ are $\tau_1,\tau_2,\cdots,\tau_m$, where $\mathcal{F}$ denotes the standard fundamental domain for the action of $S L_2(\mathbb{Z})$ on $\mathbb{H}$. Then $f\in K(j(\tau_1),\cdots,j(\tau_m))(j)$ i.e., $f$ is a rational function in $j$ over $K(j(\tau_1),\cdots,j(\tau_m))$.
\end{theorem}
\begin{proof}
First, we consider the function 
$$
g(\tau):=\left(\prod_{\tau_0\in\mathcal{F}} \left(j(\tau)-j(\tau_0)\right)^{\operatorname{ord}_{\tau_0}f}\right)f(\tau),
$$
which is holomorphic on $\mathbb{H}$. Since $f$ is meromorphic at $i\infty$, there exists $c>0$ such that the region $\{z\in\mathcal{F} : \operatorname{Im}z>c\}\cup \{i\infty\}$ contains no pole of $f$ except possibly $i\infty$. Again, $f$ is meromorphic on $\mathbb{H}$ and the region $\{z\in \mathcal{F}: \operatorname{Im}z\leq c\}$ is compact, and so it contains only a finite number of poles of $f$. Thus, $f$ can have only finitely many poles in $\mathcal{F}$ and hence above product is finite.\\

Suppose $g$ has pole of order $M$ at $i\infty$. Then the Fourier expansion of
$g(\tau)$ at $i\infty$ has the form
$$
g(\tau)=\sum_{n=-M}^\infty c_n q^n, \, \text{ where } c_{-M}\neq 0.
$$
Note that the modular function $g(\tau)-c_{-M}\,j(\tau)^M$ is holomorphic on $\mathbb{H}$ and its $q$-series starts with at most a polar term of order $M-1$. Repeating this process, we can subtract a polynomial of $j(\tau)$ to get a holomorphic modular function that vanishes at $i\infty$, and hence identically zero. Thus, $g(\tau)$ is a polynomial of $j(\tau)$ over $K$ and so $f(\tau)$ is a rational function of $j(\tau)$ over $K(j(\tau_1),\cdots,j(\tau_m))$.
\end{proof}
The above proof proves the following :
\begin{corollary}\label{rationalfunctions}
If $f(\tau)$ is a modular function of level one which is holomorphic on $\mathbb{H}$ with the following $q$-expansion at $i\infty$
$$
f(\tau)=\sum_{n=-M}^\infty c_n q^n, \ c_{-M}\neq 0
$$
with $K=\mathbb{Q}(\text{Fourier coefficients of }f)$, then $f(\tau)$ is a polynomial in $j(\tau)$ of degree $M$ with coefficients in $K$.
\end{corollary}
\begin{theorem} \label{mdft}
Let $f$ be a modular function with respect to a congruence subgroup $\Gamma$ and
$$
K:=\mathbb{Q}\left(\text{Values of }j\text{ at poles of }f\text{ in }\mathcal{F}, \text{ Fourier coefficients of }f\text{ at all cusps}\right).
$$
Then there exists a monic polynomial $P(X)\in K(j)[X]$ such that $P(f)=0$. 
\end{theorem}
\begin{proof}
Let $[S L_2(\mathbb{Z}):\Gamma]=r$ and $\{ \gamma_1=I,\gamma_2,\cdots,\gamma_r\}$ be a complete set of right coset representatives so that $$
S L_2(\mathbb{Z})=\bigsqcup_{i=1}^r\Gamma \gamma_i.
$$
For all $1\leq i \leq r$, define the functions $$
f_i(\tau):=f(\gamma_i\tau) \quad \text{for } \tau \in \mathbb{H}.
$$ Clearly, each $f_i$ is independent of the choice of coset representatives as $f$ is $\Gamma$-invariant. Again, for any $\gamma\in S L_2(\mathbb{Z})$, we have $f_i(\gamma \tau)=f(\gamma_i\gamma\tau)=f_j(\tau)$ if $\gamma_i\gamma\in \Gamma\gamma_j$, i.e., the set $\{f_1,f_2,...,f_r\}$ gets permuted under the action of $S L_2(\mathbb{Z}).$ This observation implies that any elementary symmetric polynomial of $f_1,f_2,\cdots,f_r$ is a modular function for $S L_2(\mathbb{Z})$, and hence they are in $K(j)$ by Theorem \ref{level1}. Note that the polynomial $$
P(X)=\prod_{i=1}^r(X-f_i)$$is satisfied by $f$ as $f=f_1$. Moreover, the coefficients of $P$ are elementary symmetric polynomials in $f_1,f_2,\cdots,f_r$, and hence $P(X)$ is the required polynomial.
\end{proof}
\begin{corollary}\label{useful}
 If the modular function $f$ is holomorphic on $\mathbb{H}$ and $$K := \mathbb{Q}(\text{Fourier coefficients of }f\text{ at all cusps}),$$
 then the monic polynomial in Theorem \ref{mdft} lies in $K[j][X]$.
\end{corollary}
\begin{proof}
   This is clear from above proof and Corollary \ref{rationalfunctions}, since coefficients of $P[X]$ constructed above are modular functions of level one, holomorphic on $\mathbb{H}$ with Fourier coefficients in $K$.  
\end{proof}
\bigskip

For $N>1$ and $a=(a_1,a_2)\in N^{-1}\mathbb{Z}^2\setminus\mathbb{Z}^2$, consider the function
$$
f_{a}(\tau):=\frac{g_2(L_\tau)\,g_3(L_\tau)}{\Delta_0(L_\tau)}\,\wp_{L_\tau} (a_1\tau+a_2),
$$
which is holomorphic on $\mathbb{H}$. Moreover, It is known that each such $f_a$ is a modular function of level $N$ with algebraic Fourier coefficients with respect to $e^{\frac{2\pi i\tau}{N}}$. For an elaborate discussion on these functions, the reader is referred to (\cite{shimura}, Chapter 6).\\

Let $F_{N,\overline{\mathbb{Q}}}$ denote the field of modular functions of level $N$ with algebraic Fourier coefficients with respect to $e^{\frac{2\pi i\tau}{N}}$, and $S:=\{\left(\frac{r}{N},\frac{s}{N}\right): 0\leq r,s\leq N-1 \text{ and }(r,s)\neq(0,0)\}$. Thus, $f_a\in F_{N,\overline{\mathbb{Q}}}$ for all $a\in S$.
\begin{theorem}\label{mdft..}
$F_{N,\overline{\mathbb{Q}}}$ is a finite Galois extension of $\overline{\mathbb{Q}}(j)$ and $F_{N,\overline{\mathbb{Q}}}=\overline{\mathbb{Q}}(j,f_a |a\in S)$.
\end{theorem}
For a proof of Theorem \ref{mdft..}, we refer to (\cite{shimura}, Proposition $6.9.(1)$), where it is proved for modular functions with Fourier coefficients in $\mathbb{Q}\big(e^{\frac{2\pi i}{N}}\big)$. But, one can use the same argument for the field $\overline{\mathbb{Q}}$ instead of $\mathbb{Q}\big(e^{\frac{2\pi i}{N}}\big)$.
\bigskip

\section{\bf Proof of Main Results}
\bigskip

\subsection{Values of modular functions when $j(\tau)$ is transcendental}
\begin{lemma} \label{mdfz}
    Let $g$ be a non-constant modular function of level $N(\geq 1)$ with algebraic Fourier coefficients, i.e., $g\in  F_{N,\overline{\mathbb{Q}}}$. If $\tau\in \mathbb{H}$ such that $j(\tau)\notin\overline{\mathbb{Q}}$, then $g(\tau)$ is a non-zero number.
\end{lemma}
\begin{proof}
By Theorem \ref{mdft..}, $g$ has the minimal polynomial in $\overline{\mathbb{Q}}(j)[X]$, and we can clear denominators of it's coefficients to get an irreducible polynomial 
$$
P(X)=\sum_{r=0}^m c_r(j)\,X^r \in \overline{\mathbb{Q}}[j][X], \quad m=\deg P(X)
$$
such that
$$
P(g)=0.
$$
Here $$c_r(j)=\sum_{s=0}^{d_r}c_{r s}\,j^s\in \overline{\mathbb{Q}}[j]\quad \forall\ 0\leq r\leq m, \quad d_r=\deg c_r(j).$$
Also, $c_0(j)$ is not identically zero polynomial because $P(X)$ is irreducible over $\overline{\mathbb{Q}}[j]$. Thus, we get
\begin{equation}\label{funzero}
\sum_{r=0}^m\sum_{s=0}^{d_r}c_{r s}\,j(\tau)^s\, g(\tau)^r=0\quad \forall\, \tau\in\mathbb{H}.
\end{equation}
Now we consider $\tau\in\mathbb{H}$ such that $j(\tau)\notin\overline{\mathbb{Q}}$. We claim that $g(\tau)\neq 0$. If not, from \eqref{funzero} we get $c_0(j(\tau))=0$ and hence, $j(\tau)\in\overline{\mathbb{Q}}$, which contradicts our assumption. This completes the proof that $g(\tau)$ is non-zero. 
\end{proof}
\begin{proof}[\bf Proof of Theorem \ref{zero-pole}]
If $g(\tau)=0$, then it follows from the proof of Lemma \ref{mdfz} that $j(\tau)\in \overline{\mathbb{Q}}$. If $\tau$ is a pole of $g$, we consider the modular function
$$
h(z):=\frac{1}{g(z)}\in F_{N,\overline{\mathbb{Q}}},
$$
which vanishes at $z=\tau$, and so $j(\tau)\in\overline{\mathbb{Q}}$ by Lemma \ref{mdfz}.    
\end{proof}
The following theorem shows that the non-zero number $g(\tau)$ is indeed transcendental.
\begin{theorem}\label{tranmf}
 Let $g\in F_{N,\overline{\mathbb{Q}}}$ and $\tau\in \mathbb{H}$ such that it is not a pole of $g$. If $j(\tau)\notin\overline{\mathbb{Q}}$, then $g(\tau)$ is transcendental.
\end{theorem}
\begin{proof}
If possible, let $g(\tau)\in\overline{\mathbb{Q}}$. Then the following modular function $$h(z):=g(z)-g(\tau)\quad \forall z\in\mathbb{H}$$ belongs to $ F_{N,\overline{\mathbb{Q}}}$ but vanishes at $\tau$, which contradicts Lemma \ref{mdfz}. Therefore, $g(\tau)$ is a transcendental number.
\end{proof}
\subsection{Values of modular functions when $j(\tau)\in\overline{\mathbb{Q}}$}
\quad \bigskip

From Theorem \ref{mdft..}, we have $$ F_{N,\overline{\mathbb{Q}}}=  \overline{\mathbb{Q}}\left(\{j,f_a |a\in S\}\right), \text{ where }S=\left\{\left(\frac{r}{N},\frac{s}{N}\right): 0\leq r, s\leq N-1 \text{ and }(r,s)\neq (0,0)\right\}.$$
Hence, for any $g\in F_{N,\overline{\mathbb{Q}}}$, the algebraic nature of $g(\tau)$ is completely determined by the numbers in the set $\{f_a(\tau):a\in S\}$.\\

Note that for $N=1$, the set $S$ is empty. This shows that for any $g\in F_{1,\overline{\mathbb{Q}}}$, the number $g(\tau)\in\overline{\mathbb{Q}}$ when $j(\tau)\in\overline{\mathbb{Q}}$. For $N>1$, we have to study the numbers $f_a(\tau)$ for all $a\in S$. We start by proving the following important lemma.
\begin{lemma}\label{omegatau}
 If $\tau\in \mathbb{H}$ such that $j(\tau)$ is algebraic, then there exists a unique (up to an algebraic multiple) transcendental number $\omega_\tau$ for which $g_2(\omega_\tau L_\tau)$ and $g_3(\omega_\tau L_\tau)$ are both algebraic numbers.   
\end{lemma}
\begin{proof}
Note that for any $\omega_\tau\in\mathbb{C}^\times $,
\begin{equation}\label{omegatau.}
    j(\tau)=\frac{E_4(\tau)^3}{\Delta(\tau)}=\frac{1728\,  g_2(L_\tau)^3}{\Delta_0(L_\tau)}=\frac{1728\, g_2\left(\omega_\tau L_\tau\right)^3}{\Delta_0\left(\omega_\tau L_\tau\right)},
\end{equation}
where $L_\tau=\tau\,\mathbb{Z}\oplus\mathbb{Z}$ and $\omega_\tau L_\tau=\tau\,\omega_\tau\,\mathbb{Z}\oplus\omega_\tau\mathbb{Z}$. Here the last equality follows from the fact that \begin{equation}\label{homogenity}
g_2(\omega_\tau L_\tau)=\omega_\tau^{-4}g_2(L_\tau).
\end{equation} If we choose $\omega_\tau$ such that $\omega_\tau^4=g_2(L_\tau)$, then $g_2(\omega_\tau L_\tau)=1$. Given that $j(\tau)$ is algebraic, and so \eqref{omegatau.} shows that $\Delta_0(\omega_\tau L_\tau)$ is also algebraic. Thus, the numbers $g_2(\omega_\tau L_\tau)$ and $g_3(\omega_\tau L_\tau)$ both are algebraic. Moreover,  $\omega_\tau$ is a period of $\omega_\tau L_\tau$ and both $g_2(\omega_\tau L_\tau)$ and $g_3(\omega_\tau L_\tau)$ are algebraic. By Corollary \ref{tranp} we conclude that $\omega_\tau$ is transcendental.\\

To prove uniqueness, consider an arbitrary $\omega_\tau'\in\mathbb{C}^\times$ such that both $g_2(\omega_\tau' L_\tau)$ and $g_3(\omega_\tau' L_\tau)$ are algebraic. From \eqref{homogenity} we have 
$$
 \quad g_2(\omega_\tau' L_\tau)= (\omega_\tau')^{-4}\,g_2(L_\tau)=\alpha, (\text{say}).
$$
Therefore, we get $g_2(L_\tau)=\alpha\cdot (\omega_\tau')^{4}$ and hence $\omega_\tau^4=\alpha\cdot (\omega_\tau')^{4}.$
This completes the proof that $\omega_\tau$ is unique up to an algebraic multiple.
\end{proof}

Recall that
$$
f_a(\tau):=\frac{g_2(L_\tau)\,g_3(L_\tau)}{\Delta_0(L_\tau)}\,\wp_{L_\tau}(a_1\tau+a_2),
$$ 
where $L_\tau=\tau\,\mathbb{Z}\oplus\mathbb{Z}$ and $a=(a_1,a_2)\in S$. We can rewrite it as follows:
$$
f_a(\tau)=\frac{g_2(\omega_\tau L_\tau)\, g_3(\omega_\tau L_\tau)}{\Delta_0(\omega_\tau L_\tau)}\,\wp_{\omega_\tau L_\tau}(a_1\omega_\tau\tau+a_2\omega_\tau).
$$
\begin{lemma}\label{algfa}
    If $\tau\in\mathbb{H}$ such that $j(\tau)\in\overline{\mathbb{Q}}$, then $f_a(\tau)\in\overline{\mathbb{Q}}$ for all $a\in S$.
\end{lemma}
\begin{proof}
    Let $a=(\frac{r}{N},\frac{s}{N})\in S$. Then we have
$$
f_a(\tau)=\frac{g_2(\omega_\tau L_\tau)\, g_3(\omega_\tau L_\tau)}{\Delta_0(\omega_\tau L_\tau)}\,\wp_{\omega_\tau L_\tau}\left(\frac{r\,\omega_\tau\,\tau}{N}+\frac{s\,\omega_\tau}{N}\right),
$$
where $\omega_\tau$ is chosen as in Lemma \ref{omegatau} so that $g_2(\omega_\tau L_\tau),\, g_3(\omega_\tau L_\tau)\in\overline{\mathbb{Q}}$.\\

Recall that both $\wp_{\omega_\tau L_\tau}\left(\frac{\omega_\tau\, \tau}{N} \right)$ and $\wp_{\omega_\tau L_\tau}\left(\frac{\omega_\tau}{N} \right)$  are algebraic by Corollary \ref{tranp}. Moreover, by the addition formula for Weierstrass $\wp$-function in Theorem \ref{addf}, we get that $$\wp_{\omega_\tau L_\tau}\left(\frac{r\,\omega_\tau\, \tau}{N}+\frac{s\,\omega_\tau}{N}\right)\in\overline{\mathbb{Q}}$$ for all $a=\left(\frac{r}{N},\frac{s}{N}\right)\in S$. Thus, the number $f_a(\tau)$ is algebraic for all $a\in S$.
\end{proof}
\begin{theorem}\label{algemf}
Let $g\in F_{N,\overline{\mathbb{{Q}}}}$ and $\tau\in\mathbb{H}$ such that $j(\tau)\in \overline{\mathbb{Q}}$. If $\tau$ is not a pole of $g$, then $g(\tau)$ is algebraic.
\end{theorem}
\begin{proof}
 From Lemma \ref{algfa}, we have   $f_a(\tau)\in\overline{\mathbb{Q}}$ for all $a\in S$. Then we use the following fact $$g\in  F_{N,\overline{\mathbb{Q}}}=  \overline{\mathbb{Q}}\left(\{j,f_a |a\in S\}\right)$$  to conclude that $g(\tau)$ is algebraic. 
\end{proof}
\bigskip

\begin{proof}[\bf Proof of Theorem \ref{suma}]
     Combining Theorem \ref{tranmf} with Theorem \ref{algemf} we get (i), and (ii) follows from Theorem \ref{mahl} and (i).
\end{proof}
\begin{proof}[\bf Proof of Theorem \ref{zero-pole-value}]
   We first construct the modular function 
   $$
h(\tau):=\frac{f(\tau)^{12}}{\Delta(\tau)^k}\quad \forall \tau\in\mathbb{H},
   $$
   which belong to $F_{N,\overline{\mathbb{Q}}}$ because Fourier coefficients of $f$ and $\Delta$ are algebraic. Moreover, any point $\tau\in\mathbb{H}$ is a zero (or pole) of $h$ if and only if it is a zero (or pole) of $f$ because $\Delta$ is non-vanishing and holomorphic on $\mathbb{H}$. Thus, if $\tau\in\mathbb{H}$ is a zero or pole of $f$, from Theorem \ref{zero-pole} we get that $j(\tau)\in\overline{\mathbb{Q}}$, and hence $g(\tau)\in\overline{\mathbb{Q}}$ by Theorem \ref{suma}.
\end{proof}
\bigskip

\subsection{Values of modular forms}
\quad\bigskip

Recall that for positive integers $k$ and $N$, we denote $M_{k,N,\overline{\mathbb{Q}}}$, $M_{k,N,\overline{\mathbb{Q}}}^w$ and $M_{k,N,\overline{\mathbb{Q}}}^m$ as the set of holomorphic, weakly holomorphic and meromorphic modular forms respectively with algebraic Fourier coefficients. Here, we restate Theorem \ref{algem''} for meromorphic modular forms and give a proof following the technique from \cite{hamieh-murty}. First, we prove the following lemma:
\begin{lemma}\label{det}
    If $n$ and $m$ are positive integers such that $n>2m$, then the determinant of the matrix
    $$
A=\left[\begin{array}{cccccc}
    {{n}\choose{0}} & {{n-1}\choose{0}} & {{n-2}\choose{0}} & \ldots & {{n-m+1}\choose{0}} & {{n-m}\choose{0}} \\\\
    {{n}\choose{1}} & {{n-1}\choose{1}} & {{n-2}\choose{1}} & \ldots & {{n-m+1}\choose{1}} & {{n-m}\choose{1}}\\\\
    {{n}\choose{2}} & {{n-1}\choose{2}} & {{n-2}\choose{2}} & \ldots & {{n-m+1}\choose{2}} & {{n-m}\choose{2}}\\\\
    \vdots & \vdots & \vdots & \ldots & \vdots & \vdots \\\\
    {{n}\choose{m}} & {{n-1}\choose{m}} & {{n-2}\choose{m}} & \ldots & {{n-m+1}\choose{m}} & {{n-m}\choose{m}}
\end{array} \right]_{(m+1)\times (m+1)}
  $$ 
  is given by $(-1)^{\frac{m(m+1)}{2}}$.
\end{lemma}
\begin{proof}
  To find the determinant of $A$, we apply elementary column operations successively on $A$. By using the identity ${n\choose{k}}={n-1\choose{k-1}}+{n-1\choose{k}}$, we rewrite the first column of $A$ as follows: 
$$C_1=\left[\begin{array}{c}
     {n\choose{0}} \\\\
     {n-1\choose{0}}+{n-1\choose{1}}\\\\
     {n-1\choose{1}}+{n-1\choose{2}}\\\\
     \vdots\\\\
     {n-1\choose{m-1}}+{n-1\choose{m}}
\end{array} \right]_{(m+1)\times 1},$$ and then subtract the second column $C_2$ from $C_1$ (i.e.,  $C_1'=C_1-C_2$) to get
$$C_1'=\left[\begin{array}{c}
     0 \\\\
     {n-1\choose{0}}\\\\
     {n-1\choose{1}}\\\\
     \vdots\\\\
     {n-1\choose{m-1}}
\end{array} \right]_{(m+1)\times 1}.$$
We continue this process for each column $C_i$ and and replace $C_i$ by $C_i-C_{i+1}$ for all $1\leq i\leq m$ to obtain the following matrix  
$$
A'=\left[\begin{array}{cccccc}
    0 & 0 & 0 & \ldots & 0 & {{n-m}\choose{0}} \\\\
    {{n-1}\choose{0}} & {{n-2}\choose{0}} & {{n-3}\choose{0}} & \ldots & {{n-m}\choose{0}} & {{n-m}\choose{1}}\\\\
    {{n-1}\choose{1}} & {{n-2}\choose{1}} & {{n-3}\choose{1}} & \ldots & {{n-m}\choose{1}} & {{n-m}\choose{2}}\\\\
    \vdots & \vdots & \vdots & \ldots & \vdots & \vdots \\\\
    {{n-1}\choose{m-1}} & {{n-2}\choose{m-1}} & {{n-3}\choose{m-1}} & \ldots & {{n-m}\choose{m-1}} & {{n-m}\choose{m}}
\end{array} \right]_{(m+1)\times (m+1)}.
$$
 Again, we start from the first column of $A'$ and repeat the same process to reduce the matrix in the form 
 $$
B=\left[\begin{array}{cccccc}
    0 & 0 & 0 & \ldots & 0 & {{n-m}\choose{0}} \\\\
    0 & 0 & 0 & \ldots & {{n-m}\choose{0}} & *\\\\
    0 & 0 & {{n-m}\choose{0}} & \iddots & * & *\\\\
    \vdots & \iddots & \vdots & \ldots & \vdots & \vdots \\\\
    {{n-m}\choose{0}} & * & * & \ldots & * & *
\end{array} \right]_{(m+1)\times (m+1)}.
 $$
 Since $B$ is obtained from $A$ by elementary column operations, we conclude that $\det(A)=\det(B)=(-1)^{\frac{m(m+1)}{2}}$.
\end{proof}
\begin{theorem}\label{algem'}
   Let $f$ be a non-constant meromorphic modular form with algebraic Fourier coefficients. If $\tau\in\mathbb{H}$ such that $e^{2\pi i\tau}$ is algebraic, then $f(\tau)$ is transcendental. 
\end{theorem}
\begin{proof}
By the same argument as before, we construct the modular function
$$g(\tau):=\frac{f(\tau)^{12}}{\Delta(\tau)^k}\quad \forall\, \tau\in\mathbb{H},$$
which lies in $F_{N,\overline{\mathbb{Q}}}$. From Theorem \ref{mdft..}, we know that $g$ has the minimal polynomial $P[X]\in \overline{\mathbb{Q}}(j)[X]$, and if needed, we can clear denominators of the coefficients of $P[X]$ to get a polynomial (use the same notation)
$$P(X)=\sum_{r=0}^m c_r(j)\,X^r\in \overline{\mathbb{Q}}[j][X],\quad \deg P(X)=m$$
such that
$P(g(\tau))=0$ for any $\tau\in \mathbb{H}$. Here, we have
$$
c_r(j)=\sum_{s=0}^{d_r}c_{r,s}\,j^s\in \overline{\mathbb{Q}}[j]\quad \forall\ 0 \leq r\leq m
$$
and $d_r=\deg c_r(j)$ so that $c_{r,d_r}\neq 0$ for all $0\leq r\leq m$. Moreover, Corollary \ref{useful} shows that if $f$ is holomorphic on $\mathbb{H}$, then the minimal polynomial $P[X]$ of $g$ is monic and belongs to $\overline{\mathbb{Q}}[j][X]$, and hence $c_m(j)=1$.\\

For any $\tau\in\mathbb{H}$, we have $$P\left(\frac{f(\tau)^{12}}{\Delta(\tau)^k}\right)=0,$$
and hence for any positive integer $l$, 
$$
\begin{aligned}
& \Delta(\tau)^l\, P\left(\frac{f(\tau)^{12}}{\Delta(\tau)^k}\right)=0\\
\Rightarrow &\sum_{r=0}^m\sum_{s=0}^{d_r}c_{r,s}\,j(\tau)^s f(\tau)^{12 r}\Delta(\tau)^{l-k r}=0\\
\Rightarrow &\sum_{r=0}^m\sum_{s=0}^{d_r}c_{r,s}\,E_4(\tau)^{3 s}f(\tau)^{12\, r}\Delta(\tau)^{l-k r-s}=0.
\end{aligned}
$$
We choose $l> k m +2\cdot\operatorname{max}\{d_0,d_1,\ldots,d_m\}$ so that the powers of $\Delta(\tau)$ are always positive integers. Then we put $$\Delta(\tau)=\frac{E_4(\tau)^3-E_6(\tau)^2}{1728}$$ to obtain \begin{equation} \label{trane}
\sum_{r=0}^m \sum_{s=0}^{d_r} \sum_{t=0}^{l-k r-s} \frac{(-1)^t c_{r, s}}{1728^{l-k r-s}}\left(\begin{array}{c}
l-k\,r-s \\
t
\end{array}\right) f(\tau)^{12\, r} E_4(\tau)^{3(l-k\, r-t)} E_6(\tau)^{2\, t}=0.   
\end{equation}
Now we consider $\tau\in \mathbb{H}$ such that $e^{2\pi i\tau}$ is algebraic. By Theorem \ref{schneider}, the number $j(\tau)$ is transcendental, and hence Theorem \ref{zero-pole} guarantees that $\tau$ is neither a zero nor a pole of $f$. Moreover, we claim that the number $f(\tau)$ is not algebraic. For otherwise, \eqref{trane} becomes a polynomial in two variables over $\overline{\mathbb{Q}}$, satisfied by $E_4(\tau)$  and $E_6(\tau)$. If we can prove that this is a non-zero polynomial of $E_4(\tau)$ and $E_6(\tau)$, then we would get a contradiction with Nesterenko's Theorem \ref{nest}. Hence, it would complete the proof that $f(\tau)$ is transcendental.\\

To show that the left side of \eqref{trane} becomes a non-zero polynomial of $E_4(\tau)$ and $E_6(\tau)$ when $e^{2\pi i\tau}\in\overline{\mathbb{Q}}$, we consider the following cases:\\

\emph{Case-$1$: } If $f\in M_{k,N,\overline{\mathbb{Q}}}\,\cup M_{k,N,\overline{\mathbb{Q}}}^w$, then $c_m(j)=1$, and hence the coefficient of $E_4(\tau)^{3(l-k m)}\,E_6(\tau)^{0}$ (which appears in \eqref{trane} only for $r=m,\, t=0$) is given by $$\frac{1}{1728^{l-k m}}\cdot 1\cdot f(\tau)^{12 m}\neq 0,$$
which proves that the polynomial is non-zero.\\

\emph{Case-$2$: } If $f\in M_{k,N,\overline{\mathbb{Q}}}^m\setminus \left(M_{k,N,\overline{\mathbb{Q}}}\,\cup M_{k,N,\overline{\mathbb{Q}}}^w\right)$, then the leading coefficient of $P[X]$ is given by $$c_m(j)=\sum_{s=0}^{d_m}c_{m,s}\,j^s.$$
Moreover, it cannot be \emph{identically zero polynomial} because the degree of $P[X]$ is same as the degree of the minimal polynomial of $g$. Since $l>k\, m+2\, d_m$, the following terms
$$
E_4(\tau)^{3(l-k\, m-t)}\, E_6(\tau)^{2\, t} \quad \forall\, 0\leq t\leq d_m
$$
with coefficients
$$
\frac{(-1)^t\, f(\tau)^{12 m}}{1728^{l-km}}\sum_{s=0}^{d_m}\left(\begin{array}{c}
l-k\,m-s \\
t
\end{array}\right) (1728)^s\, c_{m,s} \quad \forall\, 0\leq t\leq d_m
$$
appear in \eqref{trane}. We claim that at least one of these terms has non-zero coefficient. If not, we consider the following system of equations: $ A\, X=0$, where 
$$
A=\left[\begin{array}{ccccc}
    {{l-km}\choose{0}} & {{l-km-1}\choose{0}} & {{l-km-2}\choose{0}} & \ldots & {{l-km-d_m}\choose{0}} \\\\
    {{l-km}\choose{1}} & {{l-km-1}\choose{1}} & {{l-km-2}\choose{1}} & \ldots & {{l-km-d_m}\choose{1}}\\\\
    {{l-km}\choose{2}} & {{l-km-1}\choose{2}} & {{l-km-2}\choose{2}} & \ldots & {{l-km-d_m}\choose{2}}\\\\
    \vdots & \vdots & \vdots & \ldots & \vdots \\\\
    {{l-km}\choose{d_m}} & {{l-km-1}\choose{d_m}} & {{l-km-2}\choose{d_m}} & \ldots & {{l-km-d_m}\choose{d_m}}
\end{array} \right]_{(d_m+1)\times (d_m+1)}  X=\left[\begin{array}{c}
   x_0  \\
    x_1 \\
    x_2 \\
    \vdots\\
    x_{d_m}
\end{array}\right]_{d_m+1.}
$$
Then this system of equations is satisfied by the non-trivial vector $$\left[\begin{array}{c}
   (1728)^0\, c_{m,0} \\
    (1728)^1\, c_{m,1} \\
    (1728)^2\, c_{m,2} \\
    \vdots\\
    (1728)^{d_m}\, c_{m,d_m}
\end{array}\right]_{(d_m+1)\times 1}.$$ From Lemma \ref{det}, we get that $\det(A)\neq 0$, which contradicts the fact that the above system has only trivial solution. Thus, it completes the proof that the polynomial relation of $E_4(\tau)$ and $E_6(\tau)$ in \eqref{trane} is non-trivial.
\end{proof}
\begin{corollary}\label{algi}
If $f\in M_{k,N,\overline{\mathbb{Q}}}\,\cup M_{k,N,\overline{\mathbb{Q}}}^w$, then $f$ is algebraically dependent with $E_4$ and $E_6$ over $\overline{\mathbb{Q}}$. In particular, If $f\in M_{k,N,\overline{\mathbb{Q}}}^m$ and $\tau\in\mathbb{H}$ is not a pole of $f$, then $f(\tau)$ is algebraic over $\overline{\mathbb{Q}}(E_4(\tau),E_6(\tau))$. 
\end{corollary}
\begin{proof}
    It directly follows from the equation \eqref{trane} because the left-hand side is a non-zero polynomial in $f, E_4$ and $E_6$ over $\overline{\mathbb{Q}}.$
\end{proof}
\bigskip

Now we consider the points $\tau\in\mathbb{H}$ such that $j(\tau)\in\overline{\mathbb{Q}}$. From Lemma \ref{omegatau}, we get a transcendental number $\omega_\tau$ such that $g_2(\omega_\tau L_\tau)$ and $g_3(\omega_\tau L_\tau)$ are both algebraic. Moreover, we have the formulae
$$
\begin{aligned}
E_4(\tau)&=\frac{3}{4\pi^4}\, g_2(L_\tau)=\frac{3}{4}\left(\frac{\omega_\tau}{\pi} \right)^4 g_2(\omega_\tau L_\tau),\\ E_6(\tau)&=\frac{27}{8\pi^6}\,g_3(L_\tau)=\frac{27}{8}\left(\frac{\omega_\tau}{\pi}\right)^6g_3(\omega_\tau L_\tau).
\end{aligned}
$$
From Nesterenko's Theorem \ref{nest}, we know that at most one of $E_4(\tau)$ and $E_6(\tau)$ is algebraic. The above formulae imply that if $\frac{\omega_\tau}{\pi}\in\overline{\mathbb{Q}}$, then both $E_4(\tau)$ and $E_6(\tau)$ are algebraic, which is a contradiction. Hence, $\frac{\omega_\tau}{\pi}$ is a transcendental number. Besides, we have 
$$
\begin{aligned}
\Delta(\tau)=\frac{E_4(\tau)^3-E_6(\tau)^2}{1728}&=\frac{1}{4^6}\left(\frac{\omega_\tau}{\pi}\right)^{12}\left(g_2(\omega_\tau L_\tau)^3-27 g_3(\omega_\tau L_\tau)^2 \right)\\
&=\frac{1}{4^6}\left(\frac{\omega_\tau}{\pi}\right)^{12}\Delta_0(\omega_\tau L_\tau).  
\end{aligned}
$$
As $g_2(\omega_\tau L_\tau)$ and $g_3(\omega_\tau L_\tau)$ both are algebraic, the number $\Delta_0(\omega_\tau L_\tau)\in\overline{\mathbb{Q}}\setminus \{0\}$. Since $\frac{\omega_\tau}{\pi}$ is transcendental and $\Delta(\tau)$ is a non-zero algebraic multiple of $\left(\frac{\omega_\tau}{\pi}\right)^{12}$, we deduce that $\Delta(\tau)$ is transcendental. We use this observation to prove the following theorem.
\begin{theorem}\label{algem}
Let $f$ be a non-constant meromorphic modular form with algebraic Fourier coefficients. If $\tau\in\mathbb{H}$ such that  $j(\tau)\in \overline{\mathbb{Q}}$ and $\tau$ is not a pole of $f$, then there exists a transcendental number $\omega_\tau$ (which depends only on $\tau$ and is $\overline{\mathbb{Q}}$-linearly independent with $\pi$) such that $\left(\frac{\pi}{\omega_\tau}\right)^k f(\tau)\in \overline{\mathbb{Q}}$. Therefore, $f(\tau)$ is either zero or transcendental.
\end{theorem}
\begin{proof}
     We first construct the function $$g(\tau):=\frac{f(\tau)^{12}}{\Delta(\tau)^k}\quad \forall \tau\in\mathbb{H},$$ which lies in $ F_{N,\overline{\mathbb{Q}}}$ by the same argument as in Theorem \ref{algem'}. Since $j(\tau)\in\overline{\mathbb{Q}}$, Theorem \ref{algemf} implies that $g(\tau)$ is an algebraic number $c$, (say). Thus, we get  \begin{equation}\label{forms}
f(\tau)^{12}=c.\Delta(\tau)^k=\frac{c}{4^6}.\Delta_0(\omega_\tau L_\tau)^k\left(\frac{\omega_\tau}{\pi} \right)^{12 k}. 
\end{equation}
This shows that $f(\tau)$ is zero when $c=0$, and otherwise, it a non-zero algebraic multiple of $\left(\frac{\omega_\tau}{\pi} \right)^{k}$ and hence transcendental.
\end{proof}
\bigskip

\subsection{Proof of Theorem \ref{deg.}}\, \bigskip

 By theorem \ref{mdft..}, the numbers $g(\tau)$ and $j(\tau)$ are algebraically dependent over $\overline{\mathbb{Q}}$. Moreover, $j(\tau)$ is algebraic if and only if $g(\tau)\in\overline{\mathbb{Q}}$ by Theorem \ref{suma}. These facts together give the first equality below:
$$
\begin{aligned}
 trdeg\,\overline{\mathbb{Q}}\left(e^{2\pi i\tau},\, f(\tau),\, g(\tau) \right)&= trdeg \,\overline{\mathbb{Q}}\left(e^{2\pi i\tau},\, f(\tau),\, j(\tau) \right)\\
 &= trdeg \,\overline{\mathbb{Q}}\left(e^{2\pi i\tau},\, \frac{f(\tau)^{12}}{\Delta(\tau)^k},\, f(\tau),\, j(\tau) \right)\\
 &= trdeg \,\overline{\mathbb{Q}}\left(e^{2\pi i\tau},\, \Delta(\tau)^k,\, f(\tau),\, j(\tau) \right)\\
 &= trdeg \,\overline{\mathbb{Q}}\left(e^{2\pi i\tau},\, \Delta(\tau),\, f(\tau),\, j(\tau) \right) \\
 &= trdeg \,\overline{\mathbb{Q}}\left(e^{2\pi i\tau},\, E_4(\tau)^3,\, E_6(\tau)^2,\, f(\tau),\, j(\tau) \right)\\
 &=trdeg \,\overline{\mathbb{Q}}\left(e^{2\pi i\tau},\, E_4(\tau),\, E_6(\tau),\, f(\tau),\, j(\tau) \right)\\
 &= trdeg \,\overline{\mathbb{Q}}\left(e^{2\pi i\tau},\, E_4(\tau),\, E_6(\tau),\, f(\tau) \right)\\
 &= trdeg \,\overline{\mathbb{Q}}\left(e^{2\pi i\tau},\, E_4(\tau),\, E_6(\tau)\right).
\end{aligned}
$$ 
Here, the second equality follows from the fact that $f^{12}/\Delta^k\in F_{N,\overline{\mathbb{Q}}}$ and  $f(\tau)^{12}/\Delta(\tau)^k$ is algebraically dependent with $j(\tau)$ by Theorem \ref{mdft..}. The third and the seventh equality hold as the fields are identical. To prove the equality of fields in the third step, one needs the assumption that $f(\tau)\neq 0$.  
The final equality holds because $f(\tau)$ is algebraic over $\overline{\mathbb{Q}}\left(E_4(\tau), E_6(\tau)\right)$ by Corollary \ref{algi}. Then we apply Nesterenko's Theorem \ref{nest} on the right side of the final equality to conclude that the transcendence degree of the field $\overline{\mathbb{Q}}\left(e^{2\pi i\tau},\, f(\tau),\, g(\tau) \right)$ over $\overline{\mathbb{Q}}$ is at least two. 
\bigskip

\subsection{Values of quasi-modular forms}\,\bigskip

Let $\widetilde{M}_{k,\overline{\mathbb{Q}}}^{(p)}(\Gamma)$ denote the set of all quasi-modular forms of weight $k$ and depth $p(>0)$ for $\Gamma$ with algebraic Fourier coefficients. We study their values at the points $\tau\in\mathbb{H}$, where one of $e^{2\pi i\tau}$ and $j(\tau)$ is algebraic.
\begin{theorem}
Let $\Gamma$ be a congruence subgroup of level $N(\geq 1)$ and $\widetilde{f}\in \widetilde{M}_{k,\overline{\mathbb{Q}}}^{(p)}(\Gamma)$. If $\tau\in\mathbb{H}$ such that $e^{2\pi i\tau}$ is algebraic, then $f(\tau)$ is transcendental. 
\end{theorem}
\begin{proof}
   By Theorem \ref{quasistructure}, we can write $\widetilde{f}$ in the form $$\widetilde{f}=\sum_{r=0}^p g_r E_2^r,\quad g_r\in M_{k,N,\overline{\mathbb{Q}}}.$$ From Corollary \ref{algi}, we know that each number $g_r(\tau)$ is algebraically dependent with $E_4(\tau), E_6(\tau)$. Therefore, we get the following equality:
 $$ trdeg\,\overline{\mathbb{Q}}(g_0(\tau), g_2(\tau),\ldots, g_p(\tau), E_4(\tau), E_6(\tau))=trdeg\,\overline{\mathbb{Q}}( E_4(\tau), E_6(\tau)).$$
If possible, let both $e^{2\pi i\tau}, \widetilde f(\tau)\in \overline{\mathbb{Q}}$. Since 
$\widetilde f(\tau)=\sum_{r=0}^p g_r(\tau) E_2(\tau)^r$ and $g_r(\tau)\neq 0$ (by Theorem \ref{algem'}), the number $E_2(\tau)$ is algebraic over $\overline{\mathbb{Q}}\left(g_0(\tau), g_1(\tau),\ldots, g_p(\tau), E_4(\tau), E_6(\tau)\right)$ and so
$$
\begin{aligned}
&\quad trdeg\, \overline{\mathbb{Q}}\left(E_2(\tau), g_0(\tau), g_1(\tau),\ldots, g_p(\tau), E_4(\tau), E_6(\tau)\right)\\ &=trdeg\,\overline{\mathbb{Q}}(g_0(\tau), g_1(\tau),\ldots, g_p(\tau), E_4(\tau), E_6(\tau))\\ 
&=trdeg\,\overline{\mathbb{Q}}(E_4(\tau), E_6(\tau)).
\end{aligned}
$$
Since $e^{2\pi i\tau}$ is an algebraic number, the following field containment holds
$$
\overline{\mathbb{Q}}(e^{2\pi i\tau},\, E_2(\tau),\,E_4(\tau), E_6(\tau))\subseteq \overline{\mathbb{Q}}\left(E_2(\tau), g_0(\tau), g_1(\tau),\ldots, g_p(\tau), E_4(\tau), E_6(\tau)\right).
$$
Therefore, we get 
$$
\begin{aligned}
 &\quad trdeg\, \overline{\mathbb{Q}}(e^{2\pi i\tau},\, E_2(\tau),\,E_4(\tau),\, E_6(\tau))\\ &\leq trdeg\,\overline{\mathbb{Q}}\left(E_2(\tau),\, g_0(\tau),\, g_1(\tau),\ldots, g_p(\tau), E_4(\tau),\, E_6(\tau)\right)\\
  &= trdeg\,\overline{\mathbb{Q}}(E_4(\tau),\, E_6(\tau))\\
  &\leq 2,
\end{aligned}
$$
which contradicts Nesterenko's Theorem \ref{nest}. Hence, $\widetilde{f}(\tau)$ is a transcendental number. 
\end{proof}
To study values of quasi-modular forms in the complementary case (i.e., for $j(\tau)\in\overline{\mathbb{Q}}$), we need the following lemma.
\begin{lemma}\label{peri}
 If $\tau\in \mathbb{H}$ such that $j(\tau)$ is algebraic, then there exists a unique (up to algebraic multiple) transcendental number $\omega_\tau$ for which $\frac{\omega_\tau}{\pi}$ and $\frac{\eta_2}{\pi}$ are algebraically independent over $\overline{\mathbb{Q}}$. Here, $\eta_2$ is the second quasi-period with respect to the lattice $\omega_\tau L_\tau$ 
\end{lemma}
\begin{proof}
    We have already proved existence of $\omega_\tau$ in Lemma \ref{omegatau}.  From the definition of the Weierstrass zeta-function, one gets
 $   
Z_{\omega_\tau L_\tau}(z)= \frac{1}{\omega_\tau}\cdot Z_{L_\tau}\left(\frac{z}{\omega_\tau}\right).
$ Then we use the definition of quasi-period and the identity \eqref{quasi-period} obtain
$$
\eta_2=Z_{\omega_\tau L_\tau}(z+\omega_\tau)-Z_{\omega_\tau L_\tau}(z)= \frac{1}{\omega_\tau}\,\eta_2(L_\tau)=\frac{1}{\omega_\tau}\,G_2(\tau)= \frac{1}{3}\cdot\frac{\pi^2}{ \omega_\tau}\,E_2(\tau).
$$
Thus, we get the following formulae 
$$
E_2(\tau)=3\,\frac{\omega_\tau}{\pi}\frac{\eta_2}{\pi}, \quad E_4(\tau)=\frac{3}{4}\left(\frac{\omega_\tau}{\pi} \right)^4 g_2(\omega_\tau L_\tau),\quad E_6(\tau)=\frac{27}{8}\left(\frac{\omega_\tau}{\pi}\right)^6g_3(\omega_\tau L_\tau).$$
The above formulae imply that $E_2(\tau)$, $E_4(\tau)$ and $E_6(\tau)$ are algebraic over $\overline{\mathbb{Q}}\left(\frac{\omega_\tau}{\pi}, \frac{\eta_2}{\pi}\right)$. If possible, let $\frac{\omega_\tau}{\pi}, \frac{\eta_2}{\pi}$ be algebraically dependent. Then we get that $$trdeg\,\overline{\mathbb{Q}}(E_2(\tau),E_4(\tau),E_6(\tau))=1,$$ which contradicts Nesterenko's Theorem \ref{nest}. Hence, the numbers $\frac{\omega_\tau}{\pi}$ and $\frac{\eta_2}{\pi}$ are algebraically independent over $\overline{\mathbb{Q}}$.
\end{proof}
Lemma \ref{peri} and Theorem \ref{quasistructure} together allow us to describe values of quasi-modular forms explicitly at the points $\tau\in\mathbb{H}$ where $j$-function is algebraic.
\begin{theorem}\label{tranq}
 Let $\widetilde{f}\in \widetilde{M}_{k,\overline{\mathbb{Q}}}^{(p)}(\Gamma)$ and $\tau\in\mathbb{H}$ such that $j(\tau)\in\overline{\mathbb{Q}}$, then $f(\tau)$ is a polynomial in $\frac{\omega_\tau}{\pi}$ and $\frac{\eta_2}{\pi}$ with algebraic coefficients. Therefore, $f(\tau)$ is either zero or transcendental.   
\end{theorem}
\begin{proof}
    From Theorem \ref{quasistructure}, we have the expression $$\widetilde{f}(\tau)=\sum_{r=0}^p g_r(\tau) E_2(\tau)^r,\quad \text{where } g_r\in M_{k,N,\overline{\mathbb{Q}}}.$$
Then we apply Theorem \ref{algem} for each coefficient $g_r$, and the above formula for $E_2$ such that 
\begin{equation}\label{last}
\widetilde{f}(\tau)=\sum_{r=0}^p c_r\left(\frac{\omega_\tau}{\pi}\right)^{k-2 r}\left(\frac{\omega_\tau}{\pi}\cdot\frac{\eta_2}{\pi}\right)^r=\sum_{r=0}^p\, c_r\, \left(\frac{\omega_\tau}{\pi}\right)^{k-r}\, \left(\frac{\eta_2}{\pi}\right)^r,
 \end{equation}
where each $c_r$ is an algebraic constant. If $\widetilde{f}(\tau)\in\overline{\mathbb{Q}}\setminus \{0\}$, then \eqref{last} shows that $\frac{\omega_\tau}{\pi}$ and $\frac{\eta_2}{\pi}$ are algebraically dependent, which contradicts Lemma \ref{peri}. Therefore, the number $\widetilde{f}(\tau)$ is either zero or transcendental.     
\end{proof}
\bigskip

\subsection{Proof of Theorem \ref{deg..}}
\,\bigskip

By the same argument as in Theorem \ref{deg.}, we get
\begin{equation}\label{final.}
trdeg\,\overline{\mathbb{Q}}\left( e^{2\pi i\tau},\, f(\tau),\, g(\tau),\, \widetilde f(\tau)\right)= trdeg\,\overline{\mathbb{Q}}\left( e^{2\pi i\tau},\, E_4(\tau),\, E_6(\tau),\, \widetilde f(\tau)\right).
\end{equation}
Moreover, we can apply Theorem \ref{quasistructure} to write
$$\widetilde{f}(\tau)=\sum_{r=0}^p g_r(\tau) E_2(\tau)^r,\quad \text{where } g_r\in M_{k,N,\overline{\mathbb{Q}}}.$$
Thus, we obtain 
  $$
\begin{aligned}
&trdeg\,\overline{\mathbb{Q}}\left( e^{2\pi i\tau},\, E_4(\tau),\, E_6(\tau),\, \widetilde f(\tau)\right)\\
&=trdeg\,\overline{\mathbb{Q}}\left( e^{2\pi i\tau},\, E_4(\tau),\, E_6(\tau),\,\sum_{r=0}^p g_r(\tau)\, E_2(\tau)^r \right)\\
&=trdeg\,\overline{\mathbb{Q}}\left( e^{2\pi i\tau}, E_4(\tau), E_6(\tau), g_0(\tau),\ldots,g_p(\tau),\sum_{r=0}^p g_r(\tau)\, E_2(\tau)^r \right)\\
&=trdeg\,\overline{\mathbb{Q}}\left( e^{2\pi i\tau},\, E_4(\tau),\, E_6(\tau),\, g_0(\tau),\,\ldots,\,g_p(\tau),\,E_2(\tau) \right)\\
&=trdeg\,\overline{\mathbb{Q}}\left( e^{2\pi i\tau},\, E_4(\tau),\, E_6(\tau),\,E_2(\tau) \right),
\end{aligned}
 $$
where the second and fourth equalities follow from the fact that each coefficient $g_r(\tau)$ is algebraic over $\overline{\mathbb{Q}}\left(E_4(\tau), E_6(\tau)\right)$ by Corollary \ref{algi}. Moreover, Nesterenko's theorem \ref{nest} implies that the right side of the final equality is at least $3$, and then we use the equality \eqref{final.} to conclude the proof. 












\begin{thebibliography}{0}

\bibitem{b-d-g-p}
K. Barr\'e-Sirieix, G. Diaz, F. Gramain, G. Philibert, \textit{Une preuve de la conjecture de Mahler-Manin}, Invent. Math., 124 (1996), 1-9.

\bibitem{bertrand}
D. Bertrand, \textit{S\'eries d’Eisenstein et transcendance.} Bull. Soc. Math. France
$104(3)$, $309–321 (1976)$

\bibitem{quasi}
    C. Y. Chang, \textit{Transcendence of special values of quasi-modular forms}, Forum Math. $24$ (2012) $539-551$.


\bibitem{chudnovsky}
    G.V. Chudnovsky, \textit{Algebraic independence of constants connected with the exponential and the elliptic functions}, Dokl.
Akad. Nauk Ukrain. SSR Ser. A $8$ $(1976)$ $698–701$.


\bibitem{stromberg}
H. Cohen and F. Stromberg, \textit{Modular Forms, A Classical Approach}, Graduate Studies In Matheatics $179$. 

\bibitem{jeon-kang-kim}
D Jeon, SY Kang, CH Kim, \textit{On values of weakly holomorphic modular functions at divisors of meromorphic modular forms}, J. Number Theory 239, 183-206.

\bibitem{hamieh-murty}
Alia Hamieh and M. Ram Murty, {A note on $q$-analogues of Dirichlet L-functions}, International Journal of Number Theory, Vol. 12 (2016) 765-773.

\bibitem{lang}
Serge Lang, \textit{Elliptic Functions}, Addision-Wesley Publishing Company, Inc.


\bibitem{murty-dewar-graves}
M. R. Murty, M. Dewar and H. Graves, \textit{Problems in 
the Theory of 
Modular
Forms}, Hindustan Book Agency.

\bibitem{r-m}
M. Ram Murty and P. Rath, \textit{Transcendental Numbers} (Springer-Verlag, $2014$).

\bibitem{r-m-g}
S. Gun, M. Ram Murty, P. Rath, \textit{Algebraic independence of values of
modular forms}. Int. J. Number Theory 7(4), 1065–1074 (2011).


\bibitem{nesterenko}
Y. Nesterenko, \textit{Modular functions and transcendence questions}, Mat. Sb. $187$ $(9)$ $(1996)$ $65–96$ (in Russian); translation in:
Sb. Math. $187$ $(9)$ $(1996)$ $1319–1348$.


\bibitem{schneider}
T. Schneider, \textit{Arithmetische Untersuchungen elliptischer Integrale}, Math. Ann.$ 113$ $(1973)$ $1–13$.

\bibitem{shimura}
 G. Shimura, \textit{Introduction to the arithmetic theory of automorphic functions}, Kano Memorial Lectures $1$.

\bibitem{1-2-3}
D. Zagier, \textit{Elliptic Modular Forms and Their Applications}, Lectures at a Summer School
in Nordfjordeid, Norway.
\end{thebibliography}

\end{document}


\documentclass[11pt, reqno]{amsart}
\usepackage[utf8]{inputenc}
\usepackage{amsmath}
\usepackage{amsthm}
\usepackage{amsfonts}
\usepackage{amssymb}
\usepackage{mathrsfs, mathtools, mathdots}
\usepackage{diagbox}
\DeclarePairedDelimiter \floor {\lfloor} {\rfloor}
\DeclareMathOperator{\Tr}{Tr}

\usepackage[top=35mm,bottom=35mm,left=27mm,right=27mm]{geometry}
\newtheorem{theorem}{Theorem}[section]
\newtheorem{conjecture}{Conjecture}
\newtheorem{corollary}{Corollary}
\newtheorem{problem}{Problem}
\newtheorem*{theorem*}{Theorem}
\newtheorem*{remark*}{Remark}
\newtheorem*{problem*}{Problem}
\newtheorem*{conjecture*}{Conjecture}
\newtheorem{lemma}[theorem]{Lemma}
\newtheorem{prop}[theorem]{Proposition}


\begin{document}

\title[Transcendence of special values of functions]{A note on transcendence of special values \\ of functions related to modularity}

\author[T. Bhowmik]{Tapas Bhowmik}

\author[S. Pathak]{Siddhi Pathak}

\address{Department of Mathematics, University of South Carolina, 1523 Greene Street, LeConte College, Columbia, South Carolina, USA 29208.}

\address{Chennai Mathematical Institute, H-1 SIPCOT IT Park, Siruseri, Kelambakkam, Tamil Nadu, India 603103.}

\email{tbhowmik@email.sc.edu}

\email{siddhi@cmi.ac.in}

\subjclass[2010]{11J91, 11F03, 11F11}

\keywords{Values of modular functions, values of modular forms, Nesterenko's theorem, Schneider's theorem}

\thanks{Research of the second author was partially supported by an INSPIRE fellowship.}



\begin{abstract}
We give a comprehensive treatment on the arithmetic nature of values of modular functions, meromorphic modular forms and quasi-modular forms with respect to arbitrary congruence subgroups, that have algebraic Fourier coefficients at all cusps. This approach unifies many of the known results, which we use to prove generalizations of the theorems of Schneider, Nesterenko and Soon Yi-Kang, et al. 
\end{abstract}

\maketitle

\section{\bf Introduction}
\bigskip

Studying the arithmetic nature of values of special transcendental functions at algebraic arguments is a well-established theme in number theory. Continuing in this spirit, this note focuses on the transcendental nature and algebraic independence of values of functions arising in the modular world, such as modular functions, modular forms and quasi-modular forms. The genesis of this study can be traced back to a 1937 theorem of Schneider \cite{schneider}, namely,
\begin{theorem}[Schneider]\label{schneider}
    If $\tau\in \mathbb{H}$ is algebraic but not imaginary quadratic, then $j(\tau)$ is transcendental.
\end{theorem}
It is known from the theory of complex multiplication that if $\tau\in\mathbb{H}$ generates an imaginary quadratic field ($\tau$ is a CM point), then $j(\tau)$ is an algebraic number. Therefore, Schneider's theorem translates to the statement: if $\tau$ is algebraic, then $j(\tau)$ is algebraic if and only if $\tau$ is CM. A further conjecture by Mahler in this regard, proved by Barr\'e-sirieix, Diaz, Gramain and Philibert \cite{b-d-g-p} states that
\begin{theorem}[Barr\'e-sirieix, Diaz, Gramain, Philibert]\label{mahl}
For any $\tau\in \mathbb{H}$, at least one of the two numbers $e^{2\pi i\tau}$ and $j(\tau)$ is transcendental.
\end{theorem}
This can be derived as a consequence of a remarkable theorem of Nesterenko \cite{nesterenko} with several application.
\begin{theorem}[Nesterenko] \label{nest}
If $\tau\in \mathbb{H}$, then at least three of the numbers
$$
e^{2 \pi i \tau},\quad E_2(\tau),\quad E_4(\tau),\quad E_6(\tau)
$$
are algebraically independent over $\overline{\mathbb{Q}}$.
\end{theorem}

Although the above theorems are about specific functions of ``level $1$", it is the aim of this note to highlight that they are sufficient to deduce the corresponding general results for functions associated with arbitrary congruence subgroups. It seems to the authors that this fact may be known among experts, but nonetheless, is not as well documented. Often, the congruence subgroup in question is restricted to be the group $\Gamma_0(N)$. For instance, in a recent paper, D. Jeon, S.-Y. Kang and C. H. Kim \cite[Theorem 2.4]{jeon-kang-kim} prove the following. Let $N \in \mathbb{N}$, $g_0(N) = $ genus of $X_0(N)$, the modular curve obtained as the quotient of the extended upper half plane by $\Gamma_0(N)$. Suppose that $g_0(N) > 0$ and $m \geq g+1 $ is an integer. Let $\mathfrak{f}_{N,m}$ denote the unique modular function with respect to $\Gamma_0(N)$ constructed in \cite{j-k-k-2} such that $\mathfrak{f}_{N,m}$ is holomorphic on the upper half plane and
\begin{equation*}
    \mathfrak{f}_{N,m}(q) = \frac{1}{q^m} + \sum_{l=1}^g a_N(m,-l) \frac{1}{q^l} + O(q),
\end{equation*}
with the coefficients of powers of $q$ being algebraic numbers. Then, they show
\begin{theorem}[Jeon, Kang, Kim]\label{jjk-thm}
   Let $f$ be a non-zero meromorphic modular form with respect to $\Gamma_0(N)$ with algebraic Fourier coefficients. If $\tau$ is either a zero of a pole of $f$, then $\mathfrak{f}_{N,m}(\tau)$ is algebraic for all $m$. 
\end{theorem}
As a corollary, they deduce that any zero or a pole of such a $f$ should be either transcendental or imaginary quadratic.\\

A detailed investigation of algebraic independence of values of modular forms and quasi-modular forms was carried out by S. Gun, M. R. Murty and P. Rath \cite{r-m-g} in 2011. Their results on values of modular forms were further elaborated upon for higher level in \cite{hamieh-murty} by A. Hamieh and M. R. Murty. These theorems will follow from our discussion later. We also correct a minor error in the statement of one of the theorems in \cite{hamieh-murty}. In the context of quasimodular forms, it was proven independently by Gun, Murty and Rath \cite[Theorem 7]{r-m-g} and C. Y. Chang \cite{chang} that
\begin{theorem}[Gun-Murty-Rath and Chang]
    If $f$ is a quasi-modular form of non-zero weight for $SL_2(\mathbb{Z})$ with algebraic Fourier coefficients and $\tau\in\mathbb{H}$ is such that $j(\tau)\in\overline{\mathbb{Q}}$, then either $f(\tau) = 0$ or $f(\tau)$ is transcendental.  
\end{theorem}
The above statement if also proved for quasimodular forms with respect to $\Gamma_0(N)$ in \cite{r-m-g}.\\

In this note, we first give an exposition of the algebraic structure of modular functions of higher level, following Shimura \cite{shimura}. Building upon this and using Schneider's theorem, we prove
\begin{theorem}\label{mdf-values-thm}
Let $g$ be a non-constant modular function with respect to a congruence subgroup $\Gamma$ of level $N$, with algebraic Fourier coefficients at all cusps. 
\begin{enumerate}
    \item[(a)] If $\tau\in\mathbb{H}$ is either a zero or pole of $g$, then $j(\tau)$ is algebraic, and hence, $\tau$ is either CM or transcendental.
    \item If $\tau\in\mathbb{H}$ is such that $\tau$ is not a pole of $g$, then $j(\tau)\in\overline{\mathbb{Q}}\Leftrightarrow g(\tau)\in\overline{\mathbb{Q}}$. Thus, at least one of $g(\tau)$ and $e^{2\pi i\tau}$ is transcendental for any $\tau \in \mathbb{H}$.
\end{enumerate}
\end{theorem}

A meromorphic modular form of weight $k \in \mathbb{Z}$, with respect to a congruence subgroup $\Gamma$, is a meromorphic function $f: \mathbb{H} \rightarrow \mathbb{C}$ such that 
\begin{enumerate}
    \item[(i)] $f \big|_k \gamma = f$ for all $\gamma \in \Gamma$
    \item[(ii)] $f \big|_k \alpha$ is meromorphic at $i \infty$ for all $\alpha \in SL_2(\mathbb{Z})$.
\end{enumerate}
Here $f \big|_k \left( \begin{smallmatrix}
    a & b \\
    c& d
\end{smallmatrix} \right) (z) = {(cz+d)}^{-k} f \left( \frac{az+b}{cz+d} \right)$. For example, if $f$ is a non-zero holomorphic modular form of weight $k > 0$, then $1/f$ is a meromorphic modular form of weight $-k$. \\

In this context, we first establish a generalization of Theorem \ref{jjk-thm} to arbitrary modular forms and arbitrary modular functions. 
\begin{theorem}\label{zero-pole-value}
Let $f$ be a non-zero, non-constant meromorphic modular form of weight $k \in \mathbb{Z}$ with respect to a congruence subgroup of level $N$ and $g$ be a non-constant modular function with respect to a congruence subgroup of level $M$. Suppose that both $f$ and $g$ have algebraic Fourier coefficients at all cusps. Let $\tau \in \mathbb{H}$ be a zero or a pole of $f$. Then either $g(z)$ has a pole at $z=\tau$ or $g(\tau)$ is algebraic.
\end{theorem}


Furthermore, we generalize the theorems in \cite{r-m-g} and \cite{hamieh-murty} to the setting of meromorphic modular forms. In particular, we show the following. 
\begin{theorem}\label{general-mero-values}
Let $f$ be a non-constant meromorphic modular form with respect to a congruence subgroup $\Gamma$ of level $N$. Suppose that $f$ has algebraic Fourier coefficients at all cusps. Suppose that $\tau$ is not a pole of $f$. 
  \begin{enumerate}
      \item If $\tau\in\mathbb{H}$ is such that $e^{2\pi i\tau}$ is algebraic, then $f(\tau)$ is transcendental.
      \item If $\tau\in\mathbb{H}$ is such that $j(\tau)\in \overline{\mathbb{Q}}$, then there exists a transcendental number $\omega_\tau$ which depends only on $\tau$ and is $\overline{\mathbb{Q}}$-linearly independent with $\pi$, such that $\left(\frac{\pi}{\omega_\tau}\right)^k f(\tau)\in \overline{\mathbb{Q}}$. Therefore, $f(\tau)$ is either zero or transcendental.
  \end{enumerate}  
\end{theorem}


With regard to quasi-modular forms, we extend the previously known results to meromorphic quasi-modular forms, namely, polynomials in the weight $2$ normalized Eisenstein series, $E_2(z)$ with coefficients that are meromorphic modular forms. A concise formulation of our results is the following. 
\begin{theorem}\label{deg..}
 Let $f$ and $g$ be as in Theorem \ref{general-mero-values}. Let $\widetilde{f}$ be a non-zero quasi-modular form of depth at least $1$ with respect to $\Gamma$, having algebraic Fourier coefficients at all cusps. If $\tau\in\mathbb{H}$ is such that $f(\tau)\neq 0$, $\widetilde{f}(\tau)\neq 0$ and $\tau$ is not a pole of $f$ and $g$, then 
    $$
trdeg\,\overline{\mathbb{Q}}\left( e^{2\pi i\tau},\, f(\tau),\, g(\tau),\, \widetilde f(\tau)\right)\geq 3. 
    $$
\end{theorem}
One can recover Theorem \ref{nest} from the above statement by choosing
$$
f=E_4, \qquad  g=j, \qquad  \widetilde f= E_6+E_4\, E_2.
$$
Note however that, the proof of Theorem \ref{deg..} uses Theorem \ref{nest}.



\bigskip

\section{\bf Preliminaries}
\bigskip

This aim of this section is to study the algebraic structure of modular functions and to record the required results from transcendental number theory. For the sake of completeness and clarity of exposition, a brief account of the proofs is included, and appropriate references are given.


\subsection{Modular and quasi-modular forms}
For each $N\in \mathbb{N}$, $$
\Gamma(N):=\left\{  \left( 
\begin{array}{cc}
  a & b \\
  c & d
\end{array}
\right)\in SL_2(\mathbb{Z}):\left( 
\begin{array}{cc}
  a & b \\
  c & d
\end{array}
\right)\equiv \left( 
\begin{array}{cc}
  1 & 0 \\
  0 & 1
\end{array}
\right) (\operatorname{mod}N) \right\}
$$ is called the \emph{principal congruence subgroup of level N}, with $\Gamma(1)=SL_2(\mathbb{Z})$. A subgroup $\Gamma$ of $SL_2(\mathbb{Z})$ is said to be congruence subgroup if there exists $N\in \mathbb{N}$ such that $\Gamma(N)\subseteq \Gamma$. The smallest such $N$ is said to be the level of $\Gamma$.\\

A \emph{holomorphic modular form} of integer weight $k \geq 0$ with respect to a congruence subgroup $\Gamma$ is a holomorphic function on the upper half-plane $\mathbb{H}$ which satisfies 
\begin{enumerate}
    \item[(i)] $f \big|_k \gamma = f$ for all $\gamma \in \Gamma$
    \item[(ii)] $f \big|_k \alpha$ is holomorphic at $i\infty$ for all $\alpha \in SL_2(\mathbb{Z})$.
\end{enumerate}
The function $f$ is said to be \emph{weakly holomorphic modular form} if $f \big| \alpha$ is allowed to have poles at $i\infty$, that is, $f$ is \emph{meromorphic} at the cusps. More generally, $f$ is called \emph{meromorphic modular form} if it is meromorphic on $\mathbb{H}$ and also at cusps. \\

We will say that $f$ has algebraic Fourier coefficients if the Fourier coefficients of $f \big|_k \alpha$ at $i \infty$ are algebraic numbers. Denote the space of all holomorphic, weakly holomorphic and meromorphic modular forms, with algebraic Fourier coefficients by $M_{k,\overline{\mathbb{Q}}}(\Gamma)$, $M_{k,\overline{\mathbb{Q}}}^w(\Gamma)$ and $M_{k,\overline{\mathbb{Q}}}^m(\Gamma)$ respectively. Clearly, $M_{k,\overline{\mathbb{Q}}}(\Gamma)\subset M_{k,\overline{\mathbb{Q}}}^w(\Gamma)\subset M_{k,\overline{\mathbb{Q}}}^m(\Gamma)$. \\

For even integer $k\geq 2$, define the normalized Eisenstein series of weight $k$ for $S L_2(\mathbb{Z})$ by
$$
E_{k}(\tau) = 1-\frac{2 k}{B_{k}} \sum_{n=1}^{\infty} \sigma_{k-1}(n) q^{n}, \text{   where }q=e^{2\pi i\tau} \text{ and } \sigma_s(n)=\sum_{\substack{d|n\\ d>0}}d^s.
$$ 
Here, $B_k$ is the $k$-th Bernoulli number. Define
\begin{equation*}
    \Delta(\tau):=\frac{E_4(\tau)^3-E_6(\tau)^2}{1728}.
\end{equation*}
For $k\geq 4$, the function $E_k\in M_{k,\overline{\mathbb{Q}}}(S L_2(\mathbb{Z}))$ and $\Delta \in M_{12,\overline{\mathbb{Q}}}(S L_2(\mathbb{Z}))$. But $E_2(\tau)$ is not a modular form (see \cite{murty-dewar-graves}, Chapter $5$), as
\begin{equation*}
    E_2 \left( \frac{-1}{\tau} \right) = {\tau}^2 E_2(\tau) + \frac{6}{i\pi} \tau.
\end{equation*}
This motivates the definition of a quasi-modular form, of which there are several equivalent formulations. We use the following characterization \cite[Proposition 20]{1-2-3}.
\begin{theorem}\label{quasistructure}
Every (holomorphic) quasi-modular form for a congruence subgroup $\Gamma$ is a polynomial in $E_2$ with modular coefficients. More precisely, if $f$ is a quasi-modular form of weight $k$ and depth $p$ on $\Gamma$, then $f$ can be uniquely written as $f=\sum_{r=0}^p g_r\, E_2^r$, where $g_r\in M_{k-2 r}(\Gamma)$ for all $0\leq r\leq p$ and $g_p\neq 0$. 
\end{theorem}
A quasi-modular form $f$ is said to have algebraic Fourier coefficients if all modular coefficients in the above expression of $f$ have algebraic Fourier coefficients. \\

\subsection{The Weierstrass $\wp$-function}
 Let $L=\omega_1 \mathbb{Z} \,  \oplus \, \omega_2 \mathbb{Z}$ be a two-dimensional lattice in $\mathbb{C}$ with $\omega_1/\omega_2 \in \mathbb{H}$. The Weierstrass $\wp$-function associated to $L$, given by
\begin{equation}\label{wpfunction}
\wp (z)=\frac{1}{z^2}+\sum_{\substack{\omega\in L\\
\omega\neq 0}}\left(\frac{1}{(z-\omega)^2}-\frac{1}{\omega^2} \right), \quad \text{for }z\in \mathbb{C}\setminus L,
\end{equation}
defines a meromorphic function on $\mathbb{C}$ with double poles at each lattice points. It satisfies the differential equation
\begin{equation} \label{dell}
 \wp'(z)^2=4\wp(z)^3-g_2(L)\wp(z)-g_3(L),   
\end{equation}
where 
\begin{equation*}
    g_2(L)=60\, \sum_{\substack{\omega \in L\\ \omega\neq 0}} \frac{1}{\omega^{4}} \qquad \text{and} \qquad g_3(L)=140\, \sum_{\substack{\omega \in L\\ \omega\neq 0}} \frac{1}{\omega^{6}}.
\end{equation*}
For the lattice $L_\tau=\tau\,\mathbb{Z}\oplus \mathbb{Z}$ with $\tau\in\mathbb{H}$, 
\begin{equation*}
    g_2(L_\tau)=\frac{4\pi^4}{3}E_4(\tau) \qquad \text{and} \qquad g_3(L_\tau)=\frac{8\pi^6}{27}E_6(\tau).
\end{equation*}
Define the discriminant of a lattice, $$\Delta_0(L):=g_2(L)^3-27g_3(L)^2\neq 0, \text{ for any two dimensional lattice }L.$$ In particular, we have $\Delta_0(L_\tau)=(2\pi)^{12}\,\Delta(\tau)$ for all $\tau\in\mathbb{H}$.\\

The Weierstrass zeta-function associated to $L=\omega_1\,\mathbb{Z}\oplus \omega_2\,\mathbb{Z}$ is defined by
$$
\zeta_{L}(z)=\frac{1}{z}+\sum_{\substack{\omega\in L\\ \omega\neq 0}}\left(\frac{1}{z-\omega}+\frac{1}{\omega}+\frac{z}{\omega^2} \right)\quad \text{for } z\in \mathbb{C}\setminus L.
$$
Note that $\zeta_{L}'(z)=-\wp_L(z)$ is a periodic function with each point of $L$ as a period. Hence, the functions $$\zeta_{L}(z+\omega_1)-\zeta_{L}(z) \text{  and  } \zeta_{L}(z+\omega_2)-\zeta_{L}(z)$$ are constants. These constants are denoted by $\eta_1(L)$ and $ \eta_2(L)$ respectively and are called \emph{quasi-periods}. Moreover, it is known that \begin{equation}\label{quasi-period}
\eta_2(L_\tau)=G_2(\tau)=\frac{\pi^2}{3}\, E_2(\tau) \text{ 
  for all }\tau\in\mathbb{H}. 
\end{equation}
The reader is referred to (\cite{lang}, chapter $18.3$) for proof of these results.\\

The Weierstrass $\wp$-function satisfies the following addition formula (\cite{r-m}, chapter $11$).
 \begin{theorem} \label{addf}
 For $z_1, z_2\in \mathbb{C}$ such that $z_1\pm z_2\notin L$ we have
$$
\wp\left(z_1+z_2\right)=-\wp\left(z_1\right)-\wp\left(z_2\right)+\frac{1}{4}\left(\frac{\wp^{\prime}\left(z_1\right)-\wp^{\prime}\left(z_2\right)}{\wp\left(z_1\right)-\wp\left(z_2\right)}\right)^2 .
$$
\end{theorem}
Using the addition formula of the Weierstrass $\wp$-function, Schneider \cite{schneider.} proved that 
\begin{theorem}\label{tranw}
    Let $L=\omega_1\,\mathbb{Z}\oplus\omega_2\,\mathbb{Z}$ be a lattice such that both $g_2(L)$, $g_3(L)$ are algebraic. If $\alpha$ is an algebraic number with $\alpha\notin L$, then $\wp(\alpha)$ is transcendental.
\end{theorem}
The addition formula also implies the following important proposition (see \cite{r-m}). Two generators $\omega_1$ and $\omega_2$ of a lattice $L$ are said to be primitive if both have minimal absolute value among all generators of $L$.
\begin{prop}\label{tranp}
Let $L=\omega_1\,\mathbb{Z}\oplus\omega_2\,\mathbb{Z}$ be such that both $g_2(L)$ and $g_3(L)$ are algebraic. Assume that $\omega_1$ and $\omega_2$ are primitive generators of $L$. Then, for any natural number $n>1$, the numbers $\wp(\frac{\omega_1}{n})$ and $\wp(\frac{\omega_2}{n})$ are algebraic. Moreover, any non-zero period of $L$ is necessarily transcendental.\\
\end{prop}

\subsection{Modular functions}

A meromorphic function $g$ on $\mathbb{H}$ is said to be a modular function if it satisfies
\begin{equation}
g\left(\frac{a\,\tau+b}{c\,\tau+d}\right)=g(\tau)\quad\text{for all }
\gamma=\left(\begin{smallmatrix}
   a  & b \\
   c  & d
\end{smallmatrix} \right)\in \Gamma,
\end{equation}
that is, $g \big|_0 \gamma = g$, and is also meromorphic at all the cusps. In particular, we call a modular function on the congruence subgroup $\Gamma(N)$ to a be modular function of level $N$. Note that if $g$ is a modular function with respect to $\Gamma$, which is of level $N$, then $g$ is a modular function of level $N$.\\

An example of modular function of level one is given by 
$$
j(\tau):=\frac{E_4^3(\tau)}{\Delta(\tau)} \text{ for all } \tau\in\mathbb{H},
$$
which has the following Fourier expansion at $i\infty$
$$
j(\tau)=\frac{1}{q}+744+196884q+\cdots,$$
where $q=e^{2\pi i\tau}$. \\


More specifically, for any modular function (or form) $f$, define 
\begin{equation*}
    \mathbb{Q}(f) := \mathbb{Q} \left( \left\{ \text{Fourier coefficients of $f$ at all $\Gamma$-in-equivalent cusps} \right\} \right)
\end{equation*}
The $j$-function is a canonical example of a modular function of level one, and governs properties of all modular functions of higher levels as well. This is made precise in the following series of propositions.
\begin{prop} \label{level1}
Let $g$ be a non-constant modular function of level one and $K := \mathbb{Q}(g)$. Let $\mathcal{F}$ denote the standard fundamental domain for the action of $S L_2(\mathbb{Z})$ on $\mathbb{H}$. Suppose that the poles of $g$ in $\mathcal{F}$ are $\tau_1,\tau_2,\cdots,\tau_m$. Then $g(\tau)$ is a rational function in $j(\tau)$ with coefficients in the field $K \left( j(\tau_1),\cdots,j(\tau_m) \right)$.
\end{prop}
\begin{proof}
Consider the function 
$$
h(\tau):=\left(\prod_{j=1}^m \left(j(\tau)-j(\tau_j)\right)^{-\operatorname{ord}_{\tau_j}(g)}\right)g(\tau),
$$
where $\operatorname{ord}_{\tau_j}(g) = - $ order of the pole of $g$ at $\tau_j$. By the compactness of $\mathcal{F} \cup i\infty$, we know that $g$ has only finitely many poles in $\mathcal{F}$. Hence, the above product is finite and $h$ is holomorphic on $\mathbb{H}$.\\

Suppose $h$ has pole of order $M \geq 1$ at $i\infty$. Then the Fourier expansion of
$h(\tau)$ at $i\infty$ has the form
$$
h(\tau)=\sum_{n=-M}^\infty c_n q^n, \, \text{ where } c_{-M}\neq 0.
$$
Note that the modular function $h(\tau)-c_{-M}\,j(\tau)^M$ is holomorphic on $\mathbb{H}$ and its Fourier expansion starts with at most a polar term of order $M-1$. Iterating this process, we can subtract a polynomial in $j(\tau)$ to get a holomorphic modular function that vanishes at $i\infty$, and hence identically zero. Thus, $h(\tau)$ is a polynomial in $j(\tau)$ over $K$, and so $g(\tau)$ is a rational function of $j(\tau)$ over $K(j(\tau_1),\cdots,j(\tau_m))$.
\end{proof}

One can also conclude the following important fact from the above proof.
\begin{corollary}\label{rationalfunctions}
If $g$ is a modular function of level one which is holomorphic on $\mathbb{H}$ with a pole of order $M$ at $i\infty$, then $g(\tau)$ is a polynomial in $j(\tau)$ of degree $M$ with coefficients in $\mathbb{Q}(g)$.
\end{corollary}

The $j$-function is sufficient to `generate' all higher level modular functions as well. This is proved below.
\begin{theorem} \label{mdft}
Let $g$ be a modular function with respect to a congruence subgroup $\Gamma$ and let $\tau_1,\, \tau_2, \cdots, \, \tau_m$ be the poles of $g$ in $\mathcal{F}$. Set
$$
K:=\mathbb{Q}(g)\left(\left\{j(\tau_1), \, j(\tau_2), \cdots, \, j(\tau_m) \right\}\right).
$$
Then there exists a \textbf{monic} polynomial $P_g(X)\in K(j)[X]$ such that $P(f)=0$. 
\end{theorem}
\begin{proof}
Let $[S L_2(\mathbb{Z}):\Gamma]=r$ and $\{ \gamma_1=I,\gamma_2,\cdots,\gamma_r\}$ be a complete set of right coset representatives so that $$
S L_2(\mathbb{Z})=\bigsqcup_{i=1}^r\Gamma \gamma_i.
$$
For all $1\leq i \leq r$, define the functions $$g_i(\tau):=g(\gamma_i\tau) \quad \text{for } \tau \in \mathbb{H}.
$$ Each $g_i$ is independent of the choice of coset representatives as $g$ is $\Gamma$-invariant. Moreover, each $g_i$ is a modular function with respect to $\Gamma$ and the Fourier expansion of $g_i$ at $i \infty$ is precisely the expansion of $g$ at the cusp $g_i(i\infty)$. For any $\gamma\in S L_2(\mathbb{Z})$, we have $g_i(\gamma \tau)=g(\gamma_i\gamma\tau)=g_j(\tau)$ for $1 \leq j \leq r$ such that $\gamma_i\gamma\in \Gamma\gamma_j$, i.e., the set $\{g_1,g_2,...,g_r\}$ gets permuted under the action of $S L_2(\mathbb{Z}).$ This observation implies that any elementary symmetric polynomial in $g_1,g_2,\cdots,g_r$ is a modular function of level one, and is in $K(j)$ by Theorem \ref{level1}. Note that the polynomial $$
P(X)=\prod_{i=1}^r(X-g_i)$$is satisfied by $g$ as $g=g_1$ and has coefficients that are elementary symmetric polynomials in $g_1,g_2,\cdots,g_r$. This proves the theorem.
\end{proof}



\begin{corollary}\label{useful}
Let $g$ be a modular function with respect to a congruence subgroup $\Gamma$ which is holomorphic on $\mathbb{H}$ and $K = \mathbb{Q}(g)$. Then the \textbf{monic} polynomial $P_g(X)$ satisfied by $g$ has coefficients in $K[j][X]$.
\end{corollary}
\begin{proof}
The coefficients of $P_g(X)$ constructed in the proof above are modular functions of level one, holomorphic on $\mathbb{H}$ with Fourier coefficients in $K$. The result now follows from Corollary \ref{rationalfunctions}. 
\end{proof}


An important remark here is that if $g$ is a modular function with respect $\Gamma$ that has algebraic Fourier coefficients, Theorem \ref{mdft} does not guarantee that the coefficients of $P_g$ are algebraic, unless $g$ is holomorphic in $\mathbb{H}$. For general modular functions, one needs to delve deeper into the structure of the field generated by modular functions of higher level in order to assert that a modular function of higher level with algebraic Fourier coefficients satisfies a polynomial over $\overline{\mathbb{Q}}(j)$.



\bigskip






To that effect, we now study the structure of the field of modular functions of a fixed level $N>1$. For any field $\mathbb{F} \subseteq \mathbb{C}$, let
\begin{equation*}
    \mathcal{F}_{N, \mathbb{F}} := \left\{ \text{Modular functions of level $N$ whose Fourier coefficients at $i \infty$ are in $\mathbb{F}$} \right\}.
\end{equation*}
Theorem \ref{mdft} implies that $\mathcal{F}_{N,\mathbb{C}}$ is an algebraic extension of $\mathbb{C}(j)$. Furthermore, we have
\begin{prop}\label{fnc}
The field $\mathcal{F}_{N,\mathbb{C}}$ is a finite Galois extension of $\mathbb{C}(j)$ and the Galois group is isomorphic to $S L_2(\mathbb{Z})/\{\pm I\}\Gamma(N)$.
\end{prop}
\begin{proof}
From Theorem \ref{mdft}, we know that every $g\in \mathcal{F}_{N,\mathbb{C}}$ satisfies $P_g(X)$. The other roots of $P_g(X)$ are the modular functions $g_i$, which are also of level $N$. Hence $\mathcal{F}_{N,\mathbb{C}}/\mathbb{C}(j)$ is a Galois extension. To find the Galois group, note that any element $\gamma\in S L_2(\mathbb{Z})$ defines an automophism of $\mathcal{F}_{N,\mathbb{C}}$ by $g \mapsto g\circ\gamma $, where $(g\circ\gamma )(\tau):=g(\gamma \tau)$ for all $\tau\in\mathbb{H}$. Moreover, $\gamma$ fixes $\mathbb{C}(j)$ as $j$ is modular function of level $1$. Define a map $$
\begin{aligned}
\Psi : \,\,  & S L_2(\mathbb{Z}) \longrightarrow \text{Gal}(\mathcal{F}_{N,\mathbb{C}}/\mathbb{C}(j))\\
& \quad\gamma \longmapsto \Psi_\gamma : \mathcal{F}_{N,\mathbb{C}} \longrightarrow \mathcal{F}_{N, \mathbb{C}}\\
&\quad \quad\quad\quad \quad \quad h \longmapsto h\circ\gamma.
\end{aligned}
$$
Then $\Psi$ is a group homomorphism with kernel $\{\pm I\}\Gamma(N)$. By the primitive element theorem, $\mathcal{F}_{N,\mathbb{C}}=\mathbb{C}(\widetilde g)$ for some $\widetilde{g} \in \mathcal{F}_{N,\mathbb{C}}$. Note that any $\sigma\in Gal(\mathcal{F}_{N,\mathbb{C}}/\mathbb{C}(j)) $ sends $\widetilde{g}$ to one of its conjugates, which is of the form $\widetilde g_i=\widetilde g\circ\gamma_i$. Then we get
$$
\sigma(h)=h\circ\gamma_i\quad \forall\ h\in \mathcal{F}_{N,\mathbb{C}},
$$
and so $\sigma=\Psi_{\gamma_i}$. This shows that $\Psi$ is onto and hence, $\text{Gal}(\mathcal{F}_N/\mathbb{C}(j)) \simeq SL_2(\mathbb{Z})/\{\pm I\}\Gamma(N).$
\end{proof}






% For $N>1$, let $F_{N,\overline{\mathbb{Q}}}$ denote the field of modular functions of level $N$ with algebraic Fourier coefficients at all the $\Gamma(N)$-in-equivalent cusps. Theorem \ref{mdft} implies that $F_{N,\overline{\mathbb{Q}}}$ is an algebraic extension of $\overline{\mathbb{Q}}(j)$.


Following \cite[Chapter 6]{shimura}, we consider explicit modular function of level $N$ whose Fourier coefficients have good rationality properties. Let $\boldsymbol{a}=(a_1,a_2)\in \frac{1}{N}\mathbb{Z}^2 \,\setminus \,\mathbb{Z}^2$, consider the function
$$
f_{\boldsymbol{a}}(\tau):=\frac{g_2(L_\tau)\,g_3(L_\tau)}{\Delta_0(L_\tau)}\, \wp_{L_\tau} (a_1\tau+a_2),
$$
which is holomorphic on $\mathbb{H}$. The following properties can be checked routinely, and we leave the proof to the reader.
\begin{prop}
Let $S$ denotes the set $\{\left(\frac{r}{N},\frac{s}{N}\right): 0\leq r,s\leq N-1 \text{ and }(r,s)\neq(0,0)\}$. Let $\boldsymbol{a}$, $\boldsymbol{b} \in \frac{1}{N}\mathbb{Z}^2 \,\setminus \,\mathbb{Z}^2$ and $f_{\boldsymbol{a}}$ be as defined above. For $\gamma = \big(\begin{smallmatrix}
p & q\\
r & s
\end{smallmatrix}\big)$ and $\boldsymbol{a} = (a_1,a_2)$, let ${\boldsymbol{a\gamma}=(pa_1+ra_2, \, qa_1+sa_2)}$. Then
\begin{enumerate}
    \item[(a)] $f_{\boldsymbol{a}} = f_{\boldsymbol{b}}$ if and only if $\boldsymbol{a} \equiv \boldsymbol{b} \bmod \mathbb{Z}^2 $,
    \item[(b)] for $\gamma \in SL_2(\mathbb{Z})$, $f_{\boldsymbol{a}} (\gamma \tau) = f_{\boldsymbol{a\gamma}}(\tau)$ for all $\tau \in \mathbb{H}$.
\end{enumerate}
Therefore, all elements in $\{f_{\boldsymbol{a}}: \boldsymbol{a} \in S\}$ satisfy modularity with respect to $\Gamma(N)$.
\end{prop}

In order to conclude that the functions $f_{\boldsymbol{a}}$ are modular functions, we study their Fourier expansions at the cusp. But first, we need to understand the behaviour of the Weierstrass $\wp$-function.
\begin{lemma}\label{weifou}
Fix $\tau \in \mathbb{H}$ and denote $\wp_{\tau} = \wp_{L_{\tau}}$. For $z \not\in L_{\tau}$, we have
\begin{equation}\label{weifou}
    (2\pi i)^{-2}\wp_\tau(z)=\frac{1}{12}+\frac{1}{\xi^{-1}-2+\xi}+\sum_{n=1}^\infty c_n q^n, 
\end{equation}
where 
$$
\begin{aligned}
  \xi=e^{2\pi i z},\, q=e^{2\pi i\tau} \text{ and }
    c_n=\sum_{d|n}d\left(\xi^{-d}-2+\xi^d\right)\ \forall\ n\geq 1.
\end{aligned}
$$
\end{lemma}
\begin{proof}
    From the definition of the $\wp$-function, we have 
$$
\begin{aligned}
\wp_\tau(z)&=\frac{1}{z^2}+\sum_{\substack{(m,n)\in \mathbb{Z}^2\\ (m,n)\neq (0,0)}}\left[\frac{1}{(z-m\tau-n)^2}-\frac{1}{(m\tau+n)^2} \right]\\
&=\frac{1}{z^2}+\sum_{n\in \mathbb{Z}\setminus \{0\}}\frac{1}{(z+n)^2}-\sum_{n\in \mathbb{Z}\setminus \{0\}}\frac{1}{n^2}+\sum_{m\neq 0}\sum_{n\in \mathbb{Z}}\frac{1}{(z-m\tau-n)^2}-\sum_{m\neq 0}\sum_{n\in \mathbb{Z}}\frac{1}{(m\tau+n)^2}\\
&=-2\zeta(2)+\sum_{n\in\mathbb{Z}}\frac{1}{(z+n)^2}-2\sum_{m=1}^\infty\sum_{n\in \mathbb{Z}}\frac{1}{(m\tau+n)^2}\\
& \quad \quad \quad \quad \quad \quad \quad \quad \quad \quad \quad +\sum_{m=1}^\infty\sum_{n\in \mathbb{Z}}\left[\frac{1}{(-z+m\tau+n)^2}+\frac{1}{(z+m\tau+n)^2} \right].
\end{aligned}
$$
Recall the Lipschitz summation formula, which states 
\begin{equation*} \label{lipc}
\sum_{n \in \mathbb{Z}} \frac{1}{(z+n)^{k}}=\frac{(-2 \pi i)^{k}}{(k-1) !} \sum_{n=1}^{\infty} n^{k-1} q^{n}.
\end{equation*}
For a proof, see \cite[Theorem 4.2.2]{murty-dewar-graves}. Applying this, we obtain
$$
\begin{aligned}
& \wp_\tau(z)=-\frac{\pi^2}{3}+(-2\pi i)^2\sum_{n=1}^\infty ne^{2\pi i n z}\\
& \quad \quad \quad \quad \quad\quad\quad\quad +(-2\pi i)^2\sum_{m=1}^\infty\sum_{n=1}^\infty n\left[e^{2\pi i n(-z+m\tau)}+e^{2\pi i n(z+m\tau)}-2e^{2\pi i m n \tau} \right] \\
&\Rightarrow (2\pi i)^{-2}\wp_\tau(z)=\frac{1}{12}+\sum_{n=1}^\infty ne^{2\pi i n z}+\sum_{m,n=1}^\infty nq^{mn}\left[e^{2\pi i n z}+e^{-2 \pi i n z}-2 \right]
  \end{aligned}
$$
This implies the result.
\end{proof}

Using the above expansion, the interpretation of $g_2(L_{\tau)}$, $g_3(L_{\tau})$ and $\Delta_{0}(L_{\tau})$ in terms of Eisenstein series, we get for $\boldsymbol{a}= \left( \frac{r}{N}, \frac{s}{N} \right) \in S$,
\begin{align}
& f_{\boldsymbol{a}}(\tau) \nonumber \\
& =\frac{-1}{2592}\cdot\frac{E_4(\tau)E_6(\tau)}{\Delta(\tau)}\left[\frac{1}{12} -\sum_{n=1}^\infty ne^{\frac{2\pi i n s}{N}} q_N^{n r}+\sum_{m,n=1}^\infty n q^{m n}\left[e^{\frac{2\pi i n s}{N}}q_N^{n r}+e^{\frac{-2\pi i n s}{N}}q_N^{-n r}-2 \right]\right]. \label{foufa}
\end{align}
Since the Fourier series of $\Delta(\tau)$ begins with $q=q_N^N$, the Fourier expansion of $f_{\boldsymbol{a}(\tau)}$  begins with a rational multiple of $q_N^{-N}$. Thus, $f_{\boldsymbol{a}}$ has pole of order $N$ at $i\infty$ for all $\boldsymbol{a}\in S$. If $s$ is any other cusp, then there exists $\gamma\in S L_2(\mathbb{Z})$ such that $s=\gamma(i\infty)$. Since $f_{\boldsymbol{a}}(\gamma\tau)=f_{\boldsymbol{a\gamma}}(\tau)$, which has pole of order $N$ at $i\infty$, we conclude that $f_{\boldsymbol{a}}$ is meromorphic at all cusps, with pole of order $N$. Thus, $f_{\boldsymbol{a}}$ is a modular function of level $N$ for all $ \boldsymbol{a} \in S$. \\

This helps us to conclude the following.
\begin{theorem}\label{primiext}
For all $\boldsymbol{a} \in S$, the Fourier coefficients of $f_{\boldsymbol{a}}$ with respect to all cusps belong to $\mathbb{Q}(\mu_N)$  $\forall \, \boldsymbol{a} \in S$, where $\mu_N=e^{2\pi i/N}$. 
\end{theorem}
\begin{proof}
   Recall that Fourier coefficients of $E_4(\tau), E_6(\tau)$ and $\Delta(\tau)$ are integers. Hence, from \eqref{foufa}, it follows that the Fourier coefficients of $f_{\boldsymbol{a}}$ at $i\infty$ lie in $\mathbb{Q}(\mu_N)$ for all $\boldsymbol{a}\in S$. If $s$ is any other cusp, then there exists $\gamma\in S L_2(\mathbb{Z})$ such $s=\gamma(i\infty)$. But $f_{\boldsymbol{a}}(\gamma\tau)=f_{\boldsymbol{a\gamma}}(\tau)$ also has Fourier coefficients in $\mathbb{Q}(\mu_N)$ with respect to $i\infty$. This completes the proof.
\end{proof}

These modular functions, together with th $j$-function serve to generate all modular functions of level $N$. That is,
\begin{theorem}\label{mdft..}
We have $\mathcal{F}_{N, \mathbb{C}} = \mathbb{C} \left(j, \, \{ f_{\boldsymbol{a}} |{\boldsymbol{a}}\in S \} \right)$, $\mathcal{F}_{N,\overline{\mathbb{Q}}}=\overline{\mathbb{Q}}\left( j, \, \{ f_{\boldsymbol{a}} |{\boldsymbol{a}}\in S\} \right)$ and $\mathcal{F}_{N,\overline{\mathbb{Q}}}$ is a finite Galois extension of $\overline{\mathbb{Q}}(j)$.
\end{theorem}
\begin{proof}
Let $E_{N,\mathbb{C}}$ denote the field $\mathbb{C}\left(j, \, \{ f_{\boldsymbol{a}} | {\boldsymbol{a}}\in S\}\right)$. Then we have containment $\mathbb{C}(j)\subseteq E_{N,\mathbb{C}} \subseteq \mathcal{F}_{N,\mathbb{C}}.$ To complete the proof, it is enough to show $[\mathcal{F}_{N,\mathbb{C}}:E_{N,\mathbb{C}}]=1$, i.e., any element of $\text{Gal}(\mathcal{F}_{N,\mathbb{C}}/\mathbb{C}(j))$ that fixes $E_{N,\mathbb{C}}$ must be trivial. Recall that $\text{Gal}(\mathcal{F}_{N,\mathbb{C}}/\mathbb{C}(j))=S L_2(\mathbb{Z})/\{\pm I\}\Gamma(N)$. Let $\gamma=\big(\begin{smallmatrix}
p & q\\
r & s\end{smallmatrix}\big)\in S L_2(\mathbb{Z}) $ fix $E_{N,\mathbb{C}}$, that is, $$f_{\boldsymbol{a}}\circ\gamma=f_{\boldsymbol{a}}\quad \text{for all }{\boldsymbol{a}}\in S.$$ Equivalently,
$$
 \quad f_{\boldsymbol{a\gamma}}=f_{\boldsymbol{a}} \quad \text{for all } {\boldsymbol{a}} \in S
$$
and hence 
$$
 {\boldsymbol{a}\gamma}\equiv \pm {\boldsymbol{a}}\pmod{\mathbb{Z}^2} \quad \text{ for all } {\boldsymbol{a}}\in S.
$$
In particular, taking two choices of ${\boldsymbol{a}}= (\frac{1}{N},0)$ and $(0,\frac{1}{N})$ we conclude $\gamma \in \{\pm I\}\Gamma(N).$ This completes the proof. \\

To prove the second assertion, we proceed exactly as above. First, we show that $E_{N,\overline{\mathbb{Q}}}:= \overline{\mathbb{Q}}\left(j, \, \{ f_{\boldsymbol{a}} | {\boldsymbol{a}}\in S\}\right)$ is a Galois extension of $\overline{\mathbb{Q}}(j)$. Note that for each $a\in S$, the modular function $f_a$ is holomorphic on $\mathbb{H}$ with algebraic Fourier coefficients at all cusps. Hence, by Corollary \ref{useful} we get a polynomial $P(X)\in\overline{\mathbb{Q}}(j)[X]$, which is satisfied by $f_a$. Thus, $E_{N,\overline{\mathbb{Q}}}/\overline{\mathbb{Q}}(j)$ is an algebraic extension. Moreover, this extension is normal because each conjugate of $f_a$ is of the form $f_{a\gamma_i}$, which is equal to $f_b$ for some $b\in S$. Since $E_{N,\overline{\mathbb{Q}}}$ is obtained from $\overline{\mathbb{Q}}(j)$ by adjoining finite number of elements, $E_{N,\overline{\mathbb{Q}}}/\overline{\mathbb{Q}}(j)$ is a finite Galois extension.\\

Now we show that $E_{N,\overline{\mathbb{Q}}}=F_{N,\overline{\mathbb{Q}}}$. To begin with, observe that $\mathbb{C}$ and $F_{N,\overline{\mathbb{Q}}}$ are linearly disjoint over $\overline{\mathbb{Q}}$. Indeed, suppose $\{c_1,c_2,\ldots,c_r\}\subset \mathbb{C}$ is an arbitrary $\overline{\mathbb{Q}}$-linearly independent subset of $\mathbb{C}$. If there exists $g_i \in \mathcal{F}_{N,\overline{\mathbb{Q}}}$ such that $\sum_{i=1}^r g_i(\tau) \, c_i = 0$ for all $\tau \in \mathbb{H}$ with
$$g_i(\tau)=\sum_n d_{i n} \, \, q_N^n, \qquad d_{in \in \overline{\mathbb{Q}}} \quad \text{for }0\leq i\leq r,$$
then $$\sum_{i=1}^rc_i\,d_{i n}=0\quad \text{ for all } n.$$ The linear independence of $c_i$ over $\overline{\mathbb{Q}}$ implies that all $d_{i n}=0$ for all $0\leq i\leq r$, and hence, $f_1=f_2=\cdots=f_r=0$. Thus, we have $E_{N,\overline{\mathbb{Q}}}\subseteq \mathcal{F}_{N,\overline{\mathbb{Q}}} \subseteq \mathbb{C}E_{N,\overline{\mathbb{Q}}}.$ If possible, let $f\in \mathcal{F}_{N,\overline{\mathbb{Q}}}\setminus E_{N,\overline{\mathbb{Q}}}$. Since $f \in \mathbb{C}E_{N,\overline{\mathbb{Q}}}$, we get a $\overline{\mathbb{Q}}$-linearly independent subset $\{f_1,f_2,\ldots,f_m\}\subseteq E_{N,\overline{\mathbb{Q}}}$ such that 
\begin{equation}\label{rep.}
  f=\sum_{i=1}^m \alpha_i f_i, \quad \alpha_i\in \mathbb{C},  
\end{equation}
with at least one of the $\alpha_i \in \mathbb{C} \setminus \overline{\mathbb{Q}}$. Since $E_{N,\overline{\mathbb{Q}}}\subseteq \mathcal{F}_{N,\overline{\mathbb{Q}}}$, the set $\{f, \,f_1,f_2,\ldots,f_m\}$ is $\overline{\mathbb{Q}}$-linearly independent subset of $\mathcal{F}_{N,\overline{\mathbb{Q}}}$ and hence, $\mathbb{C}$-linearly independent. This contradicts \eqref{rep.}. Therefore, $E_{N,\overline{\mathbb{Q}}}= \mathcal{F}_{N,\overline{\mathbb{Q}}}$ and $\mathcal{F}_{N,\overline{\mathbb{Q}}}/\overline{\mathbb{Q}}(j)$ is a finite Galois extension.
\end{proof}

Thus, we have finally established that a modular function with algebraic Fourier coefficients satisfies a polynomial over $\overline{\mathbb{Q}}(j)$.

\bigskip

\section{\bf Proof of Main Results}
\bigskip

\subsection{Values of modular functions when $j(\tau) \not\in \overline{\mathbb{Q}}$}
Before proceeding with the proof of the main theorems, we establish the following intermediary observation.
\begin{lemma} \label{mdfz}
    Let $g$ be a non-constant modular function of level $N(\geq 1)$ with algebraic Fourier coefficients, i.e., $g\in  \mathcal{F}_{N,\overline{\mathbb{Q}}}$. If $\tau\in \mathbb{H}$ such that $j(\tau)\notin\overline{\mathbb{Q}}$, then $g(\tau)$ is non-zero.
\end{lemma}
\begin{proof}
By Theorem \ref{mdft..}, we know that $g$ satisfies a non-trivial polynomial over $\overline{\mathbb{Q}}(j)$. On clearing denominators, we can assume that $g$ satisfies the irreducible polynomial
$$
P(X)=\sum_{r=0}^m \bigg( \sum_{s=0}^{d_r}c_{r s}\,j^s \bigg) \,X^r \in \overline{\mathbb{Q}}[j][X], \quad m=\deg P(X), \, \, d_r=\deg c_r(j),
$$
Note that $c_0(j) \in \overline{\mathbb{Q}}[j]$ is not the identically zero polynomial. Now suppose $\tau \in \mathbb{H}$ is such that $j(\tau)$ is transcendental. If $g(\tau)$ is zero, then we get that $c_0(j(\tau))=0$ which contradicts the transcendence of $j(\tau)$. This completes the proof. 
\end{proof}

\begin{proof}[\bf Proof of Theorem \ref{mdf-values-thm}(a)]
If $g(\tau)=0$, then it follows from Lemma \ref{mdfz} that $j(\tau)\in \overline{\mathbb{Q}}$. If $\tau$ is a pole of $g$, we consider the modular function
$$
h(z):=\frac{1}{g(z)}\in \mathcal{F}_{N,\overline{\mathbb{Q}}},
$$
which vanishes at $z=\tau$, and so $j(\tau)\in\overline{\mathbb{Q}}$ by Lemma \ref{mdfz}.    
\end{proof}

The proof of Theorem \ref{mdf-values-thm}(b) is established in two parts. We first show that if $j(\tau)$ is transcendental, then $g(\tau)$ is as well. The converse implication is proved later.
\begin{prop}\label{tranmf}
 Let $g\in \mathcal{F}_{N,\overline{\mathbb{Q}}}$ and $\tau\in \mathbb{H}$ be such that it is not a pole of $g$. If $j(\tau)$ is transcendental, then $g(\tau)$ is transcendental.
\end{prop}
\begin{proof}
Suppose that $g(\tau)\in\overline{\mathbb{Q}}$. Then the modular function $$h(z):=g(z)-g(\tau)$$ belongs to $ \mathcal{F}_{N,\overline{\mathbb{Q}}}$ and vanishes at $\tau$. This contradicts Lemma \ref{mdfz}. Therefore, $g(\tau)$ is transcendental.
\end{proof}


\subsection{Values of modular functions when $j(\tau)\in\overline{\mathbb{Q}}$}
From Theorem \ref{mdft..}, we have $$ \mathcal{F}_{N,\overline{\mathbb{Q}}}=  \overline{\mathbb{Q}}\left(\{j,f_{\boldsymbol{a}} : {\boldsymbol{a}}\in S\}\right), \text{ where }S=\left\{\left(\frac{r}{N},\frac{s}{N}\right): 0\leq r, s\leq N-1 \text{ and }(r,s)\neq (0,0)\right\}.$$
Hence, for any $g\in \mathcal{F}_{N,\overline{\mathbb{Q}}}$, the algebraic nature of $g(\tau)$ is determined by the numbers in the set $\{f_{\boldsymbol{a}}(\tau):{\boldsymbol{a}}\in S\}$.\\

For $N=1$, $\mathcal{F}_{1,\overline{\mathbb{Q}}} = \overline{\mathbb{Q}}(j)$. Hence, for any $g \in \mathcal{F}_{1,\overline{\mathbb{Q}}}$, $j(\tau) \in \overline{\mathbb{Q}}$ implies that $g(\tau) \in \overline{\mathbb{Q}}$. For $N>1$, we have to study the nature of the values $f_{\boldsymbol{a}}(\tau)$ for all ${\boldsymbol{a}}\in S$. We start by proving the following important lemma.
\begin{lemma}\label{omegatau}
 If $\tau\in \mathbb{H}$ such that $j(\tau)$ is algebraic, then there exists a unique transcendental number (up to an algebraic multiples) $\omega_\tau$ for which $g_2(\omega_\tau L_\tau)$ and $g_3(\omega_\tau L_\tau)$ are both algebraic numbers.   
\end{lemma}
\begin{proof}
Note that for any $\alpha \in\mathbb{C}^\times $,
\begin{equation}\label{omegatau.}
    j(\tau)=\frac{E_4(\tau)^3}{\Delta(\tau)}=\frac{1728\,  g_2(L_\tau)^3}{\Delta_0(L_\tau)}=\frac{1728\, g_2\left(\alpha L_\tau\right)^3}{\Delta_0\left(\alpha L_\tau\right)},
\end{equation}
because of the homogeneity properties of the $g_2$, $g_3$ functions. Here $L_\tau=\tau\,\mathbb{Z}\oplus\mathbb{Z}$ and $\alpha L_\tau= \alpha \, \tau\,\mathbb{Z}\oplus\alpha \, \mathbb{Z}$. If we choose $\omega_\tau$ such that $\omega_\tau^4=g_2(L_\tau)$, then $g_2(\omega_\tau L_\tau)=1$. Given that $j(\tau)$ is algebraic, from \eqref{omegatau.}, we get that $\Delta_0(\omega_\tau L_\tau)$ is also algebraic. Thus, the numbers $g_2(\omega_\tau L_\tau)$ and $g_3(\omega_\tau L_\tau)$ are both algebraic. Moreover, $\omega_\tau$ is a period of $\omega_\tau L_\tau$ and both $g_2(\omega_\tau L_\tau)$ and $g_3(\omega_\tau L_\tau)$ are algebraic. By Corollary \ref{tranp} we conclude that $\omega_\tau$ is transcendental.\\

To prove uniqueness, consider an arbitrary $\omega_\tau'\in\mathbb{C}^\times$ such that both $g_2(\omega_\tau' L_\tau)$ and $g_3(\omega_\tau' L_\tau)$ are algebraic. By homogeneity of $g_2$, we have 
$$
 \quad g_2(\omega_\tau' L_\tau)= (\omega_\tau')^{-4}\,g_2(L_\tau)=:\beta \in \overline{\mathbb{Q} }.
$$
Then, we get $g_2(L_\tau)=\beta \cdot (\omega_\tau')^{4}$ and hence $\omega_\tau^4=\beta \cdot (\omega_\tau')^{4}.$
This completes the proof that $\omega_\tau$ is unique up to algebraic multiples.
\end{proof}

Now recall that
$$
f_{\boldsymbol{a}}(\tau):=\frac{g_2(L_\tau)\,g_3(L_\tau)}{\Delta_0(L_\tau)}\,\wp_{L_\tau}(a_1\tau+a_2),
$$ 
where $L_\tau=\tau\,\mathbb{Z}\oplus\mathbb{Z}$ and ${\boldsymbol{a}}=(a_1,a_2)\in S$. Using homogeneity properties of the functions involved, we can rewrite this as follows:
$$
f_{\boldsymbol{a}}(\tau)=\frac{g_2(\omega_\tau L_\tau)\, g_3(\omega_\tau L_\tau)}{\Delta_0(\omega_\tau L_\tau)}\,\wp_{\omega_\tau L_\tau}(a_1\omega_\tau\tau+a_2\omega_\tau).
$$
\begin{lemma}\label{algfa}
    If $\tau\in\mathbb{H}$ is such that $j(\tau)\in\overline{\mathbb{Q}}$, then $f_{\boldsymbol{a}}(\tau)\in\overline{\mathbb{Q}}$ for all ${\boldsymbol{a}}\in S$.
\end{lemma}
\begin{proof}
    Let ${\boldsymbol{a}}=\left( \frac{r}{N},\frac{s}{N} \right)\in S$. Then we have
$$
f_{\boldsymbol{a}}(\tau)=\frac{g_2(\omega_\tau L_\tau)\, g_3(\omega_\tau L_\tau)}{\Delta_0(\omega_\tau L_\tau)}\,\wp_{\omega_\tau L_\tau}\left(\frac{r\,\omega_\tau\,\tau}{N}+\frac{s\,\omega_\tau}{N}\right),
$$
where $\omega_\tau$ is chosen as in Lemma \ref{omegatau} so that $g_2(\omega_\tau L_\tau),\, g_3(\omega_\tau L_\tau)\in\overline{\mathbb{Q}}$.\\

Recall that both $\wp_{\omega_\tau L_\tau}\left(\frac{\omega_\tau\, \tau}{N} \right)$ and $\wp_{\omega_\tau L_\tau}\left(\frac{\omega_\tau}{N} \right)$  are algebraic by Corollary \ref{tranp}. Moreover, by the addition formula for Weierstrass $\wp$-function in Theorem \ref{addf}, we get that $$\wp_{\omega_\tau L_\tau}\left(\frac{r\,\omega_\tau\, \tau}{N}+\frac{s\,\omega_\tau}{N}\right)\in\overline{\mathbb{Q}}$$ for all ${\boldsymbol{a}}=\left(\frac{r}{N},\frac{s}{N}\right)\in S$. Thus, the number $f_{\boldsymbol{a}}(\tau)$ is algebraic for all ${\boldsymbol{a}}\in S$.
\end{proof}
\begin{prop}\label{algemf}
Let $g\in \mathcal{F}_{N,\overline{\mathbb{{Q}}}}$ and $\tau\in\mathbb{H}$ such that $j(\tau)\in \overline{\mathbb{Q}}$. If $\tau$ is not a pole of $g$, then $g(\tau)$ is algebraic.
\end{prop}
\begin{proof}
 From Lemma \ref{algfa}, we have   $f_{\boldsymbol{a}}(\tau)\in\overline{\mathbb{Q}}$ for all ${\boldsymbol{a}}\in S$. Since 
 $$g \in \mathcal{F}_{N,\overline{\mathbb{Q}}}=  \overline{\mathbb{Q}}\left(\{j,f_{\boldsymbol{a}} |{\boldsymbol{a}}\in S\}\right)$$
 by Theorem \ref{mdft..}, $g(\tau)$ is also algebraic. 
\end{proof}
\bigskip

Propositions \ref{tranmf} and \ref{algemf} together complete the proof of Theorem \ref{mdf-values-thm}(b). We now prove the generalization of \cite[Theorem 2.4]{jeon-kang-kim}.
\begin{proof}[\bf Proof of Theorem \ref{zero-pole-value}]
Let $f$ be a meromorphic modular form of weight $k \in \mathbb{Z}$. Define 
   $$
h(\tau):=\frac{f(\tau)^{12}}{\Delta(\tau)^k}\quad \forall \tau\in\mathbb{H}.
   $$
The function $h \in \mathcal{F}_{N,\overline{\mathbb{Q}}}$ because the Fourier coefficients of $f$ and $\Delta$ are algebraic. Moreover, any $\tau\in\mathbb{H}$ is a zero or a pole of $h$ if and only if it is a zero or a pole of $f$ because $\Delta$ is non-vanishing and holomorphic on $\mathbb{H}$. Thus, if $\tau\in\mathbb{H}$ is a zero or pole of $f$, then $h$ is non-constant and by Theorem \ref{mdf-values-thm}(a), we deduce that $j(\tau)\in\overline{\mathbb{Q}}$, and by Theorem \ref{mdf-values-thm}(b), that either $g(z)$ has a pole at $z = \tau$ or that $g(\tau)\in\overline{\mathbb{Q}}$.
\end{proof}
\bigskip

\subsection{Values of modular forms} Recall that for any integer $k$ and $N$, we denote $M_{k,N,\overline{\mathbb{Q}}}$, $M_{k,N,\overline{\mathbb{Q}}}^w$ and $M_{k,N,\overline{\mathbb{Q}}}^m$ to be the set of holomorphic, weakly holomorphic and meromorphic modular forms respectively with algebraic Fourier coefficients at all cusps. We first establish a lemma which is crucial to the proof of Theorem \ref{general-mero-values}.

\begin{lemma}\label{det}
Let $n$ and $m$ be positive integers such that $n \geq 2m$, then the determinant of the matrix $A = [a_{ij}]_{(m+1) \times (m+1)}$ with $a_{ij} = {n+1-j \choose i-1}$
%     $$
% A=\left[\begin{array}{cccccc}
%     {{n}\choose{0}} & {{n-1}\choose{0}} & {{n-2}\choose{0}} & \ldots & {{n-m+1}\choose{0}} & {{n-m}\choose{0}} \\\\
%     {{n}\choose{1}} & {{n-1}\choose{1}} & {{n-2}\choose{1}} & \ldots & {{n-m+1}\choose{1}} & {{n-m}\choose{1}}\\\\
%     {{n}\choose{2}} & {{n-1}\choose{2}} & {{n-2}\choose{2}} & \ldots & {{n-m+1}\choose{2}} & {{n-m}\choose{2}}\\\\
%     \vdots & \vdots & \vdots & \ldots & \vdots & \vdots \\\\
%     {{n}\choose{m}} & {{n-1}\choose{m}} & {{n-2}\choose{m}} & \ldots & {{n-m+1}\choose{m}} & {{n-m}\choose{m}}
% \end{array} \right]_{(m+1)\times (m+1)}
%   $$ 
  is given by $(-1)^{\frac{m(m+1)}{2}}$.
\end{lemma}
\begin{proof}
We will use the identity ${n\choose{i}} - {n-1\choose{i}} ={n-1\choose{i-1}}$ for $1 \leq j \leq m$. On applying the column operation $C_j \mapsto C_j - C_{j+1}$ successively for $1 \leq j \leq m$, the matrix $A$ transforms into 
$$
A'=\left[\begin{array}{cccccc}
    0 & 0 & 0 & \ldots & 0 & {{n-m}\choose{0}} \\\\
    {{n-1}\choose{0}} & {{n-2}\choose{0}} & {{n-3}\choose{0}} & \ldots & {{n-m}\choose{0}} & {{n-m}\choose{1}}\\\\
    {{n-1}\choose{1}} & {{n-2}\choose{1}} & {{n-3}\choose{1}} & \ldots & {{n-m}\choose{1}} & {{n-m}\choose{2}}\\\\
    \vdots & \vdots & \vdots & \ldots & \vdots & \vdots \\\\
    {{n-1}\choose{m-1}} & {{n-2}\choose{m-1}} & {{n-3}\choose{m-1}} & \ldots & {{n-m}\choose{m-1}} & {{n-m}\choose{m}}
\end{array} \right]_{(m+1)\times (m+1)}.
$$
We can now repeat the same column operations as above, repeatedly, to reduce the matrix to a lower triangular form 
 $$
B=\left[\begin{array}{cccccc}
    0 & 0 & 0 & \ldots & 0 & {{n-m}\choose{0}} \\\\
    0 & 0 & 0 & \ldots & {{n-m}\choose{0}} & *\\\\
    0 & 0 & {{n-m}\choose{0}} & \iddots & * & *\\\\
    \vdots & \iddots & \vdots & \ldots & \vdots & \vdots \\\\
    {{n-m}\choose{0}} & * & * & \ldots & * & *
\end{array} \right]_{(m+1)\times (m+1)}.
 $$
Since $B$ is obtained from $A$ by elementary column operations that do not change the determinant, we conclude that $\det(A)=\det(B)=(-1)^{\frac{m(m+1)}{2}}$.
\end{proof}



\begin{prop}\label{algem'}
   Let $f$ be a non-constant meromorphic modular form with algebraic Fourier coefficients. If $\tau\in\mathbb{H}$ such that $e^{2\pi i\tau}$ is algebraic, then $f(\tau)$ is transcendental. 
\end{prop}
\begin{proof}[Proof of Theorem \ref{general-mero-values}(a)]
Consider the modular function
$$g(\tau):=\frac{f(\tau)^{12}}{\Delta(\tau)^k}\quad \forall\, \tau\in\mathbb{H},$$
which lies in $\mathcal{F}_{N,\overline{\mathbb{Q}}}$. By Theorem \ref{mdft..}, we know that $g$ has the minimal polynomial $P[X]\in \overline{\mathbb{Q}}(j)[X]$, and we can clear denominators of the coefficients of $P[X]$ to get a polynomial 
$$P(X)=\sum_{r=0}^m c_r(j)\,X^r\in \overline{\mathbb{Q}}[j][X],\quad \deg P(X)=m$$
such that
$P(g(\tau))=0$ for any $\tau\in \mathbb{H}$. Here, we have
$$
c_r(j)=\sum_{s=0}^{d_r}c_{r,s}\,j^s\in \overline{\mathbb{Q}}[j]\quad \forall\ 0 \leq r\leq m
$$
and $d_r=\deg c_r(j)$ so that $c_{r,d_r}\neq 0$ for all $0\leq r\leq m$. Moreover, Corollary \ref{useful} shows that if $f$ is holomorphic on $\mathbb{H}$, then the minimal polynomial $P[X]$ of $g$ is monic and belongs to $\overline{\mathbb{Q}}[j][X]$, and hence $c_m(j)=1$.\\

For any $\tau\in\mathbb{H}$, we have $$P\left(\frac{f(\tau)^{12}}{\Delta(\tau)^k}\right)=0,$$
and hence for any positive integer $l$, 
$$
\begin{aligned}
& \Delta(\tau)^l\, P\left(\frac{f(\tau)^{12}}{\Delta(\tau)^k}\right)=0\\
\Rightarrow &\sum_{r=0}^m\sum_{s=0}^{d_r}c_{r,s}\,j(\tau)^s f(\tau)^{12 r}\Delta(\tau)^{l-k r}=0\\
\Rightarrow &\sum_{r=0}^m\sum_{s=0}^{d_r}c_{r,s}\,E_4(\tau)^{3 s}f(\tau)^{12\, r}\Delta(\tau)^{l-k r-s}=0.
\end{aligned}
$$
We choose $l> k m +2\cdot\operatorname{max}\{d_0,d_1,\ldots,d_m\}$ so that the powers of $\Delta(\tau)$ are always positive integers. Then we put $$\Delta(\tau)=\frac{E_4(\tau)^3-E_6(\tau)^2}{1728}$$ to obtain \begin{equation} \label{trane}
\sum_{r=0}^m \sum_{s=0}^{d_r} \sum_{t=0}^{l-k r-s} \frac{(-1)^t c_{r, s}}{1728^{l-k r-s}}\left(\begin{array}{c}
l-k\,r-s \\
t
\end{array}\right) f(\tau)^{12\, r} E_4(\tau)^{3(l-k\, r-t)} E_6(\tau)^{2\, t}=0.   
\end{equation}
Now we consider $\tau\in \mathbb{H}$ such that $e^{2\pi i\tau}$ is algebraic. By Theorem \ref{schneider}, the number $j(\tau)$ is transcendental, and hence Theorem \ref{zero-pole} guarantees that $\tau$ is neither a zero nor a pole of $f$. Moreover, we claim that the number $f(\tau)$ is not algebraic. For otherwise, \eqref{trane} becomes a polynomial in two variables over $\overline{\mathbb{Q}}$, satisfied by $E_4(\tau)$  and $E_6(\tau)$. If we can prove that this is a non-zero polynomial of $E_4(\tau)$ and $E_6(\tau)$, then we would get a contradiction with Nesterenko's Theorem \ref{nest}. Hence, it would complete the proof that $f(\tau)$ is transcendental.\\

To show that the left side of \eqref{trane} becomes a non-zero polynomial of $E_4(\tau)$ and $E_6(\tau)$ when $e^{2\pi i\tau}\in\overline{\mathbb{Q}}$, we consider the following cases:\\

\emph{Case-$1$: } If $f\in M_{k,N,\overline{\mathbb{Q}}}\,\cup M_{k,N,\overline{\mathbb{Q}}}^w$, then $c_m(j)=1$, and hence the coefficient of $E_4(\tau)^{3(l-k m)}\,E_6(\tau)^{0}$ (which appears in \eqref{trane} only for $r=m,\, t=0$) is given by $$\frac{1}{1728^{l-k m}}\cdot 1\cdot f(\tau)^{12 m}\neq 0,$$
which proves that the polynomial is non-zero.\\

\emph{Case-$2$: } If $f\in M_{k,N,\overline{\mathbb{Q}}}^m\setminus \left(M_{k,N,\overline{\mathbb{Q}}}\,\cup M_{k,N,\overline{\mathbb{Q}}}^w\right)$, then the leading coefficient of $P[X]$ is given by $$c_m(j)=\sum_{s=0}^{d_m}c_{m,s}\,j^s.$$
Moreover, it cannot be \emph{identically zero polynomial} because the degree of $P[X]$ is same as the degree of the minimal polynomial of $g$. Since $l>k\, m+2\, d_m$, the following terms
$$
E_4(\tau)^{3(l-k\, m-t)}\, E_6(\tau)^{2\, t} \quad \forall\, 0\leq t\leq d_m
$$
with coefficients
$$
\frac{(-1)^t\, f(\tau)^{12 m}}{1728^{l-km}}\sum_{s=0}^{d_m}\left(\begin{array}{c}
l-k\,m-s \\
t
\end{array}\right) (1728)^s\, c_{m,s} \quad \forall\, 0\leq t\leq d_m
$$
appear in \eqref{trane}. We claim that at least one of these terms has non-zero coefficient. If not, we consider the following system of equations: $ A\, X=0$, where 
$$
A=\left[\begin{array}{ccccc}
    {{l-km}\choose{0}} & {{l-km-1}\choose{0}} & {{l-km-2}\choose{0}} & \ldots & {{l-km-d_m}\choose{0}} \\\\
    {{l-km}\choose{1}} & {{l-km-1}\choose{1}} & {{l-km-2}\choose{1}} & \ldots & {{l-km-d_m}\choose{1}}\\\\
    {{l-km}\choose{2}} & {{l-km-1}\choose{2}} & {{l-km-2}\choose{2}} & \ldots & {{l-km-d_m}\choose{2}}\\\\
    \vdots & \vdots & \vdots & \ldots & \vdots \\\\
    {{l-km}\choose{d_m}} & {{l-km-1}\choose{d_m}} & {{l-km-2}\choose{d_m}} & \ldots & {{l-km-d_m}\choose{d_m}}
\end{array} \right]_{(d_m+1)\times (d_m+1)}  X=\left[\begin{array}{c}
   x_0  \\
    x_1 \\
    x_2 \\
    \vdots\\
    x_{d_m}
\end{array}\right]_{d_m+1.}
$$
Then this system of equations is satisfied by the non-trivial vector $$\left[\begin{array}{c}
   (1728)^0\, c_{m,0} \\
    (1728)^1\, c_{m,1} \\
    (1728)^2\, c_{m,2} \\
    \vdots\\
    (1728)^{d_m}\, c_{m,d_m}
\end{array}\right]_{(d_m+1)\times 1}.$$ From Lemma \ref{det}, we get that $\det(A)\neq 0$, which contradicts the fact that the above system has only trivial solution. Thus, it completes the proof that the polynomial relation of $E_4(\tau)$ and $E_6(\tau)$ in \eqref{trane} is non-trivial.
\end{proof}
\begin{corollary}\label{algi}
If $f\in M_{k,N,\overline{\mathbb{Q}}}\,\cup M_{k,N,\overline{\mathbb{Q}}}^w$, then $f$ is algebraically dependent with $E_4$ and $E_6$ over $\overline{\mathbb{Q}}$. In particular, If $f\in M_{k,N,\overline{\mathbb{Q}}}^m$ and $\tau\in\mathbb{H}$ is not a pole of $f$, then $f(\tau)$ is algebraic over $\overline{\mathbb{Q}}(E_4(\tau),E_6(\tau))$. 
\end{corollary}
\begin{proof}
    It directly follows from the equation \eqref{trane} because the left-hand side is a non-zero polynomial in $f, E_4$ and $E_6$ over $\overline{\mathbb{Q}}.$
\end{proof}
\bigskip

Now we consider the points $\tau\in\mathbb{H}$ such that $j(\tau)\in\overline{\mathbb{Q}}$. From Lemma \ref{omegatau}, we get a transcendental number $\omega_\tau$ such that $g_2(\omega_\tau L_\tau)$ and $g_3(\omega_\tau L_\tau)$ are both algebraic. Moreover, we have the formulae
$$
\begin{aligned}
E_4(\tau)&=\frac{3}{4\pi^4}\, g_2(L_\tau)=\frac{3}{4}\left(\frac{\omega_\tau}{\pi} \right)^4 g_2(\omega_\tau L_\tau),\\ E_6(\tau)&=\frac{27}{8\pi^6}\,g_3(L_\tau)=\frac{27}{8}\left(\frac{\omega_\tau}{\pi}\right)^6g_3(\omega_\tau L_\tau).
\end{aligned}
$$
From Nesterenko's Theorem \ref{nest}, we know that at most one of $E_4(\tau)$ and $E_6(\tau)$ is algebraic. The above formulae imply that if $\frac{\omega_\tau}{\pi}\in\overline{\mathbb{Q}}$, then both $E_4(\tau)$ and $E_6(\tau)$ are algebraic, which is a contradiction. Hence, $\frac{\omega_\tau}{\pi}$ is a transcendental number. Besides, we have 
$$
\begin{aligned}
\Delta(\tau)=\frac{E_4(\tau)^3-E_6(\tau)^2}{1728}&=\frac{1}{4^6}\left(\frac{\omega_\tau}{\pi}\right)^{12}\left(g_2(\omega_\tau L_\tau)^3-27 g_3(\omega_\tau L_\tau)^2 \right)\\
&=\frac{1}{4^6}\left(\frac{\omega_\tau}{\pi}\right)^{12}\Delta_0(\omega_\tau L_\tau).  
\end{aligned}
$$
As $g_2(\omega_\tau L_\tau)$ and $g_3(\omega_\tau L_\tau)$ both are algebraic, the number $\Delta_0(\omega_\tau L_\tau)\in\overline{\mathbb{Q}}\setminus \{0\}$. Since $\frac{\omega_\tau}{\pi}$ is transcendental and $\Delta(\tau)$ is a non-zero algebraic multiple of $\left(\frac{\omega_\tau}{\pi}\right)^{12}$, we deduce that $\Delta(\tau)$ is transcendental. We use this observation to prove the following theorem.
\begin{theorem}\label{algem}
Let $f$ be a non-constant meromorphic modular form with algebraic Fourier coefficients. If $\tau\in\mathbb{H}$ such that  $j(\tau)\in \overline{\mathbb{Q}}$ and $\tau$ is not a pole of $f$, then there exists a transcendental number $\omega_\tau$ (which depends only on $\tau$ and is $\overline{\mathbb{Q}}$-linearly independent with $\pi$) such that $\left(\frac{\pi}{\omega_\tau}\right)^k f(\tau)\in \overline{\mathbb{Q}}$. Therefore, $f(\tau)$ is either zero or transcendental.
\end{theorem}
\begin{proof}
     We first construct the function $$g(\tau):=\frac{f(\tau)^{12}}{\Delta(\tau)^k}\quad \forall \tau\in\mathbb{H},$$ which lies in $ F_{N,\overline{\mathbb{Q}}}$ by the same argument as in Theorem \ref{algem'}. Since $j(\tau)\in\overline{\mathbb{Q}}$, Theorem \ref{algemf} implies that $g(\tau)$ is an algebraic number $c$, (say). Thus, we get  \begin{equation}\label{forms}
f(\tau)^{12}=c.\Delta(\tau)^k=\frac{c}{4^6}.\Delta_0(\omega_\tau L_\tau)^k\left(\frac{\omega_\tau}{\pi} \right)^{12 k}. 
\end{equation}
This shows that $f(\tau)$ is zero when $c=0$, and otherwise, it a non-zero algebraic multiple of $\left(\frac{\omega_\tau}{\pi} \right)^{k}$ and hence transcendental.
\end{proof}
\bigskip

\subsection{Proof of Theorem \ref{deg.}}\, \bigskip

 By theorem \ref{mdft..}, the numbers $g(\tau)$ and $j(\tau)$ are algebraically dependent over $\overline{\mathbb{Q}}$. Moreover, $j(\tau)$ is algebraic if and only if $g(\tau)\in\overline{\mathbb{Q}}$ by Theorem \ref{suma}. These facts together give the first equality below:
$$
\begin{aligned}
 trdeg\,\overline{\mathbb{Q}}\left(e^{2\pi i\tau},\, f(\tau),\, g(\tau) \right)&= trdeg \,\overline{\mathbb{Q}}\left(e^{2\pi i\tau},\, f(\tau),\, j(\tau) \right)\\
 &= trdeg \,\overline{\mathbb{Q}}\left(e^{2\pi i\tau},\, \frac{f(\tau)^{12}}{\Delta(\tau)^k},\, f(\tau),\, j(\tau) \right)\\
 &= trdeg \,\overline{\mathbb{Q}}\left(e^{2\pi i\tau},\, \Delta(\tau)^k,\, f(\tau),\, j(\tau) \right)\\
 &= trdeg \,\overline{\mathbb{Q}}\left(e^{2\pi i\tau},\, \Delta(\tau),\, f(\tau),\, j(\tau) \right) \\
 &= trdeg \,\overline{\mathbb{Q}}\left(e^{2\pi i\tau},\, E_4(\tau)^3,\, E_6(\tau)^2,\, f(\tau),\, j(\tau) \right)\\
 &=trdeg \,\overline{\mathbb{Q}}\left(e^{2\pi i\tau},\, E_4(\tau),\, E_6(\tau),\, f(\tau),\, j(\tau) \right)\\
 &= trdeg \,\overline{\mathbb{Q}}\left(e^{2\pi i\tau},\, E_4(\tau),\, E_6(\tau),\, f(\tau) \right)\\
 &= trdeg \,\overline{\mathbb{Q}}\left(e^{2\pi i\tau},\, E_4(\tau),\, E_6(\tau)\right).
\end{aligned}
$$ 
Here, the second equality follows from the fact that $f^{12}/\Delta^k\in F_{N,\overline{\mathbb{Q}}}$ and  $f(\tau)^{12}/\Delta(\tau)^k$ is algebraically dependent with $j(\tau)$ by Theorem \ref{mdft..}. The third and the seventh equality hold as the fields are identical. To prove the equality of fields in the third step, one needs the assumption that $f(\tau)\neq 0$.  
The final equality holds because $f(\tau)$ is algebraic over $\overline{\mathbb{Q}}\left(E_4(\tau), E_6(\tau)\right)$ by Corollary \ref{algi}. Then we apply Nesterenko's Theorem \ref{nest} on the right side of the final equality to conclude that the transcendence degree of the field $\overline{\mathbb{Q}}\left(e^{2\pi i\tau},\, f(\tau),\, g(\tau) \right)$ over $\overline{\mathbb{Q}}$ is at least two. 
\bigskip

\subsection{Values of quasi-modular forms}\,\bigskip

Let $\widetilde{M}_{k,\overline{\mathbb{Q}}}^{(p)}(\Gamma)$ denote the set of all quasi-modular forms of weight $k$ and depth $p(>0)$ for $\Gamma$ with algebraic Fourier coefficients. We study their values at the points $\tau\in\mathbb{H}$, where one of $e^{2\pi i\tau}$ and $j(\tau)$ is algebraic.
\begin{theorem}
Let $\Gamma$ be a congruence subgroup of level $N(\geq 1)$ and $\widetilde{f}\in \widetilde{M}_{k,\overline{\mathbb{Q}}}^{(p)}(\Gamma)$. If $\tau\in\mathbb{H}$ such that $e^{2\pi i\tau}$ is algebraic, then $f(\tau)$ is transcendental. 
\end{theorem}
\begin{proof}
   By Theorem \ref{quasistructure}, we can write $\widetilde{f}$ in the form $$\widetilde{f}=\sum_{r=0}^p g_r E_2^r,\quad g_r\in M_{k,N,\overline{\mathbb{Q}}}.$$ From Corollary \ref{algi}, we know that each number $g_r(\tau)$ is algebraically dependent with $E_4(\tau), E_6(\tau)$. Therefore, we get the following equality:
 $$ trdeg\,\overline{\mathbb{Q}}(g_0(\tau), g_2(\tau),\ldots, g_p(\tau), E_4(\tau), E_6(\tau))=trdeg\,\overline{\mathbb{Q}}( E_4(\tau), E_6(\tau)).$$
If possible, let both $e^{2\pi i\tau}, \widetilde f(\tau)\in \overline{\mathbb{Q}}$. Since 
$\widetilde f(\tau)=\sum_{r=0}^p g_r(\tau) E_2(\tau)^r$ and $g_r(\tau)\neq 0$ (by Theorem \ref{algem'}), the number $E_2(\tau)$ is algebraic over $\overline{\mathbb{Q}}\left(g_0(\tau), g_1(\tau),\ldots, g_p(\tau), E_4(\tau), E_6(\tau)\right)$ and so
$$
\begin{aligned}
&\quad trdeg\, \overline{\mathbb{Q}}\left(E_2(\tau), g_0(\tau), g_1(\tau),\ldots, g_p(\tau), E_4(\tau), E_6(\tau)\right)\\ &=trdeg\,\overline{\mathbb{Q}}(g_0(\tau), g_1(\tau),\ldots, g_p(\tau), E_4(\tau), E_6(\tau))\\ 
&=trdeg\,\overline{\mathbb{Q}}(E_4(\tau), E_6(\tau)).
\end{aligned}
$$
Since $e^{2\pi i\tau}$ is an algebraic number, the following field containment holds
$$
\overline{\mathbb{Q}}(e^{2\pi i\tau},\, E_2(\tau),\,E_4(\tau), E_6(\tau))\subseteq \overline{\mathbb{Q}}\left(E_2(\tau), g_0(\tau), g_1(\tau),\ldots, g_p(\tau), E_4(\tau), E_6(\tau)\right).
$$
Therefore, we get 
$$
\begin{aligned}
 &\quad trdeg\, \overline{\mathbb{Q}}(e^{2\pi i\tau},\, E_2(\tau),\,E_4(\tau),\, E_6(\tau))\\ &\leq trdeg\,\overline{\mathbb{Q}}\left(E_2(\tau),\, g_0(\tau),\, g_1(\tau),\ldots, g_p(\tau), E_4(\tau),\, E_6(\tau)\right)\\
  &= trdeg\,\overline{\mathbb{Q}}(E_4(\tau),\, E_6(\tau))\\
  &\leq 2,
\end{aligned}
$$
which contradicts Nesterenko's Theorem \ref{nest}. Hence, $\widetilde{f}(\tau)$ is a transcendental number. 
\end{proof}
To study values of quasi-modular forms in the complementary case (i.e., for $j(\tau)\in\overline{\mathbb{Q}}$), we need the following lemma.
\begin{lemma}\label{peri}
 If $\tau\in \mathbb{H}$ such that $j(\tau)$ is algebraic, then there exists a unique (up to algebraic multiple) transcendental number $\omega_\tau$ for which $\frac{\omega_\tau}{\pi}$ and $\frac{\eta_2}{\pi}$ are algebraically independent over $\overline{\mathbb{Q}}$. Here, $\eta_2$ is the second quasi-period with respect to the lattice $\omega_\tau L_\tau$ 
\end{lemma}
\begin{proof}
    We have already proved existence of $\omega_\tau$ in Lemma \ref{omegatau}.  From the definition of the Weierstrass zeta-function, one gets
 $   
Z_{\omega_\tau L_\tau}(z)= \frac{1}{\omega_\tau}\cdot Z_{L_\tau}\left(\frac{z}{\omega_\tau}\right).
$ Then we use the definition of quasi-period and the identity \eqref{quasi-period} obtain
$$
\eta_2=Z_{\omega_\tau L_\tau}(z+\omega_\tau)-Z_{\omega_\tau L_\tau}(z)= \frac{1}{\omega_\tau}\,\eta_2(L_\tau)=\frac{1}{\omega_\tau}\,G_2(\tau)= \frac{1}{3}\cdot\frac{\pi^2}{ \omega_\tau}\,E_2(\tau).
$$
Thus, we get the following formulae 
$$
E_2(\tau)=3\,\frac{\omega_\tau}{\pi}\frac{\eta_2}{\pi}, \quad E_4(\tau)=\frac{3}{4}\left(\frac{\omega_\tau}{\pi} \right)^4 g_2(\omega_\tau L_\tau),\quad E_6(\tau)=\frac{27}{8}\left(\frac{\omega_\tau}{\pi}\right)^6g_3(\omega_\tau L_\tau).$$
The above formulae imply that $E_2(\tau)$, $E_4(\tau)$ and $E_6(\tau)$ are algebraic over $\overline{\mathbb{Q}}\left(\frac{\omega_\tau}{\pi}, \frac{\eta_2}{\pi}\right)$. If possible, let $\frac{\omega_\tau}{\pi}, \frac{\eta_2}{\pi}$ be algebraically dependent. Then we get that $$trdeg\,\overline{\mathbb{Q}}(E_2(\tau),E_4(\tau),E_6(\tau))=1,$$ which contradicts Nesterenko's Theorem \ref{nest}. Hence, the numbers $\frac{\omega_\tau}{\pi}$ and $\frac{\eta_2}{\pi}$ are algebraically independent over $\overline{\mathbb{Q}}$.
\end{proof}
Lemma \ref{peri} and Theorem \ref{quasistructure} together allow us to describe values of quasi-modular forms explicitly at the points $\tau\in\mathbb{H}$ where $j$-function is algebraic.
\begin{theorem}\label{tranq}
 Let $\widetilde{f}\in \widetilde{M}_{k,\overline{\mathbb{Q}}}^{(p)}(\Gamma)$ and $\tau\in\mathbb{H}$ such that $j(\tau)\in\overline{\mathbb{Q}}$, then $f(\tau)$ is a polynomial in $\frac{\omega_\tau}{\pi}$ and $\frac{\eta_2}{\pi}$ with algebraic coefficients. Therefore, $f(\tau)$ is either zero or transcendental.   
\end{theorem}
\begin{proof}
    From Theorem \ref{quasistructure}, we have the expression $$\widetilde{f}(\tau)=\sum_{r=0}^p g_r(\tau) E_2(\tau)^r,\quad \text{where } g_r\in M_{k,N,\overline{\mathbb{Q}}}.$$
Then we apply Theorem \ref{algem} for each coefficient $g_r$, and the above formula for $E_2$ such that 
\begin{equation}\label{last}
\widetilde{f}(\tau)=\sum_{r=0}^p c_r\left(\frac{\omega_\tau}{\pi}\right)^{k-2 r}\left(\frac{\omega_\tau}{\pi}\cdot\frac{\eta_2}{\pi}\right)^r=\sum_{r=0}^p\, c_r\, \left(\frac{\omega_\tau}{\pi}\right)^{k-r}\, \left(\frac{\eta_2}{\pi}\right)^r,
 \end{equation}
where each $c_r$ is an algebraic constant. If $\widetilde{f}(\tau)\in\overline{\mathbb{Q}}\setminus \{0\}$, then \eqref{last} shows that $\frac{\omega_\tau}{\pi}$ and $\frac{\eta_2}{\pi}$ are algebraically dependent, which contradicts Lemma \ref{peri}. Therefore, the number $\widetilde{f}(\tau)$ is either zero or transcendental.     
\end{proof}
\bigskip

\subsection{Proof of Theorem \ref{deg..}}
\,\bigskip

By the same argument as in Theorem \ref{deg.}, we get
\begin{equation}\label{final.}
trdeg\,\overline{\mathbb{Q}}\left( e^{2\pi i\tau},\, f(\tau),\, g(\tau),\, \widetilde f(\tau)\right)= trdeg\,\overline{\mathbb{Q}}\left( e^{2\pi i\tau},\, E_4(\tau),\, E_6(\tau),\, \widetilde f(\tau)\right).
\end{equation}
Moreover, we can apply Theorem \ref{quasistructure} to write
$$\widetilde{f}(\tau)=\sum_{r=0}^p g_r(\tau) E_2(\tau)^r,\quad \text{where } g_r\in M_{k,N,\overline{\mathbb{Q}}}.$$
Thus, we obtain 
  $$
\begin{aligned}
&trdeg\,\overline{\mathbb{Q}}\left( e^{2\pi i\tau},\, E_4(\tau),\, E_6(\tau),\, \widetilde f(\tau)\right)\\
&=trdeg\,\overline{\mathbb{Q}}\left( e^{2\pi i\tau},\, E_4(\tau),\, E_6(\tau),\,\sum_{r=0}^p g_r(\tau)\, E_2(\tau)^r \right)\\
&=trdeg\,\overline{\mathbb{Q}}\left( e^{2\pi i\tau}, E_4(\tau), E_6(\tau), g_0(\tau),\ldots,g_p(\tau),\sum_{r=0}^p g_r(\tau)\, E_2(\tau)^r \right)\\
&=trdeg\,\overline{\mathbb{Q}}\left( e^{2\pi i\tau},\, E_4(\tau),\, E_6(\tau),\, g_0(\tau),\,\ldots,\,g_p(\tau),\,E_2(\tau) \right)\\
&=trdeg\,\overline{\mathbb{Q}}\left( e^{2\pi i\tau},\, E_4(\tau),\, E_6(\tau),\,E_2(\tau) \right),
\end{aligned}
 $$
where the second and fourth equalities follow from the fact that each coefficient $g_r(\tau)$ is algebraic over $\overline{\mathbb{Q}}\left(E_4(\tau), E_6(\tau)\right)$ by Corollary \ref{algi}. Moreover, Nesterenko's theorem \ref{nest} implies that the right side of the final equality is at least $3$, and then we use the equality \eqref{final.} to conclude the proof. 












\begin{thebibliography}{0}

\bibitem{b-d-g-p}
K. Barr\'e-Sirieix, G. Diaz, F. Gramain, G. Philibert, \textit{Une preuve de la conjecture de Mahler-Manin}, Invent. Math., 124 (1996), 1-9.

% \bibitem{bertrand}
% D. Bertrand, \textit{S\'eries d’Eisenstein et transcendance.} Bull. Soc. Math. France
% $104(3)$, $309–321 (1976)$

\bibitem{chang}
C. Y. Chang, \textit{Transcendence of special values of quasi-modular forms}, Forum Math. 24 (2012) 539-551.


% \bibitem{chudnovsky}
%     G.V. Chudnovsky, \textit{Algebraic independence of constants connected with the exponential and the elliptic functions}, Dokl.
% Akad. Nauk Ukrain. SSR Ser. A $8$ $(1976)$ $698–701$.


\bibitem{stromberg}
H. Cohen and F. Stromberg, \textit{Modular Forms, A Classical Approach}, Graduate Studies In Mathematics 179. 

\bibitem{jeon-kang-kim}
D. Jeon, S.-Y. Kang, C.H. Kim, \textit{On values of weakly holomorphic modular functions at divisors of meromorphic modular forms}, J. Number Theory 239, 183-206.

\bibitem{j-k-k-2}
D. Jeon, S.-Y. Kang, C.H. Kim, \textit{Hecke system of harmonic Maass functions and its applications to modular curves of higher genera}, Ramanujan J., Vol. 62, (2023) 675-717.  


\bibitem{hamieh-murty}
Alia Hamieh and M. Ram Murty, {A note on $q$-analogues of Dirichlet L-functions}, International Journal of Number Theory, Vol. 12 (2016) 765-773.

\bibitem{lang}
Serge Lang, \textit{Elliptic Functions}, Addision-Wesley Publishing Company, Inc.


\bibitem{murty-dewar-graves}
M. R. Murty, M. Dewar and H. Graves, \textit{Problems in 
the Theory of 
Modular
Forms}, Hindustan Book Agency.

\bibitem{r-m}
M. Ram Murty and P. Rath, \textit{Transcendental Numbers} Springer-Verlag, (2014).

\bibitem{r-m-g}
S. Gun, M. Ram Murty, P. Rath, \textit{Algebraic independence of values of
modular forms}. Int. J. Number Theory 7(4), 1065–1074 (2011).


\bibitem{nesterenko}
Y. Nesterenko, \textit{Modular functions and transcendence questions}, Mat. Sb. 187 (9)  (1996) 65–96 (in Russian); translation in: Sb. Math. 187 (9) (1996) 1319–1348.

\bibitem{schneider}
T. Schneider, \textit{Arithmetische Untersuchungen elliptischer Integrale}, Math. Ann. 113 (1937) 1–13.

\bibitem{schneider.}
T. Schneider, \textit{Ein Satz über ganzwertige Funktionen als Prinzip für Transzendenzbeweise}, Mathematische Annalen, 121 (1949): 131–140. 

\bibitem{shimura}
 G. Shimura, \textit{Introduction to the arithmetic theory of automorphic functions}, Kano Memorial Lectures 1.

\bibitem{1-2-3}
D. Zagier, \textit{Elliptic Modular Forms and Their Applications}, Lectures at a Summer School
in Nordfjordeid, Norway.


\end{thebibliography}

\end{document}


\documentclass[11pt, reqno]{amsart}
\usepackage[utf8]{inputenc}
\usepackage{amsmath}
\usepackage{amsthm}
\usepackage{amsfonts}
\usepackage{amssymb}
\usepackage{mathrsfs, mathtools, mathdots}
\usepackage{diagbox}
\DeclarePairedDelimiter \floor {\lfloor} {\rfloor}
\DeclareMathOperator{\Tr}{Tr}

\usepackage[top=35mm,bottom=35mm,left=35mm,right=35mm]{geometry}
\newtheorem{theorem}{Theorem}[section]
\newtheorem{conjecture}{Conjecture}
\newtheorem{corollary}{Corollary}
\newtheorem{problem}{Problem}
\newtheorem*{theorem*}{Theorem}
\newtheorem*{remark*}{Remark}
\newtheorem*{problem*}{Problem}
\newtheorem*{conjecture*}{Conjecture}
\newtheorem{lemma}[theorem]{Lemma}
\newtheorem{prop}[theorem]{Proposition}


\begin{document}

\title[Transcendence]{Modular functions and Transcendence}

\author{[Author]}

\address{Department of Mathematics and Statistics, Queen's University, Kingston, Canada, ON K7L 3N6.}
\email{murty@mast.queensu.ca}
\address{Department of Mathematics and Statistics, Queen's University, Kingston, Canada, ON K7L 3N6.}



\email{siddhi@mast.queensu.ca}

\subjclass[2010]{-------}

\keywords{Algebraic independence, Transcendence degree, modular forms}
\thanks{Research ----}
\begin{abstract}

\end{abstract}

\maketitle

\section{\bf Introduction}
\bigskip

--------- ------------ ------------

\section{\bf Notations and Preliminary Results}
\bigskip

In this section, we fix some notations that will be used later. The set of all $2\times 2$ matrices with integer entries and determinant one forms a group under usual matrix multiplication, which is called (full) \emph{modular group} and denoted by $SL_2(\mathbb{Z})$. For each $N\in \mathbb{N}$, the set $$
\left\{  \left( 
\begin{array}{cc}
  a & b \\
  c & d
\end{array}
\right)\in SL_2(\mathbb{Z}):\left( 
\begin{array}{cc}
  a & b \\
  c & d
\end{array}
\right)\equiv \left( 
\begin{array}{cc}
  1 & 0 \\
  0 & 1
\end{array}
\right) (\operatorname{mod}N) \right\}
$$ forms a subgroup of $SL_2(\mathbb{Z})$, which is called the \emph{principal congruence subgroup of level N} and denoted by $\Gamma(N)$.
 Note that $\Gamma(1)=SL_2(\mathbb{Z})$. A subgroup $\Gamma$ of $SL_2(\mathbb{Z})$ is said to be congruence subgroup if there exists $N\in \mathbb{N}$ such that $\Gamma(N)\subseteq \Gamma$, and the smallest such $N$ is  called the level of $\Gamma$.\\

 Let $\mathbb{H}$ denote the upper half-plane,
$$
\mathbb{H}=\{\tau \in \mathbb{C}: \operatorname{Im}(\tau)>0\}.
$$
 The group $SL_2(\mathbb{Z})$ acts
on $\mathbb{H}$ via the formula
$$
\left(\begin{array}{ll}
a & b \\
c & d
\end{array}\right)\cdot \tau=\frac{a \tau+b}{c \tau+d} .
$$
Define the extended upper half-plane $\mathbb{H}^\star$ as 
$$
\mathbb{H}^{\star}=\mathbb{H} \cup \mathbb{Q} \cup\{i \infty\}, .
$$
where the rational numbers and $i\infty$ are called cusps. The notation $i\infty$ is used to emphasize that the point lies far in the imaginary direction. If we identify $i\infty$ with $\frac{1}{0}$, then the action of $SL_2(\mathbb{Z})$ extends from $\mathbb{H}$ to $\mathbb{H}^{*}$ via the formula \begin{equation} \label{eact}
\left(\begin{array}{cc}
    a & b \\
    c & d
\end{array}\right)\cdot \frac{r}{s}=\frac{ar+bs}{cr+ds}.
\end{equation}
This is indeed a group action where all the cusps are $SL_2(\mathbb{Z})$ equivalent to $i\infty$. 
\bigskip

\subsection{Modular forms}
\quad \bigskip

Let $\Gamma$ be a congruence subgroup of level $N$. Here we look at holomorphic functions $f: \mathbb{H}\rightarrow \mathbb{C}$ such that
\begin{equation} \label{hmod}
f(\gamma \tau)=(c\tau+d)^kf(\tau)\quad\text{for all }
\gamma=\left(\begin{smallmatrix}
   a  & b \\
   c  & d
\end{smallmatrix} \right)\in \Gamma,\end{equation}
where $k$ is a unique integer which is called the weight of $f$. Let $h$ be the least positive integer such that $\big(\begin{smallmatrix} 
1 & h \\
0 & 1
\end{smallmatrix}\big)\in \Gamma$. Obviously, $h\leq N$ and from \eqref{hmod}, we get $f(\tau+h)=f(\tau)$. Therefore, $f$ has a series expansion of the form 
$$
f(\tau)=\sum_{n\in\mathbb{Z}} a_n\ e^{2\pi i n\tau/h},
 $$
which is called the Fourier expansion of $f$ at $i\infty$, and $a_n$'s are called Fourier coefficients of $f$ at $i\infty$. We also denote the order of $f$ at $i\infty$ as $v_{i\infty}(f)=\operatorname{inf}\{n\in \mathbb{Z}:a_n\neq 0 \}$. Moreover, we say $f$ is holomorphic at $i\infty$ if $v_{i\infty}(f)\geq 0$, and meromorphic if $-\infty<v_{i\infty}(f)<0$.   \\

Sometimes we use {\it slash} notation which is defined with respect to the weight $k$ as follows:
$$
(f \mid_k \gamma)(\tau):=(\operatorname{det} \gamma)^{k / 2} j(\gamma, \tau)^{-k} f(\gamma \tau),
$$
where  $\gamma=\big(\begin{smallmatrix}
a & b \\
c & d
\end{smallmatrix}\big)\in GL_2^+(\mathbb{R})$ and   $j(\gamma,z)=cz+d$.
Then the modularity condition \eqref{hmod} is equivalent to
$$
f \mid_k \gamma=f, \quad \forall \gamma \in \Gamma.
$$
Let $s$ be an arbitrary cusp and $\gamma\in SL_2(\mathbb{Z})$ such that $\gamma(i\infty)=s$. Then the holomorphy of $f$ at $s$ determined by the holomorphy of $f\mid_k\gamma$ at $i\infty$.\\

 A \emph{modular form of weight} $k$ on a congruence subgroup $\Gamma$ is a holomorphic function $f:\mathbb{H}\rightarrow \mathbb{C}$ that satisfies \eqref{hmod} and is holomorphic at all cusps. A modular form $f$ is said to be a {\it cusp form} if it vanishes at all cusps. The set of all modular forms of weight $k$ with respect to $\Gamma$ forms a vector space over $\mathbb{C}$ which is denoted by $M_k(\Gamma)$. Moreover, the use $M_{k,\overline{\mathbb{Q}}}(\Gamma)$ to denote the set of all non-zero modular forms with algebraic Fourier coefficients at $i\infty$.\\

 For all even integers $k\geq 4$, we consider $$
G_{k}(\tau)=\sideset{}{'}\sum_{(m, n)\in \mathbb{Z}^2} \frac{1}{(m \tau+n)^{k}}, \quad \tau \in \mathbb{H},
$$
where the dashed summation means that it runs over all $(m, n) \in \mathbb{Z}^{2}$, with $(m, n) \neq(0,0)$. The function $G_k(\tau)$ is indeed a non-zero modular form of weight $k$ and level one. It has the following Fourier expansion at $i\infty$
\begin{equation}\label{fourier}
    G_k(\tau)=2 \zeta(k)+2 \frac{(-2 \pi i)^{k}}{(k-1) !} \sum_{n=1}^{\infty} \sigma_{k-1}(n)q^n, \text{ where }q=e^{2\pi i\tau} \text{ and } \sigma_s(n)=\sum_{\substack{d|n\\ d>0}}d^s.
\end{equation}
Dividing both sides of \eqref{fourier} by $2\zeta(k)$, we obtain
$$
E_{k}(\tau)=1-\frac{2 k}{B_{k}} \sum_{n=1}^{\infty} \sigma_{k-1}(n) q^{n},
$$
which is called the {\it normalized Eisenstein} series of level one and weight $k$. Here, $B_k$ is known as the $k$-th Bernoulli number, defined by the following power series $$\frac{x}{e^x-1}=\sum_{k=0}^\infty \frac{B_kx^k}{k!}.$$ It is easy to see from series expansion of $e^x-1$ that $B_k$ is a rational number. For $k=4$ and $6$, we have the following normalized Eisenstein series: 
$$
E_{4}(\tau)=1+240 \sum_{n=1}^{\infty} \sigma_{3}(n) q^{n} \quad \text{and} \quad
E_{6}(\tau)=1-504 \sum_{n=1}^{\infty} \sigma_{5}(n) q^{n}.
$$
Note that for two non-negative integers $\alpha$ and $\beta$, the function $E_4^\alpha E_6^\beta\in M_{4\alpha+6\beta}(S L_2(\mathbb{Z}))$. This is made precise in the following theorem.
\begin{theorem}\label{basis}
The set of functions
$$
\{E_4^\alpha E_6^\beta: \alpha,\beta\in\mathbb{Z}_{\geq 0}\ \text{and} \ 4\alpha+6\beta=k\}
$$
forms a basis for $M_k(S L_2(\mathbb{Z}))$.
\end{theorem}
For proof of Theorem \ref{basis}, we refer to (\cite{murty-dewar-graves}, Corollary $4.4.4$).
\begin{corollary}\label{weightzero}
Every modular form of weight $0$ on $S L_2(\mathbb{Z})$ is constant, i.e., $M_0(S L_2(\mathbb{Z}))=\mathbb{C}$.
\end{corollary}
\begin{proof}
    From Theorem \ref{basis}, we get that the constant function $1$ forms a basis for $M_0(S L_2(\mathbb{Z}))$ and hence, $M_0(S L_2(\mathbb{Z}))=\mathbb{C}$. 
\end{proof}
\bigskip

For $k=2$, the Eisenstein series is defined by the double series \begin{equation}\label{G2}
G_2(\tau)=\sum_{c\in\mathbb{Z}}\sum_{d\in\mathbb{Z}_{c}'}\frac{1}{(c \tau+d)^2},
\end{equation}
where $\mathbb{Z}_{c}'=\mathbb{Z}\setminus\{0\}$ if $c=0$ and $\mathbb{Z}_c'=\mathbb{Z}$ otherwise. It satisfies the following transformation property
$$
G_2\left(\frac{-1}{\tau}\right)=\tau^2\, G_2(\tau)-2\pi i\tau.
$$
Thus $G_2(\tau)$ is not a modular form, but the normalized $G_2(\tau)$ has similar Fourier expansion
$$E_2(\tau)=1-24\sum_{n=1}^\infty\sigma_1(n)q^n, \text{ where }q=e^{2\pi i\tau} \text{ and }\sigma_1(n)=\sum_{\substack{d|n\\ d>0}}d.$$\\


Define a function $\Delta$ on $\mathbb{H}$ as follows $$
\Delta(\tau): = \frac{E_4(\tau)^3 - E_6(\tau)^2}{1728},
$$
where $E_4^3, E_6^2\in M_{12}(S L_2(\mathbb{Z}))$. Since $M_{12}(SL_2(\mathbb{Z}))$ is $\mathbb{C}$-vector space, $\Delta$ is a modular form of weight 12. It has the following Fourier series  $$\Delta(\tau)=q-24q^2+252q^3+ \cdots,$$ which implies that $\Delta$ is indeed a cusp form having a simple zero at $i\infty$. This function is known as \emph{Ramanujan's $\Delta$-function}.
\bigskip

\subsection{The Weierstrass $\wp$-function}
\quad\bigskip

Here we introduce the Weierstrass $\wp$-function and review some of its standard properties. The Weierstrass $\wp$-function associated with a two dimensional
lattice $L=\omega_1 \mathbb{Z} \,  \oplus \, \omega_2 \mathbb{Z}$ is defined as follows
\begin{equation}\label{wpfunction}
\wp_L (z)=\frac{1}{z^2}+\sum_{\substack{\omega\in L\\
\omega\neq 0}}\left(\frac{1}{(z-\omega)^2}-\frac{1}{\omega^2} \right), \quad \text{for }z\in \mathbb{C}\setminus L. 
\end{equation}
If the corresponding lattice $L$ is clear from the context, we simply write $\wp(z)$ instead of $\wp_L(z)$. This is a meromorphic function on $\mathbb{C}$ with double poles at each lattice points. Note that $\wp_L(z)$ is an even function and each point of $L$ is a period of $\wp_L(z)$, and so $L$ is called the {\it period module} of $L$. We refer to $\frac{\omega_1}{2},\frac{\omega_2}{2},\frac{\omega_1+\omega_2}{2}$ and their $L$-equivalent points as \emph{half-lattice} points. One can show that the zeros of $\wp_L'(z)$ are all distinct and these are the half-lattice points. Moreover, it is known that the numbers $\wp(\omega_1/2)$, $\wp(\omega_2/2)$ and $\wp((\omega_1+\omega_2)/2)$ are all distinct. For further properties of Weierstrass $\wp$-function, we refer the reader to \cite{ahlfors}.\\

The Weierstrass $\wp$-function satisfies the following differential equation
\begin{equation} \label{dell}
 \wp'(z)^2=4\wp(z)^3-g_2(L)\wp(z)-g_3(L),   
\end{equation}
where $$g_2(L)=60\, \sum_{\substack{\omega \in L\\ \omega\neq 0}} \frac{1}{\omega^{4}} \text{ and } g_3(L)=140\, \sum_{\substack{\omega \in L\\ \omega\neq 0}} \frac{1}{\omega^{6}}.$$
In particular, if we take $L_\tau=\tau\,\mathbb{Z}\oplus \mathbb{Z}$ for $\tau\in\mathbb{H}$, then $$g_2(L_\tau)=60\,G_4(\tau)=\frac{4\pi^4}{3}E_4(\tau) \text{ and } g_3(L_\tau)=140\,G_6(\tau)=\frac{8\,\pi^6}{27}E_6(\tau).$$\\

We now consider the polynomial $f(x)=4 x^3-g_2(L)x-g_3(L)$. From the zeros of $\wp_L'(z)$ and \eqref{dell}, we get that $\wp(\omega_1/2)$, $\wp(\omega_2/2)$ and $\wp((\omega_1+\omega_2)/2)$ are the roots of $f(x)$ and these are all distinct. Therefore, the discriminant of $f$ $$\Delta_0(L):=g_2(L)^3-27g_3(L)^2\neq 0, \text{ for any lattice }L.$$
In particular, For any $\tau\in\mathbb{H}$ we get
$$
\Delta_0(L_\tau)=(2\pi)^{12}\Delta(\tau),
$$
where $\Delta$ is the Ramanujan cusp form. This shows that $\Delta$ is non-vanishing on $\mathbb{H}$.\\

The Weierstrass $\wp$-function has a Fourier expansion which is given by 
\begin{equation}\label{weifou.}
(2\pi i)^{-2}\,\wp_{L_\tau(z)}=\frac{1}{12}+\sum_{n=1}^\infty n\,e^{2\pi i n z}+\sum_{m,n=1}^\infty n\,q^{mn}\left[e^{2\pi i n z}+e^{-2 \pi i n z}-2 \right].
\end{equation}
This Fourier expansion of the $\wp$-function is used later. For details about this Fourier expansion, we refer to (\cite{lang}, Chapter $4$). 
\bigskip

\subsection{\bf The Weierstrass zeta function}
\quad\bigskip

The Weierstrass zeta-function associated with a two dimensional lattice $L=\omega_1\,\mathbb{Z}\oplus \omega_2\,\mathbb{Z}$ is defined as follows:
$$
Z_{L}(z)=\frac{1}{z}+\sum_{\substack{\omega\in L\\ \omega\neq 0}}\left(\frac{1}{z-\omega}+\frac{1}{\omega}+\frac{z}{\omega^2} \right)\quad \text{for } z\in \mathbb{C}\setminus L.
$$
One can show that this double sum converges uniformly and absolutely on any compact set of $\mathbb{C}$, not containing lattice points and hence, this sum defines a holomorphic function on $\mathbb{C}\setminus L$ with simple pole at each lattice points. Then the term by term differentiation of $Z_{L}$ shows that $$
Z_{L}'(z)=-\left[\frac{1}{z^2}+\sum_{\substack{\omega\in L\\
\omega\neq 0}}\left(\frac{1}{(z-\omega)^2}-\frac{1}{\omega^2} \right)\right]=-\wp_L(z),
$$
which is a periodic function with each lattice point of $L$ as a period. Hence, the functions $$Z_{L}(z+\omega_1)-Z_{L}(z) \text{  and  } Z_{L}(z+\omega_2)-Z_{L}(z)$$ are constants. These constants are denoted by $\eta_1(L), \eta_2(L)$ respectively and called quasi-periods.\\

Now we consider a lattice of the form $L_\tau=\tau\,\mathbb{Z}\oplus\mathbb{Z}$, where $\tau\in\mathbb{H}$. The Weierstrass zeta-function has the following
$q$-series expansion (\cite{lang}, chapter $18.3$):
$$
Z_{L_\tau}(z)=\eta_2(L_\tau)z-\pi i\frac{1+e^{2\pi i z}}{1-e^{2\pi i z}}-2\pi i \sum_{n=1}^\infty \left(\frac{e^{2\pi i z}q^n}{1-e^{2\pi i z}q^n}-\frac{e^{-2\pi i z}q^n}{1-e^{-2\pi i z}q^n}\right), \text{ where }q=e^{2\pi i \tau}.
$$
Note that the function $Z_{L_\tau}$ is odd, which follows from the fact that $Z_{L_\tau}'(z)=-\wp(z)$ is an even function. Hence, the coefficient of $z$ in the $q$-expansion of $Z_{L_\tau}$ is equal to $0$, and this gives
\begin{equation}\label{quasi-period}
\eta_2(L_\tau)=G_2(\tau)=\frac{\pi^2}{3}\, E_2(\tau), 
\end{equation}
where $E_2$ is the normalized Eisenstein series of weight $2$. The identity \eqref{quasi-period} is used later to study special values of quasi-modular forms.


\subsection{Modular functions}
\quad\bigskip

Let $\Gamma$ be a congruence subgroup of $S L_2(\mathbb{Z})$ and $f$ be a function on $\mathbb{H}$. Then $f$ is said to be a \emph{modular function} with respect to $\Gamma$ if it is \emph{meromorphic} on $\mathbb{H^*}$ and satisfies the modular transformation formula \eqref{hmod} for $k=0$, that is
$$
f\left(\gamma \tau \right)=f(\tau) \quad \forall\ \gamma \in \Gamma.
$$
In particular, a modular function on the congruence subgroup $\Gamma(N)$ is said to a be modular function of level $N$. Since a modular function is meromorphic on $\mathbb{H^*}$, the set of all modular functions with respect to a congruence subgroup $\Gamma$ forms a field which is denoted by $F_{\Gamma,\mathbb{C}}$. Here we analyze the structure of the field of modular functions for congruence subgroups of arbitrary level and illustrate their connection to modular functions of level one.\\

An example of a modular function of level $1$ is the $j$-function, defined as $$
j(\tau)=\frac{E_4^3(\tau)}{\Delta(\tau)} \text{ for } \tau\in\mathbb{H}.
$$ The $j$-function is holomorphic on $\mathbb{H}$ as $E_4$ and $\Delta$ are modular forms and $\Delta$ is non-vanishing on $\mathbb{H}$. The $j$-function has the following Fourier expansion at $i\infty$ $$
j(\tau)=\frac{1}{q}+744+196884q+\cdots,$$
where $q=e^{2\pi i\tau}$.
The $j$-function is a canonical example of modular function of level $1$. This is made precise in the following theorem.
\begin{theorem} \label{level1}
Let $f$ be a modular function of level one and let $$K := \mathbb{Q}(\text{Fourier coefficients of }f)\subset\mathbb{C}.$$ Suppose that the poles of $f$ in $\mathcal{F}$ are $\tau_1,\tau_2,\cdots,\tau_m$, where $\mathcal{F}$ denotes the standard fundamental domain for the action of $S L_2(\mathbb{Z})$ on $\mathbb{H}$. Then $f\in K(j(\tau_1),\cdots,j(\tau_m))(j)$ i.e., $f$ is a rational function in $j$ over $K(j(\tau_1),\cdots,j(\tau_m))$.
\end{theorem}
\begin{proof}
First, we consider the function 
$$
g(\tau):=\left(\prod_{\tau_0\in\mathcal{F}} \left(j(\tau)-j(\tau_0)\right)^{\operatorname{ord}_{\tau_0}f}\right)f(\tau),
$$
which is holomorphic on $\mathbb{H}$. Since $f$ is meromorphic at $i\infty$, there exists $c>0$ such that the region $\{z\in\mathcal{F} : \operatorname{Im}z>c\}\cup \{i\infty\}$ contains no pole of $f$ except possibly $i\infty$. Again, $f$ is meromorphic on $\mathbb{H}$ and the region $\{z\in \mathcal{F}: \operatorname{Im}z\leq c\}$ is compact, and so it contains only a finite number of poles of $f$. Thus, $f$ can have only finitely many poles in $\mathcal{F}$ and hence above product is finite.\\

Suppose $g$ has pole of order $M$ at $i\infty$. Then the Fourier expansion of
$g(\tau)$ at $i\infty$ has the form
$$
g(\tau)=\sum_{n=-M}^\infty c_n q^n, \, \text{ where } c_{-M}\neq 0.
$$
Note that the modular function $g(\tau)-c_{-M}j(\tau)^M$ is holomorphic on $\mathbb{H}$ and its $q$-series starts with at most a polar term of order $M-1$. Repeating this process, we can subtract a polynomial of $j(\tau)$ to get a holomorphic modular function that vanishes at $i\infty$, and hence identically zero by Corollary \ref{weightzero}. Thus, $g(\tau)$ is a polynomial of $j(\tau)$ over $K$ and so $f(\tau)$ is a rational function of $j(\tau)$ over $K(j(\tau_1),\cdots,j(\tau_m))$.
\end{proof}
The above proof proves the following :
\begin{corollary}\label{rationalfunctions}
If $f(\tau)$ is a modular function of level one which is holomorphic on $\mathbb{H}$ with the following $q$-expansion at $i\infty$
$$
f(\tau)=\sum_{n=-M}^\infty c_n q^n, \ c_{-M}\neq 0
$$
with $K=\mathbb{Q}(\text{Fourier coefficients of }f)$, then $f(\tau)$ is a polynomial in $j(\tau)$ of degree $M$ with coefficients in $K$.
\end{corollary}
By using Theorem \ref{level1}, we prove the following result, which describes the connection between $j$-function and modular functions of arbitrary level.
\begin{theorem} \label{mdft}
Let $f$ be a modular function with respect to a congruence subgroup $\Gamma$ and
$$
K:=\mathbb{Q}\left(\text{Values of }j\text{ at poles of }f\text{ in }\mathcal{F}, \text{ Fourier coefficients of }f\text{ at all cusps}\right).
$$
Then there exists a monic polynomial $P(X)\in K(j)[X]$ such that $P(f)=0$. 
\end{theorem}
\begin{proof}
We construct a polynomial satisfied by $f$ with coefficients that are rational functions of $j(\tau)$. Theorem \ref{level1} suggests that these coefficients are modular functions of level $1$.\\

Let $[S L_2(\mathbb{Z}):\Gamma]=r$ and $\{ \gamma_1=I,\gamma_2,\cdots,\gamma_r\}$ be a complete set of right coset representatives so that $$
S L_2(\mathbb{Z})=\bigsqcup_{i=1}^r\Gamma \gamma_i.
$$
For all $1\leq i \leq r$, define the functions $$
f_i(\tau):=f(\gamma_i\tau) \quad \text{for } \tau \in \mathbb{H}.
$$ Clearly, each $f_i$ is independent of the choice of coset representatives as $f$ is $\Gamma$-invariant. Again, for any $\gamma\in S L_2(\mathbb{Z})$, we have $f_i(\gamma \tau)=f(\gamma_i\gamma\tau)=f_j(\tau)$ if $\gamma_i\gamma\in \Gamma\gamma_j$, i.e., the set $\{f_1,f_2,...,f_r\}$ gets permuted under the action of $S L_2(\mathbb{Z}).$ This observation implies that any elementary symmetric polynomial of $f_1,f_2,\cdots,f_r$ is a modular function with respect to $S L_2(\mathbb{Z}).$\\

Consider the polynomial $$
P(X)=\prod_{i=1}^r(X-f_i),$$ which is satisfied by $f$ as $f=f_1$. Note that the coefficients of $P$ are elementary symmetric polynomials in $f_1,f_2,\cdots,f_r$, and hence, modular functions with respect to $S L_2(\mathbb{Z})$ with Fourier coefficients in $K$. Thus by Theorem \ref{level1}, the coefficients of $P$ are in $K(j)$ and so $P(X)$ is the required polynomial.
\end{proof}
\begin{corollary}\label{useful}
 If the modular function $f$ is holomorphic on $\mathbb{H}$ and $$K := \mathbb{Q}(\text{Fourier coefficients of }f\text{ at all cusps}),$$
 then the monic polynomial in Theorem \ref{mdft} lies in $K[j][X]$.
\end{corollary}
\begin{proof}
   This is clear from above proof and Corollary \ref{rationalfunctions}, since coefficients of $P$ constructed above are modular functions of level one, holomorphic on $\mathbb{H}$ with Fourier coefficients in $K$.  
\end{proof}
\bigskip

For a positive integer $N$, the field of all modular functions of level $N$ is denoted by $F_{N,\mathbb{C}}$.
In the following theorem, we prove an algebraic property of the field 
 $F_{N,\mathbb{C}}$.
\begin{theorem}\label{fnc}
The field $F_{N,\mathbb{C}}$ is a finite Galois extension of $\mathbb{C}(j)$ and the  Galois group is given by $S L_2(\mathbb{Z})/\{\pm I\}\Gamma(N)$.
\end{theorem}
\begin{proof}
From the Theorem \ref{mdft}, we have an explicit form of  polynomial satisfied by $f$ for each $f\in F_{N,\mathbb{C}}$.The normality of the extension $F_{N,\mathbb{C}}/\mathbb{C}(j)$, is also clear because all the conjugates of $f$, denoted by $f_i$'s are also inside $F_{N,\mathbb{C}}$. Therefore, the field extension $F_{N,\mathbb{C}}/\mathbb{C}(j)$ is a Galois extension. To find the Galois group, firstly,  note that any element $\gamma\in S L_2(\mathbb{Z})$ defines an automophism of $F_{N,\mathbb{C}}$ by $f \mapsto f\circ\gamma $, where $(f\circ\gamma )(\tau):=f(\gamma \tau)$ for all $\tau\in\mathbb{H}$. Moreover, $\gamma$ fixes the base field as $j$ is modular function of level one.
Now, define a map $$
\begin{aligned}
\Psi :\, & S L_2(\mathbb{Z}) \longrightarrow Gal(F_{N,\mathbb{C}}/\mathbb{C}(j))\\
& \quad\gamma \longmapsto \Psi_\gamma : F_{N,\mathbb{C}} \longrightarrow F_N\\
&\quad \quad\quad\quad \quad \quad f \longmapsto f\circ\gamma.
\end{aligned}
$$
Then $\Psi$ is clearly a group homomorphism with kernel $\{\pm I\}\Gamma(N)$. From primitive element theorem, we have $F_{N,\mathbb{C}}=\mathbb{C}(\widetilde f)$ for some primitive element. Note that any $\sigma\in Gal(F_{N,\mathbb{C}}/\mathbb{C}(j)) $ sends $f$ to one of the conjugates, which is of the form $\widetilde f_i=\widetilde f\circ\gamma_i$. Then we get
$$
\Psi(f)=f\circ\gamma_i\quad \forall\ f\in F_{N,\mathbb{C}},
$$
and so $\sigma=\Psi_{\gamma_i}$. This shows that $\Psi$ is onto. Therefore, by the First Isomorphism theorem we obtain $Gal(F_N/\mathbb{C}(j))=SL _2(\mathbb{Z})/\{\pm I\}\Gamma(N).$
\end{proof}
\bigskip

From Theorem \ref{fnc}, we already know that $F_{N,\mathbb{C}}$ is finitely generated. Here we identify a set of generators for $F_{N,\mathbb{C}}$. For $a=(a_1,a_2)\in \mathbb{Q}^2\setminus\mathbb{Z}^2$, consider the function
$$
f_{a}(\tau):=\frac{g_2(L_\tau)\,g_3(L_\tau)}{\Delta_0(L_\tau)}\,\wp_\tau (a_1\tau+a_2),
$$
which is holomorphic on $\mathbb{H}$ because $g_2, g_3$ are holomorphic functions, $\Delta_0$ is non-vanishing on $\mathbb{H}$ and $a_1\tau+a_2\notin L_\tau$.\\

If $\gamma=\big(\begin{smallmatrix}
p & q\\
r & s
\end{smallmatrix}\big)\in S L_2(\mathbb{Z})$, then we have
$$
\begin{aligned}
f_a(\gamma\tau)&=\frac{g_2(L_{\gamma\tau})\,g_3(L_{\gamma\tau})}{\Delta_0(L_{\gamma\tau})}\,\wp_{\gamma \tau}(a_1\gamma\tau+a_2)\\
&=\frac{g_2(L_\tau)\,g_3(L_\tau)}{\Delta_0(L_\tau)}\,\wp_\tau((a_1p+a_2r)\tau+(a_1q+a_2s))\\
&=f_{a\gamma}(\tau)\quad \forall\ \tau\in\mathbb{H}.
\end{aligned}
$$
For $a=(a_1,a_2),b=(b_1,b_2)\in \mathbb{Q}^2\setminus\mathbb{Z}^2$, one gets  $$
 f_a=f_b\Leftrightarrow \wp(a_1\tau+a_2;L_\tau)=\wp(b_1\tau+b_2;L_\tau) \Leftrightarrow a\equiv \pm b \pmod{\mathbb{Z}^2}.
$$
Thus, the function $f_a$ is completely determined by $a \bmod \mathbb{Z}^2$.\\

Moreover, if $N$ is a positive integer such that $Na\in \mathbb{Z}^2$, we claim that $$
f_a(\gamma \tau)=f_a(\tau) \quad \forall\, \gamma \in \{\pm I\}\Gamma(N).
$$
Indeed, if $a=(a_1,a_2)$ and $\gamma=\big(\begin{smallmatrix}
p & q\\
r & s
\end{smallmatrix}\big)\in \{\pm I\} \Gamma(N)$, then we calculate
$$
\begin{aligned}
    a\gamma &=(a_1p+a_2 r, a_1q+a_2 s)\\
    &\equiv (a_1\,p, a_2\, s)\pmod{\mathbb{Z}^2} \quad [\text{ As } q,r\equiv 0\pmod{N} \text{ and }(Na_1,N a_2)\in\mathbb{Z}^2]\\
    &\equiv\pm (a_1, a_2)\pmod{\mathbb{Z}^2}\\
    &\equiv\pm a\pmod{\mathbb{Z}^2}.
\end{aligned}
$$
Thus, we conclude that for any $a\in N^{-1}\mathbb{Z}^2\setminus\mathbb{Z}^2$, 
$$
f_a(\gamma \tau)=f_{a\gamma}(\tau)=f_a(\tau) \quad \forall\, \gamma \in \{\pm I\}\Gamma(N).
$$
This shows that $f_a$ satisfies the modularity property \eqref{hmod} for $k=0$ with respect to $\Gamma(N)$.\\

 Let $S$ denote the set $\{\left(\frac{r}{N},\frac{s}{N}\right): 0\leq r,s\leq N-1 \text{ and }(r,s)\neq(0,0)\}$. Since for $a,b\in\mathbb{Q}^2\setminus\mathbb{Z}^2$
    $$
f_a=f_b\Leftrightarrow a\equiv \pm b \pmod{\mathbb{Z}^2},
$$
the following equality of two sets holds
$$
\{f_a: a\in N^{-1}\mathbb{Z}^2\setminus\mathbb{Z}^2\}=\{f_a: a\in S\}.
$$
Therefore, it is enough to study only the functions $f_a$ for $a\in S$. \\

Now we analyze behaviour of $f_a(\tau)$ at $i\infty$ for all $a\in S$. For $a=(\frac{r}{N},\frac{s}{N})\in S$, we have 
 $$
 \begin{aligned}
f_a(\tau)&=\frac{g_2(L_\tau)\,g_3(L_\tau)}{\Delta_0(L_\tau)}\,\wp_\tau (a_1\tau+a_2)\\
&=\frac{(2\pi)^{-2}}{2592}\cdot \frac{E_4(\tau)\,E_6(\tau)}{\Delta(\tau)}\,\wp_\tau\left(\frac{r\tau+s}{N}\right).
\end{aligned}
$$
Taking $z=\frac{r\tau+s}{N}$ in the Fourier expansion of Weierstrass $\wp$-function \eqref{weifou.}, we get
$$
(2\pi i)^{-2}\wp_\tau\left(\frac{r\tau+s}{N}\right)=\frac{1}{12} -\sum_{n=1}^\infty n e^{\frac{2\pi i n s}{N}} q_N^{n r}+\sum_{m,n=1}^\infty n q^{m n}\left[e^{\frac{2\pi i ns}{N}}q_N^{n r}+e^{\frac{-2\pi i n s}{N}}q_N^{-n r}-2 \right],
$$
where $q_N=e^{2\pi i\tau/N}$ has always non-negative power.
Thus, we obtain
\begin{equation}\label{foufa}
f_a(\tau)=-\frac{E_4(\tau)E_6(\tau)}{2592\,\Delta(\tau)}\left[\frac{1}{12} -\sum_{n=1}^\infty ne^{\frac{2\pi i n s}{N}} q_N^{n r}+\sum_{m,n=1}^\infty n q^{m n}\left[e^{\frac{2\pi i n s}{N}}q_N^{n r}+e^{\frac{-2\pi i n s}{N}}q_N^{-n r}-2 \right]\right].
\end{equation}
Since the Fourier series of $\Delta(\tau)$ begins with $q=q_N^N$, the Fourier expansion of $f_a(\tau)$  begins with a rational multiple of $q_N^{-N}$. Thus, $f_a$
has pole of order $N$ at $i\infty$ for all $a\in S$.\\

If $s$ is any other cusp, then there exists $\gamma\in S L_2(\mathbb{Z})$ such that $s=\gamma(i\infty)$. Since $f_a(\gamma\tau)=f_{a\gamma}(\tau)$, which has pole of order $N$ at $i\infty$, we conclude that $f_a$ is meromorphic at all cusps, with pole of order $N$. Thus, $f_a$ is a modular function of level $N$ for all $a\in S$.
\begin{lemma}\label{primitex}
    The Fourier coefficients of $f_a$ belong to $\mathbb{Q}(\mu_N)$ with respect to all cusps $\forall a\in S$, where $\mu_N=e^{2\pi i/N}$.
\end{lemma}
\begin{proof}
   Recall that Fourier coefficients of $E_4(\tau), E_6(\tau)$ and $\Delta(\tau)$ are integers. Hence, from \eqref{foufa}, it follows that the Fourier coefficients of $f_a$ at $i\infty$ lie in $\mathbb{Q}(\mu_N)$ for all $a\in S$. If $s$ is any other cusp, then there exists $\gamma\in S L_2(\mathbb{Z})$ such $s=\gamma(i\infty)$. But $f_a(\gamma\tau)=f_{a\gamma}(\tau)$ has Fourier coefficients in $\mathbb{Q}(\mu_N)$ with respect to $i\infty$. This completes the proof.
\end{proof}
\begin{theorem}
The field of all modular functions of level $N$, that is,  $$F_{N,\mathbb{C}}=\mathbb{C}\left(\{j, f_a | a\in S\}\right).$$
\end{theorem}
\begin{proof}
Let $E_{N,\mathbb{C}}$ denote the field $\mathbb{C}\left(\{j, f_a | a\in S\}\right)$. Then we have the following containment $\mathbb{C}(j)\subseteq E_{N,\mathbb{C}} \subseteq F_{N,\mathbb{C}}.$ To complete the proof, it is enough to show $[F_{N,\mathbb{C}}:E_{N,\mathbb{C}}]=1$, i.e., any element of $Gal(F_{N,\mathbb{C}}/\mathbb{C}(j))$ that fixes $E_{N,\mathbb{C}}$ must be trivial.\\

Recall that $Gal(F_{N,\mathbb{C}}/\mathbb{C}(j))=S L_2(\mathbb{Z})/\{\pm I\}\Gamma(N)$. Let $\gamma=\big(\begin{smallmatrix}
p & q\\
r & s\end{smallmatrix}\big)\in S L_2(\mathbb{Z}) $ fix $E_{N,\mathbb{C}}$, that is, $$f_a\circ\gamma=f_a\quad \text{for all } a\in S.$$ Equivalently,
$$
 \quad f_{a\gamma}=f_a\quad \text{for all } a\in S
$$
and hence 
$$
 a\gamma\equiv \pm a \pmod{\mathbb{Z}^2} \quad \text{ for all } a\in S.
$$
In particular, taking two choices of $a= (\frac{1}{N},0)$ and $(0,\frac{1}{N})$ we conclude $\gamma \in \{\pm I\}\Gamma(N).$ This completes the proof.
\end{proof}
\bigskip

If we denote the field of modular functions of level $N$ with algebraic Fourier coefficients at $i\infty$ by  $F_{N,\overline{\mathbb{Q}}}$, then Lemma \ref{primitex} implies that $f_a\in F_{N,\overline{\mathbb{Q}}}$ for all $a\in S$. The following theorem proves that these functions are indeed all generators of $F_{N,\overline{\mathbb{Q}}}$.
\bigskip

\begin{theorem}\label{mdft..}
$F_{N,\overline{\mathbb{Q}}}$ is a finite Galois extension of $\overline{\mathbb{Q}}(j)$ and $F_{N,\overline{\mathbb{Q}}}=\overline{\mathbb{Q}}(j,f_a |a\in S)$.
\end{theorem}
\bigskip

Before proceeding with the proof, recall that if $L/K,L'/K$ are two field extensions such that both $L,L'$ are contained in an extension of $K$, then the fields $L$ and $L'$ are called \emph{linearly disjoint} over $K$ if any $K$-linearly independent subset of $L$ is also $L'$-linearly independent when the elements are considered as members of the larger field.

\begin{proof}
    Let $E_{N,\overline{\mathbb{Q}}}$ denote the field $\overline{\mathbb{Q}}(j,f_a |a\in S)$. First, we show that $E_{N,\overline{\mathbb{Q}}}/\overline{\mathbb{Q}}(j)$ is a Galois extension. Note that for each $a\in S$, the modular function $f_a$ is holomorphic on $\mathbb{H}$ and it has algebraic Fourier coefficients at all cusps by Lemma \ref{primitex}. Hence, by Corollary \ref{useful}, we get a polynomial $P(X)\in\overline{\mathbb{Q}}(j)[X]$, which is satisfied by $f_a$. Thus, $E_{N,\overline{\mathbb{Q}}}/\overline{\mathbb{Q}}(j)$ is an algebraic extension. Moreover, this extension is normal because each conjugate of $f_a$ is of the form $f_{a\gamma_i}$, which is equal to $f_b$ for some $b\in S$. Since $E_{N,\overline{\mathbb{Q}}}$ is obtained from $\overline{\mathbb{Q}}(j)$ by adjoining finite number of elements, $E_{N,\overline{\mathbb{Q}}}/\overline{\mathbb{Q}}(j)$ is a finite Galois extension.\\

Now we prove that $E_{N,\overline{\mathbb{Q}}}=F_{N,\overline{\mathbb{Q}}}$. We claim that $\mathbb{C}$ and $F_{N,\overline{\mathbb{Q}}}$ are linearly disjoint over $\overline{\mathbb{Q}}$. To prove this, consider an arbitrary $\overline{\mathbb{Q}}$-linearly independent set $\{c_1,c_2,\ldots,c_r\}\subset \mathbb{C}$. If there exists $$f_i(\tau)=\sum_n d_{i n}\,q_N^n \in F_{N,\overline{\mathbb{Q}}}\quad \text{for }0\leq i\leq r$$ such that $\sum_{i=1}^rc_i\,f_i\equiv 0$, then $$\sum_{i=1}^rc_i\,d_{i n}=0\quad \text{ for all } n,$$ which implies that all $d_{i n}=0$ for all $0\leq i\leq r$, and hence, $f_1=f_2=\cdots=f_r=0$.\\

Again, we have $E_{N,\overline{\mathbb{Q}}}\subseteq F_{N,\overline{\mathbb{Q}}} \subseteq \mathbb{C}E_{N,\overline{\mathbb{Q}}}.$ If possible, let $f\in F_{N,\overline{\mathbb{Q}}}\setminus E_{N,\overline{\mathbb{Q}}}$. Then we get a $\overline{\mathbb{Q}}$-linearly independent subset $\{f_1,f_2,\ldots,f_m\}\subseteq E_{N,\overline{\mathbb{Q}}}$ such that 
\begin{equation}\label{rep.}
  f=\sum_{i=1}^m c_i f_i, \quad c_i\in \mathbb{C}.  
\end{equation}
Since $E_{N,\overline{\mathbb{Q}}}\subseteq F_{N,\overline{\mathbb{Q}}}$, the set $\{f_1,f_2,\ldots,f_m\}$ is $\overline{\mathbb{Q}}$-linearly independent subset of $F_{N,\overline{\mathbb{Q}}}$ and hence, $\mathbb{C}$-linearly independent, which contradicts \eqref{rep.}. Therefore, $E_{N,\overline{\mathbb{Q}}}=F_{N,\overline{\mathbb{Q}}}$ and $F_{N,\overline{\mathbb{Q}}}/\overline{\mathbb{Q}}(j)$ is a finite Galois extension.
\end{proof}
\begin{corollary}\label{coefficients}
    If a modular function $f$ has algebraic Fourier coefficients at $i\infty$ (i.e. $f\in F_{N,\overline{\mathbb{Q}}}$ ), then $f$ has algebraic Fourier coefficients at all cusps.
\end{corollary}
\begin{proof}
    It follows from the fact that $f\in F_{N,\overline{\mathbb{Q}}}=\overline{\mathbb{Q}}(j,f_a |a\in S)$ and each $f_a$ has algebraic Fourier coefficients at all cusps by Lemma \ref{primitex}.
\end{proof}
\begin{corollary}
    If a modular form $f$ of level $N$ and weight $k$ has algebraic Fourier coefficients at $i\infty$, then $f$ has algebraic Fourier coefficients at all cusps. 
\end{corollary}
\begin{proof}
    We first construct the modular function 
$$g(\tau):=\frac{f(\tau)^{12}}{\Delta(\tau)^k}\quad \forall \tau\in\mathbb{H}.$$
Since $\Delta(\tau)$ has algebraic Fourier coefficients, the modular function $g$ belongs to $F_{N,\overline{\mathbb{Q}}}$, and so we can apply Corollary \ref{coefficients} to conclude the proof.
\end{proof}
\bigskip

\subsection{Quasi-modular forms} \quad \bigskip

A quasi-modular form $f$ of weight $k$ and depth $p$ with respect to a congruence subgroup $\Gamma$ is a holomorphic function on $\mathbb{H}$ satisfying the followings:
\begin{enumerate}
    \item There exists component functions $f_0=f,\, f_1,\,\ldots, f_p$ such that each $f_i$ is holomorphic on $\mathbb{H}^{*}$ and $f_p\neq 0$.
    \item The function $f$ satisfies the following transformation property $$f\mid_k \gamma\, (\tau)=\sum_{0\leq i\leq p} f_i(\tau)\,\left(\frac{c}{c\tau+d}\right)^i \text{ for all }\gamma=\left(\begin{smallmatrix}
        a & b\\
        c & d
\end{smallmatrix}\right)\in \Gamma. $$ 
\end{enumerate}
We use the notation $\widetilde{M}_k^{(\leq p)}(\Gamma)$ to denote the space of \emph{quasi-modular forms} of \emph{weight} $k$ and \emph{depth} $\leq p$ with respect to $\Gamma$. There are many other characterizations for quasi-modular forms, but we use the following.
\begin{theorem}\label{quasistructure}
Every quasi-modular form on $\Gamma$ is a polynomial in $E_2$ with modular coefficients. More precisely, if $f\in \widetilde{M}_k^{(\leq p)}(\Gamma)$, then $f$ can be uniquely written as $f=\sum_{r=0}^p g_r E_2^r$, where $g_r\in M_{k-2 r}(\Gamma)$ for all $0\leq r\leq p$. 
\end{theorem}
If all the modular coefficients of $f$ has algebraic Fourier coefficients, then we say that the quasi-modular form $f$ has algebraic Fourier coefficients. For a detailed theory of quasi-modular forms, we refer the reader to (\cite{stromberg}, Chapter $5.1$). 
\bigskip

\section{\bf Some Known Results on Transcendence}
\bigskip

Form the last section, we have the following normalized Eisenstein series on $\mathbb{H}$ 
$$
\begin{aligned}
E_2(\tau)&=1-24 \sum_{n=1}^{\infty} \sigma_1(n) e^{2 \pi i n \tau},\\
E_4(\tau)&=1+240 \sum_{n=1}^{\infty} \sigma_3(n) e^{2 \pi i n \tau},\\
E_6(\tau)&=1-504 \sum_{n=1}^{\infty} \sigma_5(n) e^{2 \pi i n \tau}, \text{ where }\sigma_k(n)=\sum_{\substack{d|n\\ d>0}}d^k.
\end{aligned}
$$
The values of these functions on $\mathbb{H}$ is a rich source of transcendental numbers. An earlier result regarding the transcendental values of $E_4$ and $E_6$ was given by D. Bertrand \cite{bertrand}. He proved that for $\tau\in\mathbb{H}$, at least one of the numbers $E_4(\tau), E_6(\tau)$ is transcendental. Later, Chudnovsky \cite{chudnovsky} showed that if $\tau\in\mathbb{H}$, then at least two of the numbers $E_2(\tau), E_4(\tau), E_6(\tau)$ are algebraically independent over $\overline{\mathbb{Q}}$. In $1996$, Nesterenko made a significant advancement in this direction by proving the following theorem \cite{nesterenko}. 
\begin{theorem}[Nesterenko] \label{nest}
If $\tau\in \mathbb{H}$, then at least three of the numbers
$$
e^{2 \pi i \tau},\quad E_2(\tau),\quad E_4(\tau),\quad E_6(\tau)
$$
are algebraically independent over $\overline{\mathbb{Q}}$.
\end{theorem}
 

Recall that the $j$-function is given by
$$
j(\tau)=\frac{E_4(\tau)^3}{\Delta(\tau)}=\frac{1728\,E_4(\tau)^3}{E_4(\tau)^3-E_6(\tau)^2}.
$$
An important consequence of Theorem \ref{nest} is the following result which was conjectured earlier by Mahler and proved in a joint work of Barr\'e-sirieix, Diaz, Gramain and Philibert \cite{b-d-g-p}. 
\begin{theorem}[Barr\'e-sirieix et al]\label{mahl}
For any $\tau\in \mathbb{H}$, at least one of the two numbers $e^{2\pi i\tau}$ and $j(\tau)$ is transcendental.
\end{theorem}
\bigskip

An element $\tau\in\mathbb{H}$ is called a CM point if it generates
a quadratic extension over the field of rational numbers. It is known, from theory of complex multiplication, that if $\tau$ is a CM point, then $j(\tau)$ is an algebraic number. For non-CM algebraic points in $\mathbb{H}$, Schneider \cite{schneider} proved the following result:
\begin{theorem}[Schneider]\label{schneider}
    If $\tau\in\mathbb{H}$ is algebraic, then $j(\tau)$ is algebraic if and only if $\tau$ is CM.
\end{theorem}
Note that the theorem of Schneider does not follow from Nesterenko's Theorem \ref{nest}. Nesterenko pointed out that both Theorem \ref{nest} and Theorem \ref{schneider} will follow from the following conjecture.
\begin{conjecture*}
If $\tau\in\mathbb{H}$ such that at most three of the numbers 
$$
\tau,\quad e^{2\pi i\tau},\quad E_2(\tau),\quad E_4(\tau),\quad E_6(\tau)
$$
are algebraically independent, then $\tau$ is necessarily a CM point.
\end{conjecture*}
\bigskip
\subsection{Special values of modular forms}\quad \bigskip

 In 2010, an extensive study on algebraic independencece of values of modular forms was done by M. R. Murty, P. Rath and S. Gun \cite{r-m-g}.  Here we mention some of their results.
 \begin{theorem}\label{algem}
     Let $f$ be a non-zero modular form of weight $k$ for $SL_2(\mathbb{Z})$. If the Fourier coefficients of $f$ are algebraic, then
     \begin{enumerate}
         \item any zero of $f$ is either CM or transcendental.
         \item $f(\tau)$ is algebraically independent with $e^{2\pi i\tau}$ except when $f(\tau)=0$.
     \end{enumerate}
 \end{theorem}
In another paper of A. Hamieh and M. R. Murty \cite{hamieh-murty}, they pointed out that similar results hold for higher level modular forms. 
\begin{theorem}\label{algem'}
  Let $f$ be a non-zero modular form of weight $k$ and level $N$ with algebraic Fourier coefficients.\begin{enumerate}
      \item If $\tau\in\mathbb{
      H
      }$ such that $e^{2\pi i\tau}$ is algebraic, then $f(\tau)$ is transcendental.
      \item If $\tau\in\mathbb{H}$ such that  such that $j(\tau)\in \overline{\mathbb{Q}}$, then there exists a transcendental number $\omega_\tau$ which depends only on $\tau$ and is $\overline{\mathbb{Q}}$-linearly independent with $\pi$ such that $\left(\frac{\pi}{\omega_\tau}\right)^k f(\tau)\in \overline{\mathbb{Q}}$. Therefore, $f(\tau)$ is either zero or transcendental.
  \end{enumerate}  
\end{theorem}
\begin{proof}
   (1) We proceed through the following two steps:\\

\textbf{Step 1} : If $N=1$, then $f$ is a modular form of weight $k$ with respect to $S L_2(\mathbb{Z})$. By Theorem \ref{basis}, $f$ can be expressed as follows :
\begin{equation}\label{level1}
f(\tau)=\sum_{\substack{4\alpha+6\beta=k\\ \alpha,\, \beta\geq 0}}c_{\alpha\beta}\,E_4(\tau)^\alpha\, E_6(\tau)^\beta\quad \text{for all } \tau\in\mathbb{H}
.\end{equation}
Since $f$ is non-zero modular form with algebraic Fourier coefficients and the Fourier coefficients of $E_4,\, E_6$ are integers, all the coefficients $c_{\alpha\beta}\in\overline{\mathbb{Q}}$, and at least one coefficient is non-zero. Given that $e^{2\pi i\tau}\in\overline{\mathbb{Q}}$, and so the numbers $E_4(\tau)$ and $E_6(\tau)$ are algebraically independent by Nesterenko's Theorem \ref{nest}. Therefore, the above expression implies that $f(\tau)$ is transcendental.\\

\textbf{Step 2} : For $N>1$, we use a different technique. We first construct the function $$g(\tau):=\frac{f(\tau)^{12}}{\Delta(\tau)^k}\quad \forall \tau\in\mathbb{H},$$ which is modular invariant under action of $\Gamma(N)$ because $f\in M_k(\Gamma(N))$. Since $f(\tau)$ is holomorphic on  $\mathbb{H^*}$ and $\Delta(\tau)$ is non-vanishing on $\mathbb{H}$ with a simple zero at $i\infty$, the function $g(\tau)$ is holomorphic on $\mathbb{H}$ and meromorphic at all cusps. Moreover, $g(\tau)$ has algebraic Fourier coefficients because the Fourier coefficients of $\Delta(\tau)$ and $f(\tau)$ are algebraic. Thus, we conclude that $g$ belongs to $F_{N,\overline{\mathbb{Q}}}$.\\

From Theorem \ref{mdft}, we know that the minimal polynomial of $g$ is of the form: $$P(X)=\sum_{r=0}^m c_r(j)X^r\in \overline{\mathbb{Q}}(j)[X],\quad \deg P(X)=m.$$ Moreover, Corollary \ref{useful} shows that $P(X)$ is a monic polynomial and belongs to $\overline{\mathbb{Q}}[j][X]$ i.e., $$
\begin{aligned}
c_r(j)&=\sum_{s=0}^{d_r}c_{r s}j^s\in \overline{\mathbb{Q}}[j]\quad \forall\ 0 \leq r\leq m-1  \quad \text{with }d_r=\deg c_r(j),\\  
   c_m(j)&=c_{m 0}.
\end{aligned}
$$ Here $c_0(j)$ cannot be the zero polynomial of $j$ because $P(X)$ is irreducible over $\overline{\mathbb{Q}}(j)$.\\

Now we have $$P\left(\frac{f(\tau)^{12}}{\Delta(\tau)^k}\right)=0\quad \text{for all } \tau\in\mathbb{H}.$$
For any positive integer $l$, we get
$$
\begin{aligned}
& \Delta(\tau)^l P\left(\frac{f(\tau)^{12}}{\Delta(\tau)^k}\right)=0\\
\Rightarrow &\sum_{r=0}^m\sum_{s=0}^{d_r}c_{r,s}\,j(\tau)^s f(\tau)^{12 r}\Delta(\tau)^{l-k r}=0\\
\Rightarrow &\sum_{r=0}^m\sum_{s=0}^{d_r}c_{r,s}\,E_4(\tau)^{3 s}f(\tau)^{12 r}\Delta(\tau)^{l-k r-s}=0.
\end{aligned}
$$
We choose $l> k m +\operatorname{max}\{d_0,d_1,\ldots,d_m\}$ so that the powers of $\Delta(\tau)$ are always positive integers. Then we use $$\Delta(\tau)=\frac{E_4(\tau)^3-E_6(\tau)^2}{1728}$$ to obtain \begin{equation} \label{trane}
\sum_{r=0}^m \sum_{s=0}^{d_r} \sum_{t=0}^{l-k r-s} \frac{(-1)^t c_{r, s}}{1728^{l-k r-s}}\left(\begin{array}{c}
l-k r-s \\
t
\end{array}\right) f(\tau)^{12 r} E_4(\tau)^{3(l-k r-t)} E_6(\tau)^{2 t}=0.   
\end{equation}
Note that the equation \eqref{trane} holds for all $\tau\in\mathbb{H}$. Now we consider $\tau\in \mathbb{H}$ such that $e^{2\pi i\tau}$ is algebraic. Then we have $f(\tau)\neq 0$, as otherwise $$c_0(j(\tau))=P\left(\frac{f(\tau)^{12}}{\Delta(\tau)^k}\right)=0.$$ This implies that $j(\tau)\in\overline{\mathbb{Q}}$, which contradicts Theorem \ref{mahl}. Moreover, we claim that the number $f(\tau)$ is not algebraic. For otherwise, \eqref{trane} becomes a polynomial in two variables over $\overline{\mathbb{Q}}$, satisfied by $E_4(\tau)$  and $E_6(\tau)$. This polynomial is non-zero because the coefficient of the term $E_4(\tau)^{3(l-k m)}\,E_6(\tau)^{0}$ (which appears in \eqref{trane} only for $r=m,\, t=0$) is given by $$\frac{1}{1728^{l-k m}}\cdot c_{m 0}\cdot f(\tau)^{12 m}\neq 0.$$ This shows that if $f(\tau)$ is algebraic, the equation \eqref{trane} gives  an algebraic dependence relation between $E_4(\tau)$ and $E_6(\tau)$, which contradicts Theorem \ref{nest}. Thus, for any $\tau\in \mathbb{H}$ such that $e^{2 \pi i\tau}$ is algebraic, the number $f(\tau)$ is transcendental.
\bigskip

(2) We deduce this part in the next section 
as a corollary of another result regarding values of modular functions.
\end{proof}
\begin{corollary}\label{algi}
If a modular form $f\in M_k(\Gamma(N))$ has algebraic Fourier coefficients, then $f$ is algebraically dependent with $E_4$ and $E_6$ over $\overline{\mathbb{Q}}$.
\begin{proof}
For $N=1$, it directly follows from the equation \eqref{level1}. For $N>1$, the assertion follows from the equation \eqref{trane} because the left-hand side is a non-zero polynomial in $f, E_4$ and $E_6$ over $\overline{\mathbb{Q}}.$
\end{proof}  
\end{corollary}
\bigskip

 Investigation on transcendence on special values of quasi-modular forms for $SL_2(\mathbb{Z})$ also has been carried out by C. Y. Chang \cite{quasi}.
 \begin{theorem}
 Let $f$ be a quasi-modular form of non-zero weight for $SL_2(\mathbb{Z})$ with algebraic Fourier coefficients. If $\tau\in\mathbb{H}$ such that $j(\tau)\in\overline{\mathbb{Q}}$ and $f(\tau)\neq 0$, then $f(\tau)$ is transcendental.    
 \end{theorem}
We shall prove all these results including higher level cases in the next section. 
\bigskip

\subsection{\bf Special values of the Weierstrass $\wp$-function}
\quad \bigskip

Here we mention some results regarding the transcendence of special values of the Weierstrass $\wp$-function and deduce important corollaries which are used later. At first, we note that the $\wp$-function satisfies the following addition formula (\cite{r-m}, chapter $11$).
 \begin{theorem} \label{addf}
 For $z_1, z_2\in \mathbb{C}$ such that $z_1\pm z_2\notin L$ we have
$$
\wp\left(z_1+z_2\right)=-\wp\left(z_1\right)-\wp\left(z_2\right)+\frac{1}{4}\left(\frac{\wp^{\prime}\left(z_1\right)-\wp^{\prime}\left(z_2\right)}{\wp\left(z_1\right)-\wp\left(z_2\right)}\right)^2 .
$$
When $z_1=z_2=z$, we have
$$
\wp(2 z)=-2 \wp(z)+\frac{1}{4}\left(\frac{\wp^{\prime \prime}(z)}{\wp^{\prime}(z)}\right)^2 .
$$
 \end{theorem}
By using this addition formula together with a remarkable theorem known as Schneider–Lang theorem, one can prove the following result (\cite{r-m}, Theorem $11.2$).
\begin{theorem}\label{tranw}
    Let $L=\omega_1\,\mathbb{Z}\oplus\omega_2\,\mathbb{Z}$ be a lattice and both $g_2(L)$, $g_3(L)$ are algebraic. If $\alpha$ is an algebraic number such that $\alpha\notin L$, then $\wp(\alpha)$ is transcendental.
\end{theorem}

We now recall the differential equation \eqref{dell}
\begin{equation}
\begin{aligned}
 \wp'(z)^2=4\wp(z)^3-g_2(L)\wp(z)-g_3(L)
 \end{aligned}
\end{equation}
and the fact that 

$$
\wp'\left(\frac{\omega_1}{2}\right)=\wp'\left(\frac{\omega_2}{2}\right)=\wp'\left(\frac{\omega_1+\omega_2}{2}\right)=0.
$$
This shows that if $g_2(L)$ and $g_3(L)$ are algebraic, then the three numbers
$$
\wp\left(\frac{\omega_1}{2}\right),\ \wp\left(\frac{\omega_2}{2}\right),\ \wp\left(\frac{\omega_1+\omega_2}{2}\right)
$$
are also algebraic. Thus, we can deduce the following corollaries of Theorem \ref{tranw} :
\begin{corollary}\label{tranp}
    Let the Weierstrass $\wp$ function be associated with $L=\omega_1\,\mathbb{Z}\oplus\omega_2\,\mathbb{Z}$ such that both $g_2(L)$ and $g_3(L)$ are algebraic. Then, for any natural number $n$, the numbers $\wp(\frac{\omega_1}{n})$ and $\wp(\frac{\omega_2}{n})$ are algebraic. 
\end{corollary}
\begin{proof}
    By the addition formula from Theorem \ref{addf} together with the fact that the numbers  $\wp\left(\frac{\omega_1}{2}\right), \wp\left(\frac{\omega_2}{2}\right), \wp\left(\frac{\omega_1+\omega_2}{2}\right)$ are algebraic, we get this result.
\end{proof}
\begin{corollary}\label{tranp.}
   Let the Weierstrass $\wp$ function associated with $L=\omega_1\,\mathbb{Z}\oplus\omega_2\,\mathbb{Z}$ such that both $g_2(L)$ and $g_3(L)$ are algebraic. Then, any non-zero period of $\wp(z)$ is necessarily transcendental. 
\end{corollary}
\begin{proof}
   Since $\wp\left(\frac{\omega_1}{2}\right), \wp\left(\frac{\omega_2}{2}\right), \wp\left(\frac{\omega_1+\omega_2}{2}\right)$ are algebraic numbers,  Theorem \ref{tranw} implies that $\omega_1/2$, $\omega_2/2$ and $(\omega_1+\omega_2)/2$ are transcendental. Then we can divide any non-zero period by suitable power of $2$ and use transcendence of these three numbers to obtain this result. 
\end{proof}

























\section{\bf Main Results}
\bigskip

Recall that for a positive integer $k$ and a congruence subgroup $\Gamma$ we have 
$$
\begin{aligned}
M_{k}(\Gamma)&:=\{\text{ The space of all modular forms of weight }k\text{ with respect to }\Gamma\}\\
M_{k,N,\overline{\mathbb{Q}}}&:=\{ f\in M_k(\Gamma(N)): f\text{ is a non-zero and it has algebraic Fourier coefficients}\}\\
\widetilde{M}_{k}^{(\leq p)}(\Gamma)&:=\{\text{ The space of all quasi-modular forms of weight k and depth } \leq p \text{ for }\Gamma\}\\
F_{N,\overline{\mathbb{Q}}}&:=\{\text{ The field of modular functions of level $N$ with algebraic Fourier coefficients} \}\\
&\,\,=\overline{\mathbb{Q}}\left(\left\{j,f_a| a\in S\right\}\right).
\end{aligned}
$$
Moreover, we use $f,\, \widetilde{f}$ and $g$ to denote arbitrary \emph{non-constant} elements of $M_{k,N,\overline{\mathbb{Q}}},\,\widetilde{M}_{k}^{(\leq p)}(\Gamma(N))$ and $F_{N,\overline{\mathbb{Q}}}$ respectively.
\bigskip

\subsection{\bf Values of modular functions when $j(\tau)$ is transcendental}\bigskip

\begin{lemma} \label{mdfz}
    Let $g\in  F_{N,\overline{\mathbb{Q}}}$, and let $\tau\in \mathbb{H}$ such that it is not a pole of $g$. If $j(\tau)\notin\overline{\mathbb{Q}}$, then $g(\tau)\neq 0$.
\end{lemma}
\begin{proof}
By Theorem \ref{mdft..}, we have an irreducible polynomial 
$$
P(X)=\sum_{r=0}^m c_r(j)\,X^r \in \overline{\mathbb{Q}}[j][X], \quad m=\deg P(X)
$$
such that
$$
P(g)=0.
$$
Here $$c_r(j)=\sum_{s=0}^{d_r}c_{r s}\,j^s\in \overline{\mathbb{Q}}[j]\quad \forall\ 0\leq r\leq m, \quad d_r=\deg c_r(j).$$
Also, $c_0(j)$ is not identically zero polynomial because $P(X)$ is irreducible over $\overline{\mathbb{Q}}[j]$. Thus, we get
\begin{equation}\label{funzero}
\sum_{r=0}^m\sum_{s=0}^{d_r}c_{r s}\,j(\tau)^s\, g(\tau)^r=0\quad \forall \tau\in\mathbb{H}.
\end{equation}
Now we consider $\tau\in\mathbb{H}$ such that $e^{2\pi i\tau}$ is algebraic. We claim that $g(\tau)\neq 0$. If not, from \eqref{funzero} we get $c_0(j(\tau))=0$ and hence, $j(\tau)\in\overline{\mathbb{Q}}$, which contradicts Theorem \ref{mahl}. This completes the proof that $g(\tau)$ is non-zero. 
\end{proof}
The following theorem shows that the non-zero number $g(\tau)$ is indeed transcendental.
\begin{theorem}\label{tranmf}
 Let $g\in F_{N,\overline{\mathbb{Q}}}$ and $\tau\in \mathbb{H}$ such that it is not a pole of $g$. If $j(\tau)\notin\overline{\mathbb{Q}}$, then $g(\tau)$ is transcendental.
\end{theorem}
\begin{proof}
If possible, let $g(\tau)\in\overline{\mathbb{Q}}$. Then the following modular function $$h(z):=g(z)-g(\tau)\quad \forall z\in\mathbb{H}$$ belongs to $ F_{N,\overline{\mathbb{Q}}}$ but vanishes at $\tau$, which contradicts Lemma \ref{mdfz}. Therefore, $g(\tau)$ is a transcendental number.
\end{proof}
\begin{corollary}
If $g\in F_{N,\overline{\mathbb{Q}}}$ and $\tau_0$ is a zero or a pole of $g$, then $j(\tau_0)\in \overline{\mathbb{Q}}$. 
\end{corollary}
\begin{proof}
If $g(\tau_0)=0$, then Lemma \ref{mdfz} implies that $j(\tau_0)\in \overline{\mathbb{Q}}$. If $\tau_0$ is a pole of $g$, we consider the function
$$
h(\tau):=\frac{1}{g(\tau)}\in F_{N,\overline{\mathbb{Q}}},
$$
which satisfies $h(\tau_0)=0$, and hence $j(\tau_0)\in\overline{\mathbb{Q}}$ by Lemma \ref{mdfz}. 
\end{proof}
\bigskip

\subsection{Values of modular functions when $j(\tau)\in\overline{\mathbb{Q}}$}
\quad \bigskip

Here we study values of modular functions at the points on $\mathbb{H}$ where $j$-function takes algebraic values. From Theorem \ref{mdft..}, we already know that $$ F_{N,\overline{\mathbb{Q}}}\subset  \overline{\mathbb{Q}}\left(\{j,f_a |a\in S\}\right), \text{ where }S=\left\{\left(\frac{r}{N},\frac{s}{N}\right): 0\leq r, s\leq N-1 \text{ and }(r,s)\neq (0,0)\right\}.$$
Hence, for any $g\in F_{N,\overline{\mathbb{Q}}}$, the algebraic nature of $g(\tau)$ is completely determined by the numbers in the set $\{f_a(\tau):a\in S\}$.\\

Note that for $N=1$, the set $S$ is empty. This shows that for any $g\in F_{1,\overline{\mathbb{Q}}}$, the number $g(\tau)\in\overline{\mathbb{Q}}$ as $j(\tau)\in\overline{\mathbb{Q}}$. For $N>1$, we have to study the numbers $f_a(\tau)$ for all $a\in S$. We start by proving the following important lemma.\bigskip

\begin{lemma}\label{omegatau}
 If $\tau\in \mathbb{H}$ such that $j(\tau)$ is algebraic, then there exists a unique (up to an algebraic multiple) transcendental number $\omega_\tau$ for which $g_2(\omega_\tau L_\tau)$ and $g_3(\omega_\tau L_\tau)$ are both algebraic numbers.   
\end{lemma}
\begin{proof}
Note that for any $\omega_\tau\in\mathbb{C}^\times $,
\begin{equation}\label{omegatau.}
    j(\tau)=\frac{E_4(\tau)^3}{\Delta(\tau)}=\frac{1728\,  g_2(L_\tau)^3}{\Delta_0(L_\tau)}=\frac{1728\, g_2\left(\omega_\tau L_\tau\right)^3}{\Delta_0\left(\omega_\tau L_\tau\right)},
\end{equation}
where $L_\tau=\tau\,\mathbb{Z}\oplus\mathbb{Z}$ and $\omega_\tau L_\tau=\tau\,\omega_\tau\,\mathbb{Z}\oplus\omega_\tau\mathbb{Z}$. Here the last equality follows from the fact that \begin{equation}\label{homogenity}
g_2(\omega_\tau L_\tau)=\omega_\tau^{-4}g_2(L_\tau).
\end{equation} If we choose $\omega_\tau$ such that $\omega_\tau^4=g_2(L_\tau)$, then $g_2(\omega_\tau L_\tau)=1$. Given that $j(\tau)$ is algebraic, and so \eqref{omegatau.} shows that $\Delta_0(\omega_\tau L_\tau)$ is also algebraic. Thus, the numbers $g_2(\omega_\tau L_\tau)$ and $g_3(\omega_\tau L_\tau)$ both are algebraic. Moreover,  $\omega_\tau$ is a period of $\omega_\tau L_\tau$ and $g_2(\omega_\tau L_\tau),\, g_3(\omega_\tau L_\tau)\in\overline{\mathbb{Q}}$. By Corollary \ref{tranp.} we conclude that $\omega_\tau$ is transcendental.\\

To prove uniqueness, consider an arbitrary $\omega_\tau'\in\mathbb{C}^\times$ such that both $g_2(\omega_\tau' L_\tau),\, g_3(\omega_\tau' L_\tau)\in\overline{\mathbb{Q}}$. From \eqref{homogenity} we have 
$$
 \quad g_2(\omega_\tau' L_\tau)= (\omega_\tau')^{-4}\,g_2(L_\tau)=\alpha, (\text{say}).
$$
Therefore, we get $g_2(L_\tau)=\alpha\cdot (\omega_\tau')^{4}$ and hence $\omega_\tau^4=\alpha\cdot (\omega_\tau')^{4}.$
This completes the proof that $\omega_\tau$ is unique up to an algebraic multiple.
\end{proof}
\bigskip

Recall that for $\tau\in\mathbb{H}$
$$
f_a(\tau):=\frac{g_2(L_\tau)\,g_3(L_\tau)}{\Delta_0(L_\tau)}\wp_{L_\tau}(a_1\tau+a_2),
$$ 
where $L_\tau=<\tau,1>=\tau\,\mathbb{Z}\oplus\mathbb{Z}$ and $a=(a_1,a_2)\in S$. We can rewrite it as follows:
$$
\begin{aligned}
    f_a(\tau)&=\frac{\omega_\tau^{-4}\,g_2(L_\tau)\cdot \omega_\tau^{-6}\, g_3(L_\tau)}{\omega_\tau^{-12}\,\Delta_0(L_\tau)}\, \omega_\tau^{-2}\,\wp_{L_\tau}(a_1\tau+a_2)\\
    &=\frac{g_2(\omega_\tau L_\tau)\, g_3(\omega_\tau L_\tau)}{\Delta_0(\omega_\tau L_\tau)}\,\omega_\tau^{-2}\,\wp_{L_\tau}(a_1\tau+a_2). 
\end{aligned}
$$
Note that 
$$
\begin{aligned}
    \wp_{\omega_\tau L_\tau}(a_1\omega_\tau \tau+a_2\,\omega_\tau)&=\frac{1}{(a_1\,\omega_\tau\,\tau+a_2\,\omega_\tau)^2}\\
    & +\sideset{}{'}\sum_{(m,n)\in \mathbb{Z}^2}\left[\frac{1}{\left(a_1\omega_\tau \tau+a_2\omega_\tau-m\omega_\tau\tau-n\omega_\tau\right)^2}-\frac{1}{(m\omega_\tau\tau+n\omega_\tau)^2}\right]\\
    &=\frac{1}{\omega_\tau^2}\left[\frac{1}{(a_1\tau+a_2)^2}+\sideset{}{'}\sum_{(m,n)\in \mathbb{Z}^2}\frac{1}{(a_1\,\tau+a_2-m\tau-n)^2}-\frac{1}{(m\tau+n)^2}\right]\\
    &=\frac{1}{\omega_\tau^2}\,\wp_{L_\tau}(a_1\tau+a_2).
\end{aligned}
$$
Thus, we get 
$$
f_a(\tau)=\frac{g_2(\omega_\tau L_\tau)\, g_3(\omega_\tau L_\tau)}{\Delta_0(\omega_\tau L_\tau)}\,\wp_{\omega_\tau L_\tau}(a_1\omega_\tau\tau+a_2\omega_\tau).
$$
\bigskip

\begin{theorem}\label{algfa}
    If $\tau\in\mathbb{H}$ such that $j(\tau)\in\overline{\mathbb{Q}}$, then $f_a(\tau)\in\overline{\mathbb{Q}}$ for all $a\in S$.
\end{theorem}
\begin{proof}
    Let $a=(\frac{r}{N},\frac{s}{N})\in S$. Then we have
$$
f_a(\tau)=\frac{g_2(\omega_\tau L_\tau)\, g_3(\omega_\tau L_\tau)}{\Delta_0(\omega_\tau L_\tau)}\,\wp_{\omega_\tau L_\tau}\left(\frac{r\,\omega_\tau\,\tau}{N}+\frac{s\,\omega_\tau}{N}\right),
$$
where $\omega_\tau$ is chosen as in Lemma \ref{omegatau} so that $g_2(\omega_\tau L_\tau),\, g_3(\omega_\tau L_\tau)\in\overline{\mathbb{Q}}$.\\

Recall that both $\wp_{\omega_\tau L_\tau}\left(\frac{\omega_\tau\, \tau}{N} \right)$ and $\wp_{\omega_\tau L_\tau}\left(\frac{\omega_\tau}{N} \right)$  are algebraic by Lemma \ref{tranp}. Moreover, by the addition formula for Weierstrass $\wp$-function in Theorem \ref{addf}, we get that $$\wp_{\omega_\tau L_\tau}\left(\frac{r\,\omega_\tau\, \tau}{N}+\frac{s\,\omega_\tau}{N}\right)\in\overline{\mathbb{Q}}$$ for all $a=\left(\frac{r}{N},\frac{s}{N}\right)\in S$. Thus, the number $f_a(\tau)$ is algebraic for all $a\in S$.
\end{proof}
\begin{theorem}\label{algemf}
Let $g\in F_{N,\overline{\mathbb{{Q}}}}$ and $\tau\in\mathbb{H}$ such that $j(\tau)\in \overline{\mathbb{Q}}$. If $\tau$ is not a pole of $g$, then $g(\tau)$ is algebraic.
\end{theorem}
\begin{proof}
 From Theorem \ref{algfa}, we have   $f_a(\tau)\in\overline{\mathbb{Q}}$ for all $a\in S$. Then we use the following fact $$g\in  F_{N,\overline{\mathbb{Q}}}=  \overline{\mathbb{Q}}\left(\{j,f_a |a\in S\}\right)$$  to conclude that $g(\tau)$ is algebraic. 
\end{proof}
\begin{corollary}\label{suma}
Let $g\in F_{N,\overline{\mathbb{Q}}}$ and $\tau\in\mathbb{H}$ such that $\tau$ is not a pole of $g$. Then

(i) $g(\tau)\in\overline{\mathbb{Q}} \Leftrightarrow j(\tau)\in\overline{\mathbb{Q}}$.

(ii) At least one of $e^{2\pi i\tau}$ and $g(\tau)$ is transcendental.
\begin{proof}
    Combining Theorem \ref{tranmf} and Theorem \ref{algemf} we get (i), and (ii) follows from Theorem \ref{mahl} and (i).
\end{proof}
\end{corollary}
\bigskip

\subsection{\bf Values of modular forms} \quad \bigskip

For an arbitrary element $f\in M_{k,N,\overline{\mathbb{Q}}}$, we study transcendence of $f(\tau)$ for two cases: one if $e^{2\pi i\tau}\in\overline{\mathbb{Q}}$, and another if $j(\tau)\in\overline{\mathbb{Q}}$. Unfortunately, the current techniques do not allow us to deal with the case when both $e^{2\pi i\tau},\, j(\tau)\notin\overline{\mathbb{Q}}$.\\

Recall that in the previous section, we have already proved the first part of  Theorem \ref{algem'}, following the technique from \cite{hamieh-murty}, which states the following:
\begin{theorem}\label{tranm}
  If $f\in M_{k,N,\overline{\mathbb{Q}}}$ and $\tau\in \mathbb{H}$ such that $e^{2\pi i\tau}$ is algebraic, then $f(\tau)$ is transcendental. 
\end{theorem}
 Here we shall prove the second part of Theorem \ref{algem'}.  Recall that if $\tau\in\mathbb{H}$ such that $j(\tau)\in\overline{\mathbb{Q}}$, then there exists a unique (up to an algebraic multiple) transcendental number $\omega_\tau$ for which the numbers $g_2(\omega_\tau L_\tau)$ and $g_3(\omega_\tau L_\tau)$ are both algebraic. Moreover, we have the formulae
$$
\begin{aligned}
E_4(\tau)&=\frac{3}{4\pi^4}\, g_2(L_\tau)=\frac{3}{4}\left(\frac{\omega_\tau}{\pi} \right)^4 g_2(\omega_\tau L_\tau),\\ E_6(\tau)&=\frac{27}{8\pi^6}\,g_3(L_\tau)=\frac{27}{8}\left(\frac{\omega_\tau}{\pi}\right)^6g_3(\omega_\tau L_\tau).
\end{aligned}
$$
From Nesterenko's Theorem \ref{nest}, we know that at most one of $E_4(\tau)$ and $E_6(\tau)$ is algebraic. The above formulae imply that if $\frac{\omega_\tau}{\pi}\in\overline{\mathbb{Q}}$, then the numbers $E_4(\tau),\, E_6(\tau)\in\overline{\mathbb{Q}}$, which is a contradiction. Hence, $\frac{\omega_\tau}{\pi}$ is a transcendental number. Besides, we have 
$$
\begin{aligned}
\Delta(\tau)=\frac{E_4(\tau)^3-E_6(\tau)^2}{1728}&=\frac{1}{4^6}\left(\frac{\omega_\tau}{\pi}\right)^{12}\left(g_2(\omega_\tau L_\tau)^3-27 g_3(\omega_\tau L_\tau)^2 \right)\\
&=\frac{1}{4^6}\left(\frac{\omega_\tau}{\pi}\right)^{12}\Delta_0(\omega_\tau L_\tau).  
\end{aligned}
$$
As $g_2(\omega_\tau L_\tau)$ and $g_3(\omega_\tau L_\tau)$ both are algebraic, the number $\Delta_0(\omega_\tau L_\tau)\in\overline{\mathbb{Q}}\setminus \{0\}$. Since $\frac{\omega_\tau}{\pi}$ is transcendental and $\Delta(\tau)$ is a non-zero algebraic multiple of $\left(\frac{\omega_\tau}{\pi}\right)^{12}$, we deduce that $\Delta(\tau)$ is transcendental. We use this observation in the following theorem to study values taken by the elements of $ M_{k,N,\overline{\mathbb{Q}}}$. \bigskip

\begin{theorem}\label{algem}
If $f\in M_{k,N,\overline{\mathbb{Q}}}$ and $\tau\in\mathbb{H}$ such that $j(\tau)\in \overline{\mathbb{Q}}$, then there exists a transcendental number $\omega_\tau$ (which depends only on $\tau$ and is $\overline{\mathbb{Q}}$-linearly independent with $\pi$) such that $\left(\frac{\pi}{\omega_\tau}\right)^k f(\tau)\in \overline{\mathbb{Q}}$. Therefore, $f(\tau)$ is either zero or transcendental.
\end{theorem}
\begin{proof}
     We first construct the function $$g(\tau):=\frac{f(\tau)^{12}}{\Delta(\tau)^k}\quad \forall \tau\in\mathbb{H},$$ which lies in $ F_{N,\overline{\mathbb{Q}}}$ by the same argument as in Theorem \ref{algem'}. Since $j(\tau)\in\overline{\mathbb{Q}}$, Theorem \ref{algemf} implies that $g(\tau)$ is an algebraic number $c$, (say). Thus, we get  \begin{equation}\label{forms}
f(\tau)^{12}=c.\Delta(\tau)^k=\frac{c}{4^6}.\Delta_0(\omega_\tau L_\tau)^k\left(\frac{\omega_\tau}{\pi} \right)^{12 k}. 
\end{equation}
This shows that $f(\tau)$ is zero when $c=0$, and otherwise, it a non-zero algebraic multiple of $\left(\frac{\omega_\tau}{\pi} \right)^{k}$ and hence transcendental.
\end{proof}
\bigskip

\subsection{Algebraic independence of values of modular forms and functions}\quad \bigskip

Let $f\in M_{k,N,\overline{\mathbb{Q}}}$ and $g\in F_{N,\overline{\mathbb{Q}}}$. By using the transcendental nature of $f(\tau)$ and $g(\tau)$ from previous theorems, we prove the following result.
\begin{theorem}\label{deg.}
    If $\tau\in\mathbb{H}$ such that $f(\tau)\neq 0$ and $\tau$ is not a pole of $g$, then the transcendence degree of the field $\overline{\mathbb{Q}}\left(e^{2\pi i\tau},\, f(\tau),\, g(\tau) \right)$ over $\overline{\mathbb{Q}}$ is at least two, that is
    $$
trdeg \,\overline{\mathbb{Q}}\left(e^{2\pi i\tau},\, f(\tau),\, g(\tau) \right) \geq 2.
$$  
\end{theorem}
\begin{proof}
    By theorem \ref{mdft..}, the numbers $g(\tau)$ and $j(\tau)$ are algebraically dependent over $\overline{\mathbb{Q}}$. Moreover, the number $j(\tau)$ is algebraic if and only if $g(\tau)\in\overline{\mathbb{Q}}$ by Corollary \ref{suma}. These two facts together give the first equality below:
$$
\begin{aligned}
 trdeg\,\overline{\mathbb{Q}}\left(e^{2\pi i\tau},\, f(\tau),\, g(\tau) \right)&= trdeg \,\overline{\mathbb{Q}}\left(e^{2\pi i\tau},\, f(\tau),\, j(\tau) \right)\\
 &= trdeg \,\overline{\mathbb{Q}}\left(e^{2\pi i\tau},\, \frac{f(\tau)^{12}}{\Delta(\tau)^k},\, f(\tau),\, j(\tau) \right)\\
 &= trdeg \,\overline{\mathbb{Q}}\left(e^{2\pi i\tau},\, \Delta(\tau)^k,\, f(\tau),\, j(\tau) \right)\\
 &= trdeg \,\overline{\mathbb{Q}}\left(e^{2\pi i\tau},\, \Delta(\tau),\, f(\tau),\, j(\tau) \right) \\
 &= trdeg \,\overline{\mathbb{Q}}\left(e^{2\pi i\tau},\, E_4(\tau)^3,\, E_6(\tau)^2,\, f(\tau),\, j(\tau) \right)\\
 &=trdeg \,\overline{\mathbb{Q}}\left(e^{2\pi i\tau},\, E_4(\tau),\, E_6(\tau),\, f(\tau),\, j(\tau) \right)\\
 &= trdeg \,\overline{\mathbb{Q}}\left(e^{2\pi i\tau},\, E_4(\tau),\, E_6(\tau),\, f(\tau) \right)\\
 &= trdeg \,\overline{\mathbb{Q}}\left(e^{2\pi i\tau},\, E_4(\tau),\, E_6(\tau)\right).
\end{aligned}
$$ 
Here the second equality follows from the fact that $f^{12}/\Delta^k$ belongs to $F_{N,\overline{\mathbb{Q}}}$ and  $f(\tau)^{12}/\Delta(\tau)^k$ is algebraically dependent with $j(\tau)$ by Theorem \ref{mdft..}. The third and the seventh equality hold as the fields are identical.  
The final equality holds because $f(\tau)$ is algebraic over $\overline{\mathbb{Q}}\left(E_4(\tau), E_6(\tau)\right)$ by Corollary \ref{algi}. Then we apply Nesterenko's Theorem \ref{nest} on the right side of the final equality to conclude that the transcendence degree of the field $\overline{\mathbb{Q}}\left(e^{2\pi i\tau},\, f(\tau),\, g(\tau) \right)$ over $\overline{\mathbb{Q}}$ is at least two.    
\end{proof}
\bigskip

\subsection{\bf Special values of quasi-modular forms} \quad \bigskip

 Here we study the values of quasi-modular forms of weight $k(>0)$, depth $p$ and level $N(\geq 1)$ with algebraic Fourier coefficients. We consider their values at the points $\tau\in\mathbb{H}$ such that one of $e^{2\pi i\tau}$ and $j(\tau)$ is algebraic.
\begin{theorem}
Let $\Gamma$ be a congruence subgroup of level $N(\geq 1)$ and $\widetilde{f}\in \widetilde{M}_{k}^{(p)}(\Gamma)$ with algebraic Fourier coefficients. If $\tau\in\mathbb{H}$ such that $e^{2\pi i\tau}$ is algebraic, then $f(\tau)$ is transcendental. 
\end{theorem}
\begin{proof}
   By Theorem \ref{quasistructure}, we can write $\widetilde{f}$ in the form $$\widetilde{f}=\sum_{r=0}^p g_r E_2^r,\quad g_r\in M_{k-2r}(\Gamma).$$ Since $\widetilde{f}$ is a non-zero quasi-modular form with algebraic Fourier coefficients, each modular coefficient $g_r$ has algebraic Fourier coefficients and at least one $g_r$ is not identically zero. From Corollary \ref{algi}, we know that each number $g_r(\tau)$ is algebraically dependent with $E_4(\tau), E_6(\tau)$. Therefore, we get
 $$ trdeg\,\overline{\mathbb{Q}}(g_0(\tau), g_2(\tau),\ldots, g_p(\tau), E_4(\tau), E_6(\tau))=trdeg\,\overline{\mathbb{Q}}( E_4(\tau), E_6(\tau)).$$
If possible, let both $e^{2\pi i\tau}, \widetilde f(\tau)\in \overline{\mathbb{Q}}$. Since 
$\widetilde f(\tau)=\sum_{r=0}^p g_r(\tau) E_2(\tau)^r$ and $g_r(\tau)\neq 0$ (by Theorem \ref{tranm}), the number $E_2(\tau)$ is algebraic over $\overline{\mathbb{Q}}\left(g_0(\tau), g_1(\tau),\ldots, g_p(\tau), E_4(\tau), E_6(\tau)\right)$ and so
$$
\begin{aligned}
&\quad trdeg\, \overline{\mathbb{Q}}\left(E_2(\tau), g_0(\tau), g_1(\tau),\ldots, g_p(\tau), E_4(\tau), E_6(\tau)\right)\\ &=trdeg\,\overline{\mathbb{Q}}(g_0(\tau), g_1(\tau),\ldots, g_p(\tau), E_4(\tau), E_6(\tau))\\ 
&=trdeg\,\overline{\mathbb{Q}}(E_4(\tau), E_6(\tau)).
\end{aligned}
$$
Since $e^{2\pi i\tau}$ is an algebraic number, the following field containment holds
$$
\overline{\mathbb{Q}}(e^{2\pi i\tau},\, E_2(\tau),\,E_4(\tau), E_6(\tau))\subseteq \overline{\mathbb{Q}}\left(E_2(\tau), g_0(\tau), g_1(\tau),\ldots, g_p(\tau), E_4(\tau), E_6(\tau)\right).
$$
Therefore, we get 
$$
\begin{aligned}
  trdeg\, \overline{\mathbb{Q}}(e^{2\pi i\tau},\, E_2(\tau),\,E_4(\tau), E_6(\tau))& \leq trdeg\,\overline{\mathbb{Q}}\left(E_2(\tau), g_0(\tau), g_1(\tau),\ldots, g_p(\tau), E_4(\tau), E_6(\tau)\right)\\
  &= trdeg\,\overline{\mathbb{Q}}(E_4(\tau), E_6(\tau))\\
  &\leq 2,
\end{aligned}
$$
which contradicts Nesterenko's Theorem \ref{nest}. Hence, $\widetilde{f}(\tau)$ is a transcendental number.
\end{proof}
To study values of quasi-modular forms in the complementary case (i.e., for $j(\tau)\in\overline{\mathbb{Q}}$), we need the following lemma.
\begin{lemma}\label{peri}
 If $\tau\in \mathbb{H}$ such that $j(\tau)$ is algebraic, then there exists a unique (up to algebraic multiple) transcendental number $\omega_\tau$ for which $\frac{\omega_\tau}{\pi}$ and $\frac{\eta_2}{\pi}$ are algebraically independent over $\overline{\mathbb{Q}}$. Here, $\eta_2$ is the second quasi-period with respect to the lattice $\omega_\tau L_\tau$ i.e., $\eta_2:=Z_{\omega_\tau L_\tau}(z+\omega_\tau)-Z_{\omega_\tau L_\tau}(z)$ for all $z\in\mathbb{C}\setminus \omega_\tau L_\tau$, where $Z_{\omega_\tau L_\tau}$ is the Weierstrass zeta function with respect to $\omega_\tau L_\tau$.
\end{lemma}
\begin{proof}
    By Lemma \ref{omegatau}, we already have a transcendental number $\omega_\tau$ such that both $g_2(\omega_\tau L_\tau)$ and $g_3(\omega_\tau L_\tau)$ are algebraic. From the definition of the Weierstrass zeta-function, one gets
    $$
Z_{\omega_\tau L_\tau}(z)= \frac{1}{\omega_\tau}\cdot Z_{L_\tau}\left(\frac{z}{\omega_\tau}\right).
$$
Now we use the definition of quasi-period and the identity \eqref{quasi-period} to obtain 
$$
\begin{aligned}
    \eta_2=\eta_2(\omega_\tau L_\tau)&=Z_{\omega_\tau L_\tau}(z+\omega_\tau)-Z_{\omega_\tau L_\tau}(z)\\
    &= \frac{1}{\omega_\tau}\left(Z_{L_\tau}(z/\omega_\tau+1)-Z_{L_\tau}(z)\right)\\
    &= \frac{1}{\omega_\tau}\,\eta_2(L_\tau)\\
    &=\frac{1}{\omega_\tau}\,G_2(\tau)\\
    &= \frac{1}{3}\cdot\frac{\pi^2}{ \omega_\tau}\,E_2(\tau).
\end{aligned}
$$
Thus, we get the following formulae 
$$
E_2(\tau)=3\,\frac{\omega_\tau}{\pi}\frac{\eta_2}{\pi}, \quad E_4(\tau)=\frac{3}{4}\left(\frac{\omega_\tau}{\pi} \right)^4 g_2(\omega_\tau L_\tau),\quad E_6(\tau)=\frac{27}{8}\left(\frac{\omega_\tau}{\pi}\right)^6g_3(\omega_\tau L_\tau).$$
Note that the numbers $E_2(\tau)$, $E_4(\tau)$ and $E_6(\tau)$ are algebraic over $\overline{\mathbb{Q}}\left(\frac{\omega_\tau}{\pi}, \frac{\eta_2}{\pi}\right)$. If possible, let $\frac{\omega_\tau}{\pi}, \frac{\eta_2}{\pi}$ be algebraically dependent. Then the above formulae imply that $trdeg\,\overline{\mathbb{Q}}(E_2(\tau),E_4(\tau),E_6(\tau))=1$, which contradicts Nesterenko's Theorem \ref{nest}. Hence, the numbers $\frac{\omega_\tau}{\pi}$ and $\frac{\eta_2}{\pi}$ are algebraically independent over $\overline{\mathbb{Q}}$.
\end{proof}

This lemma and Theorem \ref{quasistructure} together allow us to describe values of quasi-modular forms explicitly at the points $\tau\in\mathbb{H}$ where $j$-function is algebraic.
\begin{theorem}\label{tranq}
 Let $\Gamma$ be a congruence subgroup of level $N(\geq 1)$ and $\widetilde{f}\in \widetilde{M}_{k}^{(p)}(\Gamma)$ with algebraic Fourier coefficients. If $\tau\in\mathbb{H}$ such that $j(\tau)\in\overline{\mathbb{Q}}$, then $f(\tau)$ is a polynomial in $\frac{\omega_\tau}{\pi}$ and $\frac{\eta_2}{\pi}$ with algebraic coefficients. Therefore, $f(\tau)$ is either zero or transcendental.   
\end{theorem}
\begin{proof}
    From Theorem \ref{quasistructure}, we have the expression $$\widetilde{f}(\tau)=\sum_{r=0}^p g_r(\tau) E_2(\tau)^r,\quad \text{where } g_r\in M_{k-2 r}(\Gamma).$$ Since $f$ is a non-zero quasi-modular form with algebraic Fourier coefficients, each $g_r\in M_{k-2r}(\Gamma)$ has algebraic Fourier coefficients and at least one $g_r$ is not identically zero. Then we use Theorem \ref{algem} and the formula for $E_2(\tau)$ such that
\begin{equation}\label{last}
\widetilde{f}(\tau)=\sum_{r=0}^p c_r\left(\frac{\omega_\tau}{\pi}\right)^{k-2 r}\left(\frac{\omega_\tau}{\pi}\cdot\frac{\eta_2}{\pi}\right)^r=\sum_{r=0}^p\, c_r\, \left(\frac{\omega_\tau}{\pi}\right)^{k-r}\, \left(\frac{\eta_2}{\pi}\right)^r,
 \end{equation}
 where each $c_r$ is an algebraic constant. If $\widetilde{f}(\tau)\in\overline{\mathbb{Q}}\setminus \{0\}$, then the above expression of $f(\tau)$ shows that $\frac{\omega_\tau}{\pi}$ and $\frac{\eta_2}{\pi}$ are algebraically dependent, which contradicts Lemma \ref{peri}. Therefore, the number $\widetilde{f}(\tau)$ is either zero or transcendental.     
\end{proof}
\bigskip

\subsection{Algebraic independence of values of modular forms, functions and quasi-modular forms}
\quad \bigskip

In the previous theorems, we have studied the transcendental nature of values of modular functions, forms and quasi-modular forms respectively. Here we study algebraic independency of these numbers over $\overline{\mathbb{Q}}$. 

\begin{theorem}\label{deg..}
    Let $f,\, g$ and $\widetilde f$ denote modular form, modular function and quasi-modular form (of positive depth) respectively for a congruence subgroup $\Gamma$ with algebraic Fourier coefficients. If $\tau\in\mathbb{H}$ such that $f(\tau)\neq 0$ and $\tau$ is not a pole of $g$, then 
    $$
trdeg\,\overline{\mathbb{Q}}\left( e^{2\pi i\tau},\, f(\tau),\, g(\tau),\, \widetilde f(\tau)\right)\geq 3. 
    $$
\end{theorem}
\begin{proof}
 By the same argument as in Theorem \ref{deg.}, we get
\begin{equation}\label{final.}
trdeg\,\overline{\mathbb{Q}}\left( e^{2\pi i\tau},\, f(\tau),\, g(\tau),\, \widetilde f(\tau)\right)= trdeg\,\overline{\mathbb{Q}}\left( e^{2\pi i\tau},\, E_4(\tau),\, E_6(\tau),\, \widetilde f(\tau)\right).
\end{equation}
Moreover, we can apply Theorem \ref{quasistructure} to express the quasi-modular form $\widetilde f$ in the form
  $$
\widetilde{f}(\tau)=\sum_{r=0}^p g_r(\tau)\, E_2(\tau)^r,
  $$
  where each $g_r\in M_{k-2 r}(\Gamma)$ has algebraic Fourier coefficients. Thus, we obtain 
  $$
\begin{aligned}
&trdeg\,\overline{\mathbb{Q}}\left( e^{2\pi i\tau},\, E_4(\tau),\, E_6(\tau),\, \widetilde f(\tau)\right)\\
  & =trdeg\,\overline{\mathbb{Q}}\left( e^{2\pi i\tau},\, E_4(\tau),\, E_6(\tau),\,\sum_{r=0}^p g_r(\tau)\, E_2(\tau)^r \right)\\
&=trdeg\,\overline{\mathbb{Q}}\left( e^{2\pi i\tau}, E_4(\tau), E_6(\tau), g_0(\tau),\ldots,g_p(\tau),\sum_{r=0}^p g_r(\tau)\, E_2(\tau)^r \right)\\
&=trdeg\,\overline{\mathbb{Q}}\left( e^{2\pi i\tau},\, E_4(\tau),\, E_6(\tau),\, g_0(\tau),\,\ldots,\,g_p(\tau),\,E_2(\tau) \right)\\
&=trdeg\,\overline{\mathbb{Q}}\left( e^{2\pi i\tau},\, E_4(\tau),\, E_6(\tau),\,E_2(\tau) \right),
\end{aligned}
 $$
where the second and fourth equalities follow from the fact that each coefficient $g_r(\tau)$ is algebraic over $\overline{\mathbb{Q}}\left(E_4(\tau), E_6(\tau)\right)$ by Corollary \ref{algi}. Moreover, Nesterenko's theorem \ref{nest} implies that the right side of the final equality is at least $3$, and then we use the identity \eqref{final.} to conclude the proof. 
\end{proof}
\bigskip

 \section{\bf Concluding Remarks}
\bigskip

We observe that the proof of Theorem \ref{deg..} mainly uses our knowledge on transcendence of values of modular functions and Nesterenko's Theorem \ref{nest}. Again, if we choose 
$$
f=E_4, \, g=j,\, \widetilde f= E_6+E_4\, E_2,
$$
the statement of Theorem \ref{deg..} implies Nesterenko's theorem. Therefore, Theorem \ref{deg..} is essentially the Nesterenko's theorem in more general setup.
\par
$$ \quad $$
\noindent {\bf Acknowledgements.}





\begin{thebibliography}{0}

\bibitem{ahlfors}
Lars V. Ahlfors, \textit{An Introduction to the Theory of Analytic 
Functions of One Complex Variable}, Third Edition, International series in pure and applied mathematics.


\bibitem{b-d-g-p}
K. Barr\'e-Sirieix, G. Diaz, F. Gramain, G. Philibert, \textit{Une preuve de la conjecture de Mahler-Manin}, Invent. Math., 124 (1996), 1-9.

\bibitem{bertrand}
D. Bertrand, \textit{S\'eries d’Eisenstein et transcendance.} Bull. Soc. Math. France
$104(3)$, $309–321 (1976)$

\bibitem{quasi}
    C. Y. Chang, \textit{Transcendence of special values of quasi-modular forms}, Forum Math. $24$ (2012) $539-551$.


\bibitem{chudnovsky}
    G.V. Chudnovsky, \textit{Algebraic independence of constants connected with the exponential and the elliptic functions}, Dokl.
Akad. Nauk Ukrain. SSR Ser. A $8$ $(1976)$ $698–701$.


\bibitem{stromberg}
H. Cohen and F. Stromberg, \textit{Modular Forms, A Classical Approach}, Graduate Studies In Matheatics $179$.

\bibitem{cohen}
P. Cohen, \textit{On the coefficients of the transformation polynomials for the elliptic modular function}, Math. Proc. Cambridge Phil. Soc. $95$, (1984), 389-402. 

\bibitem{hamieh-murty}
Alia Hamieh and M. Ram Murty, {A note on $q$-analogues of Dirichlet L-functions}, International Journal of Number Theory, Vol. 12 (2016) 765-773.

\bibitem{lang}
Serge Lang, \textit{Elliptic Functions}, Addision-Wesley Publishing Company, Inc.


\bibitem{murty-dewar-graves}
M. R. Murty, M. Dewar and H. Graves, \textit{Problems in 
the Theory of 
Modular
Forms}, Hindustan Book Agency.

\bibitem{r-m}
M. Ram Murty and P. Rath, \textit{Transcendental Numbers} (Springer-Verlag, $2014$).

\bibitem{r-m-g}
S. Gun, M. Ram Murty, P. Rath, \textit{Algebraic independence of values of
modular forms}. Int. J. Number Theory 7(4), 1065–1074 (2011).


\bibitem{nesterenko}
Y. Nesterenko, \textit{Modular functions and transcendence questions}, Mat. Sb. $187$ $(9)$ $(1996)$ $65–96$ (in Russian); translation in:
Sb. Math. $187$ $(9)$ $(1996)$ $1319–1348$.


\bibitem{schneider}
T. Schneider, \textit{Arithmetische Untersuchungen elliptischer Integrale}, Math. Ann.$ 113$ $(1973)$ $1–13$.

\bibitem{shimura}
 G. Shimura, \textit{Introduction to the arithmetic theory of automorphic functions}, Kano Memorial Lectures $1$.


\end{thebibliography}

\end{document}

